% !TEX program = xelatex
\documentclass[12pt,a4paper,fontset=fandol]{ctexbook}

% ==================== 页面 ====================
\usepackage[top=2.5cm,bottom=2.5cm,left=2.5cm,right=2.5cm]{geometry}

% ==================== 字体 ====================
\usepackage{fontspec}
\usepackage{iffont}
% CJK 钉 fandol（两侧都取自 TeX 树里的同一份 FandolSong-Regular.otf）；
% mono 钉 DejaVu Sans Mono：CI 由 fontconfig 命中同名，本机没有该系统字体，
% 就用 MiKTeX dejavu 包里的同一版本 TTF（kpathsea 按文件名命中）。
\IfFontExistsTF{DejaVu Sans Mono}
  {\setmonofont{DejaVu Sans Mono}[Scale=MatchLowercase]}
  {\setmonofont[Extension=.ttf,
      UprightFont=DejaVuSansMono,
      BoldFont=DejaVuSansMono-Bold,
      ItalicFont=DejaVuSansMono-Oblique,
      BoldItalicFont=DejaVuSansMono-BoldOblique]
     {DejaVuSansMono}[Scale=MatchLowercase]}

% ==================== 颜色 ====================
\usepackage[dvipsnames,svgnames,x11names]{xcolor}
\definecolor{codebg}{HTML}{F5F5F5}
\definecolor{codeframe}{HTML}{E0E0E0}
\definecolor{cppkeyword}{HTML}{1A1AFF}
\definecolor{cppcomment}{HTML}{2E8B2E}
\definecolor{cppstring}{HTML}{B22222}
\definecolor{linenumgray}{HTML}{999999}
\definecolor{perfheadbg}{HTML}{35544B}
\definecolor{numgood}{HTML}{1A7F37}
\definecolor{numbad}{HTML}{B62324}
\definecolor{eventpurple}{HTML}{5B3E9D}

% ==================== listings ====================
\usepackage{listings}

% ---- 所有代码块共用的排版（先定义，各语言 style 组合它）----
\lstdefinestyle{codebase}{
  basicstyle=\small\ttfamily,
  keywordstyle=\color{cppkeyword}\bfseries,
  commentstyle=\color{cppcomment}\itshape,
  stringstyle=\color{cppstring},
  numberstyle=\tiny\color{linenumgray},
  numbers=left,
  frame=single,
  rulecolor=\color{codeframe},
  framesep=6pt,
  breaklines=true,
  tabsize=4,
  columns=flexible,
  keepspaces=true,
  backgroundcolor=\color{codebg},
  showstringspaces=false,
  escapeinside={(*@}{@*)},
}

\lstdefinestyle{cppstyle}{
  style=codebase,
  language=C++,
  morekeywords={constexpr,consteval,constinit,decltype,auto,
    concept,requires,co_await,co_yield,co_return,
    noexcept,override,final,nullptr,static_assert,
    template,typename,namespace,using,mutable,volatile,
    explicit,operator,friend,inline,virtual,
    thread_local,alignas,alignof},
  deletekeywords={int,char,bool,float,double,long,short,unsigned,signed,void,wchar_t},
  emph={size_t,int32_t,uint32_t,int64_t,uint64_t,ssize_t,
    string,vector,map,set,unique_ptr,shared_ptr,optional,variant,any,
    span,string_view,FILE,wchar_t,chrono,steady_clock,high_resolution_clock},
  emphstyle=\color{teal!80!black},
}
\lstset{style=cppstyle}

% ---- C：POSIX 计时、perf_event_open、__builtin_ 家族 ----
\lstdefinestyle{cstyle}{
  style=codebase,
  language=C,
  morekeywords={bool,FILE,size_t,restrict,inline},
  emph={size_t,int32_t,uint32_t,int64_t,uint64_t,ssize_t,clockid_t,
    timespec,timeval,itimerspec,mmap_perf_page},
  emphstyle=\color{teal!80!black},
}

% ---- 本册自建 shell 方言：perf、benchmark 与剖析器的命令行 ----
% 注意：不能借用邻册的 ShellLD，它只在 Audio Video Development 册内定义。
\lstdefinelanguage{ShellPM}{
  morekeywords={if,then,else,elif,fi,for,in,do,done,while,case,esac,
    function,return,local,export,source,echo,cd,set,exec,trap,read,
    printf,test,exit,shift,eval,alias},
  sensitive=true,
  morecomment=[l]{\#},
  morestring=[b]",
  morestring=[d]',
}
\lstdefinestyle{shellstyle}{
  style=codebase,
  language=ShellPM,
  keywordstyle=\color{cppkeyword}\bfseries,
  emph={perf,stat,record,report,annotate,mem,data,time,bench,trace,
    top,convert,buildid-cache,evlist,list,event,call-graph,period,dwarf,
    google_bench,runbench,hyperfine,wrk,tracy,vtune,amplxe,
    ncu,nvidia-smi,nsys,cuda-memcheck,valgrind,callgrind,heaptrack,
    gdb,lldb,perf2json,stackcollapse-perf,flamegraph.pl,inchworm,
    xperf,wpr,esrcli,instruments,xctrace,dtrace,osxpmem,powermetrics,
    taskset,numactl,lscpu,lstopo,ethtool,blktrace,fio,iostat,monitor,
    grep,sed,awk,xargs,jq,python3,pip,pip3,pwsh,bash,git,cmake,ninja,
    make,ccache,tar,mkdir,touch,ls,cat,wc,diff,find,sort,uniq,tee,rm,
    objdump,nm,readelf,strings,cxxflags},
  emphstyle=\color{teal!80!black}\bfseries,
}

% ---- 本册自建 CMake 方言：门禁拒收裸 language=CMake ----
\lstdefinelanguage{CMakePM}{
  morekeywords={cmake_minimum_required,project,add_executable,
    add_library,target_link_libraries,target_include_directories,
    target_compile_definitions,target_compile_options,set,if,elseif,
    else,endif,find_package,option,message,install,file,
    target_sources,add_subdirectory,enable_language,set_target_properties,
    FetchContent_Declare,FetchContent_MakeAvailable,list,string},
  sensitive=false,
  morecomment=[l]{\#},
  morestring=[b]",
}
\lstdefinestyle{cmakestyle}{
  style=codebase,
  language=CMakePM,
  keywordstyle=\color{cppkeyword}\bfseries,
}

% ---- 工具输出：perf report、火焰图折叠栈、benchmark 表格（无语言，纯文本）----
\lstdefinestyle{outstyle}{
  style=codebase,
  language={},
  numbers=none,
  frame=single,
  morecomment=[l][\color{eventpurple}\bfseries]{\#},
  emph={IPC,cache-misses,branch-misses,cycles,instructions,L1,TLB,
    p50,p90,p95,p99,p999,MIPS,GFLOPS,ns,us,ms,us/op,ns/op,B/s,iops},
  emphstyle=\color{numbad}\bfseries,
}

% ---- Python：统计与 bootstrap 脚本 ----
\lstdefinestyle{pythonstyle}{
  style=codebase,
  language=Python,
  morekeywords={self,lambda,async,await,yield,with,as,pass,del,
    global,nonlocal,is,in,not,and,or,True,False,None},
  emph={statistics,numpy,np,scipy,random,bootstrap,median,pandas,pd,
    matplotlib,plt,math,sys,os,json,time,subprocess,typing,dataclass},
  emphstyle=\color{teal!80!black},
}

% ---- 统一 diff：口径变更前后对比 ----
\lstdefinestyle{diffstyle}{
  style=codebase,
  numbers=none,
  frame=none,
  backgroundcolor=\color{white},
  morecomment=[l][\color{numbad}]{-},
  morecomment=[l][\color{numgood}]{+},
  morecomment=[l][\color{blue!70!black}\bfseries]{@@},
}

% ==================== tcolorbox ====================
\usepackage[most]{tcolorbox}

\newtcolorbox{guideline}[1][]{
  enhanced,breakable,
  colback=blue!5!white,colframe=blue!75!black,
  fonttitle=\bfseries,title=要点,#1,
  boxed title style={colback=blue!75!black},
  attach boxed title to top left={yshift=-2mm,xshift=2mm},
  arc=2mm,boxrule=0.8mm,
}
\newtcolorbox{compare}[1][]{
  enhanced,breakable,
  colback=green!3!white,colframe=green!50!black,
  fonttitle=\bfseries,title=对比,#1,
  boxed title style={colback=green!50!black},
  attach boxed title to top left={yshift=-2mm,xshift=2mm},
  arc=2mm,boxrule=0.8mm,
}
\newtcolorbox{warning}[1][]{
  enhanced,breakable,
  colback=red!5!white,colframe=red!75!black,
  fonttitle=\bfseries,title=常见陷阱,#1,
  boxed title style={colback=red!75!black},
  attach boxed title to top left={yshift=-2mm,xshift=2mm},
  arc=2mm,boxrule=0.8mm,
}
\newtcolorbox{note}[1][]{
  enhanced,breakable,
  colback=orange!5!white,colframe=orange!80!black,
  fonttitle=\bfseries,title=注意,#1,
  boxed title style={colback=orange!80!black},
  attach boxed title to top left={yshift=-2mm,xshift=2mm},
  arc=2mm,boxrule=0.8mm,
}
\newtcolorbox{bestpractice}[1][]{
  enhanced,breakable,
  colback=purple!5!white,colframe=purple!60!black,
  fonttitle=\bfseries,title=最佳实践,#1,
  boxed title style={colback=purple!60!black},
  attach boxed title to top left={yshift=-2mm,xshift=2mm},
  arc=2mm,boxrule=0.8mm,
}

% ==================== 章节标题 ====================
\usepackage{titlesec}
\usepackage{fancyhdr}
\usepackage{graphicx}
\usepackage{amssymb}
\usepackage{amsmath}
\usepackage{tabularx}
\usepackage{booktabs}
\usepackage{longtable}
\usepackage{array}
\usepackage{multirow}
\usepackage{enumitem}
\usepackage{seqsplit}
\usepackage{float}

\titleformat{\chapter}[display]
  {\normalfont\huge\bfseries\color{perfheadbg}}
  {\CTEXthechapter}{20pt}{\Huge}
\titleformat{\section}
  {\normalfont\Large\bfseries\color{blue!70!black}}
  {\thesection}{1em}{}
\titleformat{\subsection}
  {\normalfont\large\bfseries\color{blue!60!black}}
  {\thesubsection}{1em}{}

\pagestyle{fancy}
\fancyhf{}
\fancyhead[LE,RO]{\thepage}
\fancyhead[RE]{\leftmark}
\fancyhead[LO]{\rightmark}
\setlength{\headheight}{14.5pt}

\setlist[itemize]{leftmargin=2em,itemsep=2pt}
\setlist[enumerate]{leftmargin=2em,itemsep=2pt}
\emergencystretch=6em

% ==================== 超链接 ====================
\usepackage{hyperref}
\hypersetup{
  colorlinks=true,
  linkcolor=blue!70!black,
  urlcolor=blue!60!black,
  bookmarks=true,
  pdfauthor={C++ Reference},
  pdftitle={C/C++ 性能剖析与测量学参考手册：时钟、统计、剖析器与归因方法},
}

% ==================== 自定义命令 ====================
\newcommand{\cpp}[1]{C++#1}
\newcommand{\Fn}[1]{\texttt{#1()}}
\newcommand{\Tp}[1]{\texttt{#1}}
\newcommand{\Opt}[1]{\texttt{-\/-#1}}
\newcommand{\Bk}{\discretionary{}{}{}}
\newcommand{\Lib}[1]{\textbf{\texttt{#1}}}
\newcommand{\Lang}[1]{\textsf{#1}}
\newcommand{\CL}{\Lang{C}}
\newcommand{\CPP}{\Lang{C++}}
\newcommand{\PY}{\Lang{Python}}
% 口径标签：本册讲"同一次测量"时统一用这四个写法，避免同义多形。
% \MWall 墙钟时间 / \MCpu CPU 时间 / \MLat 延迟 / \MThr 吞吐
\newcommand{\MWall}{\Lang{Wall}}
\newcommand{\MCpu}{\Lang{CPU}}
\newcommand{\MLat}{\Lang{Latency}}
\newcommand{\MThr}{\Lang{Throughput}}
% 指标名一律用 \Mx{}，与 \Tp{}（代码里的标识符）区分开
\newcommand{\Mx}[1]{\textsf{#1}}

% ====================================================================
% book 双页模式默认 \flushbottom，会靠拉伸行距把每页撑满；本仓库框体与长表多，
% 撑页产生的 Underfull \vbox 是排版噪声而非缺陷，这里统一改用小底部对齐。
\raggedbottom

\begin{document}

\frontmatter

\begin{titlepage}
\centering
\vspace*{4cm}
{\Huge\bfseries C/C++ 性能剖析与测量学参考手册\par}
\vspace{1cm}
{\LARGE 时钟、统计、剖析器与归因方法\par}
\vspace{0.4cm}
{\Large perf、Google Benchmark、Tracy、VTune \\ 与 CUDA event 的跨平台实测口径\par}
\vfill
{\large \today\par}
\end{titlepage}

\tableofcontents

\mainmatter

\cleardoublepage

% ====================================================================
% 正文：五个部分、十五章加三个附录，全部就地排布（不再引外部片段）
% ====================================================================
\part{测量的地基：口径、时钟与统计}

\chapter{分界与成因：为什么``更快''必须先变成一个数}
\label{ch:pm01-why}

仓库里其它分册回答的是``怎么实现''：怎么把音频帧按时交付，怎么把渲染预算花完，怎么让
一段内核代码少拷贝一次内存。本册回答另一个问题——\textbf{当你要说``这版更快''时，你
到底凭什么说}。两者看起来相邻，实际是分界：前者是工程决策，后者是证据决策。一个补丁
完全可以``确实更快''却拿不出一个可复现的数，于是它在评审里活不过一次重新运行；也完全
可以``数字很漂亮''而负载、编译器、机器都和发布时不一样，于是它在生产里活不过第一个小时。

本册的三句地基贯穿后面每一章，这里先摆出来，后面每一章都在为其中某一句提供操作细节：
\begin{itemize}
  \item \textbf{测量先于优化}：没有一个可复现的数，``变快了''只是叙事，不是结论。
    支撑它的是重复、分布与判据（第~\ref{ch:pm03-stats} 章）。
  \item \textbf{口径先于数字}：同一个程序，\MWall、\MCpu、\Mx{IPC}、\Mx{p99} 会各自
    ``更快''，先说清你问的是哪一个。支撑它的是时钟源与计数器（第~\ref{ch:pm02-clock} 章）。
  \item \textbf{归因先于改动}：指标只能把你送到代码位置附近，从位置到原因需要另一套
    证据（布局、依赖、分支、拓扑）。把相关性当因果，是本册后半部分的主要病例。
\end{itemize}

本章不用任何工具。它把``性能问题''拆成三种真实提问并说明各自的证据形态，把四个口径的
定义与包含关系钉死，给出一个可手算的收益模型，最后用一张按症状的入口索引表收束。

\section{性能问题的三种真实提问}
\label{sec:pm01-question}

日常口语里只有一句``性能不行''，但它在证据层面是三道不同的题。答错种类比答错数值更常见，
也更贵。

\subsection*{提问一：``太慢了''}

这是一个\textbf{绝对性能}问题，隐含一句没说的判据：慢是相对于某个契约——交互要 $16$
毫秒一帧、请求要在 $100$ 毫秒内返回、批量任务要在今晚跑完。它需要的证据形态是\emph{延迟
分账}：把这 $100$ 毫秒拆成若干段（排队、准备、计算、拷贝、等待外设、收尾），指出哪一段
占大头，并说明剩下的段即使全部归零是否够。最常见的伪答案是\emph{一个函数名}：``慢在
序列化''。函数名只是位置，不是分账——序列化可能在等网络，可能在触缺页，也可能只是被误标
（内联后归属到调用者，这一点在第~\ref{ch:pm04-sampling} 章展开）。

\subsection*{提问二：``比上个版本慢''}

这是一个\textbf{差分}问题：它不关心绝对值，只关心两个二进制在同一负载下的差。它需要的
证据形态是\emph{配对与重复}：交替运行（A/B/A/B）、同样的负载、同样的构建配置、足够多的
样本，然后用区间而不是点值下结论（第~\ref{ch:pm03-stats} 章）。最常见的伪答案是\emph{各跑一次比大小}：
两次运行之差里，代码改动只占一部分，剩下的来自频率、页缓存冷热、调度邻居、后台索引任务。
在噪声达 $5\%$ 量级的机器上（这个量级必须在目标机上自己标定，方法同在第~\ref{ch:pm03-stats} 章），
``单跑对比''几乎总能给出一个方向，于是它看起来永远有效。

\subsection*{提问三：``能不能扛住 N''}

这是一个\textbf{容量与伸缩性}问题，它的证据不是平均值也不是单次成功运行，而是一条
\emph{负载扫描曲线}：并发或速率从低到高，每档记录 \MWall、\MCpu、\MThr{} 与 \Mx{p99}，
再回答三个子问题——拐点在哪、拐点之后是延迟崩还是吞吐崩、离饱和还剩多少余量。最常见的
伪答案是\emph{一次成功的压测}：``$200$ 并发跑通了''没有任何边界信息，它既可能是 $200$
已经超了而你在测自己的缓冲，也可能是队列已经很长、你把延迟搬进了队列（第~\ref{ch:pm09-concurrency} 章）。

三种提问对证据的要求完全不同，但口语里它们共用一句``性能有问题''。表
\ref{tab:pm01-question} 把三者的最小证据与常见误答并列，作为后续讨论的起点。

\begin{longtable}{>{\raggedright\arraybackslash}p{0.15\textwidth}
                  >{\raggedright\arraybackslash}p{0.20\textwidth}
                  >{\raggedright\arraybackslash}p{0.25\textwidth}
                  >{\raggedright\arraybackslash}p{0.20\textwidth}}
\caption{三种真实提问各自需要的证据形态与最常见的伪答案}\label{tab:pm01-question}\\
\toprule
\textbf{提问形态} & \textbf{它真正在问什么} & \textbf{最小证据（可交付）}
& \textbf{最常见的伪答案}\\
\midrule
\endfirsthead
\toprule
\textbf{提问形态} & \textbf{它真正在问什么} & \textbf{最小证据（可交付）}
& \textbf{最常见的伪答案}\\
\midrule
\endhead
\midrule
\multicolumn{4}{r}{\footnotesize 接下页}\\
\endfoot
\bottomrule
\endlastfoot
``太慢了'' &
延迟是否违反某个显式或隐式的契约 &
延迟分账：一段请求从上到下各段的耗时与占比，外加契约值本身 &
一个热点函数名；一个平均耗时数字；一张火焰图截图\\
``比上个版本慢'' &
同一负载下两个二进制之间是否存在超出噪声的差 &
交替重复的成对样本、差值区间、判据（什么算变差）、环境指纹 &
各跑一次比大小；只报均值；先跑新版本再跑旧版本\\
``能不能扛住 N'' &
在多大负载上仍能满足延迟目标、余量还剩多少 &
负载扫描曲线（每档给 \MWall、\MCpu、\MThr、\Mx{p99} 四个数）、拐点、饱和判据 &
一次成功的压测；``CPU 才用 $60\%$''；只看吞吐峰值\\
``哪里最值当改'' &
改动的期望收益与收益上限 &
收益模型：单次节省 $\times$ 调用次数 $\times$ 是否在关键路径 &
占比最高即最值当；把 \Mx{IPC} 低直接当成内存受限的结论\\
\end{longtable}

最后一行不属于三种提问，却是团队里最常见的第四种问法：它的证据形态是一个\textbf{算式
加一次验证}，而不是任何工具的输出。这正是第 \ref{sec:pm01-economics} 节的内容。

\section{口径先于数字：\MWall、\MCpu、\MThr{} 与 \Mx{p99}}
\label{sec:pm01-metric}

同一个程序、同一次运行，可以同时``变快''和``变慢''——只要口径选得不同。这不是诡辩，
而是四种量在测量物理上互不包含。先把定义钉死。

\textbf{\MWall{}（墙钟时间）}：调用方视角从起点到终点的时长，用单调时钟读取
（\Fn{clock{\Bk}\_gettime} 的 \Tp{CLOCK{\Bk}\_MONOTONIC}；Windows 侧对应
\Fn{Query{\Bk}Performance{\Bk}Counter}，macOS 对应 \Fn{mach{\Bk}\_absolute{\Bk}\_time}，
细节在第~\ref{ch:pm02-clock} 章）。它包含一切：真正在 CPU 上执行的指令、被调度出去的等待、缺页、锁自旋、
外设等待，以及其他线程尚未完成的部分。\textbf{用户感知的是 \MWall，不是别的}。

\textbf{\MCpu{}（CPU 时间）}：进程（或线程）实际占用 CPU 的时间总和。它可以\emph{大于}
\MWall：八个线程各跑一秒，\MCpu{} 约八秒而 \MWall{} 约一秒；也可以明显小于 \MWall：
一次以等待为主的 I/O 型负载里，两者的比值可以远低于 $0.5$。这种``利用率歧义''是第~\ref{ch:pm09-concurrency} 章
的主题，这里只需记住：\MWall{} 与 \MCpu{} 的比值既不是``效率''也不是``并行度''，它是
两者的混合，单独报它等于不报。

\textbf{\MThr{}（吞吐）}：单位时间完成的工作量，单位必须跟着负载定义走（op/s、帧/s、
GiB/s、请求/s）。吞吐隐含两个假设：请求可以并发、系统供给充足。没有饱和点信息的单点
吞吐数字，几乎无法解释。

\textbf{\MLat{} 与 \Mx{p99}}：延迟是单个操作的 \MWall{} 分布，\Mx{p99} 是它的 $0.99$
分位数。分位数不是``更严格的均值''，而是分布的一个切点：样本量不够时高尾分位数毫无意义
（要稳定估计 \Mx{p99}，样本量至少得能看见尾部若干个点，判据在第~\ref{ch:pm03-stats} 章）；多峰分布下
\Mx{p99} 甚至可能落在两峰之间的空隙里，``代表''一个并不存在的操作。

包含关系一句话说清：\MWall{} 包含本线程的 \MCpu{} 与全部等待；多线程程序的 \MWall{}
包含\emph{最慢那条关键路径}上的 \MCpu 之和，其余线程的 \MCpu{} 被遮蔽；\MThr{} 是
\MWall{} 的倒数乘以并发度，只在没有队列堆积时成立；\Mx{p99} 是 \MLat{} 分布上的一个点，
而 \MLat{} 本身是逐请求的 \MWall。

\subsection{为什么一个补丁能让 \MCpu{} 降而 \MWall{} 升}
\label{sub:pm01-diverge}

这不是玄学，至少四种机制可以单独或共同造成它。定稿任何``变快''的结论之前逐条排除：

\begin{enumerate}
  \item \textbf{并行度换了形态}。串行改多线程，总 \MCpu{} 通常\emph{增加}（同步、原子、
    每线程初始化）而 \MWall{} 减少；反过来，为减少竞争而收紧并发，\MCpu{} 减少而
    \MWall{} 上升。这类``打架''是同一件事的两面，两个口径必须同时报。
  \item \textbf{背压被搬走}。缩短一段处理时间，可能让上游更快把请求灌进下一段队列；
    下游没变则队列变长，于是单请求 \MLat{} 上升而 \MThr{} 不变甚至上升。吞吐与尾延迟
    同时``变好''和``变坏''是同一枚硬币，队列长度必须作为第三个指标一起记录（第~\ref{ch:pm09-concurrency} 章）。
  \item \textbf{运行时与分配器被改动}。减少分配次数、换内存池、关掉某个统计埋点，会降低
    当前线程的 \MCpu，却可能触发分配器后台回收、页错误、跨线程同步，把成本搬到别人身上；
    若被搬过去的那条线程正好在关键路径上，\MWall{} 就变差（分配器口径在第~\ref{ch:pm07-memory} 章）。
  \item \textbf{频率不是常数}。让 CPU 更忙（更少等待、更高指令密度）会改变 DVFS 与 turbo
    的判断，短任务与长任务的爬升时延不同；同一次改动可能让 \MCpu{} 降两三成却把执行频率
    拖低，于是 \MWall{} 升。这类效应只能靠环境指纹（governor、当前频率、温度策略）记录，
    否则下一位读者复现不出来（第~\ref{ch:pm02-clock} 章）。
\end{enumerate}

四种读数形态放在一起看更直观。下面的对比框给的是\textbf{形态}而不是数值：它告诉你哪种
组合该信哪个口径、哪种组合必须回去补证据。

\begin{compare}[title={四个口径同时变化时的五种典型形态}]
\small
\begin{tabular}{>{\raggedright\arraybackslash}p{0.17\textwidth}
                >{\raggedright\arraybackslash}p{0.30\textwidth}
                >{\raggedright\arraybackslash}p{0.28\textwidth}}
\toprule
\textbf{形态} & \textbf{四个口径的读数长什么样} & \textbf{该下什么结论}\\
\midrule
真整体变快 &
\MWall{} 降、\MCpu{} 降、\MThr{} 升、\Mx{p99} 持平或降 &
可以写``变快''，但四个数都要报，并附判据\\
并行化生效 &
\MWall{} 明显降、\MCpu{} 升、\MThr{} 升、\Mx{p99} 略升 &
两个口径一起报，说明``用核换时间''是否可接受\\
改在关键路径之外 &
\MCpu{} 降、\MWall{} 几乎不动、\Mx{p99} 不变 &
结论是``无可测收益''，不是``CPU 少了两成所以更快''\\
延迟被搬进队列 &
\MThr{} 升、\Mx{p99} 升、队列长度或积压升 &
不能只报吞吐；必须同时报尾延迟与积压\\
环境漂移 &
四个口径同向变化、与代码改动方向无关 &
先复现再下结论；查频率、邻居进程、缓存冷热\\
\bottomrule
\end{tabular}
\end{compare}

\subsection{把四个口径同时打印出来}
\label{sub:pm01-harness}

结论里的每一个数都应该由\emph{同一次运行}给出，否则你比较的是四次运行而不是四个口径。
清单 \ref{lst:pm01-harness} 是最小骨架：一次运行同时给出 \MWall、\MCpu、逐请求 \MWall{}
的 \Mx{p50}/\Mx{p99}/\Mx{max}，以及由 \MWall{} 推出的吞吐；只用标准库与时钟调用，方便
在任何平台上改。

\begin{lstlisting}[style=cppstyle,caption={四个口径同源于一次运行的最小测量骨架
（占位工作函数请换成你自己的被测代码）},label={lst:pm01-harness}]
// One run, four metrics. No external dependency on purpose.
#include <algorithm>
#include <cstdlib>
#include <ctime>
#include <cstdio>
#include <vector>

static double now_wall() {              // CLOCK_MONOTONIC: no NTP steps
  timespec ts{};
  clock_gettime(CLOCK_MONOTONIC, &ts);
  return ts.tv_sec + ts.tv_nsec * 1e-9;
}
static double now_cpu() {               // sum over all threads of the process
  timespec ts{};
  clock_gettime(CLOCK_PROCESS_CPUTIME_ID, &ts);
  return ts.tv_sec + ts.tv_nsec * 1e-9;
}

static volatile double sink = 0.0;      // keep the work from being erased
static void work_item(int i) {
  double acc = 0.0;
  for (int k = 0; k < 64; ++k) acc += (i * 31 + k) * 1.5;
  sink = acc;
}

int main(int argc, char** argv) {
  const int n = (argc > 1) ? std::atoi(argv[1]) : 200000;
  std::vector<double> lat;
  lat.reserve(static_cast<size_t>(n));

  const double w0 = now_wall(), c0 = now_cpu();
  for (int i = 0; i < n; ++i) {
    const double a = now_wall();
    work_item(i);
    lat.push_back(now_wall() - a);      // per-request wall latency
  }
  const double wall = now_wall() - w0;
  const double cpu  = now_cpu() - c0;

  std::sort(lat.begin(), lat.end());
  auto at = [&](double q) {
    size_t idx = static_cast<size_t>(q * double(lat.size() - 1));
    return lat[idx];
  };
  std::printf("wall=%.6f  cpu=%.6f  ratio=%.3f\n", wall, cpu, cpu / wall);
  std::printf("thr=%.1f  p50=%.9f  p99=%.9f  max=%.9f\n",
              n / wall, at(0.50), at(0.99), at(1.00));
  return 0;
}
\end{lstlisting}

骨架里有三处是刻意的，三处都是后面章节的伏笔：

\begin{itemize}
  \item 用 \Tp{CLOCK{\Bk}\_MONOTONIC} 而不是 \Tp{CLOCK{\Bk}\_REALTIME}：后者会被 NTP
    步进，因而测出负数耗时（六个 clock id 的语义差别列在第~\ref{ch:pm02-clock} 章的表里）。
  \item 用一个 \Tp{volatile} 的汇（sink）接住计算结果：否则编译器可能把整个循环判为死
    代码，你测到的就只是循环自身的结构（这一点在第~\ref{ch:pm15-antipattern} 章展开）。
  \item 逐请求延迟用\emph{同一次} \Fn{now\_wall} 的差值采集，但每次采集本身要花时间：
    计时调用进入被测路径会抬高 \Mx{p99}，这叫观测者效应。一次时钟读取通常在几十 ns 量级，
    取决于是否走 vDSO、内核版本与 CPU；想知道自己机器上的真实成本，把 \Fn{work\_item}
    换成空函数再跑一次，两次的差就是埋点成本（完整方法在第~\ref{ch:pm02-clock} 章）。
\end{itemize}

清单 \ref{lst:pm01-four-out} 是这个骨架的输出\emph{形状}。

\begin{lstlisting}[style=outstyle,caption={四个口径在一次运行中的读数——格式示例
（示意输出，非本仓库实测值；本册的构建与 CI 环境没有 PMU 权限、通常也没有 perf）},
label={lst:pm01-four-out}]
# Illustrative output: same workload, one machine, version B = patch applied.
# These numbers are shapes, NOT measurements taken in this repository.
[version A: baseline]
wall=12.480311  cpu=12.412003  ratio=0.995
thr=16025.3  p50=0.000061302  p99=0.000512004  max=0.003100991

[version B: after patch]
wall=11.022004  cpu=8.401772  ratio=0.762
thr=18143.8  p50=0.000048011  p99=0.000802331  max=0.004120884

# deltas A -> B:  cpu -32.9%   wall -11.7%   thr +13.2%   p99 +56.7%
# ratio dropped: the run now spends less of the wall time on cpu,
# i.e. more of it is waiting. Three metrics got better, one got worse,
# and the meaning of the ratio changed. Report all four, not a verdict.
\end{lstlisting}

这份示意输出值得逐行读：\MCpu{} 降了三成、\MWall{} 只降一成、吞吐升 $13.2\%$、而
\Mx{p99} 升了 $56.7\%$。任何单口径都能给出一个``结论''，而真实结论是：这个补丁把工作从
CPU 侧搬到了等待侧，并且牺牲了尾部。这笔交易划不划算，只有契约（提问一）能回答——这就
是口径必须先于数字说清的原因。

\section{优化的经济学：收益、占比与关键路径}
\label{sec:pm01-economics}

工具能告诉你``哪里花时间''，不会告诉你``改哪里值得''。后者是一道算术题，可以在开工前用
一支笔做完。基本式只有三个因子：

\textbf{期望收益 $=$ 单次节省 $\times$ 调用次数 $\times$ 是否在关键路径上}。

三个因子各自都能单独把收益压成零，所以三个都要有来源，三个都不能靠猜。

\subsection{先算上限：一个可手算的 Amdahl 式算例}
\label{sub:pm01-amdahl}

设一次请求的 \MWall{} 为 $1.00$ 单位，剖析显示某个函数占其中 $0.30$。把它加速到原来的
一半（$0.30$ 变 $0.15$），总时间变 $0.70 + 0.15 = 0.85$，端到端收益是
$1.00/0.85 \approx 1.18$，约 $18\%$；把它彻底删掉（$0.30$ 变 $0$），总时间 $0.70$，上限
$1.00/0.70 \approx 1.43$，约 $43\%$。为一个占 $30\%$ 的函数设计完美方案，天花板是
$43\%$，而合理方案能拿走的是那 $18\%$ 里的一部分。

反过来，一个只占 $3\%$ 的函数即使优化到零，收益上限也只有约 $3\%$——如果它还不在每轮都
要重复的关键路径上，这笔投入可能连一次复测的噪声都跨不过去。\textbf{占比给出上限，上限
决定投入的天花板}；这一条不需要任何工具权限，就能否决掉一半的``优化提议''。

\subsection{再算实际：为什么占比高的函数常常改了没用}
\label{sub:pm01-critpath}

上限是 \MCpu{} 口径的算法，实际收益是 \MWall{} 口径的算法，两者之间隔着一个关键路径。
设一个两阶段流水线：阶段二耗时 $0.50$、阶段三耗时 $0.70$，严格串行，端到端 $1.20$。把
阶段二优化到 $0.35$，端到端 $1.05$，收益 $12.5\%$；即使阶段二再降到 $0.20$，端到端
$0.90$，累计也只有 $25\%$——阶段三成了新的天花板。\textbf{瓶颈会迁移}：一次优化完成后
``第二名''接管，下一次的上限由这一次的结果决定，必须重算，不能外推。

如果阶段二与阶段三在\emph{不同线程}上并行，结论更极端：只要阶段二比阶段三少 $0.20$，
它无论怎么改，\MWall{} 一动不动；直到它降到 $0.50$ 以下、比阶段三更早完成，改动才第一次
显现在墙钟上。剖析器给的``占比''是 \MCpu{} 口径的占比，关键路径是 \MWall{} 口径的结构。
下面三类``占比高但不值得改''就是这么来的：

\begin{itemize}
  \item \textbf{被遮蔽的工作}：日志压缩、指标上报、后台序列化、分配器回收，占比不小，
    但都不在请求返回前必须完成的那条链上。
  \item \textbf{冷路径上的热点}：某函数占比高，只因为你压测时走了某个罕见分支（首次
    加载、异常处理、缓存重建）。负载定义一改，它就不在了。
  \item \textbf{被摊薄的公共函数}：\Fn{memcpy}、分配器、锁的慢路径常排在最前，它们是
    \emph{结果}而不是原因；真正要问的是谁把它们调用了这么多次（第~\ref{ch:pm06-micro} 章讲这类归因）。
\end{itemize}

\begin{guideline}[title={开工前先算的三个数}]
\small
\textbf{一、单次节省}：从剖析里读该段占比，换算成每请求的绝对时间，换算写成正文句子
（``占比 $0.30$、每请求 $1.2$ 毫秒，即约 $0.36$ 毫秒''），不要只报百分比。
\textbf{二、调用次数}：必须有来源——计数埋点、事件计数或折叠栈的命中次数（第~\ref{ch:pm05-events} 章），
不要用``估计''。次数要和负载定义绑定：每请求 $4$ 次与每秒 $4$ 次是完全不同的收益。
\textbf{三、是否在关键路径}：并发程序先画一张阶段依赖图（哪段完成后哪段才开始），标出从
请求进入到返回的最长链；不在链上的工作收益一律记零。三个数相乘，再和实现成本比较。
\end{guideline}

\section{本册的路线：定口径、拿数据、让结论活得久}
\label{sec:pm01-roadmap}

本册按三件事分块，顺序即依赖顺序：口径不清楚，数据无意义；数据不落地，归因是猜；结论不能
复现，前三步全部作废。

\begin{lstlisting}[style=outstyle,caption={本册三段式路线的框图（示意排版，不是任何
工具的输出，不含任何数值）},label={lst:pm01-route}]
  stage 1: define        stage 2: collect        stage 3: keep it alive
 +--------------------+ +--------------------+ +----------------------+
 | ch02 clocks        | | ch04 sampling      | | ch13 CI regressions  |
 | ch03 statistics    | | ch05 events, PMU   | | ch14 prod profiling  |
 | ch10 build configs | | ch06 cache, branch | | ch15 anti-patterns   |
 | metric definitions | | ch07 mem, ch08 io  | | apA three platforms  |
 | (what to measure)  | | ch09 concurrency   | | apB flags, fields    |
 +--------------------+ | ch11 gpu, ch12 bw  | +----------------------+
   metric + clock  --->  evidence  --->  a conclusion that survives review
\end{lstlisting}

各段的一句话职责：

\begin{itemize}
  \item \textbf{定口径}（第 2、3、10 章）：时钟源选错，后面每一步都在算错的东西；重复
    次数与判据说不清，``快 $3\%$''和``没有结论''是同一句话。第~\ref{ch:pm10-compiler} 章也在这里，因为它决定
    ``你测的是哪个二进制''——优化级别、断言开关、LTO 与 PGO 是口径问题，不是优化技巧。
  \item \textbf{拿数据}（第 4 至 12 章）：采样与计数两条线（第 4、5 章），归因三条线
    （第 6、7、9 章），I/O 与设备侧（第 8、11、12 章）。
  \item \textbf{让结论活得久}（第 13 至 15 章）：进 CI 当门禁、进生产当常驻剖析、进评审
    当一份可复现的结论，最后是反面清单——不可复现的结论长什么样。
\end{itemize}

\subsection{三平台与``示意输出''的声明}
\label{sub:pm01-platforms}

本册以 Linux 加 \Fn{perf} 为主线叙述，因为它的口径链条最完整可查；Windows 侧（ETW、
XPerf、VTune）与 macOS 侧（Instruments、\Tp{powermetrics}）在关键章各给一节对照，三平台
的完整并列放在附录 A。``主线是 Linux''不等于``只有 Linux 能测''：本册不写这类结论，也不
建议把只能在一个平台成立的判据写进跨平台门禁。

同样重要的一句：\textbf{本仓库的构建与发布用 CI 环境没有 PMU 权限，通常也没有 \Fn{perf}}。
因此本册所有工具输出清单——包括清单 \ref{lst:pm01-four-out}——都只是\emph{形态示意}，用来
教你怎么读字段，而不是给你一个可引用的数值。凡是 caption 标明``示意输出''或``格式示例
（非本仓库实测值）''的地方，都不要把结论建立在那组数字上。量级一律写成区间或``取决于''
的形式，并附读者在目标机上自己标定的方法；这是本册对可复现性的自我约束，不是行文懒惰。

\subsection{我该从哪一章进来}
\label{sub:pm01-entry}

表 \ref{tab:pm01-entry} 按症状索引全书，假定你手上\emph{没有}任何数据，只有一句
``不对劲''。第一列按口语写，第二列给该先读的章，第三列给今天就能做的一步。

\begin{longtable}{>{\raggedright\arraybackslash}p{0.27\textwidth}
                  >{\raggedright\arraybackslash}p{0.16\textwidth}
                  >{\raggedright\arraybackslash}p{0.34\textwidth}}
\caption{按症状的本册入口索引（手上没有数据时从哪里开始）}\label{tab:pm01-entry}\\
\toprule
\textbf{症状（口语）} & \textbf{先读} & \textbf{今天就能做的一步}\\
\midrule
\endfirsthead
\toprule
\textbf{症状（口语）} & \textbf{先读} & \textbf{今天就能做的一步}\\
\midrule
\endhead
\midrule
\multicolumn{3}{r}{\footnotesize 接下页}\\
\endfoot
\bottomrule
\endlastfoot
只知道``变慢了''，不知慢在哪个口径 &
第~\ref{ch:pm02-clock} 章 &
按清单 \ref{lst:pm01-harness} 跑一次，同时打印 \MWall、\MCpu、\Mx{p99}\\
两次跑出两成以上的差 &
第~\ref{ch:pm03-stats} 章 &
改成 A/B/A/B 交替各 $20$ 次，先看重复性再谈差异\\
想证明某个补丁确实更快 &
第~\ref{ch:pm03-stats} 章 &
先写下判据（要检出多大的差、什么算变差），再跑\\
热点里全是 \Fn{memcpy} 与分配器 &
第 4、7 章 &
先确认栈完整（帧指针还是解卷），再数调用次数\\
想看 \Mx{IPC} 与缓存未命中率 &
第 5、6 章 &
先确认本机有没有 PMU 权限；没有就明确退到软件事件\\
加了线程反而更慢 &
第~\ref{ch:pm09-concurrency} 章 &
扫线程数 $1,2,4,8$ 直到物理核数，每档同时报 \MWall{} 与 \MCpu\\
只争 \Tp{-O2} 还是 \Tp{-O3} &
第~\ref{ch:pm10-compiler} 章 &
先确认两次比较用的是\emph{不同}二进制（产物哈希进记录）\\
延迟没问题，吞吐上不去 &
第 8、9 章 &
做负载扫描找拐点，而不是在单一并发度上调参\\
磁盘读写``时快时慢'' &
第~\ref{ch:pm08-io} 章 &
把冷启动与热启动分开测，先确定页缓存有没有参与\\
GPU 段测出 $0.001$ 毫秒 &
第~\ref{ch:pm11-gpu} 章 &
检查是不是漏了同步：kernel 是排队执行，不是立即执行\\
显存或内存带宽不知够不够 &
第~\ref{ch:pm12-bandwidth} 章 &
先算算术强度，判断该跟带宽比还是跟算力比\\
CI 基准忽红忽绿 &
第~\ref{ch:pm13-ci} 章 &
环境指纹入库；确认 runner 权限能否拿到硬件事件\\
线上不能停进程，想看它在忙什么 &
第~\ref{ch:pm14-production} 章 &
先定观测预算（能接受多大比例的开销），再选手段\\
结论被一句``你环境不一样''推翻 &
第~\ref{ch:pm15-antipattern} 章 &
按第 \ref{sec:pm01-trouble} 节的起手式重写结论\\
要在 Windows 或 macOS 上做同样的事 &
附录 A &
查三平台的采样与计数能力对照、权限确认命令\\
\end{longtable}

第三列刻意只写``今天能做的一步''，不写``这样就能得到答案''：本册没有一步到位的捷径，只有
把口径、样本、归因三件事分别做对。

\section{排障与速查：三类开局错误}
\label{sec:pm01-trouble}

本章以负面清单收束——本册后面各章里绝大部分返工，根源都是下面三个开局动作之一。

\begin{warning}[title={三类一开工就把结论作废的做法}]
\small
\textbf{一、没有基线就开工。}``先改改看''意味着连差值的参考点都没有：改动之后的数既是新值
又是首值，无法与任何历史比较。正确顺序是先固定负载与口径，跑出一段基线序列（哪怕只有
$10$ 次），再动代码。\textbf{二、在调试构建上测。}\Tp{-O0} 或带断言的构建里，函数调用、
迭代器检查、边界检查都存在，热点分布与发布产物完全不同；更糟的是``调试构建里成立的直觉''
会让人把优化押在发布版根本不热的路径上。调试构建可以验证\emph{正确性}与\emph{趋势}，写进
结论的数字必须来自发布配置（第~\ref{ch:pm10-compiler} 章）。\textbf{三、拿 sanitizer 跑的数字当性能数。}地址、
线程、未定义行为这三类 sanitizer 会在内存访问或同步点插桩，开销量级远超一般优化空间，还会
改变缓存占用与线程时序，即\emph{改变被测对象本身}。sanitizer 的用法与配置不在本册重讲：
见《内存泄漏检测与防御》册；CI 里作为门禁的 clang-tidy 与 sanitizer 组合见《LSP 与代码
质量工具》第 12 章。本册只把它们当\textbf{测量对象}讨论——它们自身的开销量级，以及为什么
``开着 sanitizer 变慢了''不构成性能结论（第~\ref{ch:pm07-memory} 章）。
\end{warning}

还有一条与内联有关的开局错误值得在这里就写清，免得各章口径不一：\textbf{\Tp{inline}
首先是 ODR 与链接语义，不是提速开关}。有没有被展开由编译器启发式、优化级别、LTO、PGO
决定；\Tp{noinline} 是硬保证，而 \Fn{always\_inline} 与 \Tp{\_\_force\_inline} 只是尽力
而为（GCC 可能报错、Clang 可能警告、MSVC 的文档明言不保证）；``到底有没有展开''只能在
\emph{发布产物}上反汇编确认，在探针构建上看不算数。任何``加个 \Tp{inline} 就快''的写法都
是错的——本册与《代码风格参考手册》在这一点上口径相同，机制细节在本册
第~\ref{ch:pm10-compiler} 章。

\subsection{正确起手式}
\label{sub:pm01-open}

\begin{bestpractice}[title={开工前的八步起手式（每次性能任务重跑一遍）}]
\small
\textbf{1. 写下提问类型}：``太慢了''／``比上个版本慢''／``能不能扛住 N'' 三选一，写进
工单标题。\textbf{2. 写下口径}：结论里每个数都用第 \ref{sec:pm01-metric} 节定义的四个口径
词之一标注，不许出现裸的``性能''与``耗时''。\textbf{3. 写下判据}：要检出多大的差、重复
多少次、什么算变差、什么算无结论——先写后跑（第~\ref{ch:pm03-stats} 章）。\textbf{4. 固定负载定义}：输入、
并发度、预热策略、结束条件，写成可执行的一行命令，而不是文字描述。\textbf{5. 固定构建
配置}：确认优化级别、断言开关、目标二进制哈希；``改了代码''不等于``改了二进制''。
\textbf{6. 采环境指纹}：CPU 型号与频率策略、内核版本、编译器版本、机器名、是否独占、绑核
与 NUMA 设置（第~\ref{ch:pm02-clock} 章）。\textbf{7. 确认权限与可用性}：硬件事件能否拿到；拿不到就写明
``本轮只用软件事件''，绝不把工具报出的 ``not supported'' 连同零值当成``开销为零''。
\textbf{8. 先跑基线再改动}：基线与候选交替运行，产物、日志与判据一起归档。
\end{bestpractice}

\subsection{本章带走四句话}
\label{sub:pm01-takeaway}

\begin{enumerate}
  \item 先分清提问：绝对性能、版本差分、容量伸缩三者的证据形态互不通用，答错种类比答错
    数值更贵。
  \item 先定口径：\MWall{} 是用户感知，\MCpu{} 是投入，\MThr{} 是产出，\Mx{p99} 是分布
    切点；一次运行同时打印四个，永远不要只报其中一个。
  \item 先算经济：单次节省 $\times$ 调用次数 $\times$ 是否在关键路径，三个因子各有来源才
    构成收益估计；占比只给上限，不给最优。
  \item 先声明可复现性：环境指纹、构建哈希、判据与样本量一起进结论；缺任一项，结论在别人
    那里活不过一天。
\end{enumerate}

第~\ref{ch:pm02-clock} 章处理这四句话里最底层的一件：这些口径的数到底从哪个计数器读出来，以及为什么``读
一次时间''本身就值几十纳秒到几微秒的注意力；第~\ref{ch:pm03-stats} 章把``一个数''变成``一个可以辩护的差值''。

% 本章口径：所有清单里的数字都只是形态示意（非本仓库实测值），量级一律写成区间并给出自测方法。

\chapter{时钟源与计数器：从 clock\_gettime 到 TSC 与 GPU 时钟域}
\label{ch:pm02-clock}

上一章把``更快''拆成了 \MWall（墙钟）、\MCpu（CPU 时间）、\MThr（吞吐）与 \Mx{p99}（尾延迟）
几个互不覆盖的口径（第\ref{ch:pm01-why}章）。本章解决下一个绕不开的问题：这些数到底是用什么
量出来的。计时听起来是全书最不需要讨论的部分——调一次 \Fn{clock{\Bk}\_gettime}、做一次减法，
不就完了？但恰恰是最简单的动作里埋着最多能毁掉结论的坑：选错 clock id 会量出负数耗时；
把 \Fn{rdtsc} 当纳秒用会把一个时代的频率差异读成性能差异；在多线程程序里用 \MWall 声称
``\MCpu 变少了''；笔记本合上盖子再打开，整段测量窗口就已经作废。本章把``用什么计数、
它随什么变、一次调用多少钱、什么情况下它会说谎''四件事逐条钉死，形成第\ref{ch:pm03-stats}章
做统计之前必须备好的那把尺子。

和后面各章的接口在哪里：本章只管时钟与计数器本身。重复次数与置信区间是第\ref{ch:pm03-stats}章；
剖析器自己的计时（\Tp{task-clock}、\Tp{cpu-clock} 事件）在第\ref{ch:pm04-sampling}章与
第\ref{ch:pm05-events}章；并发下 \MCpu 与 \MWall 的歧义在第\ref{ch:pm09-concurrency}章；
GPU 侧的 event 协议在第\ref{ch:pm11-gpu}章。三平台对照在本章各留一节，完整并列见附录
\ref{ch:pmapA-platform}。

\section{单调、稳定、原始：clock id 的语义与选错后果}
\label{sec:pm02-clockid}

\CL 的 \Fn{clock{\Bk}\_gettime} 与 \Fn{clock{\Bk}\_getres}（\Tp{<time.h>}）是 Linux 用户态计时的
事实入口。关键不在函数，而在第一个参数 \Tp{clockid{\Bk}\_t}：同一个函数，换一个 id，
量的就是完全不同的物理量。先记住三条主线。

\textbf{单调性。} \Tp{CLOCK{\Bk}\_REALTIME} 是系统的``墙''，管理员改时间、NTP 的整秒步进、
虚拟机恢复时的强制校正，都会让它\emph{向后跳}。拿它做区间测量，得到的负数不是 bug，是语义。
\Tp{CLOCK{\Bk}\_MONOTONIC} 保证不回退，但它仍然被 NTP 的\emph{频率调整}（把时钟拉慢一点再拉回来）
平滑修正——它单调，但不保证``走速''恒定；而且在 Linux 上\textbf{挂起期间的时间不计入}。
\Tp{CLOCK{\Bk}\_MONOTONIC{\Bk}\_RAW} 直接挂在硬件计数上，步进、slew、调频都不影响它，代价是它和
墙钟之间的速率差会随时间漂移，不能拿来回答``现在几点''。

\textbf{对 suspend 的行为。} \Tp{CLOCK{\Bk}\_BOOTTIME} 把挂起时间也计进去，是``从开机算起''的口径；
跨笔记本合盖、嵌入式深度睡眠的测量必须用它，否则你会得到一段凭空消失的耗时。

\textbf{CPU 时间。} \Tp{CLOCK{\Bk}\_PROCESS{\Bk}\_CPUTIME{\Bk}\_ID} 与
\Tp{CLOCK{\Bk}\_THREAD{\Bk}\_CPUTIME{\Bk}\_ID} 量的不是流逝时间，而是``这个进程/线程实际占到 CPU 的时间''，
只在被调度运行时前进。它们是本章第四节 \MCpu 口径的来源之一，绝不能和 \MWall 混用。

\textbf{粗粒度家族。} \Tp{CLOCK{\Bk}\_MONOTONIC{\Bk}\_COARSE} 与 \Tp{CLOCK{\Bk}\_REALTIME{\Bk}\_COARSE}
用内核 tick 的粒度换更低的读数成本。对几十微秒级的测量它们会静默吃掉尾数；
对只要求毫秒分辨率的定时器回调，它们反而省一次高精度路径。选不选，取决于你要求的分辨率，见下一节。

表 \ref{tab:pm02-clockid} 把六组 id 汇总成一张可查的表。列口径：随什么变、会不会跳、适合测什么。

{\begin{longtable}{@{}>{\raggedright\arraybackslash}p{0.19\textwidth} >{\raggedright\arraybackslash}p{0.22\textwidth} >{\raggedright\arraybackslash}p{0.14\textwidth} >{\raggedright\arraybackslash}p{0.27\textwidth}@{}}
\caption{POSIX clock id 一览：随什么变、会不会跳、适合测什么}\label{tab:pm02-clockid}\\
\toprule
\textbf{clock id} & \textbf{随什么变} & \textbf{会不会跳} & \textbf{适合测什么；选错的后果}\\
\midrule
\endfirsthead
\toprule
\textbf{clock id} & \textbf{随什么变} & \textbf{会不会跳} & \textbf{适合测什么；选错的后果}\\
\midrule
\endhead
\midrule
\bottomrule
\endlastfoot
\Tp{CLOCK{\Bk}\_REALTIME} & 系统时间：NTP 步进与 slew、\Fn{clock{\Bk}\_settime}、时区无关（内部是 UTC） &
会，可前可后 & 与外部事件对齐时刻。测区间会得到负数或虚大；一切``耗时''都不该用它\\
\Tp{CLOCK{\Bk}\_MONOTONIC} & 单调流逝；受 NTP 频率调整（slew）；不计挂起时间 & 不会倒退 &
默认的工作负载计时。长窗口下走速被 NTP 轻微修正，通常可接受；跨 suspend 会丢时间\\
\Tp{CLOCK{\Bk}\_MONOTONIC{\Bk}\_RAW} & 纯硬件计数，不受 slew/步进/调频 & 不会（但走速恒定不代表与墙钟同速） &
校准计时器自身、长窗口测周期。拿它当``秒''会漂\\
\Tp{CLOCK{\Bk}\_BOOTTIME} & 单调流逝加挂起时间 & 不会 & 跨睡眠/深度低功耗的窗口（笔记本、嵌入式）；
忘了用它，挂起段直接消失\\
\Tp{CLOCK{\Bk}\_PROCESS{\Bk}\_CPUTIME{\Bk}\_ID} & 本进程所有线程占用的 CPU 时间 & 不会 &
回答``烧了多少 CPU''。只前进于被调度的时间，不能当 \MWall 用\\
\Tp{CLOCK{\Bk}\_THREAD{\Bk}\_CPUTIME{\Bk}\_ID} & 本线程占用的 CPU 时间 & 不会 &
每线程 CPU 账；多线程程序里对总账要自己加\\
\Tp{CLOCK{\Bk}\_MONOTONIC{\Bk}\_COARSE} 等 & 同对应精细版本，但按 tick 粒度更新 & 不会 &
毫秒分辨率够用时更便宜；对微秒级测量会静默量化掉尾数\\
\end{longtable}}

后果最常见的三种错法，各自的症状长这样：
其一，NTP 步进：用 \Tp{CLOCK{\Bk}\_REALTIME} 的常驻服务在凌晨校时瞬间算出负耗时或巨大耗时；
其二，suspend：晚上合盖、早上打开，\Tp{CLOCK{\Bk}\_MONOTONIC} 测出的``运行 6 小时''实际只跑了 20 分钟；
其三，频率调整：用 \Tp{CLOCK{\Bk}\_MONOTONIC} 跑数天窗口、再和 \Tp{CLOCK{\Bk}\_MONOTONIC{\Bk}\_RAW}
或硬件计数对账，发现速率差——这正是下一节 ``分辨率、精度、准确度三词不许混''里准确度那个问题。

先动手确认自己机器上有哪些 id、名义分辨率各是多少（清单 \ref{lst:pm02-clockprobe}）。

\begin{lstlisting}[style=cstyle, title={clock id 探针：存在性与名义分辨率}, label={lst:pm02-clockprobe}]
#include <stdio.h>
#include <stddef.h>
#include <time.h>

/* 在这台机器上逐个试探 clock id；clock_getres 失败通常意味着该 id 不存在。 */
int main(void) {
  struct { const char *name; clockid_t id; } clocks[] = {
#ifdef CLOCK_REALTIME
    {"CLOCK_REALTIME", CLOCK_REALTIME},
#endif
#ifdef CLOCK_MONOTONIC
    {"CLOCK_MONOTONIC", CLOCK_MONOTONIC},
#endif
#ifdef CLOCK_MONOTONIC_RAW
    {"CLOCK_MONOTONIC_RAW", CLOCK_MONOTONIC_RAW},
#endif
#ifdef CLOCK_BOOTTIME
    {"CLOCK_BOOTTIME", CLOCK_BOOTTIME},
#endif
    {"CLOCK_PROCESS_CPUTIME_ID", CLOCK_PROCESS_CPUTIME_ID},
    {"CLOCK_THREAD_CPUTIME_ID",  CLOCK_THREAD_CPUTIME_ID},
#ifdef CLOCK_MONOTONIC_COARSE
    {"CLOCK_MONOTONIC_COARSE", CLOCK_MONOTONIC_COARSE},
#endif
  };
  for (size_t i = 0; i < sizeof clocks / sizeof clocks[0]; ++i) {
    struct timespec res;
    if (clock_getres(clocks[i].id, &res) == 0) {
      printf("%-28s res=%ld.%09ld s\n", clocks[i].name,
             (long)res.tv_sec, (long)res.tv_nsec);
    } else {
      printf("%-28s unavailable\n", clocks[i].name);
    }
  }
  return 0;
}
\end{lstlisting}

\begin{note}[title={\CPP{} 的 \Tp{std::chrono} 映射到哪个 id：实现定义，必须自己确认}]
\Tp{steady{\Bk}\_clock} 在 Linux 主流实现上\emph{通常}落在 \Tp{CLOCK{\Bk}\_MONOTONIC}，
\Tp{system{\Bk}\_clock} 对应 \Tp{CLOCK{\Bk}\_REALTIME}，而 \Tp{high{\Bk}\_resolution{\Bk}\_clock}
只保证是两者之一（甚至可以是别名），\textbf{不保证}是 \Tp{RAW}。映射关系随标准库与版本变。
确认方法：同一瞬间各取一次 \Tp{steady{\Bk}\_clock::now()} 的纳秒值与 \Fn{clock{\Bk}\_gettime}
（\Tp{CLOCK{\Bk}\_MONOTONIC}）的读数，差值稳定说明同源，随后再对 \Tp{RAW} 做一次；
也可以直接读标准库实现。把``用了 \Tp{steady{\Bk}\_clock} 就等于单调且不受 NTP''当前提的结论，
先补上这一步再说。
\end{note}

\section{调用成本与解析度：vDSO、rdtsc，以及三词之辨}
\label{sec:pm02-cost}

\MWall 的测量误差有两半：时钟本身的性质，和\emph{取时钟这个动作}的成本。后者在被计时区间很短时
会直接混进答案。

\textbf{vDSO 路径。} 现代 Linux 把 \Fn{clock{\Bk}\_gettime} 的高频分支编进内核映射到用户态的
vDSO 页：普通线程取 \Tp{CLOCK{\Bk}\_MONOTONIC} 不需要系统调用，成本接近一次函数调用加几条访存指令——
常见形态是几十纳秒量级，具体取决于内核版本、CPU 频率、是否走 vDSO，以及有没有别的东西
（ptrace、部分剖析器、老 glibc、某些容器/seccomp 组合）把它打回真系统调用。
自查方法：\Tp{strace} \Tp{-c} 跑一次被测程序，若 \Tp{clock{\Bk}\_gettime} 出现在系统调用统计里，
你付的就是系统调用的价，量级通常高一个档。vDSO 的快还依赖内核维护的``时间基''数据：
CPU 频率变化、跨核计数不同步这类因素都会让这段逻辑退化，所以``几十纳秒''只能当作
读者自己复测前的预期，不是结论。

\textbf{\Fn{rdtsc}/\Fn{rdtscp}。} \Fn{rdtsc} 读时间戳计数器，几条指令的开销，比任何 OS API 都便宜；
但它不是函数调用语义：读出的数既没有确定的单位（见第三节），也\textbf{不保证顺序}——它会乱序完成，
可能被重排到被测代码之后执行。经典（笨重）的序列化办法是前后各插一条 \Tp{cpuid}：它把所有前序
指令执行完才返回，顺手清掉乱序窗口。代价是 \Tp{cpuid} 本身几百个周期、还会踩一堆 MSR，
把它插进微基准，你测到的就掺进了序列化的开销。更细的做法是 \Tp{lfence} 加 \Fn{rdtsc} 读前界、
\Fn{rdtscp} 加 \Tp{lfence} 读后界；\Fn{rdtscp} 保证``前序都完成''，但挡不住后序指令上移，所以后面还要补
一条栅栏。ARM 上没有 TSC：读通用定时器要用 \Tp{cntvct{\Bk}\_el0} 并遵守先读 \Tp{cntfrq} 与 \Tp{isb}
的配对，习惯别照搬。

\begin{guideline}[title={分辨率、精度、准确度：三个词不许混用}, label={box:pm02-3words}]
\textbf{分辨率}（resolution）是计时器能表示的最小刻度：\Fn{clock{\Bk}\_getres} 告诉你的是这个。
分辨率是 1\,ns 只说明``能写出 1\,ns 的数''，不说明每个 1\,ns 都可信。\newline
\textbf{精度}（precision）是单次读数对真值的随机误差，以及两次读数之间能分辨的最小\emph{区间}——
一个调用开销 30\,ns 的时钟，去量 50\,ns 的函数，有效数字所剩无几。\newline
\textbf{准确度}（accuracy）是计数速率与真值的系统偏差：NTP slew 修走的、RAW 留下的漂移、
TSC 换算常数错了、VM 迁移后频率变了，全都是准确度问题——它们不会让数字抖动，只会让数字\emph{稳定地错}。
\end{guideline}

注意一个容易上当的细节：\Fn{clock{\Bk}\_getres} 对 \Tp{CLOCK{\Bk}\_PROCESS{\Bk}\_CPUTIME{\Bk}\_ID}
一类的 CPU 时钟常报 1\,ns，但底层记账实际在调度 tick/上下文切换处更新——名义分辨率和真实更新粒度
是两回事（第四节展开）。

清单 \ref{lst:pm02-overhead} 给出``读者自己量一次调用开销''的完整骨架：累积式计时、防死代码消除、
基线扣除。关键点有三个：其一，编译器完全可以把一个不被观察的 \Tp{now()} 循环整个删掉或把调用折叠，
所以读数必须被 \Fn{Do{\Bk}Not{\Bk}Optimize} 一类的依赖屏障``压住''；其二，空循环基线（同一循环体但不取时）
要单独跑一次再相减，差值才是``一次调用''的量级——差值本身随机器、内核、glibc 变，量级请以你的
复测为准；其三，Google Benchmark 等框架把重复与时间预算交给你：\Opt{benchmark{\Bk}min{\Bk}time} 这类
选项的默认值只是保守起步，不是测量协议；重复次数与判据归第\ref{ch:pm03-stats}章。

\begin{lstlisting}[style=cppstyle, title={自测一次取时调用的开销：反死代码消除与基线扣除}, label={lst:pm02-overhead}]
// 编译参考：先 -O2，再 -O0 各跑一次看常数项是否变化。
// 注意：这里的数字不是结论，量级由你自己的机器给出。
#include <chrono>
#include <cstdint>
#include <cstdio>

// 依赖屏障：编译器既不能删除 value，也不能把相邻调用挪到它另一侧。
// GCC/Clang 的 inline asm 写法；MSVC 下用 volatile 累加加 _mm_sfence 等价物。
template <class T>
void do_not_optimize(const T& value) {
  asm volatile("" : : "r,m"(value) : "memory");
}

using my_clock = std::chrono::steady_clock;  // 映射到哪个 id 见第一节的 note

// 累 N 次取时；sink 保证每次读数都被消费。
std::uint64_t timed_sum(std::uint64_t n) {
  std::uint64_t sink = 0;
  for (std::uint64_t i = 0; i < n; ++i) {
    auto ticks = my_clock::now().time_since_epoch().count();
    do_not_optimize(ticks);
    sink += static_cast<std::uint64_t>(ticks);
  }
  return sink;
}

// 同一循环但不取时：这就是要扣掉的基线。
std::uint64_t baseline_sum(std::uint64_t n) {
  std::uint64_t sink = 0;
  for (std::uint64_t i = 0; i < n; ++i) {
    do_not_optimize(i);
    sink += i;
  }
  return sink;
}

std::uint64_t elapsed_ns(std::uint64_t (*f)(std::uint64_t), std::uint64_t n) {
  auto t0 = my_clock::now();
  std::uint64_t s = f(n);
  auto t1 = my_clock::now();
  do_not_optimize(s);
  return static_cast<std::uint64_t>(
      std::chrono::duration_cast<std::chrono::nanoseconds>(t1 - t0).count());
}

int main() {
  constexpr std::uint64_t kN = 10000000;
  (void)timed_sum(kN);            // warmup：频率爬升、分支预测、页就位
  std::uint64_t wall_timed = elapsed_ns(timed_sum, kN);
  std::uint64_t wall_base = elapsed_ns(baseline_sum, kN);
  std::printf("per-call upper bound: %llu ns\n",
              static_cast<unsigned long long>(wall_timed / kN));
  std::printf("per-call minus loop baseline: %lld ns\n",
              static_cast<long long>(wall_timed) - static_cast<long long>(wall_base));
  return 0;
}
\end{lstlisting}

清单 \ref{lst:pm02-overheadout} 是一份\emph{形态}示意：清单 \ref{lst:pm02-clockprobe} 的输出大致长什么样、
每行说明什么。数值是占位性质的格式示例，非本仓库实测值。

\begin{lstlisting}[style=outstyle, title={格式示例（非本仓库实测值）：clock id 探针的示意输出}, label={lst:pm02-overheadout}]
# 示意输出：具体数值取决于内核版本、HZ 配置与 vDSO 路径
CLOCK_REALTIME                 res=0.000000001 s
CLOCK_MONOTONIC                res=0.000000001 s
CLOCK_MONOTONIC_RAW            res=0.000000001 s
CLOCK_BOOTTIME                 res=0.000000001 s
CLOCK_PROCESS_CPUTIME_ID       res=0.000000001 s   # 名义值；真实更新粒度更粗
CLOCK_THREAD_CPUTIME_ID        res=0.000000001 s
CLOCK_MONOTONIC_COARSE         res=0.001000000 s   # 受内核 tick 频率约束
\end{lstlisting}

\section{TSC 的三件事，以及另两个平台}
\label{sec:pm02-tsc}

\Fn{rdtsc} 读的是 Time Stamp Counter：一个每核 64 位计数器。它能不能当时间用，取决于三个相互
独立的性质，一个都不能想当然。

\textbf{其一，恒定频率（\Tp{constant{\Bk}\_tsc}）。} 早期 x86 的 TSC 跟核心倍频走：降频、休眠、
turbo 都会改变每个 tick 的真实时长。现代大部分 x86 实现了``TSC 以恒定速率计数''——
计数速率与核心频率脱钩，这是 \Fn{rdtsc} 可用的前提，但不是全部平台都保证；虚拟化与老硬件
都可能破坏它。\textbf{其二，不眠（\Tp{nonstop{\Bk}\_tsc}）。} 深睡眠状态下计数器是否继续走。
Linux 用这个标志决定是否允许 \Tp{CLOCK{\Bk}\_MONOTONIC} 以 TSC 为底跨低功耗区间。
\textbf{其三，跨核同步（ART/ADSC 类机制）。} 每个核各有一个计数器。即便都恒定，
核与核之间也可能有偏移；新平台用 Always Running Timer / ADSC 一类机制让计数在低功耗下不断、
内核在 CPU 上线时校验一致性。校验失败，内核会把 TSC 标记为不可靠（\Tp{Marking TSC unstable}），
把 clocksource 换成 \Tp{hpet} 或 \Tp{acpi{\Bk}-pm}——但注意：这只影响 \Fn{clock{\Bk}\_gettime} 的底层实现，
你手写的 \Fn{rdtsc} 照样读到那个\emph{已经失步}的计数器。这是``内核告诉你它不信任 TSC，
而你未必听得见''的典型例子。

在 Linux 上确认三件事（清单 \ref{lst:pm02-tsccheck}）：标志位看 \Tp{/proc{\Bk}/cpuinfo} 的 flags 行
（\Tp{constant{\Bk}\_tsc}、\Tp{nonstop{\Bk}\_tsc}），内核是否信任看 \Tp{current{\Bk}\_clocksource}，
校验失败记录进 \Fn{dmesg} 找 \Tp{tsc}。换算常数（每 tick 多少 ns）来源更碎：有的内核暴露
\Tp{tsc{\Bk}\_khz}，新平台可用 CPUID leaf \Tp{0x15} 给出核心晶体频率比，都没有就拿
\Tp{CLOCK{\Bk}\_MONOTONIC{\Bk}\_RAW} 自校准一次。在虚拟机里三者都可能失真：宿主机调频、迁移、
嵌套虚拟化都会引入偏移，最诚实的判断方法是实测两个读数源的一致性，而不是读文档。

\begin{lstlisting}[style=shellstyle, title={Linux：TSC 三性确认命令（形态示意，路径随内核变）}, label={lst:pm02-tsccheck}]
grep -m1 'model name' /proc/cpuinfo
grep -c 'constant_tsc'  /proc/cpuinfo   # 期望等于在线 CPU 数；少了就是有的核不保证
grep -c 'nonstop_tsc'   /proc/cpuinfo
cat /sys/devices/system/clocksource/clocksource0/current_clocksource
dmesg | grep -iE 'tsc|clocksource' | tail -n 20   # Marking TSC unstable 在这里
# 换算常数：以下来源逐一试，哪个存在用哪个，都不在就自校准
ls /sys/devices/system/cpu | grep -i tsc          # tsc_khz / tsc_freq_khz，未必暴露
\end{lstlisting}

\begin{bestpractice}[title={决定直接用 \Fn{rdtsc} 之前，逐项打勾}, label={box:pm02-tsc-rules}]
（1）\Tp{constant{\Bk}\_tsc} 与 \Tp{nonstop{\Bk}\_tsc} 在\emph{每个}在线核上都存在；
（2）\Tp{current{\Bk}\_clocksource} 是 \Tp{tsc}，且 \Fn{dmesg} 无 unstable 记录；
（3）换算常数来自本机确认的来源之一，并随结论一起记录；
（4）读前界用 \Tp{lfence} 或 \Tp{cpuid}、读后界用 \Fn{rdtscp} 加 \Tp{lfence}，且栅栏开销已计入基线；
（5）多线程场景确认过跨核偏移（同一瞬间两核各读一次，差值在承诺精度内）；
（6）跨 suspend/迁移/长时间窗口的场景改用 OS API（\Tp{CLOCK{\Bk}\_MONOTONIC{\Bk}\_RAW} 或 \Tp{BOOTTIME}），
不要手写 \Fn{rdtsc}。六条全过，\Fn{rdtsc} 才配进热路径；缺任何一条，先回到 \Fn{clock{\Bk}\_gettime}。
\end{bestpractice}

\subsection{Windows：QPC 不是``周期计数''}
\label{sub:pm02-qpc}

Windows 用户态的高分辨率单调计时用 QPC：接口是
\Fn{Query{\Bk}Performance{\Bk}Counter}，返回的是 tick 计数；换算成时间，要除以
\Fn{Query{\Bk}Performance{\Bk}Frequency} 报告的频率。三个要点。
其一，这个频率由平台实现决定，不同代硬件上量级从百万到上亿 tick/秒都有，
并且和核心实际运行频率\emph{无关}——所以 QPC 读数永远不能当``CPU 周期数''用；
想要周期口径就得承认那是另一件事（PMU 事件，第\ref{ch:pm05-events}章）。
其二，现代 Windows 保证 QPC 跨线程单调，但历史上（老 HAL、多插座老平台）存在跨核不同步的问题；
在虚拟机里，后端与虚拟化增强（TSC page 一类）会让行为再变一层——确认方法是跑一段
线程迁移下的一致性测试，而不是查文章。其三，QPC 的调用成本与实现后端有关，
通常在几十纳秒到更低量级之间浮动，同样请读者用第二节的骨架自测。
墙钟侧对应 \Fn{Get{\Bk}SystemTime{\Bk}Precise{\Bk}AsFileTime}（受系统时间修改影响，语义等同
\Tp{CLOCK{\Bk}\_REALTIME} 的风险面）。剖析工具（ETW/XPerf/VTune）在第\ref{ch:pm04-sampling}章展开，
完整对照见附录。

\subsection{macOS：tick 要过 timebase 换算}
\label{sub:pm02-macos}

\Tp{mach{\Bk}\_absolute{\Bk}\_time} 给出单调 tick 计数，但它的单位换算系数\emph{不是} 1\,ns：
要用 \Tp{mach{\Bk}\_timebase{\Bk}\_info} 的 \Tp{numer}/\Tp{denom} 换算成纳秒。系数是开机时确定的，
取决于底层硬件——Intel 时代多与 TSC 类计数挂钩，Apple 硅上则是片上固定频率的通用定时器。
结论是通用的：在 macOS 上把\emph{裸 tick} 当纳秒、或把某个历史换算常数硬编码进代码，
都是把上一代机器的性质带进这一代。近年 macOS 也支持 \Tp{CLOCK{\Bk}\_UPTIME{\Bk}\_RAW}
一类 POSIX id（语义接近``开机起算''），同样用清单 \ref{lst:pm02-clockprobe} 试探本机实际支持哪些；
挂起与睡眠的计数行为请按机型实测，别照搬 Linux 的结论。频率/温度侧的观察工具
（\Tp{powermetrics}）受权限模型约束，见附录对照。

\section{别用墙钟算 CPU 事：\MCpu{} 的来处与口径}
\label{sec:pm02-cputime}

\MWall 由第一节解决；\MCpu 需要另一套来源，因为它是``被调度的时间之和''，不是``流逝的时间''。
Linux 上常用的四个来源：

\begin{itemize}
\item \Fn{getrusage}（\Tp{RUSAGE{\Bk}\_SELF}）：给出进程级 \Tp{ru{\Bk}\_utime}/\Tp{ru{\Bk}\_stime}，
即用户态与内核态分别烧掉的时间。口径陷阱：不同实现对``已退出子进程''是否计入的答案不同
（Linux 的 \Tp{RUSAGE{\Bk}\_SELF} 不计子进程，要走 \Tp{RUSAGE{\Bk}\_CHILDREN}）；粒度受内核记账约束。
\item \Tp{/proc{\Bk}/self{\Bk}/stat} 的 \Tp{utime}/\Tp{stime} 字段：单位是 clock tick，先用
\Tp{sysconf} 查询 \Tp{SC{\Bk}\_CLK{\Bk}\_TCK} 换算成秒——tick 粒度换算出的 \MCpu 有量化误差，
短任务尤其明显。
\item \Fn{times} 与 \Tp{tms}：同样的口径，更老的接口。
\item \Tp{CLOCK{\Bk}\_PROCESS{\Bk}\_CPUTIME{\Bk}\_ID} 与线程版：纳秒接口，但底层还是调度器记账，
更新发生在调度事件处——名义分辨率高不等于真实粒度细（第二节提过）。
\end{itemize}

由此得出本章对 \MCpu 的两条硬结论。

\textbf{多线程程序的 \MCpu 大于 \MWall 是正常现象，不是测量错误。} \MCpu 是多核时间的总和：
4 个线程各跑 1 秒，\MWall 是 1 秒，\MCpu 接近 4 秒。反过来说，一个补丁把 \MCpu 降了一半而 \MWall
几乎没变，可能只是``少烧了一个没人等的核''——两者根本不是同一把尺。并发下谁包含谁、
利用率的三种歧义形态，第\ref{ch:pm09-concurrency}章整章处理；这里只钉死口径：
\textbf{报告 \MCpu 时，必须同时报告 \MWall 与当时的并行度（线程数/绑核）}，否则数字无法解释。

\begin{note}[title={\MCpu 的定义，以及它不包含什么}]
\MCpu 等于该进程（或线程）在用户态与内核态实际占有 CPU 的时间之和。它\textbf{不包含}：
阻塞等待 I/O 的时间、就绪但没被调度的时间、睡眠的时间。所以 I/O 受限的负载会出现
\MWall 很大而 \MCpu 很小；被调度延迟坑掉的负载也是。想看``等了多久''，
\Tp{/proc{\Bk}/self{\Bk}/schedstat}（运行、等待、切换三段）与 cgroup 的节流计数
（\Tp{cpu.stat} 里的 \Tp{usage{\Bk}\_usec}、\Tp{nr{\Bk}\_throttled}、\Tp{throttled{\Bk}\_usec}）
是另两本账，容器与 CI 环境里必须分开记（第\ref{ch:pm13-ci}章）。
\end{note}

清单 \ref{lst:pm02-rusage} 把一次同步计算的两个口径放在一起量：循环结束时读一次 rusage，
再与墙钟对照。注意 \MCpu 的读数在进程生命周期末尾取（或确保没有并发线程在跑）才干净。

\begin{lstlisting}[style=cstyle, title={同一次运行同时量 \MWall 与 \MCpu（数值由你自己的机器给出）}, label={lst:pm02-rusage}]
#include <stdio.h>
#include <sys/resource.h>
#include <time.h>

static double ts_diff_ns(const struct timespec *a, const struct timespec *b) {
  return 1e9 * (double)(b->tv_sec - a->tv_sec)
       + (double)(b->tv_nsec - a->tv_nsec);
}
static double tv_to_s(const struct timeval *tv) {
  return (double)tv->tv_sec + 1e-6 * (double)tv->tv_usec;
}

int main(void) {
  struct timespec t0, t1;
  clock_gettime(CLOCK_MONOTONIC, &t0);

  volatile double acc = 0.0;                 /* 被测工作：单线程 CPU 密集 */
  for (long i = 1; i <= 20000000L; ++i) acc += 1.0 / (double)i;

  clock_gettime(CLOCK_MONOTONIC, &t1);
  struct rusage ru;
  getrusage(RUSAGE_SELF, &ru);

  double wall_s = ts_diff_ns(&t0, &t1) / 1e9;
  double cpu_s  = tv_to_s(&ru.ru_utime) + tv_to_s(&ru.ru_stime);
  printf("wall=%.6f s  cpu(user+sys)=%.6f s  cpu/wall=%.3f\n",
         wall_s, cpu_s, cpu_s / wall_s);
  /* 单线程纯计算：cpu/wall 应接近 1；明显偏离时先查被抢占、
     内核态占比、以及 rusage 的更新粒度，再谈别的结论。 */
  return acc > 0.0 ? 0 : 1;
}
\end{lstlisting}

\section{频率、温度与``同一台机器今天不一样''}
\label{sec:pm02-env}

第一节处理的是``时钟会不会跳''，本节处理更难缠的：时钟没有跳，但\textbf{被测程序变慢了，
因为机器不在同一个运行点}。现代 CPU 的瞬时频率由 DVFS（动态调频调压）决定：turbo 预算、
温度（热限值一到立刻降速）、供电、活跃核数、SMT 邻居的占用，全部实时影响每纳秒干多少活。
这对测量的含义是：短促爆发和持续满载跑在同一段代码上，频率可以差出一个台阶；
第一次迭代慢（冷启动、频率还没爬上去）不是噪声，是物理现象。想测``优化后的代码''而不是
``优化后的代码加上今天的室温''，就得把运行点钉住——能钉多少钉多少，钉不住的写进结论。

环境指纹就是这个动作：把``这台机器今天的运行点''变成数据的一部分。清单 \ref{lst:pm02-fingerprint}
给一份可跑的采集脚本骨架：\Tp{lscpu}、\Tp{scaling{\Bk}\_governor} 与 \Tp{scaling{\Bk}\_cur{\Bk}freq}、
turbo 开关、缓解位、内核与编译器版本。负载期间的逐核频率/温度观察，Intel 平台常用
\Tp{turbostat}（需要 root，参数与可用性随版本变，先跑 \Tp{turbostat} \Opt{help} 确认），
AMD/其他平台用各自的工具或退到``测前后各采一次均值''的粗口径——口径写进结论，别借别家平台的工具名。

\begin{lstlisting}[style=shellstyle, title={采集环境指纹（骨架：字段随发行版不同，缺项记录而不是跳过）}, label={lst:pm02-fingerprint}]
fp=perf-fingerprint; mkdir -p "$fp"
uname -a                       >  "$fp/kernel.txt"
lscpu                          >  "$fp/cpu.txt"
cat /proc/loadavg              >  "$fp/loadavg.txt"      # 有没有别的租户
cat /sys/devices/system/cpu/cpu0/cpufreq/scaling_governor > "$fp/governor.txt"
for d in /sys/devices/system/cpu/cpu*/cpufreq; do
  printf '%s %s\n' "$d" "$(cat "$d"/scaling_cur_freq 2>/dev/null)"
done > "$fp/freq.txt" 2>/dev/null                        # 缺目录=没开调频视图，记录之
cat /sys/devices/system/cpu/intel_pstate/no_turbo > "$fp/no_turbo.txt" 2>/dev/null
for v in /sys/devices/system/cpu/vulnerabilities/*; do
  printf '%s %s\n' "$(basename "$v")" "$(cat "$v")"
done > "$fp/mitigations.txt"
cc --version                   >  "$fp/compiler.txt"
# 对照实验期间可选锁定运行点；锁定的意义是"可复现"，
# 但它测的是那个运行点下的性能，结论必须写明这一点。
\end{lstlisting}

\begin{guideline}[title={环境指纹是结论的一部分，不是附录的装饰品}]
一份没有指纹的性能结论，等价于一张没有比例尺的地图。至少记：CPU 型号与核数（\Tp{lscpu}）、
内核版本（\Tp{uname} \Tp{-r}）、governor 与测时频率、turbo 状态、编译命令与编译器版本、
\textbf{测量起止时刻}。写``同一台机器今天不一样''不是借口：只要今天的指纹记下来了，
``不一样''本身就是可解释的数据。第\ref{ch:pm13-ci}章把这条纪律用在 CI 上。
\end{guideline}

最后划清本章工具与剖析器的分工：\Tp{perf} \Opt{stat} 的 \Tp{cycles}/\Tp{instructions} 事件读的是
内核经 PMU 或通用定时器给出的\emph{硬件}计数，它和你手选的 clock id 是两条独立通道，
对账（\Tp{task-clock} 对 \MCpu、cycles 对频率乘 \MWall）是第\ref{ch:pm05-events}章的内容。
本章的职责到``尺子本身可信''为止。

\section{GPU 与外设的时钟域}
\label{sec:pm02-gpu-domains}

CPU 之外的执行单元各自带独立时钟。GPU 有自己的 SM 时钟域：动态调频同样存在，且核时间
（kernel 内建的周期计数）与主机毫秒之间隔着一次换算与一次队列边界。这带来两条纪律：
其一，设备侧时间与主机时间\textbf{只在同一 stream/event 体系内部可比}——两个 CUDA event
夹出来的时间差可信，而``把主机时间戳和 kernel 时间戳排进同一张表''默认不可信，
对齐需要专门协议，误差也要专门量化；其二，比较``快慢''先问``同域吗''：
同一 kernel 换一次驱动/功耗设置，SM 频率变、周期数不变，ns 结论跟着变——这不是性能变化，
是换了尺子。具体测量协议（event 计时、同步语义、时钟锁定确认）留给第\ref{ch:pm11-gpu}章，
多卡与跨设备搬运涉及第二个域在第\ref{ch:pm12-bandwidth}章；本章只钉口径：
\textbf{跨时钟域不做减法}。

外设同理：网卡与音频设备各有硬件时钟（音频时钟域与漂移管理是《音视频开发》分册的主题），
嵌入式 SoC 的 ARM 通用定时器与 CPU 频率脱钩（见《嵌入式与硬实时》分册）。
所有``把两个不同源的计数排在一起''的需求，先回答：同源吗？偏移估过吗？漂多少量过吗？

\section{排障与速查}
\label{sec:pm02-trouble}

\begin{warning}[title={用 \Tp{CLOCK{\Bk}\_REALTIME} 测出负数耗时}, label={box:pm02-negative}]
症状：区间耗时为负，或者偶尔出现一个巨大的正数。成因按频率排：NTP 步进（尤其凌晨）、
虚拟机恢复或宿主强制校时、容器内进程自行 \Fn{clock{\Bk}\_settime}、双系统时区混乱叠加校时。
处理分三级：能改代码就改 id（\Tp{CLOCK{\Bk}\_MONOTONIC} 起步，跨 suspend 用 \Tp{BOOTTIME}）；
只能拿外部时间戳（日志、报文）就用它、但把负值\textbf{显式记录}进数据而不是丢弃或取绝对值；
常驻服务加一条自检：周期性对比 \Tp{REALTIME} 与 \Tp{MONOTONIC} 的差值，跳变即报警，
并把报警时刻盖进性能数据里——否则那一段测量全部作废都无人知道。
\end{warning}

\begin{compare}[title={三种典型场景的时钟选型}, label={box:pm02-choice}]
\begin{tabular}{@{}>{\raggedright\arraybackslash}p{0.26\textwidth} >{\raggedright\arraybackslash}p{0.24\textwidth} >{\raggedright\arraybackslash}p{0.27\textwidth}@{}}
\toprule
\textbf{场景} & \textbf{推荐} & \textbf{理由 / 要防的坑}\\
\midrule
\textbf{基准区间（单次任务）} & \Tp{CLOCK{\Bk}\_MONOTONIC} 或 \Tp{steady{\Bk}\_clock} &
步进免疫；但 slew 会轻修走速，映射与精度自己确认（第一、二节）\\
\textbf{跨睡眠的长窗口} & \Tp{CLOCK{\Bk}\_BOOTTIME} &
\Tp{MONOTONIC} 不计挂起段，整段测量会``凭空变快''（笔记本、嵌入式）\\
\textbf{热路径细粒度计时} & \Tp{CLOCK{\Bk}\_MONOTONIC{\Bk}\_RAW}，\Fn{rdtsc} 仅六条全过时 &
RAW 不被任何时钟服务改写；手写 TSC 的三性确认见第三节，缺一项回 \Fn{clock{\Bk}\_gettime}\\
\bottomrule
\end{tabular}
\end{compare}

章末速查，按症状：

\begin{itemize}
\item \textbf{耗时是负数或偶尔巨大}：先查 clock id 是不是 \Tp{REALTIME}；其次查这台机器上有没有
会改系统时间的服务（VM 工具、ntp 步进模式）。
\item \textbf{两次读数相减恒为 0 或恒为步长倍数}：区间短到低于实际更新粒度——先读
\Fn{clock{\Bk}\_getres}，再用第二节的骨架量出调用自身开销，两者都大于区间时这组测量没有意义。
\item \textbf{\Fn{clock{\Bk}\_gettime} 慢得反常（微秒级）}：vDSO 没生效，或被剖析器/ptrace 拉回
系统调用路径；\Tp{strace} \Tp{-c} 一条命令分辨。
\item \textbf{多线程下两个核的 \Fn{rdtsc} 对不上}：跨核失步。查 \Fn{dmesg} 与
\Tp{current{\Bk}\_clocksource}；确认后所有 ``\Fn{rdtsc} 结论''降级为``\Fn{clock{\Bk}\_gettime} 重测''。
\item \textbf{跑满一小时后 \Fn{rdtsc} 换算的时间开始漂}：频率假设失效（虚拟化调频、
恒定标志缺失），或换算常数本身就错；回到第三节的确认清单逐项重打勾。
\item \textbf{\MCpu 远大于 \MWall}：不是错误，是多线程求和；报数字时把并行度一起报出来，
口径歧义去第\ref{ch:pm09-concurrency}章。
\item \textbf{\MCpu 的读数只有 10\,ms 级台阶}：\Fn{getrusage}/\Tp{stat} 路径的粒度或刷新时机所致，
短区间的 \MCpu 先标成``量化占位值''，需要细账走 \Tp{schedstat} 或 PMU 的 \Tp{cpu-clock}
（第\ref{ch:pm05-events}章）。
\item \textbf{同一二进制今天比昨天慢}：先看第五节指纹——governor、温度史、邻居负载，
再谈代码。第\ref{ch:pm03-stats}章把``变慢''变成一个统计上站得住的判断。
\end{itemize}

本章把``用什么计数''钉死了：区间用单调类 id，CPU 用调度类来源，TSC 要三性俱全且换算常数
本机确认，跨时钟域不做减法。接下来把尺子读出的数变成结论——那需要重复次数与判据，见
第\ref{ch:pm03-stats}章。

\chapter{统计判据：重复次数、分布、置信区间与最小可检出差异}
\label{ch:pm03-stats}

本章解决一个朴素的问题：手上只有若干个耗时数字时，凭什么说某个补丁``快了''、``慢了''，或者``看不出来''。
第~\ref{ch:pm01-why} 章把``更快''变成了口径问题，第~\ref{ch:pm02-clock} 章解决了这些数从哪个计数器读出来；
本章处理剩下的三步：采多少个、信哪个统计量、差多大才值得写进结论。写法是\textbf{判据为主 + 可跑代码}：
结论一律落成``看到什么信号就怎么做''的表格与清单，不推公式。\PY 清单只用标准库，\CPP 清单只用
\Tp{std::{\Bk}chrono}。清单与输出样例里的数字全部是\textbf{合成数据或示意输出}，不是本仓库的实测值：
发布用的 CI runner 既没有 PMU 权限也通常没有 \Fn{perf}，具体百分比与纳秒数只能由你在自己的机器上采。

\section{一次测量的四种失败}
\label{sec:pm03-failures}

性能数字出问题的方式远多于产生它的方式，但绝大多数能归进四类。分清是哪一类很重要，因为处理动作完全不同：
有的要加重复次数，有的要加预热，有的要换机器，有的要改跑的顺序。

\textbf{单次样本}是跑一次就拿一个数比较：它没给你任何拒绝的余地，一次观测既不告诉你噪声多大，也不告诉你
这次是不是恰好撞上了一次降级。\textbf{冷启动}是第一批迭代处在与``稳定运行''完全不同的物理状态：页还没
映射进来，分支预测器与 TLB 还是空的，分配器 arena 还没建立，动态符号还没解析，频率还停在低频档；这些状态
会自己消失，所以``取第一次''和``扔掉第一次''给出两个不同的数。\textbf{系统噪声}是与负载无关的外部扰动
（邻居进程、调度迁移、后台索引与更新任务、电源与温控策略），它时间相关，同一段代码上午和下午不是一个数。
\textbf{阶跃偏置}最阴，因为它看起来就是版本差异：先跑完基线 A 再改代码跑 B，两批样本落在不同时间段上，
于是这段时间里的任何漂移都被记到 B 头上，而它的方向与真实差异无关。

\begin{longtable}{>{\raggedright\arraybackslash}p{0.16\textwidth} >{\raggedright\arraybackslash}p{0.31\textwidth} >{\raggedright\arraybackslash}p{0.34\textwidth}}
\caption{一次测量的四种失败：症状与你能做的处理}\label{tab:pm03-failure}\\
\toprule
失败形态 & 典型症状 & 你能做的处理 \\
\midrule
\endfirsthead
\toprule
失败形态 & 典型症状 & 你能做的处理 \\
\midrule
\endhead
\bottomrule
\endfoot
\endlastfoot
单次样本 & 别人复现不出来；隔天重跑结论翻转；报告里只有一个孤零零的数 &
按第~\ref{sec:pm03-steady} 节定重复次数下限；只报区间与样本量，不报单点 \\
冷启动 & 前几轮明显慢，慢的量随工作集大小变化；``跑久了自然就快了'' &
先采全量迭代时间戳，用累积均值判稳（表~\ref{tab:pm03-warmrule}）；预热段与计数段在脚本里分开写 \\
系统噪声 & 同一份二进制一天内差几十个百分点；\Mx{run{\Bk}-to{\Bk}-run variability} 很大 &
先做机器体检（第~\ref{sec:pm03-variability} 节）；记环境指纹；必要时换独占机或换时段，而不是加重复 \\
阶跃偏置 & A、B 各自内部很稳，两组之间却``恰好''差一个固定量；换顺序结论就变 &
改成交替配对，同一轮内 A/B/A/B 成对；差异只在配对差上做统计（第~\ref{sec:pm03-compare} 节） \\
\end{longtable}

四种失败有两个共同点：都不会报错，都能被``多跑几次''部分掩盖。所以判据要先立起来。

\begin{guideline}[title={开工前的三个判据}, label={box:pm03-criteria}]
按下运行键之前把这三句写进工作记录，否则采回来的数无法判读：
\begin{itemize}
  \item \textbf{要检出的最小差异是多少}（\Mx{MDE}，minimum detectable difference）。用相对量写，
        例如``低于 $5\%$ 的中位数变化这次不关心''。\Mx{MDE} 决定重复次数下限；没有它，
        ``跑多少次算够''这个问题没有答案。
  \item \textbf{主指标是哪一个口径}：\MWall、\MCpu、\MThr 还是 \Mx{p99}（定义见第~\ref{ch:pm01-why} 章）。
        多个指标同时看时先声明``结论只由主指标决定''，其余只用于归因；否则报告阶段必然滑向
        选择性报告（第~\ref{sec:pm03-multiplicity} 节）。
  \item \textbf{判据三档长什么样}：明显变好 / 无结论 / 明显变差分别对应区间的哪个位置。现在就写死，
        跑完不许改（预注册见第~\ref{sec:pm03-multiplicity} 节）。
\end{itemize}
\end{guideline}

\section{重复与稳态：预热多久、跑多少次}
\label{sec:pm03-steady}

\subsection{预热：判稳，而不是拍脑袋}
\label{sub:pm03-warmup}

预热的目的不是``让它快一点''，而是\textbf{让测量落在一个可以被重复的物理状态上}。所以预热长度不该由
``跑 $1000$ 次''这种整数决定，而该由可观测量决定：先做一整轮``每个迭代都计时''的探查跑
（清单~\ref{lst:pm03-harness} 保留逐迭代时间戳并分段打印累积均值），再问``这个序列从第几项开始不再系统地变''。
表~\ref{tab:pm03-warmrule} 的三条信号互相独立，任意一条不满足就加长预热，三条同时满足才进入计数段。

\begin{longtable}{>{\raggedright\arraybackslash}p{0.24\textwidth} >{\raggedright\arraybackslash}p{0.26\textwidth} >{\raggedright\arraybackslash}p{0.31\textwidth}}
\caption{预热的三条判稳信号：看什么、什么样算稳、不满足怎么办}\label{tab:pm03-warmrule}\\
\toprule
观察量 & 稳定时的样子 & 不稳定的处理 \\
\midrule
\endfirsthead
\toprule
观察量 & 稳定时的样子 & 不稳定的处理 \\
\midrule
\endhead
\bottomrule
\endfoot
\endlastfoot
累积均值（第 $1$ 到第 $k$ 次的均值随 $k$ 变化） & 后段近似水平；再加 $10$ 个样本也看不出整体位移 &
把预热延长到曲线转平的位置，再留一倍余量 \\
块间中位数（把序列切成等长块，比较块中位数） & 相邻块的中位数差远小于 $1$ 个 \Mx{MDE} &
还有一段缓慢漂移（通常是频率档位或缓存驻留）：继续预热，或在负载里加稳态等待 \\
尾部比例（超过中位数若干倍的迭代占多少） & 比例本身稳定，不再随轮数单调下降 &
缺页或外部后台任务还没结束；先按第~\ref{sec:pm03-variability} 节排除噪声，再谈预热 \\
\end{longtable}

\subsection{重复次数的下限由 MDE 反推}
\label{sub:pm03-repeats}

``重复 $30$ 次''本身没有意义，``为了检出 $5\%$ 的差异需要重复多少次''才有意义。筛查判据是这么一句话：
一组样本的分散程度（\Mx{MAD} 或标准差都行）除以重复次数的平方根，其量级近似于中位数差的区间半宽；
要让半宽压到 \Mx{MDE} 的一半以下，就按这个量级反解次数。平方根意味着\textbf{想把可检出差异缩小到
十分之一，样本量要放大到一百倍}。这是本章唯一需要背下来的数量关系，它解释了尾延迟与微小回归为什么贵。

这只是筛查，不是保证。真正的下限要在你自己的数据上验：清单~\ref{lst:pm03-bootstrap} 的
\Fn{confirm\_{\Bk}runs} 用重抽样回答``在这个样本量下，一个刚好等于 \Mx{MDE} 的真实改进有多大概率
被判成明显变好''。把它调到 $0.8$ 以上再去跑真版本，比套任何公式都可靠。

\begin{warning}[title={``跑得快的那一次''不是基准}, label={box:pm03-mintrap}]
把 \Mx{min} 当主指标（``取 $100$ 次里最快的一次''）在噪声大的机器上看起来很稳，实际它度量的是
``这台机器在哪一轮里被打扰得最轻''，与代码关系很弱。只看均值同理：一次抢占或一次热降频就能把均值抬起来，
而中位数无动于衷。定次数也别交给默认值：基准框架的时间预算默认值面向``毫秒级负载、每次几十微秒''的典型
场景，单个迭代几秒的负载可能只给你个位数样本，纳秒级负载却给你几百万样本；采样剖析器的频率默认值
（第~\ref{ch:pm04-sampling} 章）与模板里``只跑一轮取最小值''的写法同理。这些默认值与语义都随版本变，
以你手上那版的 \Opt{help} 为准，然后按 \Mx{MDE} 反推自己的次数，或把结论写成``无法判定''。
\end{warning}

清单~\ref{lst:pm03-harness} 是可以直接编译的采集骨架：预热与计数分开，逐迭代保留时间戳，
顺手算出中位数、四分位距、\Mx{MAD}、最近秩分位数与线性插值分位数，数值由你自己的机器与负载给出。

\begin{lstlisting}[language=C++, style=cppstyle, caption={\textbf{可跑}：预热、逐迭代采样与分位数摘要的采集骨架（清单内不含任何本仓库实测值，数值由读者自己的机器给出）}, label={lst:pm03-harness}]
// pm_harness.cpp   g++ -std=c++17 -O2 -Wall -o pm_harness pm_harness.cpp
// usage: ./pm_harness <iterations> [warmup]
#include <algorithm>
#include <chrono>
#include <cmath>
#include <cstdio>
#include <cstdlib>
#include <vector>

using clk = std::chrono::steady_clock;      // CLOCK_MONOTONIC family

// The workload under test. Replace it; keep setup out of the timed region.
static long double workload(std::size_t n) {
  long double acc = 0.0L;
  for (std::size_t i = 0; i < n; ++i) acc += 1.0L / (1.0L + static_cast<long double>(i));
  return acc;
}

// nearest rank: rank = ceil(q * N), clamped into [1, N]
static double quantile_nr(std::vector<double> xs, double q) {
  if (xs.empty()) return 0.0;
  std::sort(xs.begin(), xs.end());
  const std::size_t n = xs.size();
  std::size_t rank = static_cast<std::size_t>(std::ceil(q * static_cast<double>(n)));
  if (rank == 0) rank = 1;
  if (rank > n) rank = n;
  return xs[rank - 1];
}

// linear interpolation between neighbours (numpy's default flavour)
static double quantile_lin(const std::vector<double>& s, double q) {
  if (s.empty()) return 0.0;
  if (s.size() == 1) return s.front();
  const double pos = q * static_cast<double>(s.size() - 1);
  std::size_t lo = static_cast<std::size_t>(std::floor(pos));
  if (lo + 1 >= s.size()) return s.back();
  const double w = pos - static_cast<double>(lo);
  return s[lo] * (1.0 - w) + s[lo + 1] * w;
}

static double mad_of(const std::vector<double>& xs) {
  std::vector<double> s = xs;
  std::sort(s.begin(), s.end());
  const double m = quantile_lin(s, 0.5);
  std::vector<double> dev;
  for (double x : s) dev.push_back(std::fabs(x - m));
  std::sort(dev.begin(), dev.end());
  return quantile_lin(dev, 0.5);
}

int main(int argc, char** argv) {
  if (argc < 2) { std::printf("usage: %s <iterations> [warmup]\n", argv[0]); return 2; }
  const long iters = std::strtol(argv[1], nullptr, 10);
  const long warm  = (argc > 2) ? std::strtol(argv[2], nullptr, 10) : 0;
  const std::size_t elements = 1u << 14;

  for (long i = 0; i < warm; ++i) { volatile long double sink = workload(elements); (void)sink; }

  std::vector<double> ns;                        // one timestamp per iteration
  for (long i = 0; i < iters; ++i) {
    const auto t0 = clk::now();
    volatile long double sink = workload(elements);   // keep the work observable
    (void)sink;
    const auto t1 = clk::now();
    ns.push_back(std::chrono::duration<double, std::nano>(t1 - t0).count());
  }
  if (ns.empty()) return 1;

  std::vector<double> sorted = ns;
  std::sort(sorted.begin(), sorted.end());
  double mean = 0.0;
  for (double x : ns) mean += x;
  mean /= static_cast<double>(ns.size());

  // Warmup probe: cumulative mean at every decile. Stop the warmup where this
  // curve stops moving.
  double running = 0.0;
  const std::size_t step = ns.size() / 10 + 1;
  for (std::size_t i = 0; i < ns.size(); ++i) {
    running += ns[i];
    if ((i + 1) % step == 0)
      std::printf("  k=%6zu cum_mean=%.1f\n", i + 1, running / static_cast<double>(i + 1));
  }

  // One line per statistic so a shell pipeline can pick one out later.
  std::printf("# median=%.1f\n", quantile_lin(sorted, 0.5));
  std::printf("# n=%zu mean=%.1f mad=%.1f iqr=%.1f\n", ns.size(), mean, mad_of(ns),
              quantile_lin(sorted, 0.75) - quantile_lin(sorted, 0.25));
  std::printf("# p50/p90/p99/p999 nearest-rank: %.1f %.1f %.1f %.1f\n",
              quantile_nr(ns, 0.5), quantile_nr(ns, 0.9), quantile_nr(ns, 0.99),
              quantile_nr(ns, 0.999));
  std::printf("# p50/p99 interpolated: %.1f %.1f (say which one you used)\n",
              quantile_lin(sorted, 0.5), quantile_lin(sorted, 0.99));
  return 0;
}
\end{lstlisting}

\section{分布的形状：均值、中位数与分位数}
\label{sec:pm03-distribution}

性能数据不是正态的，原因很具体：一条迭代路径会被离散事件打断，缺页、线程被换下、缓存整体失效、
分配器回落慢路径、频率换档。这类事件让每个样本``要么差不多，要么差很多''，于是出现\textbf{长尾}
（右尾拉得很长，均值被少数样本拖走）甚至\textbf{多峰}（两个峰对应两条路径，均值落在两峰之间的空谷里，
那个值一个样本都不是）。选统计量之前先问形状：把 \Mx{MAD} 换算成与标准差可比的量（乘 $1.4826$
这个正态假设下的比例因子即可，不必深究来源）再与标准差相除，比值明显大于 $2$ 就按``均值不可信''处理。

\Mx{p99} 在小样本上不是``慢的那个百分位''，而是一句空话。最近秩算法下，样本量为 $N$ 时的 \Mx{p99}
取第 $0.99N$ 向上取整的那个顺序统计量：$N = 100$ 时它恰好是\textbf{最慢的那一个样本}，$N = 50$ 时是第二慢。
换机器、换时段、运气不好都会整体搬走它。判据是：想让某个分位数的尾部有东西撑着，尾部样本数至少要到
$20$ 以上，这是硬下限；于是 \Mx{p99} 需要几百到几千个样本，\Mx{p999} 需要上万。样本不够时宁可只报到
\Mx{p90} 并写明``尾部样本不足''，也不要报一个没有支撑的 \Mx{p99}。

更隐蔽的是：``p99'' 这个词在不同工具里算法不一致。同一批数在 \PY 标准库的 \Fn{statistics{\Bk}.quantiles}
（默认 exclusive 那套算法）、\Fn{numpy{\Bk}.percentile}（默认线性插值）、剖析器的直方图
（对数分桶，给出的是桶边界）里会算出三个不同的值。单个版本``自己跟自己比''无所谓，只要算法固定；
\textbf{跨工具比数时它就是错误的直接来源}。清单~\ref{lst:pm03-harness} 把最近秩与线性插值同时打印出来，
就是为了让你看见这个差：两个都报，并说明发布结论用的是哪一个。

\begin{lstlisting}[style=outstyle, caption={\textbf{示意输出 / 格式示例（非本仓库实测值）}：同一份 $240$ 个迭代样本，换一个统计量就换一个结论；数值是为说明分布形状而构造的}, label={lst:pm03-threeways}]
# batch: 240 retained iterations, one process, one machine (SYNTHETIC SHAPE)
# units: us per iteration

  mean        1834.2
  median      1211.0
  p90         1588.0
  p99        11402.0     <-- only 3 samples sit above it: not a tail, an accident
  mad           41.3     (1.4826 * mad = 61.2 ; sd = 2780.4 ; ratio far above 2)

  [   0,  200)   11 ###########
  [ 200,  400)  139 ########################################
  [ 400,  600)   84 ########################
  [ 600,  800)    2 ##
  [9600,11600)    3 ###

# three conclusions you could (wrongly) draw, one per number:
  "mean 1834 us"   -> "the hot loop is slow, rewrite it"
  "median 1211 us" -> "the typical path is fine, ship it"
  "p99 11402 us"   -> "there is a 10 ms stall, chase it"
# what the shape says: the bulk is bimodal and the tail has 3 samples out of 240.
\end{lstlisting}

写给别人看的数要配齐三件东西：\textbf{统计量名称 + 算法 + 样本量}，例如``中位数（最近秩，$N=812$）''。
只有均值、或只有``最好的一次''的报告（第~\ref{ch:pm15-antipattern} 章的反模式）无法判读，因为它没告诉你
形状；而形状本身（多峰、长尾、两峰比例）常常就是线索：负载有两条路径，该去做分账而不是去优化平均值。

\section{比较两个版本：配对、交替与 bootstrap 区间}
\label{sec:pm03-compare}

\subsection{顺序即变量：A/B/A/B}
\label{sub:pm03-abab}

只要 A 与 B 是``先后各跑一批''，两批样本之间就多了一个自变量：时间。消掉它要做两件事：
\textbf{交替}，同一轮里 A、B 各跑一次（A/B/A/B/\ldots），这一轮产生的一对样本共享当天的机器状态；
\textbf{随机化顺序}，每轮用随机数决定先跑谁，消掉``先跑的把缓存与频率状态改给了后跑的''这种顺序效应。
配对之后比较对象从``两组数''变成``每轮一个差值''：相减把当天漂移的共同部分抵消了，剩下的分散才是
版本间的真实不确定度。当配对不可能时（例如 A 只在那台旧机器上跑得起来），要把这件事写进结论的限制里，
而不是假装能比。

\subsection{bootstrap：不假设分布的区间}
\label{sub:pm03-bootstrap}

中位数差的置信区间不必靠正态假设。做法是重抽样：对每个版本，从它的样本里有放回地抽同样多个数、
算一个中位数；两组相减得到一个差；重复几千次，把这列差排序后取 $2.5\%$ 与 $97.5\%$ 两个位置，
就是 $95\%$ 区间。它便宜、直观，对多峰和长尾都不敏感；代价是 $N$ 太小时区间本身也是编的：
每组至少几十个点才值得跑，几百个点才谈得上稳定。清单~\ref{lst:pm03-bootstrap} 给了可跑的实现、
三档判据函数，以及两个配套工具（按 \Mx{MDE} 估最小次数、用重抽样复核这个次数够不够）。

\begin{lstlisting}[style=pythonstyle, caption={\textbf{可跑}：bootstrap 中位数差、三档判据、按 MDE 估最小重复次数与检出率复核（\PY 标准库实现，\Tp{numpy} 只用于加速；\texttt{\_\_main\_\_} 里的数据是合成的）}, label={lst:pm03-bootstrap}]
# stats_verdict.py   Python 3.9+, stdlib only (numpy optional, for speed)
import math
import random
import statistics

def quantile_nr(xs, q):
    """Nearest rank: rank = ceil(q * N), clamped into [1, N]."""
    if not xs:
        raise ValueError("empty sample")
    s = sorted(xs)
    rank = min(len(s), max(1, int(math.ceil(q * len(s)))))
    return s[rank - 1]

def mad(xs):
    m = statistics.median(xs)
    return statistics.median([abs(x - m) for x in xs])

def bootstrap_median_diff(a, b, rounds=20000, alpha=0.05, seed=1):
    """Percentile interval for median(b) - median(a); lower time means faster."""
    if len(a) < 20 or len(b) < 20:
        raise ValueError("tiny sample: this interval would be fiction")
    rng = random.Random(seed)
    diffs = []
    for _ in range(rounds):
        ma = statistics.median([a[rng.randrange(len(a))] for _ in a])
        mb = statistics.median([b[rng.randrange(len(b))] for _ in b])
        diffs.append(mb - ma)
    return quantile_nr(diffs, alpha / 2.0), quantile_nr(diffs, 1.0 - alpha / 2.0)

def verdict(a, b, mde, alpha=0.05, rounds=20000, seed=1):
    """Three tiers only; the mapping table is in the text."""
    lo, hi = bootstrap_median_diff(a, b, rounds, alpha, seed)
    point = statistics.median(b) - statistics.median(a)
    if hi < -mde:
        tier = "better"      # whole interval beyond the gain threshold
    elif lo > mde:
        tier = "worse"       # whole interval beyond the regression threshold
    else:
        tier = "unclear"     # covers "straddles zero" and "smaller than mde"
    return tier, point, lo, hi

def min_runs(spread_a, spread_b, mde, factor=3.0, cap=5000):
    """Screening bound, not a proof: shrink spread / sqrt(k) under mde / factor."""
    spread = max(spread_a, spread_b)
    k = 20
    while k < cap and spread / math.sqrt(k) > mde / factor:
        k *= 2
    return k

def confirm_runs(a, b, mde, rounds=300, inner=200, seed=7):
    """Detection rate of a true shift of exactly mde at this sample size."""
    src = random.Random(seed)
    shifted = [x - mde for x in b]         # pretend B really is mde faster
    hits = 0
    for _ in range(rounds):
        sa = [a[src.randrange(len(a))] for _ in a]
        sb = [shifted[src.randrange(len(shifted))] for _ in shifted]
        tier, _, _, _ = verdict(sa, sb, mde, rounds=inner, seed=src.randrange(1 << 30))
        hits += 1 if tier == "better" else 0
    return hits / float(rounds)

if __name__ == "__main__":
    rng = random.Random(0)                 # SYNTHETIC demo data, not a measurement
    base = [100.0 + rng.gauss(0.0, 6.0) for _ in range(300)]
    patched = [x * 0.99 for x in base]     # a true 1 percent gain
    print(verdict(base, patched, mde=1.0), min_runs(mad(base), mad(patched), mde=1.0))
    print(confirm_runs(base, patched, mde=1.0))
\end{lstlisting}

区间算出来了还要变成可执行的判定。表~\ref{tab:pm03-verdict} 就是开工前该写死的那张表：$d$ 表示
``中位数(B) 减 中位数(A)'',单位与主指标一致，越小越快。

\begin{table}[htbp]
\centering
\begin{tabular}{>{\raggedright\arraybackslash}p{0.23\textwidth} >{\raggedright\arraybackslash}p{0.12\textwidth} >{\raggedright\arraybackslash}p{0.44\textwidth}}
\toprule
区间的位置 & 档位 & 报告里该写什么 \\
\midrule
整段都在 \Mx{MDE} 的负侧之外 & 明显变好 &
``中位数低 $1.8\%$（$95\%$ bootstrap 区间 $-2.4\%$ 到 $-1.3\%$，每组 $N=600$，\Mx{MDE} $=1\%$）'' \\
整段都在 \Mx{MDE} 的正侧之外 & 明显变差 & 同样给出区间；先按第~\ref{ch:pm06-micro} 章归因，再决定改哪一段 \\
区间跨 $0$，或整段落在 \Mx{MDE} 之内 & 无结论 & 写``本轮不足以判定''；不要写成``没有变化''或``基本持平'' \\
区间宽度超过 $2$ 倍 \Mx{MDE} & 判据不成立 & 按第~\ref{sub:pm03-repeats} 节重估次数、复核检出率后重跑，而不是改判据 \\
\bottomrule
\end{tabular}
\caption{三档判据：从区间位置到结论的映射（\Mx{MDE} 必须是开工前写下的那一个）}\label{tab:pm03-verdict}
\end{table}

\begin{warning}[title={区间不跨过 $0$ 不等于``有意义''}, label={box:pm03-significance}]
区间完全在负侧只说明``差异大概率是真的''，不说明``值得改''；区间含 $0$ 也不说明``两个版本一样''，
它只说明这轮样本不足以判定。两种情况都要用 \Mx{MDE} 收口：\textbf{差异的绝对值小于 \Mx{MDE} 时，
无论区间长什么样，档位都是无结论}。写结论时把 \Mx{MDE} 一并写进去，否则一年后有人拿着这条
``显著但只快 $0.3\%$''的记录来要求合并。同理，``B 快 $0.6\%$''这种点值写法会让读者据此上线，
而那 $0.6\%$ 完全在噪声里。
\end{warning}

\section{变异性本身是指标}
\label{sec:pm03-variability}

很多``测不出来''不是方法错，是机器今天不行。所以分散度要当成一个正式指标记下来：一轮测量结束时，
除了中位数还要留下 \Mx{IQR}（四分位距，\Mx{p75} 减 \Mx{p25}，天然抗尾）与 \Mx{MAD}
（每个样本到中位数的距离再取中位数，比 \Mx{IQR} 更稳、不易被单个点带动），再用 MAD 除以中位数得到
一个无量纲的相对分散度。有了它，``这台机器现在适不适合测''就是可以自己标定的判据。
\textbf{不要抄别人的阈值}：别人的机器、内核与负载跟你的差异，往往正好就是你要测的那个量。
标定方法：挑一个与被测对象无关的\textbf{探针负载}（一个固定的、你永远不会去优化的小循环最合适），
按固定日程跑，把每天的相对分散度入库，阈值取自己历史分布的高分位（例如 $95\%$ 分位），超出就今天不测。

\begin{lstlisting}[style=shellstyle, caption={\textbf{可跑}：机器体检记录的采集与判读骨架（阈值来自你自己的历史分布，不给任何实测数值）}, label={lst:pm03-probe}]
# one probe run per slot, 20 slots, one summary line per slot into a daily log
for i in $(seq 1 20); do
  ./pm_harness 2000 50 | grep '^# median=' >> "probe-$(date +%F).csv"
  sleep 5
done
# relative spread of today, from the medians you just logged:
python3 - <<'PY'
import statistics
xs = [float(l.split('=')[1]) for l in open('probe-today.csv') if l.startswith('# median')]
med = statistics.median(xs)
mad = statistics.median([abs(x - med) for x in xs])
print('spread = mad/median = %.4f' % (mad / med))
PY
# Threshold: the 95th percentile of every daily spread you logged last month.
# Above it, stop measuring. No amount of repetition buys back a noisy machine.
\end{lstlisting}

体检不合格时的处理顺序是：先看\textbf{频率与温控}（governor、当前频率、是否正在降频，采集方法见
第~\ref{ch:pm02-clock} 章），再看\textbf{邻居}（后台索引、备份、更新任务；声称独占的机器要验证它真独占），
最后看\textbf{自己}（负载里有没有非确定性的东西：随机种子没固定、线程数跟着环境变、工作集没预热到同一规模）。
绑核与内存绑定会让``换个机器就复现不出''现形，细节在第~\ref{ch:pm09-concurrency} 章；这里只要一条：
绑与不绑都留记录，只报绑核后的数字是一种口径偷换。

\begin{warning}[title={把方差当成优化空间}, label={box:pm03-variance}]
分散度大不是一个待优化的指标。看到 MAD 很高就动手``消抖动''，绝大多数时候你优化的是\textbf{机器状态}
而不是代码：换台机器曲线立刻变平，补丁本身什么都没做。正确顺序是：\Mx{IQR} 与 \Mx{MAD} 异常时先归因
（谁在制造离散：频率、迁移、缺页、分配器、外部任务），再决定要不要测；归因完成、分散被解释之后，
剩下的差异才值得用代码去换。另一个方向同样要警惕：\Mx{MAD} 异常地小，通常意味着你在测一个全在缓存里的、
被死代码消除吃掉的、或者根本没有执行的负载（第~\ref{ch:pm15-antipattern} 章的微基准反模式）。
\end{warning}

\section{多重比较与选择性报告}
\label{sec:pm03-multiplicity}

最容易骗过自己的动作是``扫参数''：五十个编译开关、线程数、分块大小的组合各跑一轮，把最好的那一格报出来。
那一格不是性能的界，它是\textbf{噪声的一个顺序统计量}：在几十个真实性能彼此相当的组合里，最大值的期望位置
随候选数增加而稳定上移。换句话说，候选足够多时，即使代码一点没变，你也能``测出''一个更快的版本；
报出来的那个数天然偏高，而且复现不了，下一次采样会给你另一个冠军。可用的处理有三条，缺一条都会退回去：

\begin{itemize}
  \item \textbf{探索与确认分开}。扫参数那一轮只用来生成假设（``第 $12$ 格值得跟''），它的数字一律不进结论。
        要写进结论的那一次必须用\textbf{新采的样本}重跑冠军与基线的一对一对照，判据仍是
        第~\ref{sec:pm03-compare} 节的三档。
  \item \textbf{收紧单次判据}。确实要在一堆候选里直接选，就把单次门槛按候选数量同比例抬高
        （朴素做法是把允许误报的比例除以比较次数）。代价很实在：要检出的差异变大，或者样本量按
        第~\ref{sub:pm03-repeats} 节的规则放大。天下没有免费的显著性。
  \item \textbf{报全部对照}。把扫描的完整表一起给出，让读者看见冠军在整片区域里的位置。
        选择性报告的问题不在于数字错，而在于读者看不到分布。
\end{itemize}

\begin{bestpractice}[title={预注册判据与``报全部对照''的写法模板}, label={box:pm03-prereg}]
\textbf{预注册}：开跑之前把下面六行填完并存档（跟代码一起提交最好）。跑完不许改判据，改判据等于重跑。
\begin{itemize}
  \item 主指标与口径：例如``中位数 \MWall，每组 $N=600$ 次交替''
  \item \Mx{MDE} 与档位：三档各对应区间的哪个位置（表~\ref{tab:pm03-verdict}）
  \item 负载与输入：数据规模、随机种子、预热轮数与判稳依据
  \item 机器指纹：CPU 型号与频率档位、内核版本、编译器版本（记法见第~\ref{ch:pm02-clock} 章）
  \item 停止规则：``区间宽度达标或样本数到达上限''，而不是``跑到满意为止''
  \item 报告范围：扫过的每一组都列出来，骨架见下表
\end{itemize}
\textbf{报全部对照}的表头骨架（每行一组，不许只留冠军）：
\begin{center}
\begin{tabular}{>{\raggedright\arraybackslash}p{0.16\textwidth} >{\raggedright\arraybackslash}p{0.15\textwidth} >{\raggedright\arraybackslash}p{0.17\textwidth} >{\raggedright\arraybackslash}p{0.16\textwidth}}
\toprule
组（配置） & 中位数（相对基线） & 区间（$95\%$） & 档位（三档） \\
\midrule
base & $1.000$ & 不适用 & 基线 \\
cfg-07 & 填 & 填 & 填 \\
cfg-12 & 填 & 填 & 填 \\
\bottomrule
\end{tabular}
\end{center}
连续性能门禁（第~\ref{ch:pm13-ci} 章）里这条尤其重要：阈值必须来自这份预注册记录，而不是来自
``上次 CI 红了几次''。
\end{bestpractice}

\section{排障与速查}
\label{sec:pm03-trouble}

本章的判据收在表~\ref{tab:pm03-checklist} 里，按测之前、测之中、测之后三段各五条，
是本章的``我明天怎么测''版本：把它放在基准脚本目录里，每次开工前对着勾一遍。

\begin{longtable}{>{\raggedright\arraybackslash}p{0.08\textwidth} >{\raggedright\arraybackslash}p{0.24\textwidth} >{\raggedright\arraybackslash}p{0.22\textwidth} >{\raggedright\arraybackslash}p{0.25\textwidth}}
\caption{统计判据自检表：之前 / 之中 / 之后各五条}\label{tab:pm03-checklist}\\
\toprule
阶段 & 自检项 & 不合格的样子 & 处置 \\
\midrule
\endfirsthead
\toprule
阶段 & 自检项 & 不合格的样子 & 处置 \\
\midrule
\endhead
\bottomrule
\endfoot
\endlastfoot
前 & 主指标与口径已经写死 & 三个指标都想要，最后哪个都不报 & 只定一个主指标，其余降为归因用 \\
前 & \Mx{MDE} 与三档判据写成文字并存档 & ``先跑跑看再说'' & 按框~\ref{box:pm03-prereg} 做预注册 \\
前 & 重复次数下限按 \Mx{MDE} 反推过 & 沿用别人的 $10$ 次 & 先估分散度，再用重抽样复核检出率 \\
前 & 机器体检在自己的阈值内 & 阈值是抄来的，或从没标定过 & 先建立探针负载历史（第~\ref{sec:pm03-variability} 节） \\
前 & 构建配置与产物指纹已记录 & ``差不多是 Release 吧'' & 回指第~\ref{ch:pm10-compiler} 章；不确认就别开跑 \\
中 & 预热按判稳停止，不按轮数停止 & 固定写死一个预热数字 & 看累积均值与块间中位数（表~\ref{tab:pm03-warmrule}） \\
中 & 版本对照交替且顺序随机 & 先跑完 A 再跑完 B & 改成同一轮内配对（第~\ref{sec:pm03-compare} 节） \\
中 & 逐迭代时间戳入库，不只存汇总值 & 只留了一个平均数 & 重采；汇总值无法复原分布形状 \\
中 & 丢弃规则事先声明 & ``这次看着不像，删掉吧'' & 记录丢弃数量与理由，另报一组不丢弃的敏感性 \\
中 & 剖析器与 sanitizer 没在跑 & 一边剖析一边定性能数 & 关掉再测；剖析自身开销另论（第~\ref{ch:pm04-sampling} 章） \\
后 & 报区间不只报点值，并给出每组样本量 & 一个孤零零的百分比 & 换成 bootstrap 区间与三档结论 \\
后 & 统计量名称、算法、样本量三件齐全 & 只写一个分位数名字 & 补上算法；跨工具比数前先对齐算法 \\
后 & 尾部样本量够（每个尾至少 $20$ 个点） & $N=50$ 就报 \Mx{p99} & 只报到 \Mx{p90}，并写明尾部样本不足 \\
后 & 变异性一并入账 & 只报集中趋势 & 把 \Mx{IQR} 与 \Mx{MAD} 当第二个正式指标报出 \\
后 & 扫过的每一组都报了 & 只报冠军 & 扫描属探索；冠军要用新样本重验后再写结论 \\
\end{longtable}

本章四件工具各有分工：采集靠清单~\ref{lst:pm03-harness}（预热判稳、逐迭代入库），判读靠
清单~\ref{lst:pm03-threeways} 那种形状观察（先问是不是长尾多峰再选统计量），比较靠
清单~\ref{lst:pm03-bootstrap} 的交替配对加 bootstrap 区间，能不能开测靠清单~\ref{lst:pm03-probe}
的机器体检。往后走，第~\ref{ch:pm13-ci} 章把这些判据压成无人值守的门禁，
第~\ref{ch:pm15-antipattern} 章给出违反它们会长什么样的负面清单，指标名与参数速查在
附录~\ref{ch:pmapB-params}。

\part{工具链：采样、计数与火焰图}

\chapter{采样剖析：perf record 与三种调用栈展开}
\label{ch:pm04-sampling}

前三章把``测量''这件事的地基讲完了：口径（第~\ref{ch:pm01-why}~章）、时钟（第~\ref{ch:pm02-clock}~章）、
统计判据（第~\ref{ch:pm03-stats}~章）。从本章起进入工具链：先解决最古老也最常被误用的一问——
\emph{时间到底花在哪段代码里}。采样剖析器回答这个问题，但它给的是直方图，不是账单；
直方图能不能读，取决于你写代码时有没有给它留栈、符号解析时它找没找对二进制、
以及你有没有把``采样器自己吃掉的开销''记进账。本章以 Linux 上的 \Tp{perf record} 为主线，
把三种调用栈展开方式摆在一起比较，再给 Windows（XPerf/ETW、VTune）与 macOS（Instruments、\Tp{sample}）
各一节对照；完整的三平台并列放在附录~\ref{ch:pmapA-platform}。
事件计数与微架构分层不在本章，那是第~\ref{ch:pm05-events}~章的事——
采样定位到代码位置，计数与 TopDown 解释位置上的成因，两条腿缺一不可。

\section{采样剖析的物理模型}
\label{sec:pm04-model}

\subsection{样本从哪来}

采样剖析器和插桩计时（第~\ref{ch:pm02-clock}~章的 \Fn{clock{\Bk}\_gettime} 路线）的根本区别在于：
它\emph{不}在函数入口出口打点，而是在固定的时间节拍或固定的事件计数节拍上，
由内核（或平台等价的机制）抓住当前执行点，记下一个样本：指令指针、时间戳、线程号，
外加（可选的）当时的调用栈。程序跑完后，把所有样本按``栈顶落在哪个符号的地址区间''分桶计数，
除以样本总数，就是报告里每个百分比的来源。

因此``某函数占 12\%''这句话的准确读法是：\emph{落在该函数地址区间内的样本占 12\%}。
它有三个推论，后面几节全部从这里长出来：

\begin{itemize}
  \item 样本是随机的，占比是估计值——置信区间和最小可检出差异的判据原样适用于剖析报告，
        见第~\ref{ch:pm03-stats}~章；
  \item 地址区间不等于函数边界——被内联进宿主的代码、没有符号的 JIT 码、
        共享库里的同名包装，都会让样本归错桶（第~\ref{ch:pm10-compiler}~章讲内联的真实语义）；
  \item 节拍本身可以改变被测程序——每个样本都要暂停程序、抓栈、写缓冲区，
        这就是观测者效应，本章第~\ref{sec:pm04-window}~节给度量方法。
\end{itemize}

\subsection{事件驱动采样与定时采样不是一回事}

\Tp{perf} 的采样可以由两类源驱动：软件时钟（\Tp{cpu-clock}、\Tp{task-clock}，本质是定时器）
和硬件性能事件（\Tp{cycles}、\Tp{instructions}、\Tp{cache-misses} 等）。
前者近似均匀地覆盖\emph{时间}；后者把样本压在\emph{该事件发生的位置}上——
按 \Tp{cache-misses} 采样的图，热点是访存指令的密集程度，不是耗时占比。
两种图都叫``火焰图''，回答的是两个不同的问题。事件侧的完整口径（PEBS 与 \Tp{event:pp}、
精确事件把样本修正到真正肇事的指令）放在第~\ref{ch:pm05-events}~章，
本章只用它来立一条纪律：\textbf{读任何剖析报告之前，先确认样本是按什么节拍的}。

\subsection{\Tp{-F}、\Opt{freq} 与 \Opt{count}：三种指定节拍的方式}

一次采样会话的节拍有三种写法，各有代价：

\begin{itemize}
  \item \Tp{-F} \emph{数字}（长选项为 \Opt{freq}）：每秒约 \emph{数字} 个样本。
        样本数可控、跨运行可比；代价是采样覆盖对短函数同样均匀撒网，稀疏处白采。
  \item \Opt{count} \emph{N}：每累计 N 个事件采一次。周期固定，样本数正比于事件总量；
        但同样的 N 在快慢不同的机器上给出的样本密度完全不同，跨机器不可比。
        更隐蔽的是：事件节拍会把样本``偏置''到事件本身——按周期数采样近似时间，
        按缺页采样则把样本全押在分配点上。
  \item \Tp{-F} max：以内核允许的上限采样。只该用于``程序太短、低频率采不到几个样本''的场景，
        并且要预期可观的剖析自身开销。
\end{itemize}

需要强调的是：\Tp{perf record} 在不指定频率时的默认值（常见版本为 \Tp{-F 4000}，
即 4kHz 附近）\emph{是工具默认值，不是结论}。它对小任务可能采不出统计意义，
对单核紧张的设备又可能明显伤速。默认频率、可用频率集合、内核采样率上限都随版本与内核变，
落笔定基线前先在你的机器上确认：\Tp{perf record --help} 看默认，
\Tp{cat /proc/sys/kernel/perf{\Bk}\_event{\Bk}\_max{\Bk}\_sample{\Bk}\_rate} 看内核上限，
跑一次后从 perf report 头部的``Samples: N''行看实际采到多少。
低占比函数需要更长窗口是同一个统计问题：要在 99\% 置信度上分辨 1\% 的份额，
样本总量得进到千的量级——窗口长度由``预估份额 $\times$ 频率 $\times$ 秒数''倒推，
判据细节回第~\ref{ch:pm03-stats}~章。

\begin{lstlisting}[style=shellstyle, caption={最小可用的一次采样会话：先确认权限与上限，再记录，再看。}, label={lst:pm04-first}]
cat /proc/sys/kernel/perf_event_paranoid   # 值越大限制越严, 见第 3 节
cat /proc/sys/kernel/perf_event_max_sample_rate
perf record -F 399 --call-graph fp ./myapp   # 目标程序短时可加 -F 997
perf report --stdio                          # 无图形环境时用 --stdio
# 换机器、换内核后这两行确认命令都要重跑, 不要抄别人的默认值
\end{lstlisting}

\begin{guideline}[title={开工前三个问题}]
开一次剖析会话之前，先把三件事写进你的记录：
(1) 栈从哪来——用哪种展开方式（第~\ref{sec:pm04-unwind}~节），二进制是否为此重编过；
(2) 窗口多长——按目标函数的预估份额倒推样本数，而不是``跑十秒看看''；
(3) 节拍是什么——定时还是事件驱动，事件驱动的报告只回答``该事件在哪密集''，
不回答``时间在哪''。三个问题里有任何一个答不上来，报告先别读。
\end{guideline}

\section{三种调用栈展开的取舍}
\label{sec:pm04-unwind}

没有栈的采样剖析只是地址直方图：它告诉你 \Fn{Scene::cull} 里耗时，但不告诉你\emph{谁}调用了它，
而归因恰恰需要``谁调用''这一层。三种展开方式的物理机制完全不同，失真也不同。

\subsection{帧指针链：便宜，但要求你重新编译}

最古老的方式：函数序言把旧的 \Tp{rbp} 压栈、\Tp{rbp} 指向本帧，栈就是一棵可走的链表。
perf 端用 \Opt{call-graph fp}（旧写法是 \Tp{-g}）读这条链。
运行时开销最小（内核只需沿链读若干个对齐指针），但它成立的前提是链没断——
而 \Tp{-O2} 之后的 GCC/Clang 默认 \Tp{-fomit-frame-pointer}，会把 \Tp{rbp} 挪作通用寄存器，链即断。
需要显式加：

\begin{lstlisting}[style=shellstyle, caption={为帧指针展开重编发布配置：两个开关缺一不可，且要落在真正被测的那个二进制上。}, label={lst:pm04-fpbuild}]
# 与发布配置同级别, 只是加回帧指针与调试信息; 不要拿 -O0 的构建来剖析
g++ -O2 -g -fno-omit-frame-pointer -o myapp main.cpp render.cpp
# CMake 工程: 在既有 RELEASE 标志上追加, 而不是新建一个探针配置
# cmake -DCMAKE_CXX_FLAGS_RELEASE="-O2 -g -fno-omit-frame-pointer" ...
perf record -F 399 --call-graph fp ./myapp
\end{lstlisting}

失真来源有三：跨``没有帧指针的库''（发行版预编译的 \Lib{libc}、\Lib{libstdc++}）时栈断在半路；
信号处理与内联让帧``存在但撒谎''；只重编了应用、没重编静态库时，库里贡献的对象仍断。
判断标准很简单——报告里调用链是否频繁终止在共享库符号上。
要不要用发行版自带的、编了帧指针的库变体（各发行版策略不同且在变），需要读者自己查本机包源。

\subsection{DWARF 解卷：完整，但贵}

不要求帧指针：内核在每个样本处把\emph{栈顶附近的一段原始栈内存}（默认量级为 8KB，
可给 \Opt{call-graph dwarf} 追加逗号与快照字节数来调整）连同寄存器快照写进 \Tp{perf.data}，
离线时用二进制里的 \Tp{.eh{\Bk}frame} 展开信息把栈重建出来。栈的完整性最好、
不需要改编译选项，代价是：单样本负载大（几百倍于帧指针法）、perf.data 暴涨、
离线解卷慢、且\textbf{最容易触发样本丢失}（第~\ref{sec:pm04-window}~节的缓冲区问题）。
两个上限需要读者自己确认而不是抄：单样本可携带的栈字节数受 perf 与内核实现限制，
栈深度上限看 \Tp{cat /proc/sys/kernel/perf{\Bk}\_event{\Bk}\_max{\Bk}\_stack}。
深递归程序配 8KB 小快照会截断——这正是``栈看着只有一层''的常见原因之一。

\subsection{LBR：硬件给的近叶子栈}

\Opt{call-graph lbr} 不用软件展开，改用 CPU 调试寄存器里的历史分支记录
（Intel 的 LBR、AMD 的 IBS 提供各自的近亲机制，能力与深度\emph{随微架构型号变}：
常见是十几到三十几层分支对；确认方法是直接跑
\Tp{perf record --call-graph lbr}——不支持时工具会明确报错，而不是给你半截结果）。
它零重编、单样本开销近乎免费、且天然带样本\emph{之间}的分支路径；
但它只覆盖叶子附近——栈深超过记录容量，祖先帧就没了，
火焰图会呈现``根部缺失的一丛叶子''。它是``想知道这段时间走过了哪些基本块''时的正确工具，
不是完整调用树的替代品。

\begin{longtable}{>{\raggedright\arraybackslash}p{0.13\textwidth} >{\raggedright\arraybackslash}p{0.17\textwidth} >{\raggedright\arraybackslash}p{0.12\textwidth} >{\raggedright\arraybackslash}p{0.20\textwidth} >{\raggedright\arraybackslash}p{0.19\textwidth}}
\caption{三种栈展开方式对照。``需要重新编译吗''一列针对帧指针法；其余的开销与权限随内核版本浮动，按第~\ref{sec:pm04-model}~节方法自查。\label{tab:pm04-unwind}}\\
\hline
方式 & 完整性 & 需重编 & 运行时开销 & 依赖与权限 / 常见失真 \\ \hline
\endfirsthead
\hline
方式 & 完整性 & 需重编 & 运行时开销 & 依赖与权限 / 常见失真 \\ \hline
\endhead
\hline
\endfoot
帧指针链 \Tp{fp} &
只到最近的\emph{带帧}调用者；跨未重编的库会断 &
是：\Tp{-O} 同级加 \Tp{-g} \Tp{-fno-{\Bk}omit{\Bk}-frame{\Bk}-pointer} &
最小；每样本固定几十到几百字节 &
低权限即可；失真：断链后半截栈被折叠进库符号、信号帧混入 \\
DWARF 解卷 \Tp{dwarf} &
理论最完整（用 \Tp{.eh{\Bk}frame} 全量重建）&
否；但发布产物须带展开信息（多数发行版默认有）&
大：每样本携带栈快照（默认 8KB 量级），I/O 与离线解卷都进账 &
受 \Tp{perf{\Bk}\_event{\Bk}\_max{\Bk}\_stack} 与单样本大小限制；失真：栈截断、样本丢失 \\
LBR（或厂商等价机制）&
仅叶子附近若干层，深度随微架构 &
否 &
近零（硬件记录）&
依赖 PMU 与调试寄存器权限（\Tp{perf{\Bk}\_event{\Bk}\_paranoid}、虚拟化透传）；失真：无根的一丛叶子 \\
\end{longtable}

\begin{warning}[title={DWARF 的缓冲代价会被``默认值''藏起来}]
\Opt{call-graph dwarf} 在默认频率下就可能把 perf record 的环形缓冲打爆，
症状是会话结束时报出可观的丢失样本数（见第~\ref{sec:pm04-window}~节）。
丢了样本，占比就不再是无偏估计——而报告不会为此停下来说话。
用 DWARF 时，把频率降下来（例如 \Tp{-F 99}）往往比把缓冲堆到天上去更接近``可信的 12\%''。
\end{warning}

\section{符号、内联与折叠}
\label{sec:pm04-symbols}

样本里的指令指针要变成函数名，栈里的返回地址要变成调用链，靠的是符号与展开信息。
这一层的坑比采样本身更多，而且\textbf{每一个都可能静默发生}。

\subsection{内核符号与权限}

报告里满屏 \Tp{[unknown]} 的头号原因是内核地址被遮蔽：
\Tp{/proc/kallsyms} 在非特权读时按 \Tp{kernel.kptr{\Bk}\_restrict} 的值返回全零，
而 perf 能否打开内核事件由 \Tp{/proc/sys/kernel/perf{\Bk}\_event{\Bk}\_paranoid} 决定。
确认顺序：先 \Tp{cat} 这两个文件记下当前值，再决定是临时降限制、用
\Tp{perf record --kallsyms=/proc/kallsyms} 在特权时刻留一份拷贝供离线解析
（选项可用性随 perf 版本变，以 \Tp{perf record --help} 为准），还是接受``只看用户态''并写明这一口径。
容器与云主机里这几个开关可能被外部固定，属于附录~\ref{ch:pmapA-platform}~的权限模型话题。

\subsection{buildid：剖析报告绑的是哪个二进制}

\Tp{perf.data} 里每个 DSO 记录的是 build-id，报告时用\emph{当时的}文件系统去反查符号。
剖析之后重新编译、换机器看旧 \Tp{perf.data}、strip 掉符号表——都会让反查失败，
表现为整列 \Tp{[unknown]} 或``符号对但行号飘''。三条纪律：
剖析会话结束时用 \Tp{perf buildid-list -i perf.data} 把涉及的 buildid 抄进记录；
需要长期归档的报告，先用 \Tp{perf buildid-cache --add} 把二进制放进
用户主目录下的 \Tp{.debug} 缓存目录（或系统级 build-id 缓存），或用 \Opt{symfs} 指定镜像根目录；
\Opt{buildid-all} 会为所有 DSO（包括系统库）记 buildid，更稳但 \Tp{perf.data} 更大、符号解析更慢——
值不值得取决于你是否要跨机复盘，别默认全开。

\subsection{内联帧与折叠}

内联函数在样本里\emph{不存在独立返回地址}：它的指令地址落在宿主符号区间内。
于是``报告里看不到某个 \Tp{inline} 函数''\emph{不能}推出``它没被执行、没占比''——
它可能正顶着宿主的名字。部分工具链能把样本地址经调试信息展开成内联栈
（离线解析一侧常见做法是 \Fn{addr2line} 的 \Tp{-i} 选项族；perf 各子命令对内联帧的
显示与版本有关，先自查 \Tp{--help} 再下断言）。
这也解释了火焰图的一个折叠口径：同一个内联体被展开时，不同宿主里它显示在不同路径下——对；
不展开时，全部样本计入宿主——也对。两种读法必须选一个写进结论，混着读必然打架。
报告内的过滤（如 \Opt{grep} 按子串筛行；\Opt{sym-type} 一类的显示/类型选项在不同 perf 版本
含义与存在性都不同，用前先跑 \Tp{perf report --help} 过滤一遍再决定依不依赖它）
只是显示口径，不改变底层样本——但人会因为显示口径变了而忘记自己换过问题。

\subsection{火焰图：流水线与三形态}

火焰图不是 perf 的原生输出，是折叠产物。经典流水线三步：

\begin{lstlisting}[style=shellstyle, caption={perf 到火焰图的标准折叠流水线。最后一步换参数即得 icicle 形态；sunburst 与 treemap 由同生态的其他脚本或工具生成，可用性需按你所用版本自查。}, label={lst:pm04-flame}]
perf record -F 399 --call-graph dwarf,16384 -o perf.data ./myapp
perf script -i perf.data > script.out
git clone https://github.com/brendangregg/FlameGraph   # 工具仓库, 不是被测代码
./FlameGraph/stackcollapse-perf.pl script.out > folded.txt
./FlameGraph/flamegraph.pl --title "myapp" folded.txt > flame.svg
# icicle 形态（根在上、叶在下）: 给 flamegraph.pl 加 --reverse
# 折叠行格式: 分号分隔的栈帧; 制表符分隔的样本计数 —— 用分页器肉眼先过一遍
\end{lstlisting}

折叠文件（栈帧用分号串起、行尾挂样本计数）值得\emph{至少肉眼看一遍}：
它能让``栈全断''（每行只有一两帧）与``栈全同''（展开失败退化成单一根）这两种失真现形。
形态上，flame 适合看叶子热点、icicle 适合自上而下追问``这条子树为什么存在''、
sunburst 与 treemap 在深栈与窄热点上各有可读性——选形态不改变数据，
改变的是你能看见的失真类型。

\begin{note}[title={折叠与去重的两个口径陷阱}]
一是``同名不同源''：动态生成代码、多个静态库里的同名符号、不同 buildid 的同版本库，
折叠后并成一帧；二是``样本太少的宽帧''：宽度为几个百分点的根帧下挂着长尾，
眼睛会优先读宽帧，而回归往往藏在长尾里。这两条都不是绘图 bug，是阅读纪律问题，
下一节用同一份数据说明。
\end{note}

\section{读报告的纪律}
\label{sec:pm04-report}

\subsection{自身占比与总占比}

树形剖析报告至少有两种正交的排序口径：\emph{自身}（样本的栈顶落在谁的地址区间）与
\emph{总计}（栈里\emph{经过}这个帧的样本，含它的所有被调者）。
文本视图（\Tp{perf report --stdio}）默认展示的 Overhead 列是自身占比；
展开调用树后同一函数会顶起一大片子树——那才是总计。
一个纯胶水函数自身占比永远小、总计占比可以是 90\%；一个叶子工具函数反之。
两份排序回答不同问题：自身→``值不值得改这段代码本身''；总计→``值不值得砍这条调用路径''。

\begin{lstlisting}[style=outstyle, caption={示意输出（非本仓库实测值，仅为 perf report --stdio 的形态示例；数值、符号、库版本都是编造的占位）。}, label={lst:pm04-report}]
# Samples: 39k of event 'cpu-clock'
# Event count (approx.): 98070000000
#
# Overhead       Samples  Command  Shared Object      Symbol
# ........  ............  .......  .................  .....................
    18.42%          7236  render   libstdc++.so.6     [.] std::_Hash_bytes
    12.71%          4992  render   myapp              [.] Scene::cull
     9.30%          3655  render   libc.so.6          [.] memcpy_erms
     7.14%          2805  render   myapp              [.] Frame::encode
     5.02%          1971  render   [kernel.kallsyms]  [k] copy_user_enhanced
        ...
\end{lstlisting}

逐列口径：Overhead 是自身样本占比；Samples 列让你用第~\ref{sec:pm04-model}~节的
``绝对数''去校验百分比是否可信（几百个样本谈 0.5\% 的差是无效的）；
Shared Object 加 buildid（开 \Opt{dsym} 时）决定符号可信度；
\Tp{[k]} 与 \Tp{[.]} 分别标内核与用户符号——满屏 \Tp{[unknown]} 则回到第~\ref{sec:pm04-symbols}~节。

\subsection{过滤阈值制造的错觉}

\Opt{percent-limit}（交互界面里的等价过滤）把低于阈值的项\emph{从显示里删掉}。
于是出现两个经典错觉：其一，剩余项``加起来等于 100\%''是假象——被删的部分还在贡献总时长，
你只是看不见；其二，``前两名占 90\%，其余忽略''的结论在阈值过滤后的屏幕上无法证伪。
纪律：任何写进结论的百分比，先回到无过滤视图（或 \Tp{perf report --stats} 一类的样本汇总）
核对样本构成；过滤只用于交互浏览。

\begin{compare}[title={同一份数据，三种口径三种决策}]
\begin{tabular}{>{\raggedright\arraybackslash}p{0.18\textwidth} >{\raggedright\arraybackslash}p{0.30\textwidth} >{\raggedright\arraybackslash}p{0.30\textwidth}}
视图 & 看到的样子 & 适合的决策 \\ \hline
自身占比、无过滤 & \Tp{std::{\Bk}\_Hash{\Bk}\_bytes} 占 18\% 居首 & 是否值得替换哈希实现 \\
总计占比、展开树 & \Fn{Frame::encode} 子树合计 41\% & 是否值得砍这条路径或缓存结果 \\
事件驱动（缺页采样）& 分配点与页错误密集处 & 与耗时无关，交给第~\ref{ch:pm07-memory}~章的口径 \\
\end{tabular}
\end{compare}

\subsection{\Tp{memcpy\_erms} 排第一常常不是结论}

\begin{warning}[title={库顶帧效应：别和 \Lib{libc} 吵架}]
\Tp{memcpy\_erms}、\Tp{memset}、\Tp{std::{\Bk}\_Hash{\Bk}\_bytes} 这类符号常年盘踞榜首，
因为它们是\emph{很多}调用者共同的叶子。自身占比把它们排第一，只说明
``这台机器确实在搬运、在哈希''，不说明``优化 \Lib{libc} 能赢''。
正确的下一步是展开调用树找到\emph{谁}在拷、拷\emph{什么}、拷\emph{为什么拷}：
是多余的中间缓冲（布局问题，第~\ref{ch:pm06-micro}~章）、是按值传大对象（接口问题）、
还是本来就该少的数据量（算法问题）。把``memcpy 占 9\%''写进优化立项的标题，
是把相关性当因果的教科书样本（全书收束见第~\ref{ch:pm15-antipattern}~章）。
\end{warning}

\section{窗口、周期与归属}
\label{sec:pm04-window}

\subsection{目标与系统级噪声}

target 族选项（\Tp{-p} 进程、\Tp{-t} 线程、\Tp{-C} CPU 列表）把样本限定在被测负载上；
\Opt{system} 采全系统，桌面、内核线程、别的租户都进来——除非你就是要在系统级找邻居，
否则剖析自己的程序时不要开系统级，至少在读报告时按 \Tp{Command}/\Tp{pid} 列过滤干净再谈百分比。

\subsection{按线程、按 CPU 归属}

\Opt{sort} 的字段组合决定``占比是对谁算的''：默认的 symbol 排序把所有线程混成一锅；
加 \Tp{tid} 后能看到某个线程独立的热点形状；加 \Tp{cpu} 后能发现``同一个工作线程被调度到
不同物理核上表现不同''（SMT 邻居与迁移的账，第~\ref{ch:pm09-concurrency}~章）。
多角色流水线（生产者/消费者、渲染线程）尤其要求按线程归属——
混排的图里，一个 10\% 的帧可能属于两个完全不同的角色。

\subsection{采样自身的开销怎么量}

剖析器改变被剖对象的行为不是玄学：每个样本要中断、抓栈、写缓冲；
DWARF 法还额外制造 I/O 与页缓存压力。度量方法就一次对照：
\textbf{同一二进制、同一负载，剖析器关与开各跑一组，交替进行}（A/B/A/B，判据回
第~\ref{ch:pm03-stats}~章），报出 \MWall{} 与 \MCpu{} 两个口径的差；
再扫一遍频率（例如 \Tp{-F 99}、\Tp{-F 399}、\Tp{-F 997}、\Tp{-F 4000}）看开销随频率的曲线形状。
``开销''的正确报法是这个比值随负载与频率的函数关系，不是单独一个百分数——
单报 3\% 会诱导别人在你没测过的负载上照抄你的采样率。
顺带一条负面结果也是结果：关掉剖析器再测的那一组，才是性能结论的合法样本；
带剖析器的运行只能用来定位位置。

\subsection{缓冲区与样本丢失}

内核侧环形缓冲（\Opt{buffer-size}、\Opt{mmap-pages} 一族选项）溢出时样本被\emph{丢弃}，
而丢失不随机——高频率、大栈（DWARF）的时段先丢，于是占比系统性偏移。
会话结束注意 perf 汇总里的丢失计数（字段措辞随版本变，找含 lost 的汇总行；
也可用 \Tp{perf report --stats} 交叉核对样本数与报告头部是否一致），
丢了就降频率或加缓冲，二选一，别加频率硬扛。

\section{另两平台的对照}
\label{sec:pm04-other-platforms}

\subsection{Windows：ETW/XPerf 与 VTune}

Windows 的采样剖析走 ETW（Event Tracing for Windows）：实时会话（图形化工具或
XPerf 命令行）采集样本事件并要求栈展开，之后在 WPA 一类查看器里分析。展开方式与本章同构：
帧指针链，或用镜像自带的 unwind 数据（x64 的 \Tp{.pdata}/\Tp{.xdata} 段）重建——
选哪种模式、你的库带不带 unwind 信息，需要读者自己在工具的栈分析配置里确认，
两种模式的失真与 Linux 侧完全对得上。符号对应 PDB：剖析机与构建机分离时，
符号路径（如 \Tp{{\Bk}\_NT{\Bk}SYMBOL{\Bk}\_PATH}）没配好，
效果和 Linux 丢 buildid 一样，是满屏地址。
MSVC 的帧指针策略受 \Tp{/Oy}/\Tp{/Oy-} 与平台、版本影响，不要假设默认保留——确认方法
与本章一致：看展开后栈是否成片截断，必要时反汇编看序言。
VTune 在 Windows 上可用自己的采样引擎（含 LBR 类硬件记录），口径可与
\Opt{call-graph lbr} 对照阅读。协程/fiber/异步续延有一个 Linux 同款的问题：
样本落在续延入口桩上时，``逻辑调用栈''在物理栈里不存在，
火焰图会把整条异步链折叠进调度器帧——读图前先问：这是栈没了，还是本来就没有栈？

\subsection{macOS：Instruments 与 \Tp{sample}}

macOS 的免安装起手式是 \Tp{sample <pid> <N>}（\emph{N} 为采样秒数）：
按定时节拍抓目标进程的线程栈，产出文本调用树；它只覆盖用户态
（内核帧需要更高权限，且受系统完整性保护状态约束——你能不能看到
\Tp{[kernel]} 与 \Tp{libsystem} 之下的帧，以实际运行为准，别抄网上的权限步骤）。
Instruments 的 Time Profiler 模板（图形界面，或命令行 \Tp{xctrace} 一族；
命令行入口在不同 Xcode 世代之间有迁移，先跑 \Tp{xcrun xctrace --help}
确认本机形态）给出同样的样本直方图加符号化，符号依赖 dSYM：
发布产物和 dSYM 的对应关系丢了，等价于 Linux 的 buildid 断链。
与 Linux 的主要口径差：macOS 上没有 perf 事件名的直接对应物，
``定时采样近似 CPU 时间''这一条成立，但被电源管理与频率爬升扰动得更明显——
环境指纹（第~\ref{ch:pm02-clock}~章）里要带上 \Tp{powermetrics} 采到的频率曲线。
三平台能力、权限与展开方式的完整并列见附录~\ref{ch:pmapA-platform}。

\section{排障与速查}
\label{sec:pm04-trouble}

\begin{guideline}[title={第一次剖析：先问栈，再问热点}]
拿到任何一份新环境的剖析报告，阅读顺序应该是：
(1) 样本总数够不够（报告头部的 Samples 行）与有没有丢失；
(2) \Tp{[unknown]} 占比——先修符号与权限，再谈热点；
(3) 抽十行折叠栈肉眼看深度——用的是哪种展开、栈断在哪一层；
(4) 这时才读第一名的百分比。跳步的顺序读不出错才怪。
\end{guideline}

\begin{longtable}{>{\raggedright\arraybackslash}p{0.22\textwidth} >{\raggedright\arraybackslash}p{0.22\textwidth} >{\raggedright\arraybackslash}p{0.24\textwidth} >{\raggedright\arraybackslash}p{0.13\textwidth}}
\caption{按症状的处置表（``详见''一列指向本章节号，或全书章号与附录）。\label{tab:pm04-trouble}}\\
\hline
症状 & 最可能的原因 & 第一步做什么 & 详见 \\ \hline
\endfirsthead
\hline
症状 & 最可能的原因 & 第一步做什么 & 详见 \\ \hline
\endhead
\hline
\endfoot
报告满屏 \Tp{[unknown]} & 符号、权限或 buildid 断链之一 & \Tp{cat} 两个权限开关；\Tp{perf buildid-list} 核对 & 3.1/3.2 \\
调用栈只有一帧，或成片断在库符号 & 用了 \Tp{fp} 展开但二进制没留帧指针 & 重编加 \Tp{-g} \Tp{-fno-{\Bk}omit{\Bk}-frame{\Bk}-pointer}，或换 \Tp{dwarf} & 2.1/2.2 \\
深递归程序的栈到中途消失 & 栈快照小于递归深度，DWARF 截断 & 给 \Opt{call-graph dwarf} 加大快照字节数，查 \Tp{perf{\Bk}\_event{\Bk}\_max{\Bk}\_stack} & 2.2 \\
会话结束报出大量丢失样本 & 环形缓冲溢出：频率过高或栈快照过大 & 降 \Tp{-F}，必要时加 \Opt{buffer-size}；重跑 & 2 与 5.4 \\
同一程序两次跑热点排名互换 & 样本数不足，或阶跃偏置（热、频率漂移）& 看 Samples 绝对数；交替重跑取配对 & 1.1，第~\ref{ch:pm03-stats}~章 \\
按缺页采样的图与按时间的图结论相反 & 事件节拍把样本偏置到事件位置 & 先声明口径：两张图回答的是两个问题 & 1.2 \\
\Tp{memcpy\_erms} 或 \Tp{std::} 符号占大头 & 库顶帧效应，尚无归因 & 展开调用树找调用方，别先动 \Lib{libc} & 4.3 \\
异步、协程程序火焰图全挂在调度器下 & 续延入口没有逻辑栈 & 改用带任务上下文的埋点，或分段各测一段 & 6.1 \\
CI 或容器里根本没有 \Tp{perf} 或没事件权限 & 共享 runner 无 PMU（本册样例因此全是示意）& 改软件事件口径，或本地带数据复现 & 第~\ref{ch:pm13-ci}~章、附录 A \\
\end{longtable}

最后把本章放回全书的位置：采样剖析把你送到\emph{代码位置}附近，
从位置到成因需要事件计数与微架构分层（第~\ref{ch:pm05-events}~章）、
布局与分支的归因（第~\ref{ch:pm06-micro}~章）；
而把这套方法压进生产环境的低开销预算，是第~\ref{ch:pm14-production}~章的内容。
本章唯一能算``硬结论''的一句话是地基三句的展开版：\emph{没有可复现的样本构成
（节拍的类型与频率、栈的展开方式、符号与 buildid 的对应、丢失计数），
剖析报告里的任何百分比都只是叙事}。下一章从计数开始。

\chapter{事件计数与微架构分层：PMU、TopDown 与三工具对照}
\label{ch:pm05-events}

第~\ref{ch:pm04-sampling}~章给的是\textbf{位置}：采样剖析告诉你时间花在哪个函数、哪一行。
本章给的是\textbf{阶段}：事件计数（event counting）不告诉你哪一行慢，它告诉你这段时间里 CPU
是在\emph{执行有用的指令}、在\emph{猜错之后回滚}、在\emph{等指令送进来}，还是在\emph{等数据送进来}。
两种证据互补，缺一个就会得出半对的结论：只看热点函数容易把 \Tp{memcpy} 当成病因，
只看指标容易把 \Mx{IPC} 当成答案。

本章要回答三件事：一次运行能同时拿多少个事件，为什么拿到之后\textbf{不能直接拿来比}
（第~\ref{sec:pm05-counters}~节）；最少需要哪几个事件，才能让一个数变成口径明确的结论，
\Mx{IPC} 在哪些情形下会误导（第~\ref{sec:pm05-minset}~节）；指标异常时怎么用 TopDown 分层
把它翻译成下一步动作，以及\textbf{怎么在自己的机器上确认哪些事件真的存在}
（第~\ref{sec:pm05-topdown}~节）。

本章同时是三工具对照章。事件计数不是 Linux 专属能力：Windows 侧有 VTune 与 ETW 两条路，
macOS 侧有 Instruments 模板与 \Tp{powermetrics}。三者的\textbf{抽象层次、权限模型与可用性差别}
比命令名差别重要得多（第~\ref{sec:pm05-tools}~节）。完整的三平台并列速查放在
附录~\ref{ch:pmapA-platform}，本章只把判断依据讲清。

一条前置警告：本章所有清单里的数字都是\textbf{形态示意}。写作与构建环境是 Windows 加 MiKTeX，
发布用的 CI runner 没有 PMU 权限、通常也没有 \Tp{perf}（见第~\ref{ch:pm13-ci}~章），所以任何一段
输出只能告诉你形状，不能告诉你数值应该是多少。

\section{事件、物理计数器与复用}
\label{sec:pm05-counters}

\subsection*{计数器是每核的少量寄存器，不是无限槽位}

PMU（Performance Monitoring Unit）在硬件里是一小把计数器寄存器，\textbf{每个逻辑 CPU 一套}。
通用计数器的数量随 CPU 型号变：x86 上常见的架构实现是 4 个通用加 2--3 个固定计数器，
部分服务器硅片给到 8 个；AMD 与不少 Arm 核给 6 个。这些数字\textbf{不要照抄}，用两条办法自查。

\begin{itemize}
  \item 软确认：\Tp{perf stat -v} 的详细输出会打印每个事件落到了哪个计数器编号
        （形如 \Tp{assigned}），数一下不同编号的个数就是可用槽位；
  \item 硬确认：\Tp{lscpu} 的型号与特性行，以及
        \Tp{/sys/bus/{\Bk}event\_source{\Bk}/devices{\Bk}/} 下的节点与厂商事件描述文件；虚拟化环境
        还要看宿主暴露了几个虚拟 PMU 计数器，虚拟机里通常\textbf{少于裸机}，症状是复用比例很低。
\end{itemize}

还有一件事会吃掉一个通用计数器：x86 上的 \Tp{nmi{\Bk}\_watchdog}（内核用来做硬锁检测的、
基于 PMU 的看门狗）。默认开启时你会少一个槽位。

\begin{lstlisting}[style=shellstyle, title={槽位、分配、复用与分组：四步自查（Linux）}]
# 1) is nmi_watchdog taking one general counter?
cat /proc/sys/kernel/nmi_watchdog

# 2) which counter did each event land on (report goes to stderr)
perf stat -v -e cycles,ref-cycles,instructions ./workload

# 3) more events than slots -> multiplexing, ratio shown per line
perf stat -e cycles,instructions,task-clock \
  -e context-switches,page-faults,branch-misses ./workload

# 4) events that will be divided against each other share one group
perf stat -e '{cycles,instructions,branch-misses}' ./workload
\end{lstlisting}

\subsection*{复用（multiplexing）与它带来的缩放}

请求的事件多于槽位时，内核不报错，而是把事件分成若干组，在时间片上\textbf{轮转}：某个窗口只计数
其中一部分。最后每个计数都按时间比例放大回总时长，\Tp{perf stat} 会在行尾打印启用比例，
形如 \Tp{(42.31\%)}。这意味着：

\begin{itemize}
  \item 这个数是\textbf{估计值}，成立前提是事件在整个窗口内均匀发生。
        有阶段性的负载（预热、主循环、收尾写盘）会让估计系统性偏斜；
  \item 比例越低，方差越大。看到 \Tp{(10.00\%)} 这一档的行，把它当定性线索，别当数据；
  \item \textbf{两个被不同比例缩放的事件相除是不可信的}。\Mx{IPC} 恰好是这种商：
        若 \Tp{cycles} 缩放 100\%、\Tp{instructions} 缩放 42\%，商是两个不同时间窗之比。
\end{itemize}

\begin{warning}[title={缩放比例不同的两个事件相除}]
一次跑 30 个事件看起来效率高，实际是把自己送进最坏的口径：所有派生量（\Mx{IPC}、分支未命中率、
每千周期未命中数）都可能来自两个不同的缩放比。正确做法是\textbf{分批跑}：每批事件数不超过
可用槽位，或者把需要相除的事件放进同一个事件组（见下一小节）。多跑几轮的代价由
第~\ref{ch:pm03-stats}~章的重复判据承担，比拿一个说不清的数便宜。
\end{warning}

\subsection*{事件组：全有或全无的原子调度}

花括号把多个事件写成一组。组的语义是\textbf{全有或全无}：内核要么在同一时间窗里同时计数组内
全部事件，要么整组都不上。于是组内各行缩放比例相同，它们的商才有意义。代价是组更容易因为
槽位不够而整组 \Tp{not counted}——尤其在系统级 \Tp{-a} 采集、两个同核兄弟线程都要计数、
或别的剖析器已经占了计数器的时候。分组到底有没有上、哪一组被整组丢弃，用 \Tp{perf stat -v}
看分配细节；分组相关的选项名随 \Tp{perf} 版本变，先跑 \Tp{perf stat -h | grep -i group}
确认本机叫什么。

\begin{guideline}[title={一次运行的事件预算}]
\begin{itemize}
  \item 先数槽位，再决定批大小；\textbf{要相除的事件必须在同一组}。
  \item 想知道阶段差异，用间隔输出（形如 \Tp{perf stat -I 1000}，毫秒粒度，
        选项名以 \Tp{perf stat -h} 为准）或分阶段埋点，而不是靠复用把三个阶段抹平。
  \item 系统级 \Tp{-a} 与进程级默认口径不同：进程级 \Tp{perf stat ./app} 会计入被派生的子进程；
        按线程用 \Tp{-t} 或 \Tp{--per{\Bk}-thread}，跨核对比用 \Tp{--per{\Bk}-socket}、
        \Tp{--per{\Bk}-core} 一类，可用性以 \Tp{perf stat -h | grep per-} 为准。
  \item 汇报时写出用了几个事件、几个槽位、缩放比例区间，否则读者无法复现你的数。
\end{itemize}
\end{guideline}

\section{最小可用事件集与 \Mx{IPC} 的读法}
\label{sec:pm05-minset}

\subsection*{不依赖具体型号的八个量}

下面这些事件名走的是 \Tp{perf} 的\textbf{通用事件层}（generic 与 software 层），不直接绑定硅片编码，
因此跨平台面最宽：\Tp{cycles}、\Tp{instructions}、\Tp{ref{\Bk}-cycles}、\Tp{task{\Bk}-clock}、
\Tp{context{\Bk}-switches}、\Tp{page{\Bk}-faults}，再加 \Tp{branch{\Bk}-instructions} 与
\Tp{branch{\Bk}-misses}。这八个够搭出第~\ref{sec:pm05-topdown}~节之前的全部判断。
注意通用名不一定映射得到硬件计数器，映射不到就是 \Tp{<not supported>}，
\textbf{而不是 0}（见第~\ref{sec:pm05-trouble}~节）。

\begin{lstlisting}[style=outstyle, label={lst:pm05-stat},
  caption={形态示意：\Tp{perf stat} 输出的骨架（数值非本仓库实测值）}]
 Performance counter stats for './fft_bench -n 18':

      1,024.37 msec task-clock:u           0.97 CPUs utilized
          3,120 context-switches
             47 cpu-migrations
          1,982 page-faults
2,884,171,640 cycles                       2.82 GHz
1,730,490,112 instructions                 0.60  insn per cycle
  412,880,111 branch-instructions
   11,023,774 branch-misses                2.67% of all branches
      <not counted> L1-dcache-load-misses  (0.00%)

   1.055312890 seconds time elapsed
   1.010002000 seconds user
   0.041001000 sys
\end{lstlisting}

清单~\ref{lst:pm05-stat} 只是形状示意，不来自本仓库任何实测；比率列的写法与单位随版本变，
以本机输出为准（本册用空白代替原输出里的分隔符，免得与注释符号混淆）。
逐字段口径见表~\ref{tab:pm05-field}。

\begin{longtable}{>{\raggedright\arraybackslash}p{0.24\textwidth} >{\raggedright\arraybackslash}p{0.25\textwidth} >{\raggedright\arraybackslash}p{0.32\textwidth}}
\caption{事件计数输出的字段口径}\label{tab:pm05-field}\\
\hline
字段 & 是什么口径 & 读的时候注意什么 \\
\hline
\endfirsthead
\hline
字段 & 是什么口径 & 读的时候注意什么 \\
\hline
\endhead
\hline
\endfoot
\hline
\endlastfoot
\Tp{task{\Bk}-clock} & \MCpu{}：被监测工作占用的 CPU 毫秒数（软件事件，含子进程） &
  多线程下它是各线程时间之和，可以大于墙钟；与墙钟之比就是行尾的利用率 \\
\Tp{context{\Bk}-switches} & 上下文切换次数（主动与被动合并） &
  高不等于抖动；自旋与主动让出都会抬高它，也可能什么都没坏 \\
\Tp{cpu{\Bk}-migrations} & 线程换核次数 &
  对缓存与频率敏感；绑核（\Tp{taskset}）后应显著下降，第~\ref{ch:pm09-concurrency}~章要用 \\
\Tp{page{\Bk}-faults} & 缺页总数（\Tp{minor{\Bk}-faults} 与 \Tp{major{\Bk}-faults} 是它的细分） &
  预热阶段的缺页与稳态阶段的缺页含义完全不同，要分阶段看 \\
\Tp{cycles} & 核心时钟周期数 &
  随频率变：同一份代码在两种频率下 \Tp{cycles} 不同，与代码无关 \\
\Tp{instructions} & 退休指令数（retired，不是译码条数也不是微操作数） &
  Intel 与 AMD 的计数粒度不同，跨厂商不要直接比大小 \\
\Tp{ref{\Bk}-cycles} & 参考周期，一条与倍频解耦的时基 &
  与 \Tp{cycles} 相除得到平均运行倍率，是判断频率漂移最便宜的证据（第~\ref{ch:pm02-clock}~章） \\
\Tp{branch{\Bk}-misses} & 预测错误的分支数 &
  行尾百分比是相对 \Tp{branch{\Bk}-instructions} 的，两者必须同组、同缩放 \\
行尾 \Tp{(42.31\%)} & 该事件被启用的时间比例，即复用缩放依据 &
  比例低于一半就当线索；没有比例的行是 100\% \\
\Tp{seconds user} 与 \Tp{sys} & 内核给的粗粒度 CPU 时间分账 &
  与 \Tp{task-clock} 口径不同（统计时机也有差），不要混用两条去推利用率 \\
\end{longtable}

\subsection*{\Mx{IPC} 怎么读，以及它的四种误导}

\Mx{IPC} 就是 \Tp{instructions} 除以 \Tp{cycles}，读作\textbf{每周期退休几条指令}。它的价值在于
归一化掉了频率，因此可以跨两次运行比较；它的危险在于把多得多的东西压成了一个数。
四种典型误导：

\begin{enumerate}
  \item \textbf{跨工作不可比}。两个版本干的事不同，\Mx{IPC} 上升可能只是执行了更便宜的指令。
        只有在同一份输入、同一算法下，\Mx{IPC} 的差才是代码的差。
  \item \textbf{与频率耦合的错觉}。降频让 \Tp{cycles} 变少，\Mx{IPC} \emph{看起来}升高，而墙钟更慢。
        先用 \Tp{cycles} 与 \Tp{ref-cycles} 之比把这一项分离出来，再看 \Mx{IPC}。
  \item \textbf{高 \Mx{IPC} 不等于快}。已经贴着发射宽度跑的程序只能靠减少指令数获益；
        低 \Mx{IPC} 的程序也可能因为并行度更高而总吞吐更好。
  \item \textbf{同核兄弟线程下的歧义}。两个线程同时跑时各自的 \Mx{IPC} 都低于单跑，但总吞吐
        可能更高；这是共享发射资源的正常形状，展开见第~\ref{ch:pm09-concurrency}~章。
\end{enumerate}

内存受限程序的典型形状是：\Tp{cycles} 很多、\Tp{instructions} 不涨、\Mx{IPC} 明显低于这颗核标称的
发射上限（量级取决于型号与配置，读者可用 \Tp{perf list metric} 里厂商给的峰值指标自查），同时
分支错误占比正常。这时把 \Mx{IPC} 当成优化目标是错的——该去查访存层次与数据布局，
第~\ref{ch:pm06-micro}~章接手。只用软件事件起手，拿到的是 \MCpu{}、切换、缺页以及与墙钟的对齐
关系，微架构层面的判断一件都做不了；八个通用事件起手，才谈得上 \Mx{IPC}、分支错误率与频率漂移。

\section{TopDown 分层与在自己的机器上确认事件}
\label{sec:pm05-topdown}

\subsection*{四桶思想}

TopDown 是一族方法的统称（Intel 的 TMA 方法学把它做成了产品，但这个思想不专属 Intel）：
把流水线看作每周期若干个\emph{槽位}（slots），每个槽位在这一拍落到四个桶之一。

\begin{itemize}
  \item \textbf{retiring}：这条指令真的退休了，槽位没白花；
  \item \textbf{bad speculation}：槽位花在了后来被丢弃的路径上（分支猜错、机器清除、恢复）；
  \item \textbf{frontend bound}：槽位饿着，因为指令没能送进取指与解码通路（一级指令缓存、指令侧
        换入缓冲、解码限流、代码布局差、恢复后重新取指）；
  \item \textbf{backend bound}：槽位饿着，因为执行资源或数据没到位（访存、端口争用、依赖链长、
        寄存器压力导致的溢出）。
\end{itemize}

第一层的四个百分比近似构成一个分配（工具会给出各桶之和偏离满值时的说明），往下还有第二、
第三层。它比单个事件更接近\textbf{该改什么}这个问题，因为一个未命中计数既可能来自前端也可能
来自后端，而分层先替你做完了归属这一步。形状见清单~\ref{lst:pm05-td}。

\begin{lstlisting}[style=outstyle, label={lst:pm05-td},
  caption={形态示意：TopDown 第一层的输出形状（非实测值，条形长度只表示相对高低）}]
  Metric group: TopdownL1
  retiring          | ##                 (low)
  bad speculation   | #
     - branch mispredicted
  frontend bound    | ##########
  backend bound     | ################
     - core bound
     - memory bound

  Level 1 only says which direction to split next.
  It never says which function, which line, which array.
\end{lstlisting}

\subsection*{不抄事件表：在自己的机器上发现可用事件}

事件名随内核版本、\Tp{perf} 版本与 CPU 型号变。抄一张网帖里的事件表，最常见的结果是一整屏
\Tp{<not supported>}。可执行的发现路径见清单~\ref{lst:pm05-discover}。

\begin{lstlisting}[style=shellstyle, label={lst:pm05-discover},
  title={事件与指标发现（Linux）}]
# list by layer: hardware / software / cache events
perf list hw
perf list sw
perf list cache

# derived metrics and metric groups: TopDown lives here
perf list metric | grep -i -E 'topdown|slots'
perf list metricgroup | grep -i -E 'Topdown|Pipeline'

# chip-level aliases: each name maps to one raw encoding
grep -r . /sys/bus/event_source/devices/cpu/events/

# the decisive step: a name existing does not mean it can count
perf stat -e cycles,instructions true
perf stat -M TopdownL1 true
\end{lstlisting}

两个要点：\Tp{perf list metric} 与 \Tp{perf list metricgroup} 里的条目名才是稳定入口，底层编码
由描述文件负责，你不必记；\Tp{true} 只跑极短，用它验证\textbf{能不能计数}就足够——看的是
\Tp{not supported}、\Tp{not counted} 或权限错误，而不是数值。

分层在不同厂商上\textbf{不是同一套}实现。Intel 侧常见入口是 \Tp{perf stat --topdown} 或 \Tp{-M}
指标组（是否支持看上一条），VTune 把它做成一个分析类型；AMD 侧没有同名的 TMA 产品，可用厂商
给的指标、AMD uProf，或退到 \Tp{stalled{\Bk}-cycles{\Bk}-frontend} 与
\Tp{stalled{\Bk}-cycles{\Bk}-backend} 这类通用名（很多平台上它们是派生量或不映射，输出
\Tp{<not supported>}）；Arm 侧同样先看 \Tp{stalled{\Bk}-cycles-*} 是否映射，核内事件编号
（形如 \Tp{r0023}）随核型号变，必须现场用 \Tp{perf list} 确认，四桶模型往往要自己用前端阻塞
与后端阻塞两类计数搭。

\begin{bestpractice}[title={把分层当成筛子，而不是结论}]
\begin{itemize}
  \item 先分层再采样：分层告诉你该查后端还是前端，采样再去那一路里找位置。
        反过来做（先看到热点再猜原因）就是第~\ref{ch:pm15-antipattern}~章列的归因偷换。
  \item 只在同一版本、同一频率状态下比较各桶百分比。频率、核型号、同核兄弟线程、别的剖析器
        都在改分母。
  \item 分层事件本身要占计数器（部分平台占专用槽，部分平台占通用槽），症状是别的事件被挤去
        复用、缩放比例暴跌；确认方法是加上分层事件后重跑 \Tp{perf stat -v} 看分配是否变了。
  \item 汇报时写清用了哪一层、哪个指标组、哪个工具版本——否则这份分层对下一个读者不可用。
\end{itemize}
\end{bestpractice}

\section{同一事件的三种拿数方式：计数、采样、精确采样}
\label{sec:pm05-modes}

\Tp{perf stat} 与 \Tp{perf record} 可以用\textbf{完全相同}的事件名，产出两种完全不同的证据。
计数只回答\textbf{一共发生了多少次}；采样回答\textbf{在哪个指令地址发生的}；精确采样还回答
\textbf{哪一次访存、带着哪个数据地址}。三者口径见表~\ref{tab:pm05-mode}。

\begin{longtable}{>{\raggedright\arraybackslash}p{0.16\textwidth} >{\raggedright\arraybackslash}p{0.26\textwidth} >{\raggedright\arraybackslash}p{0.18\textwidth} >{\raggedright\arraybackslash}p{0.22\textwidth}}
\caption{三种拿数方式的口径与代价}\label{tab:pm05-mode}\\
\hline
方式 & 拿到什么 & 开销类型 & 常见失效 \\
\hline
\endfirsthead
\hline
方式 & 拿到什么 & 开销类型 & 常见失效 \\
\hline
\endhead
\hline
\endfoot
\hline
\endlastfoot
计数（\Tp{perf stat}） & 总量与派生比（\Mx{IPC}、错误率） &
  几乎只有计数器读写开销，对行为扰动最小 & 没有时间轴也没有位置；多事件被复用缩放 \\
采样（\Tp{perf record}） & 触发事件时的指令指针与调用栈 &
  每若干个事件一次中断：保存上下文、写环形缓冲 & 周期或频率选错就漏帧或淹帧 \\
精确采样（\Tp{:pp}、\Tp{:ppp}） & 事件真正发生时的地址，而不是恢复之后的地址 &
  高于普通采样，依赖 PEBS 或 IBS 一类机制 &
  虚拟机与受限权限下常直接不可用；缓冲压力更大 \\
带数据级别（\Tp{--data}） & 附带被访问的数据地址等附加信息 &
  每个样本更重，墙钟扰动肉眼可见 & 需要硬件事件与精确采样配合，记录体积明显变大 \\
\end{longtable}

\subsection*{周期、频率与精确后缀}

指定采样疏密有三个入口：\Tp{-F}（每秒样本数，形如 \Tp{-F 997}）、\Tp{-c}（每多少个事件一个样本，
形如 \Tp{-c 100003}）、以及时间类周期选项（可用性随版本变）。\Tp{-F} 的密度按\textbf{时间}定，
负载一慢样本就变密；\Tp{-c} 的密度按\textbf{事件量}定，跨版本比较时更容易保持一致。
工具的默认值不是结论：默认频率在短负载上只采到几百个样本是常态，不足以支撑任何占比结论
（判据回到第~\ref{ch:pm03-stats}~章）。

把普通采样变成精确采样是在事件名后加后缀，原始事件写法是把标志编进事件串，两条都见
清单~\ref{lst:pm05-precise}；编码\textbf{务必}从 \Tp{perf list} 或
\Tp{/sys/bus/{\Bk}event\_source{\Bk}/devices{\Bk}/cpu{\Bk}/events/} 现场取，不要抄。
特权级过滤用同一层后缀：\Tp{:u} 只算用户态，\Tp{:k} 只算内核态——\Tp{cycles:u} 与 \Tp{cycles}
在 \Tp{perf\_event{\Bk}\_paranoid} 较严的系统上是\textbf{能不能跑}的分水岭。

\begin{lstlisting}[style=shellstyle, label={lst:pm05-precise},
  title={精确采样与原始事件写法（编码随型号变）}]
# exact sampling via the :pp suffix
perf record -e cache-misses:pp ./app

# same thing written with a raw event, plus privilege filter
perf record -e 'cpu/event=0x3c,umask=0x0,pp/' ./app
perf record -e cycles:u -c 100003 ./app

# data-address level costs more per sample
perf record --data -e cache-misses:pp ./app
\end{lstlisting}

\subsection*{软件事件与没有 PMU 权限的环境}

\Tp{perf list sw} 列出的软件事件不碰硬件计数器，因此在无 PMU 权限处也能用。八个常用项见
清单~\ref{lst:pm05-sw}，其中两个值得单拎：\Tp{alignment{\Bk}-faults} 统计非对齐访存导致内核介入
的次数（多数架构上这条路径很贵，而它是个便宜的探测器）；\Tp{emulation{\Bk}-faults} 统计指令被
软件模拟的次数，非零基本意味着你在模拟层或非法指令路径上，此时性能数没有意义。

\begin{lstlisting}[style=shellstyle, label={lst:pm05-sw},
  title={无 PMU 也可用的软件事件与权限自查}]
# software events: no PMU counters involved
perf stat -e cpu-clock,task-clock,clock-resolution ./workload
perf stat -e context-switches,cpu-migrations,page-faults ./workload
perf stat -e minor-faults,major-faults ./workload
perf stat -e alignment-faults,emulation-faults ./workload

# current strictness (meaning of each value varies by kernel version)
cat /proc/sys/kernel/perf_event_paranoid

# can hardware events actually be obtained here?
perf stat -e cycles,instructions ./workload

# container and capability hints (names vary by distro)
grep CapEff /proc/self/status
\end{lstlisting}

Linux 侧的门槛是 \Tp{perf\_event{\Bk}\_paranoid} 与能力位（较新内核走 \Tp{CAP{\Bk}\_PERFMON}，
老环境看 \Tp{CAP{\Bk}\_SYS{\Bk}\_ADMIN}）；容器还叠加 seccomp 或 AppArmor 对
\Tp{perf\_event{\Bk}\_open} 的拦截，所以\textbf{能进容器}不等于\textbf{能读 PMU}。
虚拟机默认不暴露虚拟 PMU，除非宿主显式透传。\textbf{本仓库的 CI runner 就属于这一类}：既没有
PMU 权限也常常没有 \Tp{perf}，那里能拿的只有软件事件与墙钟——这正是本章反复强调\textbf{怎么确认}
而不是抄表的原因。

附带一条口径纪律：ASan、TSan、UBSan 跑起来的进程带着巨大的影子与拦截开销，在那些配置上读到的
事件计数只能用来解释\textbf{为什么慢}，绝不能当性能数（用法回指《内存泄漏检测与防御》册，
CI 门禁回指《LSP 与代码质量工具》第 12 章）。

\section{三工具对照：\Tp{perf stat}、VTune 与 Instruments}
\label{sec:pm05-tools}

\begin{compare}[title={事件计数在三个平台上的抽象与门槛}]
\begin{tabular}{>{\raggedright\arraybackslash}p{0.14\textwidth} >{\raggedright\arraybackslash}p{0.21\textwidth} >{\raggedright\arraybackslash}p{0.21\textwidth} >{\raggedright\arraybackslash}p{0.21\textwidth}}
 & \textbf{Linux / perf} & \textbf{Windows / VTune} & \textbf{macOS / Instruments} \\
事件抽象 & 原始事件名、通用层、指标层（\Tp{perf list}） &
  分析类型加结果页里的内部指标名 & 模板加固定指标；Intel 机型曾提供计数器模板 \\
TopDown 族 & \Tp{--topdown} 或 \Tp{-M}（取决于版本） &
  微架构探索类分析直接给分层 & 无同名分层，只有系统工具给的聚合量 \\
权限模型 & \Tp{perf\_event{\Bk}\_paranoid}、\Tp{CAP{\Bk}\_PERFMON}、容器策略 &
  需装采样驱动，通常要管理员装一次，之后普通用户可用 &
  Instruments 走开发者安全开关；\Tp{dtrace} 需 root 且受完整性保护限制 \\
无 PMU 时 & 软件事件仍可用 & 退化到基于 ETW 的采样与调度证据 &
  \Tp{powermetrics} 与时间类模板仍可用 \\
常见坑 & 抄了别的型号的事件名；把复用缩放当数据 &
  驱动没装就断定采不到；符号路径没配 & 把模板给的聚合百分比当成可复现指标 \\
\end{tabular}
\end{compare}

\subsection*{Windows：两条路，先分清有没有 PMU}

VTune 这条路读的是硬件计数器，依赖它安装的采样驱动；驱动没装好的表现通常是该类分析跑不出
微架构指标。典型形态见清单~\ref{lst:pm05-win}（分析类型名各版本改过，\Tp{uarch-} 与
\Tp{microarchitecture-} 这类别名关系以现场输出为准）。

\begin{lstlisting}[style=shellstyle, label={lst:pm05-win},
  title={Windows 侧形态示意（命令行以本机 vtune -h 为准）}]
vtune -collect-help
vtune -collect microarchitecture-exploration -- app.exe
vtune -collect general-exploration -result-dir r1 -- app.exe

# ETW route: no PMU needed, gives CPU time and scheduling evidence
# usually requires an elevated session
wpr -start CPU -start Profile
wpr -stop trace.etl
xperf -on Base+Profile
xperf -stop -d trace.etl
\end{lstlisting}

ETW 一侧用 \Tp{xperf} 或 \Tp{wpr} 采集，没有退休指令计数，它给的是 Profile 采样、上下文切换、
延迟过程调用与中断服务例件事件：能把 \MCpu{} 与调度分账讲清，不能把 \Mx{IPC} 讲清。
两者互补，不是替代。栈展开机制与三种展开的对照见第~\ref{ch:pm04-sampling}~章。

\subsection*{macOS：计数器面最窄，先把口径想清楚}

\begin{lstlisting}[style=shellstyle, label={lst:pm05-mac},
  title={macOS 侧形态示意（字段名随机型与系统版本变）}]
# whether developer profiling tools may run without root
sudo DevToolsSecurity --status

# per-cluster frequency and CPI-like aggregates
sudo powermetrics --listsamplers
sudo powermetrics -i 1000 -n 3 --samplers cpu_power

# command-line capture of a time-based template
instruments -t 'Time Profiler' -o run.trace ./app
\end{lstlisting}

要点：Apple Silicon 上 \Tp{powermetrics} 给的是\textbf{按簇}（性能簇与能效簇）的聚合频率与 CPI 类
指标，不是\textbf{按核按事件}的原始计数；Instruments 的性能计数器模板在 Intel 机型上存在过，
在当前机型与当前系统上能不能用，判据是\textbf{模板列表里能不能搜到它}，而不是任何一篇文章的说法。
所以 macOS 上的分层常常要退到 \Tp{powermetrics} 的 CPI 走向加 Instruments 的时间分账，再配同机
交替对照（第~\ref{ch:pm03-stats}~章）。这条限制不是\textbf{测不了}，是\textbf{换一套口径测}；
清单~\ref{lst:pm05-mac} 的头两条命令就是先确认这两条路在这台机器上能不能走。

\begin{warning}[title={三个平台共同的假象}]
三张工具截图看起来都在说同一件事时，它们可能共用同一个错误前提：机器频率不同、输入数据不同，
或者都被剖析器自己扰动了。跨平台对照时先把 \MCpu{} 与 \MWall{} 对齐，把频率证据
（\Tp{cycles} 与 \Tp{ref-cycles} 之比、\Tp{powermetrics}、\Tp{turbostat} 任一）钉住，
再比微架构指标。工具之间能对照的是\textbf{分层的形状}，不是百分比数值。
\end{warning}

\section{排障与速查}
\label{sec:pm05-trouble}

\begin{warning}[title={\Tp{not supported} 被当成\,\textbf{零开销}}]
\Tp{<not supported>} 意思是这个事件名在你的内核与硅片组合上没映射到计数器，\textbf{不代表这件事
没发生}；\Tp{<not counted>} 意思是事件存在但这一轮没被调度上（通常是分组或槽位不够），同样
不代表 0。把两者读成 0，就会得出\textbf{分支预测没问题}、\textbf{一级缓存未命中可以忽略}这类
彻底反向的结论。看到非数值的格子，先回去确认名字与权限（第~\ref{sec:pm05-topdown}~节、
第~\ref{sec:pm05-modes}~节末），再谈解释。
\end{warning}

把本章踩过的坑按症状归拢，就是表~\ref{tab:pm05-trouble}：每一行都给\textbf{下一步}，而且下一步
几乎都是回到第~\ref{sec:pm05-topdown}~节那套现场发现流程，而不是去网上再抄一张事件表。

\begin{longtable}{>{\raggedright\arraybackslash}p{0.25\textwidth} >{\raggedright\arraybackslash}p{0.24\textwidth} >{\raggedright\arraybackslash}p{0.30\textwidth}}
\caption{事件计数排障表}\label{tab:pm05-trouble}\\
\hline
症状 & 最可能的原因 & 下一步 \\
\hline
\endfirsthead
\hline
症状 & 最可能的原因 & 下一步 \\
\hline
\endhead
\hline
\endfoot
\hline
\endlastfoot
满屏 \Tp{<not supported>} & 抄了别的 CPU 型号的事件名 &
  用 \Tp{perf list hw} 与 \Tp{/sys/bus/{\Bk}event\_source{\Bk}/devices{\Bk}/cpu{\Bk}/events/}
  现场取名字 \\
某行 \Tp{<not counted>}，旁边全是小百分比 & 事件多于槽位，整组没被调度 &
  缩小批大小；把要相除的事件放进同一组；或者分批跑 \\
\Mx{IPC} 两次差异大但墙钟没变 & 两行缩放比例不同，或频率漂移 &
  先核对行尾比例是否一致；加 \Tp{ref-cycles} 算周期之比 \\
\Mx{IPC} 明显低于该核上限 & 后端受限（访存或依赖链） &
  转分层：\Tp{perf list metricgroup} 找 TopDown 组；再回第~\ref{ch:pm06-micro}~章 \\
各桶百分比之和远离满值 & 工具用了估计层，或分层事件被挤掉 &
  读工具的说明行；把分层事件单独一批跑，别和别的事件混在同一组 \\
权限错误 & \Tp{perf\_event{\Bk}\_paranoid} 较严、容器拦截
  \Tp{perf\_event{\Bk}\_open}、虚拟机没有虚拟 PMU &
  读 \Tp{/proc/{\Bk}sys{\Bk}/kernel/{\Bk}perf\_event\_paranoid}；先用软件事件跑通，再谈提权 \\
\Tp{:pp} 报不可用 & 精确采样机制在虚拟机或受限权限下没暴露 &
  退回普通采样或纯计数；需要精确地址时换环境 \\
\Tp{task-clock} 大于墙钟 & 多线程的线程时间求和，\MCpu{} 本来可以大于 \MWall{} &
  这是口径不是错误；利用率歧义见第~\ref{ch:pm09-concurrency}~章 \\
VTune 跑完没有微架构指标 & 采样驱动没装好或没生效 &
  用 \Tp{vtune -collect-help} 与工具自带检查确认驱动，再决定是否退回 ETW 路线 \\
Instruments 找不到计数器模板 & 机型或系统不提供，或安全开关没放开 &
  \Tp{DevTools{\Bk}Security --status} 与 \Tp{powermetrics --{\Bk}listsamplers}，
  改口径而不是硬找 \\
同一命令今天与昨天差很多 & 频率、热限、后台邻居、内核更新 &
  记录环境指纹（CPU 型号、调频策略、内核与编译器版本）后再比，模板见
  附录~\ref{ch:pmapC-glossary} \\
只有本册样例对得上、自己机器对不上 & 你的事件名与槽位数本来就不同 &
  对形状，不对数；一切按第~\ref{sec:pm05-topdown}~节的现场发现流程走 \\
\end{longtable}

\begin{guideline}[title={本章的最小起手式}]
\begin{enumerate}
  \item \Tp{perf list hw} 与 \Tp{perf list metricgroup} 各跑一次，记下这台机器\textbf{有什么}；
  \item 一批跑完第~\ref{sec:pm05-minset}~节列出的那八个通用事件，批大小不超过槽位；
        命令行形态照第~\ref{sec:pm05-counters}~节的四步自查清单改事件名即可；
  \item 先取三个数：\Mx{IPC}、周期与参考周期之比、\Tp{task{\Bk}-clock} 与墙钟之比；
  \item 指标异常再用分层组（\Tp{-M} 或 \Tp{--topdown}，支持则用）确定方向；
  \item 方向定了以后，交回第~\ref{ch:pm04-sampling}~章的采样去定位位置，
        或交回第~\ref{ch:pm06-micro}~章做缓存与分支级归因；
  \item 写结论时把事件清单、槽位数、缩放比例区间、内核与工具版本一并写进去。
        参数与字段的完整口径见附录~\ref{ch:pmapB-params}。
\end{enumerate}
\end{guideline}
\part{归因：从指标回到代码位置}

\chapter{缓存、分支与指令级归因}
\label{ch:pm06-micro}

第~\ref{ch:pm05-events} 章把你送到了一堆数字面前：\Mx{cache{\Bk}-misses} 多少、\Mx{branch{\Bk}-miss} 占比、\Mx{IPC} 是 $0.7$ 还是 $2.3$。
本章处理这些数字之后的那一步——\textbf{从指标回到代码位置，再从位置回到原因}。
位置可以由采样的地址归属给出，原因只能由代码自身的形态给出：数据怎么摆、依赖链多长、
分支是否依赖输入、循环有没有被向量化、这段 \Tp{.text} 离热区有多远。
本章范围收在\textbf{同一线程之内}：所有对照都假设你绑了核、跑的是单线程热循环。
跨线程的干扰（两个线程写同一条缓存线造成的伪共享、共享计数器被打爆）留给第~\ref{ch:pm09-concurrency} 章，
那边的判据是拓扑与并发度，不是布局。构建侧的开关——优化级别、LTO、PGO、内联的真实语义——在第~\ref{ch:pm10-compiler} 章，
本章只讲\textbf{怎么观察它们的后果}，不讲怎么配。
贯穿全章的一条纪律：指标是线索，不是结论；结论要么来自一次只动一个变量的布局差分，
要么来自发布产物上的反汇编，两者都没有时，手上那个数只是相关性（第~\ref{ch:pm15-antipattern} 章把它列为归因偷换）。

\section{从 miss 率到代价：把 miss 换算成周期}
\label{sec:pm06-misscost}

\subsection{miss 率为什么单独读不出来}

命中率是一个\textbf{无量纲比例}，代价是一个\textbf{周期数}，两者之间隔着三个量：
每千条指令的访问数、被 miss 的那一层往下要走多远、以及这段代码本来每周期能干多少事。
$1\%$ 的 miss 率可以完全无关痛痒——如果循环本身是计算密集的，访存延迟会被执行单元藏起来；
$30\%$ 的 miss 率也可以不是瓶颈——如果它全落在 L1，而 L1 命中路径本来就是几个周期。
所以拿到 miss 数以后，第一个要问的不是``能不能把命中率提上去''，而是
\textbf{这些 miss 穿透到了哪一层、那一层每次要付多少周期}。

换算的思路可以写成正文句子：某一层带来的周期账，约等于落到该层的访问次数乘以该层的访问延迟；
总 stall 是各层这一项的和，再加上被乱序窗口吸收掉的那部分（吸收不掉才成为关键路径上的时间）。
粗略量级（\textbf{只作判断用，不许写进结论}）：L1 通常个位数周期、L2 几十、L3 上百、DRAM 数百，
相邻层次之间接近一个数量级的跳变而不是平滑增长。真正吞时间的从来不是 ``L1 没命中''，
而是 ``没命中之后走到哪儿停了''。要把这件事变成可用的数，你需要两样东西：各层的命中率、各层的延迟。
前者本章下面讲怎么取，后者这一节先讲怎么确认。

\subsection{层次、容量与延迟：先量自己这台机器}

不同代际、不同微架构、不同云主机的层次差别很大，\textbf{共享关系}也会变（L2 是每核独占还是 cluster 共享、
vCPU 数与实际拓扑是否一致）。抄一张网上的延迟表是最常见的错误起手式。可执行的确认路径如下。

\begin{lstlisting}[style=shellstyle,caption={在自己机器上确认缓存层次与线路径},label={lst:pm06-hier}]
# 1) 层次与容量：线大小、组数、关联度、谁共享这一层
lscpu | grep -i cache
lscpu -C                          # 部分版本给 per-level 明细
# 更细的口径直接从 sysfs 取（字段名比 lscpu 稳定）
for L in /sys/devices/system/cpu/cpu0/cache/index*; do
  echo "-- $L"
  for F in level size type ways_of_associativity coherency_line_size \
           number_of_sets shared_cpu_list; do
    printf '%-26s %s\n' "$F" "$(cat "$L/$F" 2>/dev/null)"
  done
done

# 2) 本机有哪些缓存/TLB 事件（不抄表，现场列）
perf list                # 关注 cache 族与 dTLB/iTLB 族
cat /sys/devices/cpu/events/cache/*   2>/dev/null   # raw 编码

# 3) 空载基线：本底 miss 与本底 branch-miss，后面所有差分都要减它
perf stat -e cycles,instructions,cache-references,cache-misses,\
branch-instructions,branch-misses -- taskset -c 2 sleep 5
\end{lstlisting}

层次延迟不能从这些文件里读出来，只能\textbf{测}：构造一个只依赖指针 chasing 的环形数组，
工作集分别落在 L1 / L2 / L3 / DRAM，每步跑一次 load，用第~\ref{ch:pm02-clock} 章的计时口径取每次 load 的周期数中位数；
工作集按 $2$ 的幂扫过去，曲线出现台阶的位置就是你机器上各层的\textbf{实际容量边界}，
台阶之间的高度差就是各层的\textbf{实际延迟}。这个实验的骨架与第~\ref{ch:pm07-memory} 章的局部性实验共用，
数值由你自己填——本章通篇不给任何具体周期数。

\begin{guideline}[title={要点：先换算，再谈命中率},label={box:pm06-cost}]
\begin{itemize}
\item miss 率无量纲，代价才是周期；先问 ``落到哪一层''，再谈 ``命中率好不好''。
\item 层次容量与延迟每台机器不同，必须自己测（\Tp{sysfs} 取层次、指针 chasing 取延迟），不要抄表。
\item 事件名与字段格式都随内核与 CPU 变；上面的清单给的是\textbf{探测方法}，不是可用事件表。
\item 任何差分之前先取空载基线：本底 \Tp{cache{\Bk}-misses} 与本底 \Tp{branch{\Bk}-misses} 都要减掉。
\item 换算只在 ``单线程、绑核、频率已知'' 时成立；并发与提频的影响分别在第~\ref{ch:pm09-concurrency} 章与本章末尾。
\end{itemize}
\end{guideline}

\section{局部性归因：步长、遍历方向与布局}
\label{sec:pm06-locality}

\subsection{三个可操作变量}

同一线程的 \Mx{cache{\Bk}-misses} 变化，几乎总能归因到三个变量之一。
\textbf{步长（stride）}：沿访问序列相邻两次 load 的地址距离。步长为 $1$ 时一条线被拉进来、
后面的元素直接白拿；步长跨过线路径（隔一整线取一个元素）时每条线只用一次，
有效带宽掉到 ``线大小除以步长'' 这个比例。
\textbf{遍历方向}：\CL{} 与 \CPP{} 的多维数组都是行主序，按行扫是顺序的，按列扫每一步跨一整行——
列扫描就是 ``步长等于行跨度'' 的最坏情况。
\textbf{布局}：把用到的字段聚在一起（数组结构体，SoA），还是把对象聚在一起（结构体数组，AoS）。
循环只用两个字段时，AoS 会把整条线里用不到的字段一起搬进来，等效于把步长放大；
SoA 让每个字段各自连续，代价是跨字段访问要跳数组、以及打包/转置本身的时间。

\subsection{可跑骨架：遍历顺序与转置}

下面这个骨架只做三件事：顺序扫、隔线扫、按列扫（转置），\textbf{所有数值留空，由你在自己机器上填}。
它把三个变量隔离在同一个数组、同一种访问粒度上——这样 miss 的变化才只可能来自步长和方向。

\begin{lstlisting}[style=cppstyle,caption={局部性对照骨架：步长与遍历方向（数值由读者填写）},label={lst:pm06-traverse}]
#include <cstdint>
#include <cstdio>
#include <ctime>
#include <vector>

static volatile uint64_t g_sink;      // 反死代码消除：结果必须导出

static double now_sec() {
  timespec ts;
  // 时钟口径见第 2 章：单调、不受 NTP 步进影响的时钟源
  clock_gettime(CLOCK_MONOTONIC, &ts);
  return double(ts.tv_sec) + double(ts.tv_nsec) * 1e-9;
}

// N 取到工作集明显大于 L3；元素 8 字节，一条线容纳 8 个元素
int main() {
  constexpr int N = 4096;             // 4096*4096*8B = 128 MiB，按需缩小
  std::vector<uint64_t> a(size_t(N) * N);
  for (size_t i = 0; i < a.size(); ++i) a[i] = uint64_t(i);

  uint64_t acc = 0;
  // 预热：首次跑一定慢（缺页、缓存冷、频率爬升），不要计入
  for (int i = 0; i < N; ++i) acc += a[size_t(i) * N];
  g_sink = acc;

  double t;
  // (A) 顺序扫，步长 1：一条线拉进 8 个元素全被用上
  t = now_sec(); acc = 0;
  for (size_t i = 0; i < a.size(); ++i) acc += a[i];
  t = now_sec() - t; g_sink = acc;
  std::printf("A stride=1      %12.3f ms\n", t * 1e3);

  // (B) 隔线扫：步长 = 一条线的元素数，每条线只用 1 个元素
  t = now_sec(); acc = 0;
  constexpr int EPL = 8;              // 由 coherency_line_size 推出
  for (size_t i = 0; i < a.size(); i += EPL) acc += a[i];
  t = now_sec() - t; g_sink = acc;
  std::printf("B stride=line   %12.3f ms  (工作量只有 A 的 1/8)\n", t * 1e3);

  // (C) 按列扫：步长 = 一整行，跨层 miss 的典型来源
  t = now_sec(); acc = 0;
  for (int j = 0; j < N; ++j)
    for (int i = 0; i < N; ++i) acc += a[size_t(i) * N + j];
  t = now_sec() - t; g_sink = acc;
  std::printf("C column scan   %12.3f ms  (转置后应显著变化)\n", t * 1e3);
  return 0;
}
\end{lstlisting}

(B) 那一行有个坑必须写清楚：步长扫做的\textbf{工作量}只有顺序扫的 $1/8$，时间不能直接对比。
局部性对照里正确的比法是 ``每字节的周期'' 或 ``每千次访问的 miss''——先把工作量对齐再谈快慢。
这是 ``口径先于数字''（第~\ref{ch:pm03-stats} 章）在缓存层面的具体形式。

\subsection{结构体数组还是数组结构体：先量字节}

把 AoS 换成 SoA 之前，把两件事算出来：这个循环\textbf{实际用到的字节数}占对象总字节数的比例；
对象是否 ``正好装不满一条线''（比如对象 $104$ 字节、线 $64$ 字节，那么每访问一个对象至少跨两条线，
而字段打包后同一字段可能回到一条线里）。比例越低、跨线越频繁，SoA 的收益空间越大；
如果循环本来就用到了对象里所有字段，换布局只是把同一段内存换个写法，miss 几乎不动，
反而多了转置成本。

\begin{compare}[title={同一循环、两种布局的可观察形状（只给形状，不给数）},label={box:pm06-layout}]
循环体：对每个粒子累加 \Tp{p.x} 与 \Tp{p.y}，其余字段（\Tp{id}、\Tp{tag}、\Tp{pad[8]}）不使用。
\begin{itemize}
\item AoS：每个粒子连续排布，\Tp{x} 与 \Tp{y} 之间夹着无关字段；每次访问只用到线里一小段，
  \textbf{字节利用率低}，相邻粒子的 \Tp{x} 之间隔了整个对象。
\item SoA：\Tp{x} 数组与 \Tp{y} 数组各自连续；对 \Tp{x} 的扫描是纯步长 $1$，
  线拉进来就全部用完，\textbf{字节利用率高}。
\item 正确的写法：同一份循环，两组 \Mx{L1{\Bk}-miss} 与 \Mx{LLC{\Bk}-miss}、两个 \Mx{IPC}，
  工作量对齐后用 ``每千次访问 miss'' 对照。本框只给形状，数值由读者自己填。
\item 常见误判：看到 SoA 的 miss 低就写 ``SoA 更快''。如果关键路径上还有别的 stall
  （依赖链、分支），miss 降低可以一点也换不到时间。
\end{itemize}
\end{compare}

\subsection{prefetch 与非临时访问的失效条件}

软件预取（架构相关的 \Tp{\_\_builtin{\Bk}\_prefetch} 一族，强度 $0$ 到 $3$ 对应不同缓存层）能改变的只是
\textbf{访存延迟能不能被提前}，前提是这条线本来就会被用到。所以它在\textbf{计算受限}的代码上多半只添乱：
多指令、多访存队列压力、多带宽。在访存受限、模式可预测（连续或固定步长）的循环里才可能有效。
非临时存储（\Tp{movnti} 一类）绕过缓存直写内存，对 ``写一次就不再读'' 的大块输出能减少读入污染；
如果后面还要读回来，它反而制造一次额外的 DRAM 往返。
两个失效条件要记住：\textbf{访问被打断或换出}时预取收益直接归零（第~\ref{ch:pm09-concurrency} 章的调度侧）；
\textbf{访存队列深度有限}时过于激进的预取会让真正需要的 load 排队，收益不线性叠加。
是否有效只能用差分自己量：加与不加各跑、同机绑核、交替重复，按第~\ref{ch:pm03-stats} 章的区间判据比，不看单次。

\section{冲突未命中、线路径，以及与第 9 章的分界}
\label{sec:pm06-conflict}

\subsection{容量明明够但 miss 很高：先怀疑冲突}

关联度是有限的：一个地址在某层能落在哪几组，由地址的若干位决定。工作集总量远小于容量、
命中率却仍然很差时，最常见的解释是\textbf{冲突未命中}——大量同时使用的地址被映射到同一组，彼此挤出去。
典型触发形态：数组行跨度恰好是 $2$ 的幂、哈希桶指针按固定步长排布、两个热数组刚好错开整数倍的线数。
判断依据不是 ``看着像''，而是\textbf{把跨度改成奇数（多留一行 padding）再跑一次}：
冲突型 miss 会大幅回落，真正的容量不足几乎不动。这是布局层的差分实验，与换算法无关，
也是本章 ``归因需要证据'' 里最便宜的一个样例。

\subsection{线路径与映射：不抄位段表}

地址里参与组索引的那几位（随缓存层与关联度变，不同代际还会引入地址哈希）决定了 ``谁和谁抢''。
本章不抄这张位段表——它随型号变，抄错的代价比不抄更大。可执行的做法是：从 \Tp{sysfs} 拿到
``线大小、组数、关联度''，把这几个量当成\textbf{假设}，再用上面那条 ``改跨度再跑一次'' 的实验去检验；
如果模型说 ``这个跨度会冲突'' 而实验显示不冲突，以实验为准。

\subsection{拿不到硬件事件时怎么办}

\Tp{perf{\Bk}\_c2c} 一类工具（事件名随平台与内核变，用 \Fn{perf} list 现场探测）在 Intel 上依赖
\Tp{mem{\Bk}\_load{\Bk}\_retired{\Bk}\_latency{\Bk}\_ge{\Bk}\_N} 这类带阈值、需要精确取样
（\Tp{event:pp}）的事件，普通用户经常拿不到权限。确认顺序四步：
\begin{enumerate}
\item 工具支持吗：\Fn{perf} c2c 是否存在、子命令报什么错。
\item 事件可用吗：现场列本机 \Tp{memory{\Bk}-events} 与 \Tp{cache} 族，别抄表。
\item 权限够吗：内核采样权限开关在容器与 CI runner 里通常收紧（第~\ref{ch:pm05-events} 章；完整权限模型见附录 A）。
\item 都拿不到就退回\textbf{行为差分}：加对齐填充、把行跨度改奇数、拆开两个热数组——
  一次只动一个变量，看 miss 与 \MLat{} 是否按预期动。行为差分不需要 PMU 权限，
  代价是只能给出 ``与布局有关'' 这一档结论，给不出 ``哪条线被谁抢了''。
\end{enumerate}

\begin{note}[title={同一条缓存线，两套解释：本章与第 9 章的分工},label={box:pm06-vs09}]
\begin{itemize}
\item 本章（同一线程的布局）：一条线里混装了有用与没用的字段、地址被映射到同一组、跨度踩到 $2$ 的幂。
  处置手段是重排、拆字段、改跨度、对齐；判据是\textbf{单线程绑核就能复现}。
\item 第~\ref{ch:pm09-concurrency} 章（跨线程的共享）：两个线程写同一条线的不同字节，一致性协议让线在线程间来回迁移；
  共享计数器把 \Mx{cache{\Bk}-misses} 打爆。处置手段是按线程切分数据、填充到整线；
  判据是\textbf{只在并发下出现}，单线程数据与它无关。
\item 顺序：先做本章这一半（单线程），把布局因素钉干净，再谈并发那一半。
  反过来做会把两类问题混成一个 ``加 padding 就好了'' 的故事，事后无法复现。
\end{itemize}
\end{note}

\begin{lstlisting}[style=outstyle,caption={一段 \Fn{perf} stat 输出的字段读法（格式示例，非本仓库实测值）},label={lst:pm06-stat}]
 Performance counter stats for 'taskset -c 2 ./traverse':

          1,234.56 msec task-clock          # 这段代码在该 CPU 上被计量的时间
              4,567      context-switches   # 绑核单线程仍明显非 0 => 有抢占
                xxx      cycles
                xxx      instructions      # <- 与上一行相除得 IPC，不是第三份证据
                  xx     L1-dcache-load-misses
                  xx     LLC-load-misses    # 是否支持因核而异，not supported 不等于 0
           x.xx      task clock cycles per instruction

      1.25 seconds user
\end{lstlisting}

逐格口径：\Tp{task{\Bk}-clock} 是被计量代码在该 CPU 上的时间，不是墙钟；\Tp{context{\Bk}-switches} 明显非 0 时
所有按 CPU 归属的数据都要打折；\Mx{IPC} 是两个计数相除的结果，不构成独立证据；
带 not supported 的事件计数为 $0$ 时，含义是 ``本机没给你这个事件''，不是 ``开销为 0''（第~\ref{ch:pm05-events} 章的 0 陷阱）。

\section{分支归因：\Mx{branch{\Bk}-miss} 的读法与``消分支''的代价}
\label{sec:pm06-branch}

\subsection{先分清失败数、失败率与计数口径}

\Tp{branch{\Bk}-misses} 是绝对数，读它之前先除以 \Tp{branch{\Bk}-instructions} 得到 \Mx{branch{\Bk}-miss rate}。
两个坑：\textbf{计数口径}——现代乱序 CPU 上部分实现会把分支算多次，同一程序在不同平台报出的比率
不可直接横比，跨机器比这个数之前先比口径；\textbf{可靠性提示}——比率很低时 \Fn{perf} 会打印
计数不可靠的小字，那行字属于结论的一部分，不是日志噪声。
代价侧也不是常数：分支失败后要不要回滚已经乱序执行的窗口，取决于失败点之后已经投机跑了多少工作，
短循环里的失败与大段计算之后的失败，惩罚不是一个量级。

真正值得盯的是\textbf{数据依赖分支}。预测器能学好的是模式可重复的分支（循环出口、状态机跳转、
方向恒定的判断）；学不好的是\textbf{输入决定方向}的分支（阈值判断、符号判断、命中判断），
因为它可预测性由数据分布决定——同一个分支换个数据集就能从很低跳到接近一半。
所以任何 ``这个分支很准'' 的结论都必须带上数据集描述，否则只是 ``在我手上这批输入上很准''。

\subsection{把 miss 落到具体分支}

\Fn{perf} 的逐指令归属给出 ``每条指令的采样占比''，形态如下。它的位置可信度有前提：
二进制是发布用的那一个、符号与地址对得上、没有被跨函数调度打乱。而编译器恰恰会把数据依赖分支
\textbf{向量化}掉——向量化后的代码里根本没有那条分支了。所以看到某行归属很高，
先用反汇编确认那条指令真的存在，再谈优化。

\begin{lstlisting}[style=outstyle,caption={指令级分支归属的形态（格式示例，非本仓库实测值）},label={lst:pm06-annotate}]
 # perf annotate --symbol 归一化后的形态（示意）
 Percent |  Assembly of kernel_loop
   0.00  |   cmp     %rax,%rdx
   71.3   |*  jl      .Lslow          ; 方向由数据决定，采样归到这里
    0.00  | .Lfast:
   12.4   |   movzbl  (%rdi,%rax),%ecx
 # * 标记分支本身；Percent 是它吃掉的样本占比，不是"这条指令耗时"
\end{lstlisting}

更直接的路子是\textbf{分支栈记录}：\Fn{perf} 的 last-branch 类采样（Intel 的 LBR、AMD 的对应实现，
深度与可用性都随平台变，现场确认能不能采）。它给出 ``源地址到目标地址'' 的成对记录，
可以统计 ``哪个分支跳去了哪、多久一次''，比单个失败计数更接近 ``这段代码的控制流长什么样''。

\begin{lstlisting}[style=outstyle,caption={分支栈记录聚合后的形态（格式示例，非本仓库实测值）},label={lst:pm06-brstack}]
 # perf script 带分支栈字段后的逐条记录（from:to:spec:depth:.. 形态）
 0x401a3f  0x401a72  :-38  0  0  1  1
 0x401a72  0x401a45  :-12  0  1  0  1
 # 聚合 after->before 计数 = 真实控制流边集合
 # 读出的是"哪些分支互相把关键路径拉长"，而不是单个 miss 率
\end{lstlisting}

\subsection{消分支为什么可能更慢}

把 ``取绝对值'' 或 ``按条件选值'' 写成查表、位运算、无地板技巧，或者指望编译器生成 \Tp{cmov}，
动机都是把控制冒险换成数据通路。但这不是免费的：\textbf{两条路径的副作用都执行了}，
于是它多做两件事——算了你本来不必算的那一半，load 了你本来不会碰的那张表。
在访存受限的代码上，多出来的 load 可能触发 miss，比一次可预测的分支还贵；
在计算密集的代码上，多出来的运算挤占执行端口。
不同架构的无分支手段形态也不同（条件选择/条件置位指令族、混合类向量指令），
但\textbf{到底生成了什么只能从发布产物看}，不能从 ``我写了三目运算符'' 推。

\begin{warning}[title={``把分支改成查表必然快''},label={box:pm06-brwarn}]
\begin{itemize}
\item 查表引入一条 load：表大了就跑到 L1/L2 之外，表小也照样占用取指与访存队列。
  它把控制冒险换成访存依赖，在访存受限的代码上常是净亏。
\item \Tp{cmov} 不是 ``两边都算但只取一个'' 的安全保证：有副作用的分支（解引用、函数调用）
  不能被这样改写——正确性会坏，编译器本来也不会这么做。
\item 向量化会改变分支语义：掩码执行等价于两条路径都算再选。为 ``保住分支'' 写的代码，
  可能在另一个优化级别下被向量化掉，收益方向随之翻转（第~\ref{ch:pm10-compiler} 章讲开关口径）。
\item 判据只有一条：同一台机器、同一数据集、交替重复，\MLat{} 的区间与 \Mx{branch{\Bk}-miss} 一起看。
  区间不分开就没有结论，``看着快了一点'' 不是结论（第~\ref{ch:pm03-stats} 章）。
\end{itemize}
\end{warning}

\section{指令级归因：向量化、对齐与代码布局}
\label{sec:pm06-vec}

\subsection{向量化到底发生了没有：报告一族开关}

先钉死本章立场：\textbf{报告只能告诉你 ``没向量化''，不能告诉你 ``快了还是慢了''}。
它写的是优化器 ``认为'' 的事，和这段代码在你机器上跑多快之间还隔着执行端口宽度、
迭代次数够不够填满流水线、对齐情况与访存瓶颈。判断收益方向仍然要靠差分测量。

报告开关是一族，而且\textbf{开关名随版本变，落笔前先确认你手上的编译器认哪个}：
GCC 的向量化信息/备注一族（\Opt{fopt{\Bk}-info-vec} 及其 \Tp{-optimized} / \Tp{-missed} 变体）、
Clang 的 pass 报告一族（\Tp{-Rpass} 与 \Tp{-Rpass{\Bk}-missed} 配 \Tp{loop{\Bk}-vectorize}）、
MSVC 的向量化报告级别（\Tp{/Qvec-report} 系列）、Intel 编译器的对应开关。
这些名字在不同大版本里出现、改名或合并过，本章刻意不给 ``该敲哪个'' 的唯一命令行：
以你自己版本的文档与 ``help 里 grep vec'' 的结果为准。

\begin{lstlisting}[style=shellstyle,caption={向量化报告：先确认开关存在，再读结论},label={lst:pm06-vecreport}]
# 思路：编译时让优化器逐循环报告结论，然后只看两类：
#   (a) 已向量化的：确认就是你关心的那个函数（别在别的循环上高兴）
#   (b) 未向量化的：报告一般带原因（依赖 / 规约顺序 / 步长未知 / 体内有调用）
# 开关名随版本变。先列出来再敲：
$CC --help | grep -i vec            # GCC 一族
clang --help | grep -i pass         # Clang 一族；MSVC/Intel 看各自文档
# 形态示例（不保证你的版本就是这两个名字）：
#   gcc  -O3 -c kernel.c -fopt-info-vec-optimized -fopt-info-vec-missed
#   clang -O3 -c kernel.c -Rpass=loop-vectorize -Rpass-missed=loop-vectorize
# 报告是线索：把 missed 的原因逐条对应到下面三类成因，再改代码。
\end{lstlisting}

可以改的成因三类：\textbf{依赖与别名}——循环读写同一数组而下标关系证不出来，或存在跨迭代减法依赖；
\textbf{不能安全重组}——浮点归约顺序改变结果，除非接受语义代价（那是第~\ref{ch:pm10-compiler} 章的内容，它改变结果不只是改变速度）；
\textbf{循环形态}——步长不是编译期常量、尾部长度未知、循环体里有调用或不规则访存。
第一类的处置是把信息交回给编译器（\Tp{restrict}、对齐假设一类内建），
注意 \Tp{restrict} 是 ``我保证不重叠'' 的承诺，承诺错了是\textbf{结果错}而不只是慢；
第二类的处置是决定要不要改变数值语义；第三类多半得重写循环或分块。

\subsection{最终以发布产物的反汇编确认}

\begin{lstlisting}[style=shellstyle,caption={在发布产物上确认向量化与热区形态},label={lst:pm06-disasm}]
# 纪律：在真正用于发布/基准的那个二进制上看，不要在探针构建上看。
objdump -d --no-show-raw-insn ./bench        # 定位目标函数
#   看有没有宽寄存器的 load/算术/混合指令，还是仍然是标量循环体
nm --size-sort -C ./bench                    # 函数体积排序，配合热点定位
readelf -S ./bench | grep -i text            # 段与分区情况（代码布局实验用）
# 地址或符号对不上 => 你看的不是被测的那份产物，先核对构建产物再继续。
\end{lstlisting}

反汇编还能回答报告回答不了的问题：向量化是不是只发生在探针构建上、
\Tp{memcpy} 是否被展开成向量拷贝、尾部处理被复制了几份、
循环是不是被完全展开成直线代码（展开后已经没有 ``循环'' 可供归因，
而 \Tp{inline} 的语义本来首先是 ODR 与链接问题、不是提速开关，口径与第~\ref{ch:pm10-compiler} 章一致）。

\subsection{对齐、别名与一段流水示意}

未对齐的向量 load 在部分架构上直接非法，在另一些上合法但可能跨两条线、多花一拍；
对齐的目标是 ``别跨线、别非法''，不是 ``越对齐越快''。别名是向量化最常见的第一号阻碍。

\begin{lstlisting}[style=outstyle,caption={循环被向量化后的形态示意（概念图，非实测数据）},label={lst:pm06-pipe}]
 scalar : 每元素一次 load / 一次加 / 一次 store，循环控制也在账上
   it :   1    2    3    4    5    6    7    8
 vector:  一次宽 load 覆盖若干元素，循环控制被摊薄
   vec:   [ 1  2  3  4 ]      [ 5  6  7  8 ]
 前提：步长编译期已知 / 指针对不重叠 / 归约顺序允许重排
 收益边界：循环体本来就很短，或瓶颈在访存时，
           向量化只是把瓶颈换了一个位置，不消掉它
\end{lstlisting}

被 libc 认出来的 \Tp{memcpy{\Bk}\_erms}、向量化 \Fn{memchr} 一类函数，热点地址归属常常落在你的调用行
而不是实现里（第~\ref{ch:pm04-sampling} 章讲过 ``排第一不是结论''）。这里补一句本章特有的：这些实现本身已经向量化，
在它前面再写一层 ``手工向量化'' 通常只是把同一件事干两遍。

\subsection{代码布局与指令侧资源}

取指与微指令队列这类\textbf{指令侧}（I-side）资源同样会耗尽，而消耗多少不只取决于 \Tp{.text} 多大，
还取决于热基本块在虚拟地址空间里是否彼此靠近、是否踩到页边界和取指窗口。
一个只碰 $2$ MiB \Tp{.text} 的循环，如果每次调用都要跨很多页取指令，取指带宽与 \Mx{iTLB} 都可能成瓶颈。
三条观察手段都属于 ``看''，不属于 ``改''：
\begin{itemize}
\item \textbf{热区的页分布}：把采样拿到的热地址按页聚合，看热区是集中在少数几页还是散在几十页。
\item \textbf{函数体积与顺序}：\Fn{nm} 排序找异常大的函数；跨模块的大函数被拆散会让取指链路更分散。
\item \textbf{布局变换后的差分}：\Tp{-ffunction{\Bk}-sections} 加链接器的段顺序文件、PGO 二次编译、
  BOLT 一类二进制重排都能改变 I-side 表现。它们怎么配、收益方向由什么决定在第~\ref{ch:pm10-compiler} 章；
  本章只强调一条观察纪律：\textbf{重排只搬代码、不改语义，所以数据侧的 \Tp{cache{\Bk}-misses}
  基本不该动}；如果数据 miss 也跟着变了，说明变量没隔离住（同时改了内联或代码生成）。
\end{itemize}

最后把一句账留在这一章：同一段代码，单线程下频率高（睿频没被别的核分走），满载时低，
向量指令密集的代码还会撞功耗/电流上限（频率偏置的具体行为随平台与调速器变）。
于是同一份代码在 \Tp{taskset} 单线程与满载两种条件下的 ``每周期工作量'' 与 ``每秒周期数'' 都可能不同。
本章的对照一律固定在绑核单线程并记录当时频率；频率与 DVFS 的口径在第~\ref{ch:pm02-clock} 章，并发下的解释在第~\ref{ch:pm09-concurrency} 章。

\begin{bestpractice}[title={本章归因四步起手式},label={box:pm06-start}]
\begin{itemize}
\item 一次只动一个变量：改跨度、拆字段、换布局、去分支、重排——每次挑一个，其余全钉住。
\item 同机、绑核、交替重复：A/B/A/B 跑，用第~\ref{ch:pm03-stats} 章的区间判据；换一台机器等于换一个结论。
\item 先看指标，再改代码，最后看反汇编：反汇编没确认 ``代码真的变了'' 之前，收益可能来自别处。
\item 环境指纹进结论：CPU 型号、缓存层次、内核与编译器版本、当时的频率与调速器。
  没写进结论的指纹，下次复现不出来。
\end{itemize}
\end{bestpractice}

\section{排障与速查}
\label{sec:pm06-signature}

本章的收束资产是一张\textbf{签名表}：看到 X 异常，先查 Y，别急着查 Z。
第三列写的是 ``看着相关但通常白费'' 的动作——归因阶段最容易浪费预算的就是它。
表中 ``现象'' 一列描述的是指标形态，不是本仓库的实测值。

\begin{longtable}{|>{\raggedright\arraybackslash}p{0.22\textwidth}|>{\raggedright\arraybackslash}p{0.27\textwidth}|>{\raggedright\arraybackslash}p{0.27\textwidth}|}
\caption{缓存、分支与指令级归因的签名对照（指标形态示意，非实测值）}\label{tab:pm06-signature}\\
\hline
\textbf{看到 X 异常} & \textbf{先查 Y} & \textbf{别急着查 Z}\\
\hline
\endfirsthead
\hline
\multicolumn{3}{|l|}{表~\ref{tab:pm06-signature}~~缓存、分支与指令级归因的签名对照（续）}\\
\hline
\textbf{看到 X 异常} & \textbf{先查 Y} & \textbf{别急着查 Z}\\
\hline
\endhead
\hline
\endfoot
\hline
\endlastfoot
\Mx{cache{\Bk}-misses} 高、工作集远小于 L3 & 访问步长与跨度；把行跨度改奇数做差分（冲突） & 直接换 SoA 或上对象池\\
\hline
顺序扫快、列扫慢 & 存储顺序与遍历方向；转置或分块 & 分配器局部性参数\\
\hline
AoS 改 SoA 后 miss 没降 & 循环是否真只用少数字段；对象是否本来就跨线 & 断言 ``SoA 一定快''\\
\hline
\Mx{L1{\Bk}-miss} 涨但 \Mx{LLC{\Bk}-miss} 不动 & 线内字节利用率与预取队列；宽 load 是否跨线 & 归因到 DRAM 带宽\\
\hline
\Mx{LLC{\Bk}-miss} 涨且 \Mx{IPC} 掉 & 工作集是否越过你测出的容量边界；层次延迟自己测 & 抄网上的周期数表\\
\hline
\Mx{dTLB-miss} 高而缓存计数正常 & 页大小、访问是否跨页、工作集物理连续性 & 又去改缓存布局\\
\hline
\Mx{branch{\Bk}-miss} 率高 & 该分支方向是否由输入决定；换数据集看它跟不跟着变 & 一律改查表或 \Tp{cmov}\\
\hline
单个分支归属占比极高 & 反汇编确认那条分支在发布产物里还存在 & 相信报告与源码一致\\
\hline
\Mx{branch{\Bk}-miss} 计数看着荒谬 & 本机计数口径与 ``计数不可靠'' 提示行 & 以为预测器失灵\\
\hline
\Mx{branch{\Bk}-miss} 降但 \MLat{} 升 & 无分支写法多做的工作、查表引入的 load & 继续 ``再消一个分支''\\
\hline
\Mx{IPC} 低、周期高、miss 正常 & 除法、跨迭代依赖链、串行归约 & 加更多预取\\
\hline
\Mx{IPC} 高但仍然慢 & 端口宽度、访存上限、循环迭代数太少 & 认为 ``已经到头了''\\
\hline
报告说循环已向量化的 & 发布产物反汇编里确有宽寄存器指令 & 把报告本身当收益\\
\hline
报告说未向量化（依赖） & 别名与 \Tp{restrict}、下标关系能否被证明 & 换个优化级别碰运气\\
\hline
浮点归约不向量化 & 顺序是否必须保序、能否接受结果变化 & 为速度开激进重排\\
\hline
绑核后 \Tp{context{\Bk}-switches} 明显非 0 & 是否存在抢占或迁移，归属数据要打折 & 继续算 miss 率\\
\hline
单线程快、并发慢 & 频率上限与功耗偏置；跨线程共享（第~\ref{ch:pm09-concurrency} 章） & 归因给数据布局\\
\hline
同一二进制两台机差异大 & 缓存层次、容量与微架构差异（重测本章表格） & 怀疑代码\\
\hline
事件显示 not supported 计 0 & 本机是否真的提供该事件（现场探测） & 把 0 当 ``零开销''\\
\hline
热点落在 \Tp{memcpy} 族 & 调用点的字节数与对齐；实现本身已向量化 & 自己再写一层向量拷贝\\
\hline
加了软件预取反而慢 & 代码是否本就访存受限；队列压力与带宽 & 把预取当银弹\\
\hline
换分配器后 miss 变了 & 物理连续性与页对齐变化；同线程还是跨线程 & 归因给 ``分配器更快''\\
\hline
重排 \Tp{.text} 后数据 miss 也变 & 变量没隔离（顺带改了内联或代码生成） & 认为代码布局影响数据缓存\\
\hline
\end{longtable}

\subsection{三平台这一层的观察面}

完整的能力与权限对照放在附录 A，这里只留本章用到的三条：
\begin{itemize}
\item Linux：\Fn{perf} 一族（计数、逐指令归属、分支栈记录、带地址取样）是主线，
  事件名与权限都现场确认，不抄表。
\item Windows：\Tp{XPerf} 与 VTune 一类给出 \Mx{retiring} / \Mx{bad{\Bk}-speculation} 这类分层与
  微架构指标，命名与 Linux 事件不同源；对照时先对齐口径，再对数字。
\item macOS：\Tp{Instruments} 的 CPU 时间模板与 \Tp{powermetrics}；权限模型不同，
  精确取样与地址级取样的可用性要现场确认。
\end{itemize}
三个平台都能做本章的布局差分——``只有 Linux 能测'' 是不成立的结论，
差别只在\textbf{拿不拿得到事件}，而不是\textbf{能不能改布局}。

\subsection{本章自查五条}

\begin{enumerate}
\item 我说的是 miss \textbf{率}还是 miss \textbf{数}？换算成周期了吗？层次延迟是我自己测的吗？
\item 我测的是不是只有单线程、绑核、记录了频率的那一次？跨线程因素排除了吗？
\item 这个 ``优化'' 有没有对应一次只动一个变量的差分？还是同时改了两件事（第~\ref{ch:pm15-antipattern} 章）？
\item 代码形态的判断（向量化、展开、内联）是在发布产物上反汇编确认的，还是从源码意图推的？
\item 报告与指标给我的是 ``没发生'' 还是 ``变慢了''？我有没有把编译器的说法当成收益？
\end{enumerate}

\chapter{内存与分配器：局部性、碎片与对象池的测量}
\label{ch:pm07-memory}

前两章把镜头对准指令与微架构：缓存缺失事件、分支误预测、\Mx{TopDown} 的分层。那些指标能把你送到
出问题的循环附近，但到了内存这一侧，它们只回答半个问题。真正落到工程上的提问通常是另外几种形态：
把 \CPP{} 的 \Fn{new} 换成 arena 之后到底快了多少；AoS 改成 SoA 之后为什么有的循环快了有的反而慢；
服务跑六小时后 \Mx{RSS} 涨了一大截，是泄漏、碎片，还是分配器只是没归还；换成 \Lib{ptcpmalloc} 之后
那根曲线算不算改善。难处不在于没有工具，而在于\textbf{同一次改动会同时改变四本账}：分配调用的路径
长度、访问的局部性、驻留的物理页数量、归还给内核的时机。指标只告诉你四本账的总和变了，把总和当因果，
就是本章要防的事。

本章的口径约定延续前面几章：\MWall{} 指挂钟时间，从起点到终点真实流逝的时间；\MCpu{} 指进程或线程
在 CPU 上实际被调度的时间；\MLat{} 指单次操作的耗时分布，本章关心它的尾部，因为分配路径可能突然走
系统调用；\MThr{} 指单位时间完成的工作量。四个口径在内存侧经常打架：arena 让 \MCpu{} 下降却让
\Mx{RSS} 上升，大页让 \MThr{} 上升却让小分配的 \Mx{p99} 变差，一次扩容能让平均值纹丝不动而让
\Mx{p999} 翻倍。先说清你要辩护的是哪一个，再谈数字。

接口在哪里：缓存缺失事件的选择与指令级归因（计数、\Mx{TopDown} 的 \Mx{Memory Bound} 分支）在
第~\ref{ch:pm06-micro} 章；PMU 事件名随内核与 CPU 型号变化、三个工具之间怎么对照计数在
第~\ref{ch:pm05-events} 章；线程共享缓存、伪共享与利用率歧义在 第~\ref{ch:pm09-concurrency} 章；计时源本身在
第~\ref{ch:pm02-clock} 章，重复次数与最小可检出差异在 第~\ref{ch:pm03-stats} 章；哪些指标值得进回归门禁在
第~\ref{ch:pm13-ci} 章；构建口径（优化级别、LTO、断言）在 第~\ref{ch:pm10-compiler} 章。ASan、TSan、UBSan 与
Valgrind 怎么用不在本章重讲，那是《内存泄漏检测与防御》册的内容；本章只把它们当成\textbf{测量对象}：
加了多少成本、加完之后哪个数还能看、哪个数不能看。

\section{分配开销的账目怎么分}
\label{sec:pm07-budget}

一次 \Fn{malloc} 或 \Fn{new} 的成本不是一个数，而是几段条件的和：在缓存里找到一块可用区间的成本、
改元数据（bin 链表、块头、region 位图）的成本、是否需要跨线程同步、是否需要向内核要页、以及新页
第一次被写入时的缺页成本。最后一段最容易被漏掉，因为它发生在分配调用\textbf{之后}的第一次访问上，
不在调用内部。这也是很多 ``换分配器无效'' 的实验其实有效的原因：省掉的那一段确实省掉了，只是它本来
只占总耗时的一小部分，落进了噪声。

按层拆开，才能得到可辩护的结论：你优化掉的到底是哪一层。表 \ref{tab:pm07-budget-layers} 是本章后续
所有实验的账目底本，其中第三列是\textbf{能被观测}的量、第四列是\textbf{能被改动}的旋钮，两列不能混：
\Mx{RSS} 是可观测量而非可干预变量；\Tp{M\_MMAP\_THRESHOLD} 是可干预变量，但它的影响必须通过可观测
量确认。布局层不参与分配调用，常被 ``分配优化'' 的讨论漏掉，可它决定分配之后每一次访问的成本，本册
把它算进同一本账。

\begin{bestpractice}[title={开工前的一次性判别：这次改动动的是哪一层}]
  \begin{enumerate}
    \item 先只测\textbf{数量}，不测时间：一次迭代里发生多少次分配、请求多少字节、结束时活着的字节
          是多少。分配次数除以迭代次数若小于 $10^1$，分配路径本身多半不是瓶颈，叙事应该转向布局或
          其他热点；量级到 $10^3$ 以上，才值得继续谈池与 arena。
    \item 再测\textbf{尺寸分布}：请求字节数除以次数只是均值，真正有用的是直方图。均值几十字节的负载，
          真身常常是绝大多数小请求加极少数大块，两者的干预手段完全不同，而均值把它们混成一个数。
    \item 数量与分布都确认值得做之后，才进入时间测量。顺序反了，你会在一个不存在的路径上做微优化，
          并把噪声读成收益。
  \end{enumerate}
\end{bestpractice}

{\footnotesize
\begin{longtable}{>{\raggedright\arraybackslash}p{0.15\textwidth}
                  >{\raggedright\arraybackslash}p{0.21\textwidth}
                  >{\raggedright\arraybackslash}p{0.20\textwidth}
                  >{\raggedright\arraybackslash}p{0.19\textwidth}}
\caption{内存侧开销的分层账目：每层的成本来源、你能量到什么、你能改什么（分类说明，不含实测值）}\label{tab:pm07-budget-layers}\\
\hline
层 & 成本从哪来 & 可观测量 & 可干预变量 \\
\hline
\endfirsthead
\multicolumn{4}{l}{（接下页，表 \ref{tab:pm07-budget-layers}）}\\
\hline
层 & 成本从哪来 & 可观测量 & 可干预变量 \\
\hline
\endhead
应用层（池、arena、缓存分配器） &
批量取块代替 $N$ 次细碎请求，路径退化为指针自增；顺带把生命周期相同的对象摆在一起 &
区域计数（次数、字节数）、bump 步长与已用量、每轮归还的块数 &
池上限与块大小、arena 容量、是否分代、复位时机 \\
\hline
\CPP{} 语言层（\Fn{new}、\Fn{delete}、容器） &
\Fn{bad\_alloc} 异常路径、\Fn{new} 与构造之间的边界、\Fn{vector} 与 \Fn{string} 的增长与搬迁 &
容器搬迁次数与增长因子、临时对象数量、被替换的 \Fn{new} 调用次数 &
\Fn{reserve}、\Fn{shrink{\Bk}\_to{\Bk}\_fit}、增长因子、小对象内联缓冲、自定义分配器 \\
\hline
C 分配层（\Fn{malloc}、\Fn{free}） &
bin 查找与链表操作、块头与元数据、线程 arena 的锁或无锁队列、\Fn{free} 后的合并与延迟归还 &
分配器统计（bin 直方图、region 数）、\Fn{malloc{\Bk}\_usable{\Bk}\_size}、
\Fn{malloc{\Bk}\_info} 风格输出、落在分配调用上的 \MCpu{} &
换实现、arena 数量上限、线程局部缓存开关、统计接口是否开启 \\
\hline
内核与页层（\Tp{brk}、\Tp{mmap}、缺页） &
向内核要地址区间、新区间第一次写入触发缺页、页表建立、必要时的零页填充 &
minor 与 major fault 计数、\Mx{RSS} 与 \Mx{VSZ} 的差、系统调用次数 &
预触碰（\Fn{memset} 或 \Fn{madvise}）、大页与 THP、\Tp{M\_MMAP{\Bk}\_THRESHOLD}、
\Tp{MALLOC\_TRIM{\Bk}\_THRESHOLD} \\
\hline
布局层（不参与分配调用，却决定之后的一切） &
每次访问搬运的无关字节、指针追逐造成的串行访存、对齐留下的空洞、跨缓存行的字段 &
\Mx{L1}/\Mx{L2}/\Mx{L3} 缺失、\Mx{IPC}、单位迭代搬运的估计字节数、\Mx{Memory Bound} 占比 &
AoS 与 SoA、字段顺序与打包、分块遍历、索引替代指针、缓存行对齐 \\
\hline
\end{longtable}
}

\subsection{把账目落到代码上：一个分配计数器}
\label{sub:pm07-ledger}

清单 \ref{lst:pm07-ledger} 给出最小可用的计数骨架，它有三个刻意的取舍，每一个都是测量学的而不是
风格上的。第一，计数路径本身绝不能再分配，所以内部直接走 \Fn{std::malloc} 与 \Fn{std::free}；
在被替换的 \Fn{new} 里再调用 \Fn{operator new} 会立刻递归爆栈。第二，开关是运行时布尔而不是编译期
常量，方便同一份二进制跑对照，避免把重新编译带来的代码生成差异混进结论。第三，\textbf{计数与计时
必须分成两次运行}：两个 \Fn{fetch\_add} 在单线程下只是几十纳秒量级的尾巴，在高并发下却可能比被计数
的分配更贵，贵多少取决于核数、缓存拓扑与争用程度；这个数要读者在目标机上自己标定，方法就是把开关
设为假、其余完全相同地再跑一次，报出两者 \MCpu{} 的关系。

\begin{lstlisting}[style=cppstyle, caption={分配账目计数器：替换全局 \CPP{} 分配函数并统计次数与字节数}, label={lst:pm07-ledger}]
// alloc_ledger.hpp -- count allocations inside a region; never allocate while counting.
#include <atomic>
#include <cstddef>
#include <cstdint>
#include <cstdlib>
#include <new>

struct AllocLedger {
    std::atomic<std::uint64_t> new_calls{0};
    std::atomic<std::uint64_t> new_bytes{0};
    std::atomic<std::uint64_t> delete_calls{0};
    std::atomic<std::int64_t>  live_bytes{0};
    void reset() { new_calls = 0; new_bytes = 0; delete_calls = 0; live_bytes = 0; }
};

inline AllocLedger g_ledger;
inline bool g_ledger_on = false;   // runtime switch: one binary, two runs
// Toggle around the measured region and restore the previous value, so nested
// regions and the warm-up phase never leak into the counters.

void* operator new(std::size_t n) {
    void* p = std::malloc(n == 0 ? 1 : n);       // do NOT call ::operator new here
    if (p != nullptr) {
        if (g_ledger_on) {
            g_ledger.new_calls.fetch_add(1, std::memory_order_relaxed);
            g_ledger.new_bytes.fetch_add(n, std::memory_order_relaxed);
            g_ledger.live_bytes.fetch_add(static_cast<std::int64_t>(n),
                                          std::memory_order_relaxed);
        }
        return p;
    }
    throw std::bad_alloc();
}

void operator delete(void* p) noexcept {
    if (p == nullptr) return;
    if (g_ledger_on) g_ledger.delete_calls.fetch_add(1, std::memory_order_relaxed);
    std::free(p);
}

void operator delete(void* p, std::size_t n) noexcept {  // sized delete, C++14 and later
    if (p == nullptr) return;
    if (g_ledger_on) {
        g_ledger.delete_calls.fetch_add(1, std::memory_order_relaxed);
        g_ledger.live_bytes.fetch_sub(static_cast<std::int64_t>(n),
                                      std::memory_order_relaxed);
    }
    std::free(p);
}
\end{lstlisting}

有了它，第一次跑只开计数不开计时，得到次数、请求字节数与结束时的 \Tp{live\_bytes}；第二次跑关掉计数，
只测 \MWall{} 与 \MCpu{}。两次之间不是``扣掉开销''的关系：带计数的那次连分配路径本身都被改过，所以
把不带计数的运行当唯一时间基线，带计数的运行只贡献数量；要报告观察者效应本身，再加一次全程关闭开关
的运行作为第三次对照。

这个骨架有意漏掉三件事，报告时要写成已知盲区：\CPP{} 的数组形式分配（\Tp{new[]}）、
\Fn{aligned\_new} 与 placement \Fn{new} 不经过这两个标量重载；栈对象与线程栈的增长完全不经过它；
使用自定义分配器的容器也不经过它。换句话说，它统计的是\textbf{全局 \Fn{new} 路径}，不是\textbf{程序
的内存足迹}。把这两个名字混用，是本章见到的第一类错误结论的来源。

\section{局部性改造的测量闭环}
\label{sec:pm07-locality}

局部性有四个互相独立的来源，改动时几乎总被同时碰到，值得先分开命名。空间局部性指一次访问之后很快
访问相邻地址；时间局部性指同一段数据在缓存把它逐出之前被再次用到；访问粒度指一次读取搬进缓存行的
字节里你真正用了多少；访存并行度指硬件能同时在飞多少笔加载。遍历数组时依赖链很短，而指针追逐（链表、
跳表、\Fn{unordered{\Bk}\_map} 的桶）把每个加载变成下一个加载的前置条件，于是内存级并行度退到接近 1。
前三者由布局决定，第四者由\textbf{依赖图}决定：把第四种当成布局问题去重排结构体通常毫无收益，
第~\ref{ch:pm06-micro} 章 里有它的指令级对应证据。

\begin{guideline}[title={AoS 改 SoA 之前必须存在的三个证据}]
  \begin{enumerate}
    \item 循环里\textbf{只用到结构体的部分字段}。若所有字段都被访问，AoS 与 SoA 搬运的总字节数
          相同，改造只剩向量化与对齐上的附带效果。
    \item \Mx{Memory Bound} 或缓存缺失计数确实随工作集增长而恶化，而不是在最小工作集上就恶化
          （事件与三工具对照见 第~\ref{ch:pm05-events} 章，指令级定位见 第~\ref{ch:pm06-micro} 章）。
    \item 遍历是\textbf{顺序}的。随机下标访问时 SoA 只改变你追逐的是三个数组还是一个结构体，
          串行依赖并没有被消掉。
  \end{enumerate}
  三条缺一条，就不要把 \MWall{} 的改善归因给 SoA。
\end{guideline}

\subsection{工作集扫描：不知道缓存大小也能确认主因}
\label{sub:pm07-sweep}

最强的局部性证据不是某个单点数字，而是\textbf{整条曲线}。固定算法与每元素工作量，只把元素个数 $N$
从几千扫到几千万，让工作集依次跨过各级缓存容量与 DRAM，画出每元素耗时随 $N$ 的变化。局部性改造如果
真的解决了缓存问题，它移动的是\textbf{拐点对应的 $N$}，也就是把整条曲线向右推；如果只改了指令层面，
动的是拐点\textbf{以下}那段平台的\textbf{高度}，拐点位置不动。这两件事在单点测试里看起来一模一样，
是本章最常见的归因错误来源。

\begin{lstlisting}[style=outstyle, caption={工作集扫描输出的形态示意（格式示例，非本仓库实测值；形状只用于说明读图方法，数值需读者在自己机器上重做）}, label={lst:pm07-sweep}]
$ ./locality-sweep --sizes 1e3 1e4 1e5 1e6 1e7 --layout aos
# the ns/elem column below is a RELATIVE shape, not a measured value.
   N          working set        ns/elem (AoS)   profile
   1e3        ~ few KB           ....            near-cache plateau
   1e4        ~ tens of KB       ....            near-cache plateau
   1e5        ~ few hundred KB   ..--..          knee 1 (L1 -> L2 scale)
   1e6        ~ few MB           -----           mid plateau
   1e7        ~ few tens of MB   ______/\____    knee 2 (cache -> DRAM, steep)

$ ./locality-sweep --sizes ... --layout soa
   N          working set        ns/elem (SoA)   reading
   1e3        ~ few KB           ....            same height -> not a near-cache win
   1e5        ~ few hundred KB   ...--...        knee moved RIGHT -> locality win
   1e7        ~ few tens of MB   _____/\___      slope above knee ~ unchanged

# How to read the pair:
#   plateau height moves     -> instruction-level / vectorization effect
#   knee position moves      -> locality (bytes touched per useful value)
#   knee disappears entirely -> prefetching or another limiter took over
# Sweep both variants in the same order, then repeat the sweep once backwards so
# thermal and frequency drift show up as hysteresis instead of as a win.
\end{lstlisting}

这个实验有三个污染项。第一，拐点位置本身是平台属性：页大小与 THP 设置改变 TLB 压力，会让拐点和缓存
边界一起移动，于是你把虚拟内存效应读成缓存效应；确认方法是读
\Tp{/sys{\Bk}/kernel{\Bk}/mm{\Bk}/transparent{\Bk}\_hugepage{\Bk}/enabled} 与
\Tp{/proc/{\Bk}meminfo} 里的 \Mx{AnonHugePages} 字段，并保证扫描前后一致。第二，$N$ 变大时计时调用
本身占比下降，每元素工作量太小时曲线左端会出现纯测量伪影：把每元素工作量做成可参数化的内层重复次数，
检查左端是否随重复次数线性缩放，非线性说明你还在测启动成本。第三是频率：末端的大工作集把内存控制器
压满、功耗上升，调频可能在跑完之前就把频率降下来（见 第~\ref{ch:pm02-clock} 章），于是你把降频读成了内存墙；
同时记录每元素周期数而不只看秒数能过滤掉一部分漂移，更彻底的做法是在扫描中间插入一个固定参照负载，
用参照负载自己的漂移归一化整条曲线。

\begin{warning}[title={局部性实验最常见的三个假信号}]
  \begin{enumerate}
    \item \textbf{起点没对齐}。裸数组的起点由分配器决定，两版布局可能落在不同的缓存行偏移上，于是
          你测到的是那一次对齐抖动。办法：给起点加随机偏移重复多次，先量出抖动幅度，再判断收益是否
          落在抖动之内。
    \item \textbf{循环被编译器吃掉了}。\Tp{-O2} 与 \Tp{-O3} 下，只做累加的遍历可能被完全向量化甚至
          折叠；换布局常常同时换了可口向量化的程度，两版根本不在同一层比较。办法：在发布产物上
          反汇编确认循环还在、确认访存指令形态如你所料（构建口径见 第~\ref{ch:pm10-compiler} 章）。
    \item \textbf{把 \Mx{IPC} 上升当收益}。\Mx{IPC} 与缓存缺失同时受频率、SMT 邻居与预取影响；没有
          \MWall{} 或 \MThr{} 上的对应改善，它只是解释，不是结论。
  \end{enumerate}
\end{warning}

\section{对象池与 arena 的收益怎么测}
\label{sec:pm07-pool}

池与 arena 的收益有两个来源，而且几乎总在一次改造里同时发生。一个是\textbf{分配路径变短}：一次
bump 指针自增替代一次带 bin 查找与元数据更新的 \Fn{malloc}。另一个是\textbf{布局变好}：arena 顺序
切块，生命周期相近的对象天然相邻，指针追逐可以被索引替代。如果报告只写 ``换成 arena 之后快了若干''，
读者无法知道省下的是哪一笔，而这两笔的维护含义完全不同：前者是可预期的常数，后者一旦数据布局改动
就会消失。

\subsection{三段实验：把分配与布局拆开}
\label{sub:pm07-threeway}

可辩护的做法是做三组改造而不是一组。A 组：只把 \Fn{new} 换成 arena，对象内部布局与访问顺序保持
不变（可以在 arena 上原样构造同一个 AoS 结构）。B 组：只改布局（AoS 改 SoA，或把指针换成索引），
分配仍走 \Fn{new}。C 组：两者都改。D 组是基线，什么都不改。于是有三种可能的读数：

\begin{itemize}
  \item 若 A 有收益且 C 的总收益约等于 A 加 B，两笔账独立，可以分别报告各自贡献；
  \item 若 A 几乎无收益、B 拿走全部收益、C 与 B 相同，那么收益全部来自局部性，此时
        \textbf{arena 只是把布局改好所附带的搬运成本}，池本身可以不做，省下归还策略的复杂度；
  \item 若 A 有收益但 C 明显小于 A 加 B，说明两者在争同一个瓶颈（多半是带宽）。报告里必须写明收益
        不可叠加，否则下一个人会把两个改动各算一遍，得出凭空多出来的总收益。
\end{itemize}

\begin{guideline}[title={报告 arena 收益时必须同时贴出的四个数}]
  \begin{enumerate}
    \item \textbf{分配次数的变化}：单位迭代内 \Fn{new} 与 \Fn{malloc} 调用次数前后对比，用清单
          \ref{lst:pm07-ledger} 的计数器数出来，而时间来自不开计数的那一次。
    \item \textbf{单次均摊成本}：不要拿两次运行的 \MWall{} 差直接除以次数，总耗时里除了分配还有其他
          一切。要么在分配调用处直接计时（会引入观察者效应，必须同时报出该开销的量级），要么按
          第~\ref{ch:pm03-stats} 章 把多组配对样本做成带置信区间的差值，并检查它是否超过最小可检出差异。
    \item \textbf{\Mx{RSS} 的两条线}：峰值与稳态。池的典型副作用是稳态 \Mx{RSS} 上升而峰值下降，
          把两个方向合写成一句 ``内存更省了'' 就是误导。
    \item \textbf{复位粒度}：arena 什么时候真的把内存还给内核，以及你的负载是否允许在那一刻还。
          如果答案是进程退出时，长驻服务里的收益要按驻留上涨重新估值。
  \end{enumerate}
\end{guideline}

尾延迟是池最容易被忽略的一栏。通用分配器偶发的合并、分裂与向内核要页，会让 \Mx{p99} 或 \Mx{p999}
出现少数大值；arena 通常能把这些大值搬走，代价是把它们\textbf{集中}到扩容那一刻：一次扩容可能触发
一次 \Tp{mmap} 加一批缺页，形成一个更高但更稀的尖峰。对延迟敏感的系统，正确的问法不是池让 \Mx{p99}
变好还是变差，而是尖峰被摊平了还是被前移了；要回答它，需要记录每次分配调用的耗时直方图，并按是否
首次触碰该页分开统计（口径见 第~\ref{ch:pm06-micro} 章，分位数的统计判据见 第~\ref{ch:pm03-stats} 章）。

多线程程序里池还有两个新问题：空闲链表的头指针是跨线程共享的可变状态，容易与相邻数据落在同一缓存行
上造成伪共享；无锁链表的 \Fn{CAS} 重试在高争用下会让 \MCpu{} 上升而 \MThr{} 下降，也就是利用率歧义。
两者都属于 第~\ref{ch:pm09-concurrency} 章 的范围，本章只提醒一句：池的并行版本必须与串行版本一起测。还有
一个本章特有的陷阱：bump 指针本身也是共享可变状态，而每线程 arena 会把 arena 数量放大成线程数，驻留
随之上涨，与下一节的归还延迟叠加后极难归因，所以每线程 arena 的上限要当成一个显式的被测参数扫一遍。

\section{碎片到底测什么}
\label{sec:pm07-frag}

碎片不是一个量，工程讨论里至少有四个不同的东西共用这个词，它们的测量手段互不替代。

\textbf{外部碎片}：空闲总量足够，但没有一段连续区间能满足某次请求。它只能由空闲区间尺寸的\textbf{分布}
说明，不能由总量说明。分配器自带的统计接口（\Fn{malloc{\Bk}\_info} 风格的输出，或 \Lib{ptcpmalloc}
这类实现的 \Fn{malloc{\Bk}\_stats{\Bk}\_print}）是主要来源；接口怎么开、字段怎么读属于
《内存泄漏检测与防御》册，本章只关心它们作为数据的口径。

\textbf{内部碎片}：你请求 $R$ 字节，分配器给你一块可用容量 $U$（不小于 $R$），多出的 $U-R$ 你永远
用不到，而 \Fn{malloc{\Bk}\_usable{\Bk}\_size} 给出的正是 $U$。于是内部碎片率就是 $U$ 与 $R$ 之差对
$R$ 的比例，这句话可以直接写进报告，不必上显示公式。它对\textbf{小尺寸请求}尤其敏感，因为多数分配器
把请求向上取整到某个粒度。想测它，就采样真实请求大小的序列，逐个查 $U$，报直方图而不报均值。

\textbf{驻留}：\Mx{RSS} 计的是当前映射在物理内存里的字节，同时包含活动数据、分配器留下的空闲但已
触碰过的页、代码与映射文件、以及共享库。用 \Mx{RSS} 减去分配器统计里的空闲总量，得到的差\textbf{不是}
碎片，只是那几个来源的和。

\textbf{归还延迟}：\Fn{free} 之后分配器不立刻收缩或解除映射，或者内核侧的回收有延迟。它让内存曲线
看起来只涨不跌，是泄漏与碎片之间最容易被误判的一档。

\subsection{采集清单}
\label{sub:pm07-collect}

\begin{lstlisting}[style=shellstyle, caption={内存侧时序采集清单：在同一时间轴上取驻留、缺页与分配器状态}, label={lst:pm07-collect}]
#!/bin/sh
# mem-trace.sh -- capture one time axis from several sources.
# Nothing here proves fragmentation by itself; it only gives you aligned series.
PID=$1
OUT=${2:-mem-trace}
mkdir -p "$OUT"

# (0) identity first: allocator and compiler version strings belong to the result.
ldd ./myapp > "$OUT/libs.txt" 2>&1          # which libc, which allocator
g++ --version | head -n 1 > "$OUT/compiler.txt"
uname -a > "$OUT/kernel.txt"

# (1) residency and faults, sampled on a FIXED period.
#     A variable sampling period destroys every later difference.
#     /proc field numbers differ between kernels: verify them on YOUR box first.
while kill -0 "$PID" 2>/dev/null; do
    ts=$(date +%s.%N)
    # field 10 = minflt, 12 = majflt, 23 = vsize bytes, 24 = rss in pages
    awk -v ts="$ts" '{printf "%s minflt=%s majflt=%s vsize=%s rss_pages=%s\n", \
                      ts, $10, $12, $23, $24}' "/proc/$PID/stat" >> "$OUT/stat.log"
    grep -E '^(Rss|Anonymous|Swap|THP)' "/proc/$PID/smaps_rollup" \
        >> "$OUT/smaps.log" 2>/dev/null
    sleep 1
done

# (2) page policy, captured before and after: a moving knee may be a TLB effect.
cat /sys/kernel/mm/transparent_hugepage/enabled > "$OUT/thp.txt"

# (3) allocator tunables actually in effect (names differ per allocator):
#     glibc family : MALLOC_ARENA_MAX, M_MMAP_THRESHOLD, M_TRIM_THRESHOLD
#     other family : a stats dump entry point instead of env knobs
env | grep -E '^(MALLOC_|M_)' > "$OUT/tunables.txt"

# (4) allocation counts: a SEPARATE run with the ledger enabled, never a timed run.
echo "see lst:pm07-ledger: g_ledger_on only in the counting run" >> "$OUT/notes.txt"
\end{lstlisting}

清单 \ref{lst:pm07-collect} 里按编号取 \Tp{/proc/{\Bk}stat} 字段是\textbf{脆弱}做法：不同内核版本的
字段顺序与单位会变，更稳的是按名字解析 \Tp{smaps{\Bk}\_rollup} 的键。这句话应当写在脚本注释里，比写在
结论里更容易被后来的人看到。

\begin{lstlisting}[style=outstyle, caption={分配器统计与时序对齐的输出口态示意（格式示例，非本仓库实测值；下面每一个数字都是占位符，不可引用）}, label={lst:pm07-fragfmt}]
# ---- allocator dump (shape only; field names depend on the allocator) ----
allocated :  1 234 567   # bytes requested by the program
active    :    987 654   # bytes in chunks handed out and not yet released
mapped    :  2 345 678   # bytes obtained from the OS so far
retained  :    456 789   # bytes kept in allocator pools, NOT returned to the OS
bins (small):  extent   count
   0x20     :   18 432
   0x30     :    9 001
   0x40     :      512
   ...
largest request ever issued  : 0x2000
largest free extent available: 0x18    -> cannot satisfy 0x2000
#  free TOTAL is plenty but the largest EXTENT is too small  => external fragmentation
#  usable size >> requested size across many small chunks    => internal fragmentation

# ---- time-aligned series on one axis (what mem-trace.sh gives you) ----
t=0000  rss=  80  live_from_app=  72  minflt_delta=      0
t=1500  rss= 240  live_from_app=  95  minflt_delta=  41 000   (ramp-up)
t=3000  rss= 236  live_from_app=  96  minflt_delta=   1 200
t=4500  rss= 235  live_from_app=  34  minflt_delta=       0   (released to allocator)
#  live dropped a lot while rss stayed flat  => retention, not a leak,
#  and STILL not proof of fragmentation: check the largest free extent above.
\end{lstlisting}

\begin{warning}[title={不能从一张 \Mx{RSS} 曲线得出碎片结论}]
  \textbf{第一}，要同一时间轴上的两列：分配器认为\textbf{活着}的字节，与 \Mx{RSS}。只有 \Mx{RSS} 涨
  而活着那列不涨，是空闲池增长（归还策略或外部碎片的候选）；两列一起涨，优先怀疑生命周期或泄漏，
  那是《内存泄漏检测与防御》册的问题。\textbf{第二}，单点采样看不出趋势，看不出趋势就不能说碎片。
  \textbf{第三}，外部碎片的判定需要一个具体请求做锚：只有在空闲总量够、而某个已知尺寸的最大请求拿不到
  连续区间时它才成为问题；否则 \Mx{RSS} 里的空闲只是未归还，对性能没有影响，只对内存预算有影响。
\end{warning}

换分配器（例如从 \Lib{glibc} 的 \Fn{malloc} 换成 \Lib{ptcpmalloc} 或 \Lib{jemalloc} 风格的实现）在本章
口径下不是一个变量而是三个同时动的变量：块的划分策略、跨线程同步模型、归还内核的时机，所以只报
``换了分配器快了多少'' 是不完整甚至误导的。最小实验集应当是四组：默认配置；限制 arena 数量（把
\Tp{MALLOC\_ARENA{\Bk}MAX} 设为 1 与设为核数，两个都要）；把大请求阈值扫过你的典型大块尺寸；开启与
关闭分配器统计接口各一次（它本身有成本，开着它测时间等于换了一个被测系统）。所有配对都要在
第~\ref{ch:pm03-stats} 章 的判据下比较，并写明哪些组的差异落在最小可检出差异以下：``差异不显著'' 也是结论，
它否证的是你的归因假设。

\section{sanitizer 与跟踪工具是测量对象}
\label{sec:pm07-sanitizer}

本章对 ASan 与 Valgrind 只有一句话的态度：用法、选项与报告解读属于《内存泄漏检测与防御》册，CI 里把
它们当质量门禁的做法属于《LSP 与代码质量工具》第 12 章；\textbf{它们跑出来的时间与内存数不能当性能
数}。之所以单开一节，是因为内存与分配器恰好是观察者效应最严重的一处：测量仪器改变的不只是耗时，而是
\textbf{被测路径的身份}。

以 AddressSanitizer 为例，它做了三件与本节直接相关的事。它\textbf{替换}了 \Fn{malloc} 与 \Fn{new}
路径，于是你正要优化的那一层被换成另一个实现；它给每个分配加红区并维护一份与内存成比例的影子状态，
内存占用按倍数放大；它把 \Fn{free} 之后的内存放进隔离与延迟归还机制，目的是抓 use-after-free，代价是
\Mx{RSS} 的时间形状与生产完全不同。于是同一次运行里，分配路径、驻留、归还三本账全被改过。开销具体
多少取决于负载尺寸、红区比例、线程数与选项，公开报告常见的口径是数倍到十倍以上量级，而这个数
\textbf{必须}在你的机器与你的负载上标定，标定形态见清单 \ref{lst:pm07-calib}。Valgrind 的 memcheck
更彻底：它是动态二进制翻译，指令不在原生时序上执行，所以 \Mx{IPC}、缓存缺失计数、\Mx{TopDown} 一类
指标在它下面\textbf{没有可解释的含义}（工具与事件对照见 第~\ref{ch:pm05-events} 章），开销通常还要再高一到
两个数量级。

\begin{lstlisting}[style=outstyle, caption={倍率标定与趋势门禁的形态示意（格式示例，非本仓库实测值；所有数字都是占位符，只说明记录格式与读图规则）}, label={lst:pm07-calib}]
$ ./bench-native --iters 200 ...    # build A: -O2, no instrumentation  <- timing baseline
$ ./bench-asan   --iters 200 ...    # build B: same source, -fsanitize=address
$ valgrind --tool=memcheck ./bench-native --iters 200 ...

# calibration.tsv -- record RATIOS, each tied to one machine + toolchain fingerprint
run_id   wall_ratio_asan   rss_ratio_asan   wall_ratio_vg    notes
r-001          4.3               2.9           31.2      cores=8, governor=powersave
r-002          4.1               2.9           30.8      same binaries, repeat
#  Rule 1: a sanitizer ratio is NOT a performance number.
#  Rule 2: the ratio belongs to THIS machine + compiler + flags; do not carry it to CI.
#  Rule 3: never diff a sanitizer run against a native baseline; compare findings, not timings.
#  Rule 4: counts differ too (the tool allocates its own objects and delays frees),
#          so reconcile allocation counts against build A, never against build B.

# ---- what is allowed into the regression whitelist (see ch:pm13-ci) ----
# native build only, two series on one time axis:
#   alloc_count_per_iter                          bytes_live_peak
#      |          ..\  __/\__                        |      _____________
#      |         /  \/       ___                     |     /
#      |   _____/                                    | ___/        <- flat
#      +------------ time                            +----------- time
#   repeating step every epoch   -> reset / repool is not happening
#   saw-tooth with rising floor  -> retention OR leak, still undecidable here
\end{lstlisting}

那么 sanitizer 跑出来的什么\textbf{可以}用？只有语义类信息：某条路径是否越界、是否报告泄漏、某个对象
是否被释放两次，以及这些发现\textbf{是否在原生构建里也存在}（很多越界在原生构建里是静默的）。计数类
指标要谨慎到口径级别：隔离机制与泄漏检测开启时，ASan 统计到的分配与释放次数可能不等于原生构建的次数，
因为工具自身的内部对象、被延迟的释放、异常路径上的构造都会进入统计；把它和清单
\ref{lst:pm07-ledger} 在原生构建里数出的次数直接相减，会得出一条假回归。

值得进 第~\ref{ch:pm13-ci} 章 门禁的内存指标，本章的答案是原生构建下的单位迭代分配次数、活字节峰值、稳态
\Mx{RSS} 三个：它们都是计数或字节量，不需要 PMU 权限就能复现，适合做趋势。而 \Mx{L3 miss}、\Mx{IPC}
一类需要 PMU 的事件\textbf{不适合}当门禁指标（发布用的 CI runner 通常既没有 PMU 权限也没有
\Fn{perf}），它们只属于开发机上的归因素据，写进报告但不设阈值。分配器统计接口如果自带成本，白名单里
必须注明它是开着还是关着，否则一次接口开关的变更会被读成性能回归。

\begin{warning}[title={三个不能下结论的时刻}]
  \begin{enumerate}
    \item \textbf{在 sanitizer 构建上测收益}。换 arena 后 ``内存开销变小'' 这类结论在 ASan 下毫无
          意义，因为被改的那一层已经不是你改的那一层。
    \item \textbf{在带调试断言或 \Tp{-O0} 的构建上测局部性}。未优化构建里断言与调试信息改变了结构体
          排布与循环的可向量化程度，SoA 与 AoS 的对照不成立（构建口径见 第~\ref{ch:pm10-compiler} 章）。
    \item \textbf{拿单次运行说事}。内存类指标对页缓存状态、THP 设置与邻居负载尤其敏感，重复次数、
          分布与最小可检出差异一律按 第~\ref{ch:pm03-stats} 章 给。
  \end{enumerate}
\end{warning}

\section{排障与速查}
\label{sec:pm07-troubleshoot}

\begin{compare}[title={症状、第一手指标与结论边界}]
\begin{tabular}{>{\raggedright\arraybackslash}p{0.20\textwidth}
                >{\raggedright\arraybackslash}p{0.22\textwidth}
                >{\raggedright\arraybackslash}p{0.33\textwidth}}
\hline
症状 & 先看什么 & 能支持到哪一步 \\
\hline
换 arena 后 \MWall{} 变快 & 分配次数、单次均摊、A/B/C 三段实验 &
到 ``分配路径变短''；只有 B 组也有效时才能同时声称布局收益 \\
\hline
SoA 改造后有的循环变慢 & 工作集扫描的拐点、每元素搬运字节估计 &
到 ``该循环本来就用满字段''；指令级证据回 第~\ref{ch:pm06-micro} 章 \\
\hline
\Mx{RSS} 长跑只涨不跌 & 同一时间轴的活字节、分配器 retained 字段 &
只能坐实或排除 ``未归还''；是不是泄漏走《内存泄漏检测与防御》册 \\
\hline
\Mx{p99} 出现规律尖峰 & 分配调用耗时直方图、arena 扩容事件 &
到 ``尖峰与扩容同刻''；是否被前移要看池上限与复位粒度 \\
\hline
换分配器后整体变差 & arena 数量上限、大请求阈值、统计接口开关 &
三个变量各跑一次，才能断定是哪一个而不是 ``换分配器'' \\
\hline
CI 报内存回归 & 构建指纹、sanitizer 开关、采样周期 &
先确认口径没变；工具与接口开关的变化不算被测程序的回归 \\
\hline
\end{tabular}
\end{compare}

\begin{bestpractice}[title={内存侧改动上线前的六条自检}]
  \begin{enumerate}
    \item 口径写在数字前面：这组数据是 \MWall{}、\MCpu{}、\MThr{} 还是 \Mx{p99}，谁包含谁。
    \item 计数器与计时器不在同一次运行里开；两者都用了就分别写明各来自哪一次。
    \item 归因用了三段实验（只换分配、只换布局、两者都换），并写出收益是否可叠加。
    \item 局部性结论有拐点位移的证据，而不是一个单点。
    \item 内存结论同时给出 \Mx{RSS} 与活字节两条时序，并给出池上限与归还时机。
    \item sanitizer 与 Valgrind 的倍率单独记在自己的指纹下，绝不与原生基线相减；需要 PMU 的指标
          不进 CI 门禁。
  \end{enumerate}
\end{bestpractice}

收束一句：本章所有手段都服务于同一个判据——\textbf{改动之前先能说出哪本账会被动，改动之后那本账
必须被动}。四本账（分配路径、局部性、驻留、归还）中任何一本都能独自解释一次 \MWall{} 的变化，所以
只有总和的变化不构成结论。要看线程之间的分配争用与共享缓存干扰，去 第~\ref{ch:pm09-concurrency} 章；要落到
某一行代码的缺失来源，回 第~\ref{ch:pm06-micro} 章；要把这些数变成门禁，去 第~\ref{ch:pm13-ci} 章。

\chapter{I/O 与存储测量：块设备、页缓存与延迟分账}
\label{ch:pm08-io}

\MLat{}（一次请求从发出到完成的延迟）与 \MThr{}（单位时间完成的工作量，吞吐）在 I/O 上比在
CPU 上更容易被偷换：同一段代码，冷启动的那一次与缓存已经建起来的那一次能差出数量级，
而你未必知道差在哪一层——本章把 ``跑第二次'' 当成一种\emph{口径}而不是错误，
错的是拿热运行的数去回答冷启动的问题。本章要做的就是把\emph{更快}拆开：
一次 I/O 的 \MWall{} 属于四段账——应用调用、内核（页缓存与回写）、块层与调度器、设备——
每一段用什么证据看、哪些默认值会悄悄替你写好结论。
与前面的接口：时间从哪来见 第~\ref{ch:pm02-clock} 章，重复与分布的判据见 第~\ref{ch:pm03-stats} 章，
缺页与调度事件的口径见 第~\ref{ch:pm05-events} 章；与后面的接口：跨线程的排队与背压归
第~\ref{ch:pm09-concurrency} 章，I/O 基线挂进流水线归 第~\ref{ch:pm13-ci} 章，
线上低开销地采 I/O 归 第~\ref{ch:pm14-production} 章。

\section{四段延迟分账：应用、内核、块层与设备}
\label{sec:pm08-ledger}

一次 \Fn{read} 的 \MWall{} 是四段时间的和，但只有第一段是你自己的代码。
按\emph{边界}而不是按\emph{工具}列出它们，是本章唯一需要背下来的东西。

\begin{itemize}
\item \textbf{应用段}：请求在你进程里排队、序列化、拷贝进缓冲区之前的时间。
  边界是系统调用进入内核的那一刻。这一段最常见的成因是待发请求太少而不是设备慢：
  单线程同步读、一个在飞请求都没有时，再快的设备也吃不满。
\item \textbf{内核段}：VFS、页缓存查找、缺页、回写、预读。
  边界是 ``命中缓存直接返回'' 与 ``生成一个块请求'' 的分岔——
  \textbf{命中页缓存的请求永远不到达块层}，所以块设备的任何计数器都看不见它。
\item \textbf{块层段}：请求合并、调度器排队、硬件队列深度、扇区到命令的转换。
  边界是请求被下发给驱动的那一刻。
\item \textbf{设备段}：控制器与介质、盘内缓存、RAID 条带重组与读改写；
  虚拟机里还要加上宿主机后端那一段。边界是完成中断与请求结束。
\end{itemize}

这四段在 Linux 上有直接的仪器形态：\Tp{blktrace} 给每个请求打上 \Tp{Q}（进入块层队列）、
\Tp{G}（拿到请求结构）、\Tp{I}（下发驱动）、\Tp{D}（设备完成）四类时间戳，
\Tp{blkparse} 与 \Tp{btt} 把相邻时间戳摊成三段耗时。它需要特权、要求被测设备未被独占，
而且自身开销不小——把它当\emph{形状}工具而不是\emph{数值}工具（判据见 第~\ref{ch:pm03-stats} 章）。

\begin{lstlisting}[style=shellstyle,caption={分账起手式：在同一观测窗内同时取块层与应用侧证据（命令可跑，路径请替换）},label={lst:pm08-ledger}]
WIN_START=$(date +%s)                              # 1) 先定窗口，再决定谁在窗口里跑
iostat -x -t 1 30 > /tmp/io.win.txt &               # 2) 不给间隔时 iostat 报的是开机至今均值
grep -E '^(MemAvailable|Cached|Dirty|Writeback)' /proc/meminfo > /tmp/mem.before
cat /proc/self/io                                 # 3) 应用侧提交量（换成被测进程 pid）
#    read_bytes / written_bytes 是提交给文件系统层的量，不等于设备字节
# 4) 块层内部事件（需要特权，只看形状）
#    blktrace -d /dev/nvme0n1 -o /tmp/blk -w 30 && blkparse -i /tmp/blk -d /tmp/blk.bin
#    btt -i /tmp/blk.bin     # 给出 Q2G / G2I / I2D 三段耗时
uname -r; cat /sys/block/nvme0n1/queue/scheduler   # 5) 环境指纹进结论
lsblk -d -o NAME,ROTA,TYPE,SIZE,MODEL,TRAN
\end{lstlisting}

下面这个样例把 \Tp{iostat -x} 的字段与 \Tp{fio} 报告里对应的量并排。
\textbf{数值一律是占位符}：本章不给任何看着像实测值的延迟数字，
量级要你在自己的设备上标定（第 2 节与第 3 节给方法）。

\begin{lstlisting}[style=outstyle,caption={四段分账的字段对照（格式示例，非本仓库实测值；尖括号内为待填量）},label={lst:pm08-fields}]
# ---- iostat -x 1 的表头与一行形态（字段名以你机器为准）----
Device   r/s    w/s   rkB/s   wkB/s  rrqm/s  r_await  w_await  aqu-sz  %util
nvme0n1  <r/s>  <w/s> <rkB/s> <wkB/s> <rqm/s> <t_aw>   <t_aw>   <q>     <pct>
#   r/s, w/s      : 合并后完成的请求数/秒 —— 不等于应用侧 read()/write() 次数
#   rkB/s, wkB/s  : 块层字节速率           —— 页缓存命中的读不在这里
#   rrqm/s,wrqm/s : 每秒被合并的请求数     —— 反推调用次数要把它们加回去
#   r_await       : 读请求在块层+设备的排队+服务总时间，不含应用段
#   aqu-sz        : 平均队列深度           —— 与 fio 的 iodepth 对照着看
#   %util         : 设备至少有一个在飞请求的时间占比，不等于饱和
# ---- 同一次 fio 作业报告里对应的三行（形态示意）----
#  read: IOPS=<n>  BW=<x>MiB/s  runtime=<y>msec
#    slat : <均值> <分布>     提交耗时：fio 把请求交出去花的时间
#    clat : <均值> <p50 p99>  完成耗时：块层 + 设备（与 await 同段）
#    lat  : <均值> <p50 p99>  总耗时：约等于 slat + clat
#  对照：aqu-sz 应与 iodepth 同数量级；对不上就是有一段没算进来
\end{lstlisting}

\begin{compare}[title={应用侧证据与块层证据：谁看哪一段}]
\begin{tabular}{>{\raggedright\arraybackslash}p{0.19\textwidth} >{\raggedright\arraybackslash}p{0.26\textwidth} >{\raggedright\arraybackslash}p{0.31\textwidth}}
\toprule
证据来源 & 它的边界在哪 & 看不见什么 \\
\midrule
应用计时（自己埋点） & 包住 \Fn{read} 的 \MWall{} & 分不清缓存命中与设备读 \\
\Tp{fio} 的 \Tp{clat} & 块层下发到完成 & 用户态排队与数据拷贝 \\
\Tp{iostat -x} 的 \Tp{r\_await} & 请求进块层到完成 & 页缓存命中、应用段 \\
\Tp{blktrace} 的 \Tp{I2D} & 下发驱动到设备完成 & 合并前的排队与调度 \\
\Tp{/proc/self/io} & 文件系统层提交的字节 & 设备实际写入（写放大漏掉） \\
\bottomrule
\end{tabular}
\end{compare}

\begin{warning}[title={默认值陷阱：\%util 被当成 ``设备已经满了''}]
机械盘时代的 \Tp{\%util} 接近 $100$ 确实意味着几乎没有空闲服务窗口；但现代设备普遍能并行
处理多个请求（NVMe 多队列、SATA 的 NCQ、RAID 条带、多路径、虚拟机后端），
\Tp{\%util} 只回答 ``有没有请求在飞''，不回答 ``还能不能再塞''。把它当饱和度，
会在设备远未饱和时得出 ``瓶颈在存储'' 的结论。判断饱和的正确做法是差分：
把 \Opt{iodepth}（或并发请求数）翻倍再测一次，看 \MThr{} 是否随之上升、
\Mx{clat} 的上尾是否随之上抬；两条曲线拐弯的地方才是你自己的饱和口径。
\end{warning}

\begin{guideline}[title={写一份能辩护的 I/O 结论之前，先钉死这四个边界}]
\begin{enumerate}
\item \textbf{窗口}：起止时间戳与时长，以及所有观测工具用的是不是同一个窗口；
  不同窗口的均值不能相加。
\item \textbf{边界}：你说的是应用段、块层段还是设备段？\MWall{}、\MCpu{}、
  \Mx{clat}、\Mx{await} 是四个不同的数（口径辨析见 第~\ref{ch:pm01-why} 章）。
\item \textbf{单位}：请求数、字节数还是带宽；用逻辑扇区还是应用块大小；
  \Tp{KiB} 与 \Tp{kB} 混用会让两个 ``同一个数'' 差出几个百分点。
\item \textbf{缓存}：本次运行里哪些层有缓存参与（第 2 节），
  对照组是不是同样的缓存状态。
\end{enumerate}
\end{guideline}

\section{缓存有没有参与：冷启动、\Tp{drop\_caches} 与 O\_DIRECT}
\label{sec:pm08-cache}

``先问缓存有没有参与'' 是 I/O 测量的第一句话，因为参与的层比你以为的多：
Linux 页缓存、块层预读、文件系统元数据缓存、设备内 DRAM、RAID 控制器缓存（常带
flush 语义）、虚拟机宿主的页缓存、容器 overlayfs 的上层。
其中只有前两层能用 \Tp{drop\_caches} 一类手段碰到，其余必须换设备或换路径才能排除。

\subsection{怎么知道命中了页缓存}
\label{sub:pm08-hit}

三条互相独立的证据，任一条都能单独证伪 ``这次跑打到设备了''：
\Tp{iostat -x} 在该窗口内 \Tp{r/s} 与 \Tp{rkB/s} 全程为零而应用确实读到了数据，
说明命中页缓存、块层没被惊动；\Fn{mincore}（或 \Lib{vmtouch} 这类封装）报告你关心的
区间还有多少页驻留，这是最直接的证据，能精确到偏移；\Tp{/proc/meminfo} 里 \Tp{Cached}
随读取上升，冷跑之后的增量约等于读进去的字节（被回收掉的除外）。

\subsection{三种 ``变冷'' 的口径与各自的失真}
\label{sub:pm08-cold}

\begin{enumerate}
\item \textbf{\Tp{/proc/sys/vm/drop\_caches}}：写 \Tp{1} 释放页缓存，\Tp{2} 释放
  dentry 与 inode 缓存，\Tp{3} 两者都做。三处失真：它只丢弃\emph{干净}页，
  必须先 \Fn{sync} 再用，否则你以为的冷其实是脏页还在；它是\textbf{全局}操作，
  把同机其他进程的缓存一并丢掉，共享机器上等于给所有人制造冷启动，对照组从此不可比；
  它丢掉整机缓存，于是随后第一次跑比被测程序真实需要的 ``冷'' 更冷，数字系统性偏坏。
  写这个文件需要特权，且并非所有内核都允许。
\item \textbf{只清自己：}\Fn{posix\_fadvise} 用 \Tp{POSIX\_FADV\_DONTNEED}
  逐出被测文件区间，或每次换一个新建文件。它不打扰别人，也更贴近
  ``这个数据集不在缓存里'' 这个真实问题；代价是逐出不保证瞬时完成，要再复核驻留比例。
\item \textbf{让工作集大于内存：}最稳的冷法。把测试文件大小推到明显超过可用内存，
  缓存自然帮不上忙，也顺带稀释预读的影响。代价是运行变长，且文件建立本身要计入
  （用 \Fn{fallocate} 预分配，别边扩边测）。
\end{enumerate}

\begin{warning}[title={``第二次跑当然快''：最常见的 I/O 假结论}]
一个补丁声称 ``随机读变快''，证据是第二次运行比第一次快。真正发生的事情多半是
第一次把数据搬进了页缓存，第二次压根没打设备。这一类错误有三个变体：
测试文件比内存小；预热与计时之间没有区分缓存状态；对照组与实验组不是同一次冷启动。
正确做法是把冷/热当成两个\emph{显式}口径分别编号：
``冷启动第 1 次运行'' 与 ``缓存命中运行'' 是两个结论，不是同一个结论跑两遍。
重复次数与判据见 第~\ref{ch:pm03-stats} 章；``同时改了两件事'' 的排法见
第~\ref{ch:pm15-antipattern} 章。
\end{warning}

\subsection{缓冲 I/O 与 O\_DIRECT：两个不同的问题}
\label{sub:pm08-odirect}

\Tp{O\_DIRECT}（\Tp{fio} 里的 \Opt{direct=1}）绕过页缓存，把请求直接交给块层。
它不是 ``更真实''，只是换了一个问题：\textbf{缓冲 I/O} 回答
``在这个内核、这份内存、这种访问模式下，程序会经历什么''，包含预读、包含合并、
包含命中红利，也包含内核替你做的排队；\textbf{直接 I/O} 回答
``这块设备在这种请求模式下能提供多少''，它把内核缓存策略从等式里抹掉，
因此通常让\emph{顺序读带宽}变低（没有预读）、让\emph{随机小读}更接近设备本身、
让重复读取不再有红利。

三个必须自己确认的前提。一是\textbf{对齐}：缓冲区地址、文件偏移与传输长度通常都要
按设备逻辑块大小对齐，先读 \Tp{/sys/block/<dev>/queue/logical{\Bk}\_block{\Bk}\_size}
这个值（\Tp{io\_min} 与 \Tp{io\_opt} 是它的同族视图），再用 \Fn{posix\_memalign}
一类手段取对齐缓冲。二是\textbf{短读}：直接 I/O 允许返回少于请求的长度，
测量代码必须按实际完成量记账，否则 \MThr{} 被高估。三是\textbf{元数据不归它管}：
路径解析、inode 与目录项仍走缓存，而文件在测试中被\emph{扩展}时还会额外产生元数据写，
``纯数据写'' 的口径就没了（先预分配）。此外，并非所有文件系统与挂载选项组合都支持它，
某些配置下未对齐的直接访问会静默退化成缓冲访问，此时你的 ``直接 I/O 数'' 其实是缓存数
——与 \Tp{iostat} 交叉核对是唯一可靠的办法。

\begin{note}[title={fio 会自己动缓存：\Opt{invalidate}、预读与 \Tp{fadvise}}]
\Tp{fio} 在打开文件前默认要做一次失效处理（\Tp{invalidate}，默认开启）：
对目标文件同步并把缓存里对应内容丢掉，某些引擎与文件系统组合下还会先顺序读一遍。
也就是说 ``跑第一遍就是冷的'' 未必成立，也可能是 fio 已经把它调成自己想要的冷。
同样，块层的预读窗口 \Tp{read\_ahead\_kb}（或 \Tp{blockdev --getra}，
后者以 $512$ 字节扇区为单位）是 ``顺序读带宽'' 这个数的一部分成因。
把这些默认值抄进环境指纹，别当常量。
\end{note}

\begin{lstlisting}[style=shellstyle,caption={测之前先确认设备与调度器：可跑的检查清单（全部本机自证，不抄别人的表）},label={lst:pm08-probe-cmd}]
# ---- 我是谁：设备类型与层叠关系 ----
lsblk -o NAME,ROTA,TYPE,SIZE,MODEL,TRAN,MOUNTPOINT,FSTYPE   # ROTA=0 即非旋转
findmnt -T /path/to/testfile; df -T /path/to/testfile  # 落在 tmpfs 上就没有设备参与
dmsetup ls 2>/dev/null; cat /proc/mdstat 2>/dev/null   # dm/LVM/crypt/软RAID 的层叠
# ---- 队列与调度器：决定块层段长什么样 ----
Q=/sys/block/nvme0n1/queue
cat $Q/scheduler $Q/nr_requests $Q/rotational $Q/read_ahead_kb
cat $Q/logical_block_size $Q/physical_block_size $Q/max_sectors_kb
cat $Q/optimal_io_size   # 0 表示设备未上报建议粒度
ls /sys/block/nvme0n1/mq/ 2>/dev/null | wc -l; nproc  # 硬件队列数与核数是否对上
# ---- 写缓存语义：决定写延迟归谁 ----
cat /sys/block/sda/queue/write_cache    # transient 还是 writethrough
nvme id-ctrl /dev/nvme0 2>/dev/null | grep -i vwc    # 是否暴露易失写缓存
nvme smart-log /dev/nvme0 2>/dev/null | grep -iE 'data_units|percentage_used'
# ---- 环境指纹（整段留档，写进结论）----
uname -a
sysctl vm.dirty_ratio vm.dirty_background_ratio vm.dirty_expire_centisecs 2>/dev/null
mount | sed -n '1,20p'        # barrier / sync / noatime 等挂载项
\end{lstlisting}

\begin{table}[htbp]
\centering
\begin{tabularx}{\textwidth}{>{\raggedright\arraybackslash}p{0.20\textwidth} >{\raggedright\arraybackslash}X >{\raggedright\arraybackslash}p{0.22\textwidth}}
\toprule
确认项 & 怎么确认（本机） & 忽略了会怎样 \\
\midrule
旋转还是固态 & \Tp{lsblk -o ROTA}、\Tp{queue/rotational}；再用随机读差分印证 & 把预读红利当成设备能力 \\
有没有 \Tp{dm}/RAID 层 & \Tp{lsblk} 层叠树、\Tp{dmsetup table}、\Tp{/proc/mdstat} & 分账归错层，读改写看不见 \\
调度器与队列深度 & \Tp{queue/scheduler} 的方括号项、\Tp{nr{\Bk}\_requests}、\Tp{mq} 目录 & 把排队时间当成设备时间 \\
预读窗口 & \Tp{queue/read{\Bk}\_ahead{\Bk}\_kb}、\Tp{blockdev --getra} & 顺序带宽含了内核预取 \\
写缓存语义 & \Tp{queue/write{\Bk}\_cache}、控制器自报的易失缓存 & 写带宽其实是缓存带宽 \\
路径落在哪个文件系统 & \Tp{findmnt -T}、\Tp{df -T} & tmpfs/overlay 上没有设备账 \\
对齐粒度 & \Tp{logical{\Bk}\_block{\Bk}\_size}、\Tp{optimal{\Bk}\_io{\Bk}\_size} & O\_DIRECT 报错或静默退化 \\
\bottomrule
\end{tabularx}
\caption{I/O 测量前的设备确认表（内容随内核与设备而变，一切以本机命令输出为准）}
\label{tab:pm08-probe}
\end{table}

\section{fio 的口径面：每个参数改的是哪一段账}
\label{sec:pm08-fio}

\Tp{fio} 是 ``参数即口径'' 的典型工具：它没有默认的正确测法，只有默认的参数。
纪律是\textbf{任何一组参数都要能用一句话说出它在问什么问题}，说不出就还没测。
不要相信文档里的默认值——让 \Tp{fio} 自己把生效参数打出来：
\Tp{fio --showcmd} 回显这一组作业真正生效的选项行，\Tp{--enghelp} 列出这个版本认哪些引擎，
\Tp{--cmdhelp} 列出支持的选项全集（三者都随版本变，把输出抄进结论）。

\subsection{三组典型作业}
\label{sub:pm08-fio-jobs}

\begin{lstlisting}[style=shellstyle,caption={三组典型作业骨架：每段注明它测的是哪一段账（路径与设备请自己替换）},label={lst:pm08-fio-jobs}]
FILE=/srv/fio/bigfile        # 文件口径：含文件系统与页缓存
DEV=/dev/nvme0n1             # 裸设备口径：需要特权，绕过文件系统，测的是设备
# ---- A. 顺序读带宽：这种访问模式下内核+设备能搬多快 ----
fio --name=seq-read --filename=$FILE --rw=read --bs=1M \
    --size=16G --runtime=120 --time_based --ramp_time=10 \
    --ioengine=libaio --iodepth=32 --direct=1 \
    --group_reporting --latency_log --log_avg_msec=100
#   direct=1   : 排除页缓存与预读，因此通常低于 buffered 的同款作业
#   bs=1M      : 大传输，测带宽而不是测每请求开销
#   iodepth=32 : 让足够多请求在飞，才喂得出设备并行度
#   ramp_time  : 丢掉起步瞬态（频率爬升、队列建立、后端积压）
#   注意       : ioengine=psync 这类同步引擎不支持队列深度，iodepth>1 无意义
# ---- B. 随机 4K 读 IOPS：设备能吃多少并发小请求 ----
fio --name=rand-4k-read --filename=$DEV --rw=randread --bs=4k \
    --size=16G --runtime=120 --time_based --ramp_time=10 \
    --ioengine=io_uring --iodepth=64 --numjobs=4 --direct=1 \
    --percentile_list=50:90:95:99:99.9 --group_reporting
#   numjobs=4 + iodepth=64 : 总并发是乘出来的，iodepth 本身不等于并发
#   percentile_list        : 尾延迟必须显式要，均值永远看不出上尾
#   io_uring 需内核与权限支持；退到 libaio 时语义不同但都不走页缓存
# ---- C. 混合读写（把提交语义算进来）：读写互相拖累多少 ----
fio --name=mixed-rw --filename=$FILE --rw=randrw --rwmixread=70 \
    --bs=8k --bssplit=4k/50:8k/30:64k/20 --size=8G \
    --runtime=120 --time_based --ioengine=libaio --iodepth=16 --direct=1 \
    --numjobs=4 --fsync=64 --end_group=0 --group_reporting
#   rwmixread / bssplit : 让请求分布贴近真实负载，别让最 favorable 的模式代表全部
#   fsync=64            : 每 64 次写做一次同步，把设备缓存 flush 的账加回来
#   end_group           : 组内同步时机；取值不同改的是谁在等谁
\end{lstlisting}

\begin{longtable}{>{\raggedright\arraybackslash}p{0.19\textwidth} >{\raggedright\arraybackslash}p{0.11\textwidth} >{\raggedright\arraybackslash}p{0.27\textwidth} >{\raggedright\arraybackslash}p{0.23\textwidth}}
\caption{fio 关键参数与它改变的测量含义（默认倾向以 \Opt{showcmd} 输出为准）}\label{tab:pm08-fio-param}\\
\toprule
参数 & 默认倾向 & 它改了哪一段账 & 写错的后果 \\
\midrule
\endfirsthead
\toprule
参数 & 默认倾向 & 它改了哪一段账 & 写错的后果 \\
\midrule
\endhead
\bottomrule
\endfoot
\Opt{ioengine} & 同步类 & 提交方式：能否排队、是否走内核异步路径 & 与 \Opt{iodepth} 组合无效，出现 ``加并发没变快'' \\
\Opt{iodepth} & 很浅 & 块层在飞请求数 & 浅队列喂不出并行度，把能力当上限 \\
\Opt{direct} & 缓冲 & 页缓存与预读是否参与 & 报成设备带宽，其实是缓存带宽 \\
\Opt{rw} & 顺序读 & 访问模式，决定能否吃预读与合并 & 随机业务拿顺序数辩护 \\
\Opt{bs} / \Opt{bssplit} & 单块大小 & 每请求工作量与请求数 & 只测 favorable 尺寸，忽略真实分布 \\
\Opt{size} & 取决于版本 & 工作集是否大于内存 & 数据集全塞进缓存，第二轮起全命中 \\
\Opt{numjobs} & $1$ & 进程/线程数与总并发 & 忘了 \Opt{group{\Bk}\_reporting} 时读数口径不同 \\
\Opt{time{\Bk}\_based}、\Opt{ramp{\Bk}\_time} & 跑完即止 & 计时窗口 & 建文件与收尾回写被算进稳态 \\
\Opt{fsync}、\Opt{end{\Bk}\_group} & 不同步 & 设备缓存 flush 是否入账 & 写数漂亮但没落到介质 \\
\Opt{latency{\Bk}\_log}、\Opt{percentile} & 只报摘要 & 尾延迟可见性 & 只有均值，看不到上尾 \\
\Opt{randrepeat} & 固定图形 & 两次跑是否走同一条随机序列 & 对照实验的运行之间不可比 \\
\end{longtable}

\subsection{读 fio 的输出：三行延迟不是三次测量}
\label{sub:pm08-fio-out}

\Tp{slat}、\Tp{clat}、\Tp{lat} 是三种口径的和差关系：\Tp{slat} 是提交耗时
（\Tp{fio} 把请求交出去所花的时间），\Tp{clat} 是完成耗时（块层加设备），\Tp{lat} 是总耗时。
\textbf{同步引擎下提交即阻塞}，此时 \Tp{slat} 里已经含设备时间，\Tp{clat} 的含义与异步引擎
完全不同——把两种引擎的 \Tp{slat} 并排比较就是口径偷换。另有三条纪律：
百分位是 ``这段窗口内'' 的百分位，样本量不足时 \Mx{p99.99} 只是几个样本
（判据见 第~\ref{ch:pm03-stats} 章）；\Opt{group\_reporting} 会把多作业的分布\emph{合并}，
一个慢作业的上尾被其他作业的多数样本稀释，报尾延迟前先确认合并口径；
\Mx{IOPS} 与 \Mx{BW} 是总工作量除以运行时长算出的，只要窗口里含建文件、
收尾回写或频率爬升，这两个数会一起偏。

\begin{bestpractice}[title={A/B 对照的正确跑法}]
\begin{enumerate}
\item 参数集固定成一份 job 文件并进版本库；一次只改一处，
  否则又回到 ``同时改了两件事''。
\item 冷/热各一份口径分别编号；同一 job 各跑 $N$ 次，$N$ 由要检出的差异反推。
\item 交替执行（A/B/A/B），消掉先后顺序带来的频率与温度偏置。
\item 每次运行都同时留 \Tp{iostat} 窗口与设备指纹（第 2 节清单的输出）。
\item 报分布不报单点：把 \Opt{latency\_log} 的日志留在产物里，
  结论里给分位数与区间，而不是给一个 ``平均值''。
\end{enumerate}
\end{bestpractice}

\section{同步调用、写缓存与写放大}
\label{sec:pm08-sync}

\subsection{\Fn{fsync}、\Fn{fdatasync}、\Fn{msync} 的代价怎么自己测}
\label{sub:pm08-fsync}

三者语义不同：\Fn{fsync} 把数据与 ``恢复该文件所需'' 的元数据一起推到设备可见状态；
\Fn{fdatasync} 只在元数据影响数据取回时才同步它，通常少一次元数据写；
\Fn{msync} 作用于 \Fn{mmap} 的共享映射，把脏页交回回写路径，
是否真等到设备取决于 \Tp{MS\_SYNC} 还是 \Tp{MS\_ASYNC}。
它们的耗时都不是设备的一个属性，而是\emph{当时有多少脏页}与
\emph{文件系统怎么组织提交}的函数，所以只能自己测。

\begin{lstlisting}[style=cstyle,caption={自测提交点延迟的骨架：测出来的是这台机器在这种脏页速率下的数（不含任何预期数值）},label={lst:pm08-fsync}]
// 编译：cc -O2 -std=c17 -D_GNU_SOURCE sync_probe.c
#include <fcntl.h>
#include <stdint.h>
#include <stdio.h>
#include <stdlib.h>
#include <string.h>
#include <time.h>
#include <unistd.h>
#define BLOCK (64u * 1024u)

static inline uint64_t now_ns(void) {              // 时钟选型见第 2 章
    struct timespec ts;
    clock_gettime(CLOCK_MONOTONIC, &ts);
    return (uint64_t)ts.tv_sec * 1000000000ull + (uint64_t)ts.tv_nsec;
}
static int cmp_u64(const void *a, const void *b) {
    uint64_t x = *(const uint64_t *)a, y = *(const uint64_t *)b;
    return (x > y) - (x < y);
}
int main(int argc, char **argv) {
    const char *path = argc > 1 ? argv[1] : "/var/tmp/sync_probe.bin";
    int blocks = argc > 2 ? atoi(argv[2]) : 1024;   // 工作集
    int every  = argc > 3 ? atoi(argv[3]) : 1;      // 每 every 次写做一次同步
    int rounds = argc > 4 ? atoi(argv[4]) : 30;     // 重复次数
    int fd = open(path, O_CREAT | O_RDWR | O_TRUNC, 0644);
    if (fd < 0) { perror("open"); return 1; }
    char *buf = malloc(BLOCK);                      // 换 O_DIRECT 时这里要对齐分配
    memset(buf, 0xAB, BLOCK);
    uint64_t *s = malloc(sizeof *s * (size_t)rounds);
    for (int r = 0; r < rounds; r++) {
        uint64_t t0 = now_ns();
        for (int i = 0; i < blocks; i++) {
            if (pwrite(fd, buf, BLOCK, (off_t)i * BLOCK) != (ssize_t)BLOCK) return 1;
            if ((i + 1) % every == 0 && fdatasync(fd) != 0) return 1;
        }
        if (fsync(fd) != 0) return 1;               // 收尾：把脏页压干
        s[r] = now_ns() - t0;
    }
    qsort(s, (size_t)rounds, sizeof *s, cmp_u64);
    size_t h = (size_t)(rounds * 0.5), g = (size_t)(rounds * 0.99);
    printf("every=%d rounds=%d min=%llu p50=%llu p99=%llu max=%llu ns\n",
           every, rounds, (unsigned long long)s[0], (unsigned long long)s[h],
           (unsigned long long)s[g], (unsigned long long)s[rounds - 1]);
    close(fd);                                     // 数字被打印，循环不会被整体删掉
    return 0;
}
\end{lstlisting}

读这个数时要知道四件事，它们决定了同一个 \Fn{fsync} 为什么差很多。
第一，回写由脏页阈值与超时驱动（\Tp{vm{\Bk}.dirty{\Bk}\_ratio}、
\Tp{vm{\Bk}.dirty{\Bk}\_background{\Bk}\_ratio}、
\Tp{vm{\Bk}.dirty{\Bk}\_expire{\Bk}\_centisecs} 一类参数，名称随内核版本变，
先用 \Tp{sysctl} 列出本机实际值再说话），于是刚写完就同步的孤立一次 \Fn{fsync}
与连续流里的 \Fn{fsync}（后台回写已经开始）是两个不同的数：专测提交点延迟时，
要么给连续流，要么把间隔固定写进结论。第二，日志式文件系统的 \Fn{fsync}
可能要等前一个事务，延迟呈多数很快、偶发极慢的双峰形态，这正是必须报分布的理由。
第三，带易失写缓存的盘在完成中断时数据多半还在盘内缓存里，只有 flush 或 FUA
那条命令才把介质算进来，所以不带 \Opt{fsync} 的写带宽与带 \Opt{fsync} 的写带宽
是两个问题，前者常常是缓存带宽（确认方法见第 2 节）。
第四，\Fn{mmap} 的写页只产生脏页，真正下设备是在回写或 \Fn{msync} 时，
测 ``mmap 写入快'' 要问清测的是缺页成本还是提交成本
（缺页口径见 第~\ref{ch:pm05-events} 章，分配与局部性见 第~\ref{ch:pm07-memory} 章）。

\subsection{写入放大：机制与观察点}
\label{sub:pm08-waf}

写放大（应用写多少字节、设备实际写多少字节）不是单一层的属性，至少来自五处：
文件系统的块大小与延迟分配（小写被聚成大写，反向也成立）、日志结构的两次写、
RAID 条带的读改写（部分条带写要先读旧块）、SSD 内部映射表与垃圾回收的搬移写、
以及数据库层的检查点与压实。虚拟化与容器（写时复制、精简置备）再叠一层。
正确的观察姿势是\textbf{同一次运行里取两端的字节数做比}：应用端用
\Tp{/proc/<pid>/io} 的 \Tp{written\_bytes} 或 \Tp{fio} 写的总字节，
块层端用 \Tp{iostat -x} 的 \Tp{wkB/s} 在窗口上积分，
设备端用 SMART 或 NVMe 统计里的写入单元数（单位由设备自报，
先确认它按 $1000$ 还是 $1024$ 折算再换算）。
三层两两相除给出的是不同口径，报数时要说明分的是哪两层。
压缩、克隆（reflink、\Tp{FICLONE}）与去重会让应用写的字节与设备写的字节完全脱钩——
克隆写几乎不产生设备写，这就是 ``复制大文件带宽爆表'' 的来源。

\begin{warning}[title={路径本身就能让数字失真：tmpfs、overlay、快照}]
测试文件放在 \Tp{/dev/shm} 或任何 tmpfs 上时，\Tp{iostat} 全程没有计数而 \MThr{} 高得离谱；
容器里你以为写的是宿主机磁盘，实际落在 overlayfs 上层，首次写还要付出 copy-up；
带本地快照的文件系统会让 ``写'' 变成 ``写新块''；开启压缩的挂载会拉开应用字节与设备字节。
先跑 \Tp{findmnt -T <路径>} 与 \Tp{df -T <路径>} 确认你测的是哪块设备，再讨论数字。
\end{warning}

\section{网络、管道与另两平台的对照}
\label{sec:pm08-net}

\subsection{带宽、延迟、抖动是三件事}
\label{sub:pm08-net3}

管道类测量最常见的偷换是把 \MThr{} 当 \MLat{}。链路的吞吐由最窄的一段乘时间决定，
而两侧的缓冲区（socket 缓冲、网卡环、交换队列、内核队列）足以让
``平均很快'' 与 ``某个请求等得很久'' 同时成立：缓冲足够大时 \MThr{} 好看，
代价是排队延迟搬进了尾部分布。这与 第~\ref{ch:pm09-concurrency} 章 里
``无界队列把延迟搬进队列'' 是同一件事，判据也一样：
吞吐必须与队列长度或往返时间的分位数同时报，只报其一就是把相关性当因果。

\subsection{\Tp{iperf} 一类工具的默认值陷阱}
\label{sub:pm08-iperf}

\begin{itemize}
\item \textbf{默认是 TCP}。结果含缓冲、拥塞控制、重传与自动调窗，
  它回答 ``这个 TCP 栈在这条路径上能吃多少''，不回答 ``链路容量是多少''。
\item \textbf{UDP 通常带一个保守的默认目标码率}：不给 \Opt{bitrate} 时你可能在测
  远低于线速的一点流量，却把它报成链路能力。任何 UDP 数都要显式给速率并写明。
\item \textbf{间隔报告是均值}，且看不到尾延迟：间隔越短越接近瞬时，
  但更受调度与缓冲影响；窗口长度本身也要写进结论。
\item \textbf{改窗口或缓冲就是改结论}：这类工具的窗口参数改的是 socket 缓冲，
  也就改了排队那一侧。要延迟就换能测单请求往返的工具——
  连接建立时间、往返分布与重传率不是带宽工具的输出量，反向模式也变不出它们。
\end{itemize}

\subsection{内核侧证据：连接计数、网卡统计与 NAPI}
\label{sub:pm08-nic}

\Tp{ss -ti} 给每条连接的往返时间与其方差、重传、拥塞窗口、接收队列占用，
这是单个请求的等待的侧面证据，比链路级均值更接近应用经历。
\Tp{/proc/net/dev} 的丢弃与错误类计数、\Tp{nstat} 的协议计数都\textbf{必须做差分}：
它们是开机以来的累计值，绝对值没有窗口含义。
\Tp{/proc/net/softnet{\Bk}\_stat} 每 CPU 一行的含义依次是处理数、因无预算而丢弃数、
时间片耗尽次数：最后一列持续增长说明软中断轮询被自己的预算打断，表现为
``设备没满但在丢包''，这是 NAPI 轮询模型的典型症状。
队列与中断侧用 \Tp{ethtool -l} 看通道数与核数是否对上，\Tp{ethtool -c} 看中断合并设置
（它同时改变延迟分布与 CPU 税），\Tp{/proc/interrupts} 与
\Tp{/proc/irq/<n>/smp{\Bk}\_affinity} 看中断落在哪个 CPU、\Tp{irqbalance} 会不会动它。
offload（GRO、TSO、LRO）改变 ``包数'' 的含义：统计里的一个大包在应用侧可能是若干段，
按包数换算 \MThr{} 会得出不一致的数。

\begin{lstlisting}[style=shellstyle,caption={网络侧的窗口证据：所有计数都要两次快照相减},label={lst:pm08-net}]
ss -tin state established | head -30           # 每条流的 rtt / retrans / 队列占用
ethtool -S eth0 > /tmp/nic.t0; sleep 10; ethtool -S eth0 > /tmp/nic.t1
diff /tmp/nic.t0 /tmp/nic.t1 | grep -Ei 'drop|err|fifo|miss'   # 窗口内的增量
awk 'NR<=8{print NR-1, $1, $2, $3}' /proc/net/softnet_stat     # 处理/丢弃/预算耗尽
ethtool -l eth0; ethtool -c eth0; ethtool -k eth0 | grep -Ei 'receive|segmentation'
grep -E 'eth0|virtio' /proc/interrupts         # 中断落在哪些 CPU
ss -ltn | awk 'NR>1{print $2, $3, $4}'         # 监听积压：应用来不及取，不是网络慢
\end{lstlisting}

\subsection{Windows：\Fn{ReadFile}、IOCP 与 ETW 磁盘事件}
\label{sub:pm08-win}

\Tp{O\_DIRECT} 在这里的近似物是 \Tp{FILE{\Bk}\_FLAG{\Bk}\_NO{\Bk}\_BUFFERING}：
同样要求缓冲区、偏移与长度按卷扇区对齐，同样绕过系统缓存但不消掉设备自己的写缓存；
\Tp{FILE{\Bk}\_FLAG{\Bk}\_WRITE{\Bk}\_THROUGH} 只保证不留在系统缓存，
介质与设备缓存仍可能吸收写。队列深度不由单个 \Fn{ReadFile} 的参数给出，
而由同时挂起的重叠请求个数给出，配合完成端口（IOCP）时常常只有一个待处理请求——
浅队列在 NVMe 上同样测不出并行度，这与 Linux 上 \Opt{iodepth} 的坑同构。
缓存侧的对应物是系统缓存与预取机制，第二遍读照样显著变快，
处理方式与第 2 节一致：把冷/热写成两个显式口径。

性能计数器的 \Tp{PhysicalDisk} 一组（\Tp{Avg. Disk sec/Read}、\Tp{Disk Queue Length}、
\Tp{Disk Bytes/sec}）是\emph{派生均值}，看不到尾延迟，而 ``Bytes/sec'' 是块层口径、
看不见缓存命中。要分布就走 ETW：用 Windows Performance Recorder、\Tp{xperf} 或
\Tp{logman} 采磁盘与文件 I/O 的提供程序，拿到的是逐请求耗时，可以自己摊成分位数；
Sysinternals 的 \Tp{diskmon} 适合快速看每卷速率。内核级 ETW 会话通常需要管理员权限，
符号解析需要 PDB 路径（三平台权限对照见 附录~\ref{ch:pmapA-platform}）。
两个本地干扰项要单独排除：筛选器驱动（防病毒、备份、终端管控）会插进 I/O 路径，
文件系统日志与索引也会在计时段里产生写；判据还是差分——
在能排除实时防护的环境里各跑一次，把 ``谁在路径上'' 写进结论。

\subsection{macOS：\Tp{fs\_usage}、File Activity 与 FSEvents 的边界}
\label{sub:pm08-mac}

macOS 的统一缓冲缓存受内存压力回收影响，缓存状态随整机负载变化；
\Tp{purge}（需要特权）释放文件缓存，作用与 Linux 的 \Tp{drop\_caches} 相当，
失真是同一条：整机一起变冷，共享机器上不该拿来做对照。这里的 \Tp{iostat} 字段与读法
\textbf{与 Linux 不同}（没有 \Tp{await} 与 \Tp{\%util} 这一族，速率与每请求字节
按它自己的定义给出），不要把 Linux 的判读规则搬过来。逐调用延迟看 \Tp{fs\_usage}
（需要特权，按进程给出每次系统调用的耗时）与 Instruments 的 File Activity 模板；
能耗与磁盘空闲状态看 \Tp{powermetrics}（权限模型与工具可得性见
附录~\ref{ch:pmapA-platform}）。

\Tp{FSEvents} 是\emph{目录级变更通知服务}：它会合并事件、延迟投递、只保证最终一致，
还会被索引与快照类背景活动干扰。用它测文件写入性能是口径错误——
你测的是通知链路的合并与投递策略，不是文件系统也不是设备。
APFS 的克隆与本地快照会让复制大文件几乎不产生设备写，
观察方法与 4.2 小节的写放大相同：比较应用侧字节与设备侧字节的比。

\section{排障与速查}
\label{sec:pm08-trouble}

\begin{longtable}{>{\raggedright\arraybackslash}p{0.22\textwidth} >{\raggedright\arraybackslash}p{0.26\textwidth} >{\raggedright\arraybackslash}p{0.29\textwidth}}
\caption{I/O 测量按症状排障表（先查证据，再谈结论）}\label{tab:pm08-io}\\
\toprule
症状 & 先看哪个计数器/证据 & 常见误判 \\
\midrule
\endfirsthead
\toprule
症状 & 先看哪个计数器/证据 & 常见误判 \\
\midrule
\endhead
\bottomrule
\endfoot
第二次比第一次快很多 & \Tp{iostat} 在该窗口有无 \Tp{r/s}；\Fn{mincore} 驻留比例 & 归因给补丁，其实是命中缓存 \\
顺序读带宽异常高 & \Tp{read{\Bk}\_ahead{\Bk}\_kb}、\Opt{direct} 是否为 $0$ & 把内核预取当设备能力 \\
\Tp{\%util} 接近 $100$ & 加深 \Opt{iodepth} 做差分，看 \MThr{} 是否还在涨 & 当成设备已饱和 \\
\Tp{iostat} 正常但应用慢 & 应用侧 \MWall{} 与 \Mx{clat} 的差；线程是否在等待 & 磁盘没问题当成链路没问题 \\
随机小读吞吐上不去 & \Opt{iodepth}、\Opt{numjobs}、\Tp{nr{\Bk}\_requests}、硬件队列 & 把浅队列当设备上限 \\
写数很好但掉电丢数据 & \Tp{write{\Bk}\_cache}、易失缓存、有无 \Opt{fsync} & 写延迟落实为缓存吸收 \\
偶发极慢的 \Fn{fsync} & 报分布；\Tp{dirty{\Bk}\_ratio} 与文件系统提交参数 & 用均值报同步很便宜 \\
换设备后随机读变快 & \Tp{optimal{\Bk}\_io{\Bk}\_size}、\Tp{logical{\Bk}\_block{\Bk}\_size} & 归因给介质，其实是粒度与预读 \\
容器里数字不同 & \Tp{findmnt -T}、是否 overlay/thin、配额与限流 & 默认同一块设备 \\
带宽很高但请求很慢 & \Tp{ss -ti} 的往返与队列、\Tp{softnet{\Bk}\_stat} & 用 \MThr{} 代替 \MLat{} \\
\Tp{nstat} 计数看着在涨 & 两次快照差分（它是开机累计值） & 把累计值当窗口增量 \\
fio 结果与手册不符 & \Opt{showcmd} 打印的实际生效参数 & 相信默认值，或混用同步与异步口径 \\
\Tp{tps} 对不上调用数 & \Tp{rrqm/s}、\Tp{wrqm/s} 的合并计数 & 把合并后请求数当系统调用数 \\
\end{longtable}

\begin{bestpractice}[title={把 I/O 的账与 CPU 的账分开测，但在同一次运行里对齐}]
\begin{enumerate}
\item 先定窗口：给运行打明确的起止时间戳，所有采集器覆盖同一窗口，
  窗口外的数据不参与结论。
\item 应用侧同时埋 \MWall{} 与 \MCpu{} 两条（\MWall{} 是墙钟耗时，\MCpu{} 是进程
  消耗的 CPU 时间；I/O 等待下两者会明显分家，来源见 第~\ref{ch:pm02-clock} 章）。
  \MCpu{} 低而 \MWall{} 高，说明时间花在等待而不是计算。
\item 同一窗口内取块层证据（\Tp{iostat} 或 \Tp{fio} 的 \Mx{clat}），
  另取一次 CPU 侧证据（采样与计数见 第~\ref{ch:pm04-sampling} 章 与 第~\ref{ch:pm05-events} 章），
  时间戳对齐之后，``慢在等待'' 与 ``慢在拷贝/序列化'' 才分得开。
\item 报三个量为一组：\MThr{}、\MLat{} 的中位数、\MLat{} 的上尾分位数，
  缺任何一个都可能被另一种机制解释（判据见 第~\ref{ch:pm03-stats} 章）。
\item 环境指纹进结论：设备与固件、文件系统与挂载选项、调度器与队列深度、
  预读与脏页参数、内核版本（第 2 节清单的输出整段留档）。
  换机器就重建基线，见 第~\ref{ch:pm13-ci} 章；线上采样预算见 第~\ref{ch:pm14-production} 章。
\end{enumerate}
\end{bestpractice}

\begin{guideline}[title={五问收口：一份 I/O 结论的最小体检}]
\begin{enumerate}
\item 这次运行里页缓存、预读、设备缓存各有没有参与？证据是哪一条？
\item 报的是哪一段延迟——应用、块层还是设备？工具字段与你的说法一致吗？
\item 设备、文件系统、调度器、队列深度是自己确认的还是抄来的？
\item 窗口对齐了吗？不同工具的数来自同一次运行的同一段吗？
\item 有分布吗？上尾分位数后面的样本量够支撑这句话吗？
\end{enumerate}
三平台工具与字段的完整并列见 附录~\ref{ch:pmapA-platform}，
指标与命令行参数速查见 附录~\ref{ch:pmapB-params}，按症状的术语排障见 附录~\ref{ch:pmapC-glossary}。
\end{guideline}
\part{并发、编译器与运行时}

\chapter{并发下的测量：伸缩性、伪共享与占用率歧义}
\label{ch:pm09-concurrency}

单线程世界里，\MWall 与 \MCpu 只差一个不超过 1 的系数，占用率是一个无害的旁证。线程数一过 1，
这两本账就分家：CPU 时间可以被自旋买满，墙钟时间却被一条缓存行卡住；吞吐还能涨，尾延迟先崩。
本章把并发测量拆成可操作的四步——认清占用率的三种歧义形态（第~\ref{sec:pm09-ambiguity} 节），
量化伸缩性曲线的形状与它的三个拐点（第~\ref{sec:pm09-scaling} 节），
分别处理伪共享（第~\ref{sec:pm09-fs} 节）与同步原语的测量面（第~\ref{sec:pm09-sync} 节），
收在调度迁移与异步背压（第~\ref{sec:pm09-sched} 节），最后落在一张按顺序走的排障表
（第~\ref{sec:pm09-trouble} 节）上。

与邻章的接口很明确：\textbf{同一线程内}的数据布局与自动向量化归第~\ref{ch:pm06-micro} 章，
本章只谈\textbf{跨线程}共享的行、锁与时钟域；重复次数、置信区间与最小可检出差异归第~\ref{ch:pm03-stats} 章，
本章只把分位数与 IQR 当作争用的第二证据；PMU 事件与 TopDown 分层归第~\ref{ch:pm05-events} 章，
互连带宽与 NUMA 落点归第~\ref{ch:pm12-bandwidth} 章。工具用法不在本册重讲：ThreadSanitizer（下称 TSan）
与 Helgrind 的命令、报告解读与抑制文件回指《内存泄漏检测与防御》
与《LSP 与代码质量工具》第~12 章的 sanitizer 门禁；CPU 隔离、\Tp{isolcpus}、\Tp{nohz\_full}
这类实时性手段留给《嵌入式与硬实时》。本章只保留它们的另一个身份：被测系统的一部分，以及一种开销。

\section{占用率歧义：\MWall、\MCpu 与 \MThr 在并发下分家}
\label{sec:pm09-ambiguity}

先把四个口径钉住，后文不再换词：\MWall 是墙钟时间（按下回车到结果落地的物理时间），
\MCpu 是 CPU 时间（所有核为该任务工作的时间之和），\MThr 是吞吐（单位时间的完成量），
\MLat 是延迟（单个请求的耗时，本册一律用分位数表述，如 \Mx{p99}）。
并发下两者的关系只剩一句话：\MCpu 等于 \MWall 乘以这段时间里平均占用的核数。
于是``占用率''这个数丢掉了最有用的那一半信息——它不知道有几个核在干活，更不知道干的是活还是自旋。

\begin{table}[htbp]
\centering
\small
\begin{tabular}{>{\raggedright\arraybackslash}p{0.17\textwidth}
                >{\raggedright\arraybackslash}p{0.20\textwidth}
                >{\raggedright\arraybackslash}p{0.42\textwidth}}
\toprule
形态 & 数字表现 & 第一证据与解释 \\
\midrule
一、\MCpu{} 高但更慢 & 占用率趋近 $100\%$，\MWall{} 反而变长 &
  \Mx{IPC} 下降、后端停顿占比上升、\Fn{futex} 进入次数接近零：钱花在自旋与重试上，进度为零 \\[3pt]
二、\MWall{} 短但烧核 & \MWall{} 减半，\MCpu{} 翻倍 &
  并行度换时间的正常账。它不是回退，但它是成本：容量、能耗与共存进程的抢占都按 \MCpu{} 结算 \\[3pt]
三、\MThr{} 升而 \Mx{p99} 崩 & 平均吞吐改善，尾部分位数恶化一个量级 &
  工作被搬进队列或搬进锁的等待队列。只看 \MThr{} 的报表会把排队当成优化；
  要看队列长度分位数与唤醒次数（第~\ref{sub:pm09-backpressure} 节） \\
\bottomrule
\end{tabular}
\caption{并发下三种口径分家的形态、判读与对应的第一证据（判读启发式）。}
\label{tab:pm09-forms}
\end{table}

\begin{lstlisting}[style=outstyle, label={lst:pm09-forms},
caption={同一负载扫四档线程数的示意输出与格式示例（非本仓库实测值）。}]
workload: 8 Mi items, 4 shards. strong scaling, T = threads
  T  Wall_ms  CPU_ms  busy%  Mops/s  p50us  p99us  spin%
  1    412.0   411.4   99.9    2.46      31     41     0
  4    121.7   480.2   99.0    8.34      44     63     0
  8     71.3   566.9   99.5   14.20      58    118     6
 16    168.9  2610.3  100.0    5.96     210   1043    71
# T=8  form 2: Wall halved, CPU up. Bill it as CPU_ms.
# T=16 form 1: Wall longer than at T=8, CPU 4.6x -> spin
# p99 rises while p50 barely moves: work is queued, not done
# spin% = share of scheduled time spent retrying; not a kernel metric
\end{lstlisting}

\begin{note}[title={为什么单线程基线在并发章里仍然必需}]
三件事只有 $T=1$ 的那趟才能提供。其一，加速比与伸缩模型的自变量都以单线程吞吐为分母；
没有它，``16 核跑满''既可能是成就也可能是灾难。其二，串行段的代价要能被减出来：
把共享结构换成线程本地结构再跑一趟 $T=1$，差值就是同步本身的固定税
（一次自增、一次可能被编译器折叠掉的栅栏），它决定后面所有争用数字的零点。
其三，缓存与频率的归一化基准：跨运行比较吞吐时把数字除以每周期工作量（如 ops per Mcycles），
而每周期工作量只能从干净的单线程趟次里取。
\end{note}

\subsection{跨平台对照：占用率不是同一个量}

Linux 侧最省事的读法是 \Tp{/proc/stat} 的时间字段与 cgroup 的 \Tp{cpu.stat}；
注意 \Tp{iowait} 在内核口径里算 idle，虚拟机上的 \Tp{steal} 是被宿主拿走的时间，两者都会让``忙''失真。
Windows 侧的性能计数器里，\Tp{\%{\Bk}Processor{\Bk}Time}、\Tp{\%{\Bk}Idle{\Bk}Time} 与
\Tp{\%{\Bk}Utility{\Bk}Processor} 不是同一件事，DPC 与中断另有各自的计数器；
线程级 CPU 时间走 \Fn{GetThreadTimes}，单位是 100\,ns；就绪队列与切换形态要从 ETW（\Tp{xperf}）取。
macOS 上 \Tp{powermetrics} 需要提权，并同时给频率与每核忙闲；线程时间走 \Fn{thread\_info}，
单位是 mach absolute time，换算比由 \Fn{mach\_timebase\_info} 现场取。
三个平台的``每核 $100\%$''都不等于``每核满频的周期供给''，所以结论必须连带实际频率一起记
（Linux 用 \Tp{turbostat} 或 \Tp{scaling{\Bk}\_cur{\Bk}\_freq}，另两平台用各自计数器的频率实例）。
计数器名与权限模型随版本变，先在目标机上确认再引用，别照抄本文。

\section{伸缩性曲线：定义、串行比例的来源与掉头的位置}
\label{sec:pm09-scaling}

\subsection{从加速比到三个模型}

记 $T_1$ 为单线程 \MWall，$T_N$ 为 $N$ 线程 \MWall，加速比 $S(N) = T_1 / T_N$。三个常用模型解释力不同：

\begin{itemize}
\item Amdahl：$S(N) = 1 / (\sigma + (1 - \sigma) / N)$。只承认串行份额 $\sigma$，
      曲线单调上升并渐近于 $1/\sigma$，因此结构上无法解释``线程更多反而更慢''。
\item Gustafson（弱伸缩）：固定\textbf{每线程工作量}而非总工作量，此时 $S(N) \approx N$ 近似自证。
      它回答``机器喂不喂得饱''，不回答``算法拆不拆得开''；
      把弱伸缩的数字画进强伸缩的坐标是最常见的自欺。
\item USL（通用伸缩律）：$S(N) = N / (1 + \sigma (N - 1) + \kappa N (N - 1))$。
      $\sigma$ 是串行化份额，$\kappa$ 是缓存一致性份额。最后一项随线程对数二次增长，
      正好把 Amdahl 画不出的向后弯（backwards bend）画出来，
      也让峰值位置成为可报告的量而不是拟合的副产品。
\end{itemize}

拟合先把式子线性化：$N / S(N) - 1 = \sigma (N - 1) + \kappa N (N - 1)$，
把 $(N-1)$ 与 $N(N-1)$ 当两个自变量做一次最小二乘即可，不需要非线性求解器。
$\kappa$ 项的物理来源就是下一节的伪共享与共享计数器——
这也是为什么``把曲线弯头归因于代码里的串行段''通常是错的。

\begin{warning}[title={两点定不出三个参数，也别把串行份额当成代码属性}]
只用 $N=1$ 与 $N=8$ 两档数据，Amdahl 与 USL 都能``完美解释''同一批观测，
$\kappa$ 会被随机噪声决定，置信区间宽到失去意义。线程数取 $1,2,3,4,6,8,12,16,24,32$
这类含非 2 幂的混合点，峰值邻域加密，每档至少 10 次重复并报中位数与 IQR，
并且\textbf{看残差}：残差随 $N$ 单调增大就说明模型已经不适配这台机器，该分段而不是继续调参。
另外，``串行比例''不等于代码里的临界区：全局单例、内存控制器、日志刷盘、连接池上限、
乃至整机共享的功耗预算，都会以同样形状出现在 $S(N)$ 上。
\end{warning}

\subsection{扫描骨架与曲线形状}

\begin{lstlisting}[style=cppstyle, label={lst:pm09-scan},
caption={线程数扫描骨架：每档打印 \MWall、\MCpu、\MThr{} 与 \Mx{p99}；数值列由读者填。}]
struct Row { int n; double wall, cpu, thr, p50, p99; };

// strong scaling: total work fixed. Flip to weak scaling and grow
// the work with n -- then report which one you actually ran.
std::vector<Row> sweep(const std::vector<int>& ns, int r) {
  std::vector<Row> rows;
  for (int n : ns) {
    std::vector<double> wall, cpu, lat;
    for (int k = 0; k < r; ++k) {
      auto parts = split_evenly(total_items, n);
      Pool pool(n);                   // parked on a start barrier
      const auto c0 = cpu_time();     // a process-wide cpu clock
      const auto t0 = steady_clock::now();
      for (auto& p : parts) pool.submit([q = p]{ run_shard(q); });
      pool.join();
      wall.push_back(ms(steady_clock::now() - t0));
      cpu.push_back(c0.since(cpu_time()));
      append(lat, per_op_latencies());   // keep every sample
    }
    rows.push_back({n, median(wall), median(cpu),
      units() / median(wall), pct(lat, 0.50), pct(lat, 0.99)});
  }
  return rows;   // caller prints median + IQR + min/max
}
\end{lstlisting}

\begin{lstlisting}[style=outstyle, label={lst:pm09-curve},
caption={同一二进制、每档 12 次重复的曲线与拟合片段：ASCII 形状与数值均为示意输出/格式示例（非本仓库实测值）。}]
 S(N)
   16|                       x   N* ~ 11
   12|                   x
    8|           x
    4|        x
    1|  x
      +---+---+---+---+---+---+---+---+-- N
        2   4   6   8  10  12  16  24  32
       linear knee  bend  plateau regress
# fit: sigma 0.018 (0.015..0.021)
#      kappa 0.0042 (0.0034..0.0050)  boot 10k
# resid +0.05 up to N=12, +0.21 at N=32 -> regime break
\end{lstlisting}

\subsection{曲线掉头的三类边界}

\begin{table}[htbp]
\centering
\small
\begin{tabular}{>{\raggedright\arraybackslash}p{0.15\textwidth}
                >{\raggedright\arraybackslash}p{0.24\textwidth}
                >{\raggedright\arraybackslash}p{0.40\textwidth}}
\toprule
边界 & 触发条件 & 怎么确认它在你机器上的位置 \\
\midrule
SMT 配对 & 线程数超过物理核数 &
  \Fn{lscpu} \Tp{-e} 的 \Tp{CORE}、\Tp{SOCKET} 列，或拓扑 sysfs 目录；
  跨过它时 \MThr{} 通常只涨几个百分点到几十个百分点，而每线程 \Mx{IPC} 明显下降 \\
一致性域 & 协作线程跨 LLC 分片或一致性目录 &
  \Fn{lstopo} 或 \Fn{numactl} 的硬件视图；\Tp{perf c2c} 报告里远程命中份额开始不为零 \\
内存节点 & 分配落在远端节点 &
  \Fn{numastat} 的 \Tp{-m} 与 \Tp{-p}；这一档常表现为\textbf{台阶}而非平滑下降，带宽细节归第~\ref{ch:pm12-bandwidth} 章 \\
\bottomrule
\end{tabular}
\caption{三个拐点对应的硬件边界与确认方法（Linux 为主；Windows 可用 \Fn{coreinfo}
或任务管理器的关联性视图，macOS 用 \Fn{sysctl} 的拓扑项，名称需现场确认）。}
\label{tab:pm09-boundary}
\end{table}

USL 是平滑模型，机器不是。上面三类边界都会让残差分段，
所以报告里要么给分段拟合，要么明确写``本曲线只在 $N \le 12$ 区间内可信''。

\section{伪共享：缓存行的共享与状态共享的差别}
\label{sec:pm09-fs}

\subsection{定义、行大小与代价量级}

两个逻辑上无关的变量落在同一条缓存行上时，一个核的写会让其他核的副本失效，
一致性协议被迫反复迁移该行所有权。这叫伪共享：数据没被共享，行被共享了。
代价取决于球传多远：同一 LLC 分片内、跨分片、跨节点是三个量级。本册不给点值——
那是把工具默认值当结论——只给口径：``本地通常数百周期以内，跨节点可到千余周期''，
并要求你自己测出这三个数：同一份紧凑数组与同一份填充数组各跑一趟 $T=N$，
再各跑一趟 $T=1$ 减掉固定税，差值除以自增次数即得每次行传输的周期估计。

行大小不要硬编码。\CPP{} 17 起有常量 \Tp{hardware{\Bk}\_destructive{\Bk}\_interference{\Bk}\_size}
（命名空间 \Tp{std::}），但标准库与编译器的提供时点和取值并不一致，
未提供时会给出诊断或回落到实现默认，先在自己的构建里确认它等于几。
运行时查询走 \Fn{sysconf} 与 \Tp{\_SC{\Bk}\_LEVEL1{\Bk}\_DCACHE{\Bk}\_LINESIZE}，
部分平台返回 0，或返回 L1 粒度而非实际干扰粒度。
常见架构是 64 字节，个别架构行更宽，所以``写死 64 并到处复用''是一个会一路跟你到移植阶段的错误。

\begin{lstlisting}[style=diffstyle, label={lst:pm09-pad},
caption={两种修法的形式示例（差异格式示意，非本仓库实测）。}]
@@ per-thread counters, tightly packed
-struct Counters { std::atomic<long> v[MAXT]; };
+struct Counters {
+  struct Slot { alignas(CLINE) atomic<long> v; } s[MAXT];
+};                  // CLINE from the query above, not a literal 64
@@ when padding is not enough: amortise the line transfer
+  long local = 0;    // thread-private, no atomicity required
+  for (auto it = shard.begin(); it != shard.end(); ++it)
+    local += (*it > 0);
+  if (local) drain_to_shared(local);  // per shard, not per item
\end{lstlisting}

\subsection{怎么观察：\Tp{perf c2c} 与差分}

\Tp{perf c2c} 在支持的 PMU 上按缓存行聚合访问与远程命中信息。事件名与可用性随 CPU 型号与内核版本变化，
Intel 与 AMD 走不同记录通路，ARM 侧并非所有实现都提供同类事件，所以第一步是确认可用性：
record 一趟再看报告是否为空，为空就退回差分法。差分法是自足的：同一份源码、两个构建（紧凑与填充），
同一负载、同样重复次数、同样绑核策略，吞吐差就是行争用的净贡献。
注意 \Tp{perf c2c} 本身开记录级采样，是有开销的观测，不要把它和计时跑混成同一趟。

\begin{lstlisting}[style=outstyle, label={lst:pm09-c2c},
caption={\Tp{perf c2c report} 关键片段与差分跑的示意输出/格式示例（非本仓库实测值）。}]
$ perf c2c record -o c2c.data -- ./build/app --threads 16
$ perf c2c report --stats -LLC,8
C0 0x7f2a1c4b8040   hitm    total    lcl    rmt
-----------------  ------  ------  -----  -----
  0x7f2a1c4b8040  44807   44821     14     47
  ^ ONE 64-byte line, 16 writers, offsets 0..63
  [.] Counters::s+0x0  worker_loop  hitm% 99.6
# offsets spread over the line, written by different threads
#   -> packed per-thread state (false sharing)
# one offset, many readers -> real sharing; protocol likes reads
$ ./build/app --counters=compact --threads 16   # 12 reps
 thr_mops 5.96  p99_us 1043  hitm_per_kop 3.9e3
$ ./build/app --counters=padded --threads 16    # else identical
 thr_mops 15.9  p99_us  121  hitm_per_kop 4.1e0
\end{lstlisting}

判读分三步：先确认热点集中在一行（地址列只剩一个条目），再确认该行内\textbf{多个偏移被不同线程写}，
最后确认偏移跨度接近整行。三条同时成立才动布局；只成立一条时更可能是算法层的共享。

\begin{warning}[title={padding 修不好算法层的共享}]
填充只解决``两个无关变量挤在同一行''。如果两个线程确实在更新同一个计数器、同一个队列头指针、
同一条空闲链，那么无论怎么对齐行都要迁移：此时正确动作是消除共享本身——线程本地累加后归并、
按线程 id 低几位分片、把决策移到批量提交点。检验就在同一组差分里：填充后 \MThr{} 几乎不动，
而远程命中份额仍然集中在一行，说明面对的是真共享；继续加填充只会撑大结构体、
弄坏缓存足迹与向量化条件。
另一处常见误用是以为 \Tp{alignas(64)} 会自动留出整行空隙：它只保证该对象起始地址对齐，
不保证与邻居之间有空行。放进数组或结构体时，要给槽位类型本身加对齐并让尺寸向上取整，
或显式补一个填充成员，否则第 $k$ 个槽位仍与第 $k+1$ 个共享行的尾巴。
修完必须回测单线程基准（第~\ref{ch:pm06-micro} 章的范围），确认没有把成本从跨线程账挪到单线程账。
\end{warning}

\section{锁、原子与内存序的测量面}
\label{sec:pm09-sync}

\subsection{竞争度：等待份额与重试次数是两件事}

\textbf{等待份额}回答``我有多少 CPU 时间花在没推进的循环里''，
证据是 \Tp{schedstat} 的等待时间、后端停顿占比、自研的自旋周期估计。
\textbf{重试次数}回答``我失败了多少次才成功''，证据是 CAS 失败计数、\Fn{futex} 进入与唤醒次数、
硬件事务的 abort 计数。两者正交：高频重试但每次极快的自旋锁可以有很小的等待份额；
从不重试、只睡一次的重型互斥锁可以有很大的等待份额。只报其中一个，等于把争用的两个面砍掉一个。
若要在代码里指名道姓，比较并交换的接口是 \Fn{atomic{\Bk}\_compare{\Bk}\_exchange{\Bk}\_weak}
与 \Fn{compare{\Bk}\_exchange{\Bk}\_strong}：弱形式允许虚假失败，重试计数要为它单独分栏。

\begin{table}[htbp]
\centering
\small
\begin{tabular}{>{\raggedright\arraybackslash}p{0.19\textwidth}
                >{\raggedright\arraybackslash}p{0.13\textwidth}
                >{\raggedright\arraybackslash}p{0.14\textwidth}
                >{\raggedright\arraybackslash}p{0.33\textwidth}}
\toprule
操作 & 无争用 & 争用 & 为什么值得单独测 \\
\midrule
原子读/写（同行单写者） & 一两个周期 & --- & 行留在本地，常作零点 \\
CAS 比较交换 & 数十周期 & 百至数百周期 & 天生需要行独占权 \\
\Tp{seq\_cst} 存储（x86） & 数十周期 & --- & 编成带全屏障的形式 \\
互斥锁快路径 & 数十周期 & --- & 全在用户态，不进内核 \\
互斥锁睡眠路径 & --- & 微秒级 & 陷入内核、唤醒、可能被迁移 \\
自旋等待 & 数周期/次 & 百周期/次 & 买满占用率而进度为零 \\
\bottomrule
\end{tabular}
\caption{同步操作成本阶梯：只给量级与顺序关系，不给点值。每一格都要在目标机上填数
（方法见第~\ref{sub:pm09-ladder} 节）；CPU、内核版本与 C 库版本会让不同机器上的数字互相不可比。}
\label{tab:pm09-ladder}
\end{table}

\subsection{填出自己的成本阶梯}
\label{sub:pm09-ladder}

三步法。第一步，$T=1$ 跑一次空临界区与一次真实临界区，差值是无争用固定税；
第二步，$T=N$ 让全部线程抢同一把锁，减去第一步的税再除以临界区次数，得到争用路径的每操作成本；
第三步，同一负载换成``每线程本地累加、末尾归并''，它与第二步的差就是可回收争用的上限——
这个上限决定你后续所有同步优化的天花板。循环体必须留下不可消除的副作用
（累加进一个外部可见量，或用反死代码消除的挡板），否则编译器会把整段同步折掉；
这是并发微基准最常见的假结果来源，也是``优化级别一变结论就翻''的来源。

\subsection{自旋在计数器下的指纹}

\begin{lstlisting}[style=outstyle, label={lst:pm09-spin},
caption={\Fn{perf} \Tp{stat} 与自研报告的示意输出/格式示例（非本仓库实测值）；事件名先用 \Fn{perf} \Tp{list} 确认。}]
$ perf stat -e task-clock,cycles,instructions,\
> stalled-cycles-backend,context-switches,cpu-migrations,\
> syscalls:sys_enter_futex -- ./build/app --threads 16

  158.3 s task-clock      #   9.9 CPUs utilized
  4.02e10 cycles          #   2.5 GHz (record the real freq)
  1.97e10 instructions    #   0.49 IPC
  2.31e10 stalled-cycles-backend  #  57% of cycles
    41932 context-switches
        7 syscalls:sys_enter_futex   # they spin, no sleep
# form 1 fingerprint: busy 100%, IPC down, stalls up
$ ./build/app --threads 16 --spin 0
    9.4 s task-clock  # IPC rises, futex rises: waiting moved
\end{lstlisting}

同一段程序，自旋与睡眠在数据上不是``更快与更慢''，而是两本不同的账：前者把成本记进 \MCpu{}
与停顿计数，后者把它记进等待时间与切换计数。把两种实现的 \MWall{} 直接相减当作收益，
会漏掉供给侧的排队成本。

\subsection{内存序：同一份源码在两个 ISA 上不是同一个被测系统}

\begin{itemize}
\item \textbf{重排与用 \Tp{volatile} 当同步}。非原子标志上的读可以被提到循环外，
      于是 \CPP{} 层面的停止条件永远不会被再次检查；\Tp{volatile} 只约束一部分优化，
      既不给原子性也不给顺序。这类 bug 在数据里长得像``某线程占用率恒为 $100\%$ 而吞吐不降''，
      很容易被当成性能问题去调参。
\item \textbf{顺序的代价按平台分布}。在强序平台上 \Tp{acquire} 与 \Tp{release}
      常常退化为普通装载与存储，只有全屏障形式付额外代价；弱序平台上同一行源码会编出
      带获取/释放语义的指令或显式栅栏。因此跨架构比较原子操作开销之前先比反汇编，不要先比时钟。
\item \textbf{怎么确认自己付了哪条指令}。用 \Tp{-S} 或 \Fn{objdump} 把栅栏与 \Fn{lock} 前缀
      指令数出来，写进测量元数据。这条纪律与第~\ref{ch:pm10-compiler} 章的``改一行代码前先确认改了二进制''
      是同一件事，只是并发场景下它还随目标架构改变。
\item \textbf{构建配置会改变被测对象}。优化级别、LTO 与内联决策改变临界区长度，
      临界区长度改变争用概率与持锁时间的分布。伸缩性对比必须锁定编译命令行与产物标识，
      把它们一起写进 CSV 头部。内联首先是 ODR 与链接语义而不是提速开关，
      这一口径与《代码风格参考手册》一致。
\end{itemize}

\subsection{TSan 是检测器，不是测速仪}
\label{sub:pm09-tsan}

TSan 动态维护 happens-before 关系：每次访存都要与影子状态交互，每个线程携带向量时钟，
锁与原子操作被拦截以维护同步边，插桩构建里的内联与部分优化还会被抑制。
开销是这四项叠加的结果，量级只能写成区间：CPU 时间数倍到十余倍、RSS 数倍，
并明确它由访存密度与同步事件数决定，与算法复杂度无关——
所以不存在可移植的 TSan 开销常数，只有你在自己负载上测出的那一组数。更关键的是三条不可比性：

\begin{enumerate}
\item 它改变线程交错：被拦截的同步原语会让被运气掩盖的争用出现，也会让原本的高争用路径消失。
      它给出的耗时数是插桩系统的伸缩性，不是被测程序的。
\item 向量时钟的合并成本随同步事件增长，会给一致性份额 $\kappa$
      掺进不属于被测程序的分量，于是拟合参数被工具污染。
\item 覆盖依赖实际交错。没报 race 不等于没有 race，``TSan 干净''是一次执行的证据，
      不是一个程序的证明。
\end{enumerate}

\begin{bestpractice}[title={双二进制原则}]
把\emph{计时构建}（\Tp{-O2}，无插桩）与\emph{检测构建}
（\Tp{-O1}、\Tp{-g}、\Tp{-fsanitize=thread}，注意它与 ASan 互斥）当成两个被测系统：
各出各的 CSV，共享同一份元数据（提交、编译器版本与标志、拓扑、调频策略、配额状态），
但\textbf{永不共享数据点}。TSan 的时间倍数与内存倍数只写进``检测器开销''一栏，
绝不进入加速比、拟合与回归门禁。需要复现同一交错时用确定性记录回放（如 \Lib{rr}），
别把 TSan 日志当时间线；回放自身的额外成本也要单独测。
用法、报告解读与抑制文件不在本册：见《内存泄漏检测与防御》与《LSP 与代码质量工具》第~12 章；
Helgrind 一类锁次序检查器开销再高一个量级，适合``跑一次就够''的场合。
\end{bestpractice}

\section{调度、迁移与异步背压}
\label{sec:pm09-sched}

\subsection{上下文切换计数的三种读法}

切换次数本身不是指标，它是分叉点。配合 \Tp{schedstat} 的三个字段（切换次数、累计运行时间、
累计等待时间）与自愿/非自愿之分，才有三种不同结论：切换多而等待短，是锁闪烁与唤醒风暴；
切换多且等待长，是 CPU 供给不足或被抢占（容器配额与虚拟机 \Tp{steal} 都会伪装成这一种）；
切换极少而 \MCpu{} 高，是线程在自旋，回到第~\ref{sub:pm09-ladder} 节的指纹即可确认。
读 \Tp{/proc} 里线程的 \Tp{utime} 与 \Tp{stime} 时记得它们是 tick 粒度，亚毫秒的并发段不能靠它计时；
线程级高精度计时用 \Fn{clock\_gettime} 配 \Tp{CLOCK{\Bk}\_THREAD{\Bk}\_CPUTIME{\Bk}\_ID}，
其可用性与分辨率同样要现场确认。

\begin{lstlisting}[style=shellstyle, label={lst:pm09-aff},
caption={拓扑、亲和性与迁移证据的命令骨架（格式示例，输出由读者在自己机器上得到）。}]
lscpu -e=CPU,CORE,SOCKET,NODE       # where the boundaries are
numactl -H                          # distance matrix
numactl --physcpubind=0-7 --localalloc -- ./build/app
taskset -cp 0-7 $(pgrep -f build/app)   # did the mask stick?
cat /proc/self/task/*/schedstat     # sw, run_ns, wait_ns
grep -E 'throttled' /sys/fs/cgroup/cpu.stat
\end{lstlisting}

线程内还可以用 \Fn{sched\_getcpu} 记录``这个线程实际跑在哪个 CPU''，
并把每线程的 CPU 集合与迁移次数打进输出。否则你以为在测算法，其实在测一次幸运的放置。

\subsection{把``绑了才复现''写进结论}

绑核消除迁移噪声、固定 SMT 配对与内存落点，是让曲线可重复的最省事办法；
但迁移常常正是被测现象——负载均衡、工作窃取、NUMA 感知分配都依赖它。
诚实口径是\textbf{两组都留}：一组绑定（稳定曲线、紧的置信区间），
一组默认调度器（真实抖动、更宽的分布），并写明哪一组是决策依据。
只报绑定那组，等于把复现性与生产表现混为一谈；只报默认那组，曲线形状会被邻居负载支配。
更彻底的隔离（\Tp{isolcpus}、\Tp{nohz\_full}、中断屏蔽，以及实时调度类与绑核策略的时序后果）
属于《嵌入式与硬实时》的范围，本册只提醒：这些开关会改变基线，跨这两类环境比较数字没有意义。

\subsection{异步与背压：队列长度才是指标}
\label{sub:pm09-backpressure}

生产与消费解耦之后，``变快''最常见的假象是把延迟搬进队列而不是消灭它。正确的指标组合是四条并列：
\MThr{}（完成率）、\MLat{} 的分位数（尤其 \Mx{p99}）、队列长度的分位数
（不是平均长度，平均值会把短暂挤满洗掉），以及丢弃与降级的计数口径
（必须说清一次事件是帧、任务还是请求，跨系统比较时这是最常见的不对等）。

\begin{compare}[title={无界队列的``变快''与两个必须并列的指标}]
\small
\begin{tabular}{>{\raggedright\arraybackslash}p{0.14\textwidth}
                >{\raggedright\arraybackslash}p{0.27\textwidth}
                >{\raggedright\arraybackslash}p{0.32\textwidth}}
\toprule
口径 & 无界队列（$T=8$） & 有界队列加显式降级 \\
\midrule
\MThr{} & 完成率上升：生产者不再阻塞，CPU 被填满 & 完成率略低：背压把过载挡在生产者一侧 \\
\MLat{} \Mx{p99} & 上升一个量级：等待被搬进队列 & 平稳：超预算的活被降级而不是排队 \\
队列 \Mx{p99} & 峰值可达分配上限，随后换页或失败 & 稳定在上界附近 \\
降级计数 & 没有这一项，过载被隐藏 & 有，而且它才是容量结论的依据 \\
\bottomrule
\end{tabular}
\end{compare}

\begin{lstlisting}[style=outstyle, label={lst:pm09-queue},
caption={异步流水线报告片段的示意输出/格式示例（非本仓库实测值）。}]
pipeline: capture -> decode(n=8) -> mux    60 s of load
 thr_fps  412.3     # mean over window; NOT the verdict
 lat_p50  8.4 ms  lat_p99 61.2 ms  lat_p999 188.0 ms
 queue_p50 2  queue_p99 47  queue_max 512 (bound)
 dropped 0   backpressure_events 129  degraded 12
# report lat_p99 and queue_p99 as a PAIR: throughput hides it
# queue_max == bound with dropped == 0 => saturated, not fine
\end{lstlisting}

两个细节决定这类数据能否使用。其一，采样要与负载同源：用同一时钟域取入队与出队时间戳，
跨线程只做差值、不做绝对对齐；时钟选择与 TSC 失步在第~\ref{ch:pm02-clock} 章。
其二，容器里的过载信号优先：先读 \Tp{cpu.stat} 的节流计数再谈队列形状，
因为配额节流会让``占用率 $100\%$、队列堆积、\MThr{} 断崖''三种症状同时出现，而它们的修法完全不同。

\section{排障与速查}
\label{sec:pm09-trouble}

\begin{longtable}{>{\raggedright\arraybackslash}p{0.05\textwidth}
                  >{\raggedright\arraybackslash}p{0.20\textwidth}
                  >{\raggedright\arraybackslash}p{0.25\textwidth}
                  >{\raggedright\arraybackslash}p{0.27\textwidth}}
\caption{加了线程反而慢：按顺序走的六条排查。顺序本身是方法——前两条排除``数字不成立''，后四条才谈代码。}\label{tab:pm09-trouble}\\
\toprule
序 & 检查 & 判据 & 若为真 \\
\midrule
\endfirsthead
\toprule
序 & 检查 & 判据 & 若为真 \\
\midrule
\endhead
\midrule
\endfoot
\bottomrule
\endlastfoot
1 & 配额与虚拟化 & \Tp{cpu.stat} 的 \Tp{nr{\Bk}\_throttled} 是否增长；\Tp{steal} 份额是否可观 &
  先解决供给再谈伸缩；提高配额或换裸机后重跑整条曲线 \\
2 & 基线是否同物 & 两个构建是否同提交、同标志、同拓扑；有没有混进插桩构建的数字 &
  重建对照组，执行双二进制原则（第~\ref{sub:pm09-tsan} 节） \\
3 & 是否在自旋 & \Mx{IPC} 下降、后端停顿上升、\Fn{futex} 接近零、\Tp{schedstat} 等待很短 &
  缩短临界区或改为睡眠；自旋预算按实测等待时间设定 \\
4 & 是否共享一行 & \Tp{perf c2c} 热点是否集中在一行且多偏移被多线程写 &
  填充或分片；填充后无变化则是算法层共享，回到归并与批量提交 \\
5 & 是否跨过边界 & 拐点是否与 SMT 配对、LLC 域、内存节点重合 &
  分段拟合或按拓扑放置线程与内存；带宽与落点见第~\ref{ch:pm12-bandwidth} 章 \\
6 & 是否把延迟搬进队列 & \MThr{} 升而 \Mx{p99} 与队列 \Mx{p99} 同升 &
  加背压与显式降级计数，用两个指标并列重新定义``变快'' \\
\end{longtable}

\begin{guideline}[title={并发测量的起手式}]
一趟 \MWall、一趟 \MCpu、每线程一份 \Mx{p99}，是任何并发结论的三件套；
占用率只是由此派生的第四个数字，不能单独支撑结论。线程数取含非 2 幂的一串，
每档至少 10 次重复并报 IQR；绑核与不绑核各跑一组；单线程基线与空临界区基线各跑一组；
检测器的数字进``开销''栏，不进曲线。
\end{guideline}

速查要点：加速比与两个份额（$\sigma$ 串行、$\kappa$ 一致性）用线性化式子拟合，残差有形状就分段；
伪共享先看行内偏移跨度、再看线程归属；争用用等待份额与重试次数两个正交指标描述；
TSan 的数字永远不进伸缩模型；容器里先读节流；
结论要能被``实际频率、拓扑、配额、二进制标识''这四行元数据复现。
需要三平台的工具、权限与事件对照，集中在附录 A。

\chapter{编译器与构建口径：优化级别、LTO、PGO 与内联的真实语义}
\label{ch:pm10-compiler}

前面几章把注意力放在``怎么读数''：时钟怎么选（见 第~\ref{ch:pm02-clock} 章）、样本够不够（见
第~\ref{ch:pm03-stats} 章）、指标回到代码的哪一行（见 第~\ref{ch:pm06-micro} 章）。这一章把镜头挪到读数的
\textbf{另一端}——那个被反复计时的二进制到底是怎么来的。同一段源码经过不同的优化级别、不同的链接策略、
不同的 profile 输入，产出的是\textbf{行为不同的程序}；而大多数``补丁有效''的结论，其实只是
``构建口径变了''的结论换了个名字。本章解决三个问题：第一，怎么把``测的是哪个二进制''钉成一个可以复述
的事实；第二，编译器在代码层面究竟给了哪些\textbf{可验证的}与\textbf{不可验证的}承诺（内联就是典型）；
第三，LTO、PGO、自动向量化这三类``跨翻译单元的决策''怎么观察、怎么记账。与后文的接口在
第~\ref{ch:pm13-ci} 章：构建口径决定基线能不能长期比对；与 第~\ref{ch:pm06-micro} 章 的接口在于，微架构指标只能
告诉你``结果如何''，本章的工具告诉你``编译器做了什么决定''。

\section{先把``测的是哪个二进制''钉死}
\label{sec:pm10-binary}

\subsection{构建类型不等于优化级别}

\Tp{CMAKE\_{\Bk}BUILD\_{\Bk}TYPE} 的四个常用取值各自绑定了一组编译选项，但它绑定的是\textbf{默认值}，
而默认值会随平台、随模块、随你覆盖的变量而改变。把这张表当作``起手式''而不是``结论''更合适：

\begin{itemize}
  \item \Tp{Debug}：通常是 \Tp{-O0 -g}，保留断言。它的用途是调试，不是测量。
  \item \Tp{Release}：通常是 \Tp{-O3 -DNDEBUG}（GCC/Clang 家族），断言被关掉。
  \item \Tp{RelWithDebInfo}：通常是 \Tp{-O2 -g -DNDEBUG}，做性能剖析时最常用的一档，因为它既有
        优化又有可用的符号。
  \item \Tp{MinSizeRel}：通常是 \Tp{-Os -DNDEBUG}，为体积优化，热路径形状与前几档明显不同。
\end{itemize}

关键区别是 \Tp{-DNDEBUG}：\Fn{assert} 在展开时依赖 \Tp{NDEBUG} 宏，一旦定义，\Fn{assert} 变成空操作。
也就是说，\textbf{带断言与不带断言跑出的两个数，比较的是两段不同的代码}——断言里往往夹着循环校验、
容器快照、临时分配，这些都会进 \MWall{}（墙钟时间）也进 \MCpu{}（CPU 时间）。想知道断言值多少，唯一
诚实的办法是 \Tp{-UNDEBUG} 与 \Tp{-DNDEBUG} 各构建一次、在同一个基准上交替跑（交替对照的口径见
第~\ref{ch:pm03-stats} 章）。

\subsection{优化级别改变的是``哪些承诺被允许打破''}

GCC/Clang 一族的 \Tp{-O} 族大致是这样一条阶梯；MSVC 的 \Tp{/Od}、\Tp{/O1}、\Tp{/O2}、\Tp{/Ox} 与它
不是一一对应关系，只能算方向相近（逐档语义要查你所用版本的文档，不要拿别人的表当依据）：

\begin{itemize}
  \item \Tp{-O0}：不做优化，目标是调试可用。函数调用、临时量全部落到内存。
  \item \Tp{-Og}：以``不破坏调试体验''为约束的轻量优化。它是最容易被误当成``接近发布''的一档——
        实际上它的内联、循环形状与 \Tp{-O2} 差别很大，用它做 A/B 的基线常常给出错误结论。
  \item \Tp{-O1}/\Tp{-O2}：常规发布档。\Tp{-O2} 是``优化''的默认含义，很多其它开关（例如若干内联阈值）
        以它为参照。
  \item \Tp{-O3}：在 \Tp{-O2} 之上打开更激进的循环变换、向量化与内联倾向。它\textbf{可能}让某段代码
        变慢（更大的内联体挤爆指令缓存、过度展开拉长热路径），所以``生产一定用 \Tp{-O3}''不是结论。
  \item \Tp{-Os}：以体积优先，常常在嵌入式或 I-cache 受限的场景反而更快。
\end{itemize}

\Tp{\_\_OPTIMIZE\_\_} 是 GCC/Clang 一族在 \Tp{-O} 大于 0 时定义的宏，\textbf{没有}统一数值口径，别把
它当``级别几''读；MSVC 侧并没有同名宏，只能靠 \Tp{NDEBUG} 或项目自定义宏。想要一个可靠的判断，就在
产物上确认（见本节末尾的清单），不要在源码里靠宏猜。

\begin{warning}[title={在 Debug 上测出来的数字什么也说明不了}]
一段 \Tp{-O0} 构建的剖析结果，能告诉你的只有``这段代码被执行了多少次、大概在哪''，不能告诉你任何关于
发布版本的性能结论。\Tp{-O0} 下没有内联、没有向量化、每个临时量都落内存，于是采样剖析会把大量时间归到
本不该存在的取数存数与函数序言上；反过来，\Tp{-O2} 下真实存在的瓶颈（一次向量化后的归约、一个被内联
进调用点的循环）在 \Tp{-O0} 报告里\textbf{根本不会出现}。最常见的事故是：在 Debug 上发现某个取值函数
占据百分之几，于是去``消调用开销''，而发布版本里它早就被展开掉了。正确做法：\textbf{测量一律用与发布
相同或尽量相同的优化级别与宏定义}，只在需要调试符号时补 \Tp{-g}（\Tp{-g} 只增加调试信息，正常情况下
不改变生成的代码；若某平台某版本上它改变了行为，以你自己的差分测试为准）。
\end{warning}

\subsection{开关如何落到测量上}

\begin{longtable}{>{\raggedright\arraybackslash}p{0.19\textwidth} >{\raggedright\arraybackslash}p{0.20\textwidth} >{\raggedright\arraybackslash}p{0.21\textwidth} >{\raggedright\arraybackslash}p{0.20\textwidth}}
\caption{常见构建开关与它对测量的影响（默认值随编译器与版本变，逐条以你的编译器文档为准）}
\label{tab:pm10-flag}\\
\toprule
开关 & 改变了什么 & 对测量的影响 & 怎么确认生效 \\
\midrule
\endfirsthead
\toprule
开关 & 改变了什么 & 对测量的影响 & 怎么确认生效 \\
\midrule
\endhead
\Tp{-O0} & 基本不优化 & 数字只反映执行路径，不反映发布性能 & 反汇编看是否逐语句落内存 \\
\Tp{-Og} & 以可调试为先 & 常被当发布近似，实际差距大 & 与 \Tp{-O2} 产物做体积/反汇编差分 \\
\Tp{-O2} & 常规优化 & 适合作为默认测量档 & 构建日志打印完整命令行 \\
\Tp{-O3} & 更激进的展开与内联 & 可能更快也可能更慢 & 两档各跑一次，交替对照 \\
\Tp{-Os} & 体积优先 & I-cache 受限场景可能反超 & 比较 \Tp{.text} 段大小 \\
\Tp{-DNDEBUG} & 关掉 \Fn{assert} & 断言里的校验全部消失 & \Fn{assert} 是否展开为空 \\
\Tp{-g} & 附加调试信息 & 只增大产物，正常不改代码 & 对比两次产物的代码段字节 \\
\Tp{-flto} & 跨 TU 决策 & 分开测的模块收益不再相加 & 见 第~\ref{sec:pm10-lto} 节 \\
\Tp{-fno-omit-}{\newline}\Tp{frame-{\Bk}pointer} & 保留帧指针 & 使栈可展开，自身有代价 & \Fn{perf} 报告里栈是否完整 \\
\Tp{-march=} / \Tp{-mtune=} & 目标指令集与调度模型 & 换机器可能非法指令 & 运行机与编译机是否同构 \\
\Tp{-ffast-math} & 放松浮点语义 & 结果会变，不只是速度 & 见 第~\ref{sec:pm10-vec} 节 \\
\bottomrule
\end{longtable}

最后一列是这一节的核心纪律：\textbf{每一个``我以为打开了''的开关都要有一个可执行的确认动作}。最省事的
做法是让编译器实际收到的完整命令行落到构建产物里。CMake 侧可以导出编译数据库，并让它随产物一起归档；
同时把编译器版本串编进程序，使``这个二进制是怎么来的''成为产物自身的属性（骨架见清单
\ref{lst:pm10-cmake}）。取单次构建的真实命令行，Makefile 生成器可用 \Tp{-n} 做一次只打印不执行的构建，
Ninja 侧用 \Tp{-v} 或直接读 \Tp{compile\_{\Bk}commands.json}——两条路都只是把同一件事变成可核对的文本。

\begin{lstlisting}[style=cmakestyle,title={导出编译数据库并把编译器版本写进产物},label={lst:pm10-cmake}]
set(CMAKE_EXPORT_COMPILE_COMMANDS ON)   # 生成编译数据库

execute_process(
  COMMAND ${CMAKE_CXX_COMPILER} --version
  OUTPUT_VARIABLE COMPILER_STAMP
  OUTPUT_STRIP_TRAILING_WHITESPACE)

target_compile_definitions(myapp PRIVATE
  PM_COMPILER_STAMP="${COMPILER_STAMP}"
  PM_BUILD_TYPE="${CMAKE_BUILD_TYPE}")
\end{lstlisting}

\section{内联的真实语义}
\label{sec:pm10-inline}

这是本章最容易写错、也是本册反复回指的一段。口径与《代码风格参考手册》完全一致，一句话概括：
\textbf{\Tp{inline} 首先是一个关于 ODR 与链接语义的关键字，它不是提速开关。}

\subsection{先把语义层次拆开}

\CPP{} 的 \Tp{inline} 说明符解决的是三件事，它们都不是``性能''：

\begin{enumerate}
  \item \textbf{ODR}：一个 \Tp{inline} 函数可以在多个翻译单元里各自出现一份定义，链接器把它们视为
        同一个实体。类内定义的成员函数、模板函数天然满足这条。
  \item \textbf{链接可见性}：\Tp{static} 或匿名命名空间给出内部链接，\Tp{inline} 给出的是``跨 TU 合法
        的单一实体''，两者在符号表上的形态不同（前者常成为局部符号，后者可能留下一个可被外部看到的
        弱符号，具体形态由编译器决定）。
  \item \textbf{对实现的一个提示}：标准把``是否真的展开''完全交给实现。也就是说第三层\textbf{没有}
        任何保证。
\end{enumerate}

真正决定是否展开的是优化器的启发式：函数体积、调用点热度、是否有取地址、是否跨越 TU 边界、当前
\Tp{-O} 级别、有没有 \Tp{-flto}、有没有 profile 反馈。于是会出现四种常见误解：

\begin{itemize}
  \item ``加了 \Tp{inline} 就一定会展开''——错。编译器完全可以忽略它，照样在符号表里留一个函数体并按
        调用处理。
  \item ``不加 \Tp{inline} 就不会展开''——也错。一个小函数即使没标，优化器常常自行展开（GCC/Clang 的
        \Tp{-O2} 就有``只有一个调用点就展开''这类规则，阈值随版本变）。
  \item ``\Tp{always\_{\Bk}inline} 是硬保证''——不准确。它是``尽力而为，做不到时给出诊断''：GCC 在确实
        无法展开时报错，Clang 多数情况下只给警告，MSVC 的强制内联修饰符文档明言不保证全部展开。三家
        行为并不相同，且都随版本变。
  \item ``\Tp{noinline} 只是建议''——恰恰相反，\textbf{\Tp{noinline} 才是这一族里接近硬保证的那个}：
        GCC/Clang 的属性写法与 MSVC 的声明修饰符都把``禁止展开''当成约束来实现；但边缘情形仍可能有
        例外，所以最终还是要看产物。
\end{itemize}

\begin{guideline}[title={把内联当成``需要观测的优化器决策''来写结论}]
在性能文档里，一句合法的表述只能是``在这个构建配置下，这个函数被展开/没有被展开''，并且要给出确认它的
方式（本节末尾的反汇编方法）。\textbf{非法}的表述是``我加了 \Tp{inline}，所以这里快了''——它既没有说明
观测结果，也没有说明观测对象。如果确实需要强制某种行为，正确的做法是：需要展开就用
\Tp{always\_{\Bk}inline}（并准备好接受编译器的诊断），需要保持独立函数体就用 \Tp{noinline}（这在测量里
常常更有用：把被测函数钉成一个可寻址的独立体，采样归属才干净），而不要指望裸 \Tp{inline}。
\end{guideline}

\subsection{MSVC 的 \Tp{/Ob} 一族}

MSVC 主要用命令行而不是关键字来表达展开意愿：\Tp{/Ob1}（展开标记为 \Tp{inline} 的函数）、\Tp{/Ob2}
（再加上编译器认为合适的）、\Tp{/Obin}（与 \Tp{/Os}、\Tp{/Ot} 一起隐含的那一档）。三点要注意：
第一，这些开关同样是``意愿''而不是保证，而且它们与整体优化开关 \Tp{/GL} 的组合会改变实际判定；第二，
\Tp{/Ob} 与 \Tp{/O2} 是两个正交维度，单独说``优化级别''不足以描述产物；第三，MSVC 的头文件里大量使用
带 \Tp{INLINE} 的项目宏，它们展开成什么取决于 \Tp{/Ob} 设置，因此同一份源码在两种设置下的实际代码差异
可能比你预期的大。

\subsection{唯一的确认方式：在发布产物上反汇编}

不要在一个专门打开调试开关的探针构建上看内联，也不要在 \Tp{-O0} 上看——你看到的必然是``没展开''。
要在\textbf{实际发布配置}产出的那个二进制上看。Linux/ELF 侧的一条可执行核对路径见清单
\ref{lst:pm10-inline-check}；它给的是命令形状，判读结论要落在``调用指令有没有出现''这一条硬线索上。

\begin{lstlisting}[style=shellstyle,title={在发布产物上核对展开与调用点（骨架，非本仓库实测）},label={lst:pm10-inline-check}]
# 0) 留下编译器串，作为结论的一部分
readelf -p .comment ./myapp | head -20

# 1) 这个符号还在不在？还在不代表一定被调用，不在基本说明全部展开
nm -C --defined-only libfoo.a  | grep -F 'my::helper'
nm -C -D ./myapp               | grep -F 'my::helper'

# 2) 在调用方所在的目标文件里，有没有指向它的 call
objdump -dC --no-show-raw-insn \
  build/CMakeFiles/myapp.dir/src/solver.cpp.o \
  | awk '/<my::hot_loop\(\)>:/{f=1} f&&/call/{print} /^$/{f=0}'

# 3) 展开的痕迹：被调函数的特征指令是否直接出现在调用方体内
objdump -dC ./myapp | sed -n '/<my::hot_loop>:/,/^$/p' > hot.asm
grep -c 'movsd' hot.asm

# 4) 段体积，用于判断是布局变了还是代码变了
size -A ./myapp | grep -F '.text'
\end{lstlisting}

Windows 侧对应 \Tp{dumpbin} 的反汇编模式（需要带代码的产物）与链接期的 \Tp{/MAP} 输出；macOS 侧是
\Fn{nm} 加 \Fn{otool} 的反汇编。三边共同的问题不是工具，而是\textbf{链接器与优化器可能已经改变了
函数边界}，所以判读要按``调用指令有没有出现''，而不是按符号在不在。

\subsection{内联对测量的两个副作用}

即便展开真的发生，它对读数的影响也常常是双向的：

\begin{enumerate}
  \item \textbf{归属被抹掉}。\Tp{-O} 级别越高、内联越多，采样剖析里独立函数的桶就越少，热点会向调用者
        集中。火焰图上表现为一个很宽的上层帧，此时你无法用剖析器回答``这个内联体值多少''——它已经不
        是一个可归属的单位。想要分开记，就把候选体用 \Tp{noinline} 钉住做一次\textbf{诊断专用}构建
        （它的数字不能当发布结论用）。
  \item \textbf{代码布局变差}。过度展开会拉长基本块、稀释 I-cache 与解码带宽。这在循环体很小但调用点
        很多的库代码里尤其常见。观察入口是 第~\ref{ch:pm06-micro} 章 的指令级归因与 \Mx{IPC}、
        \Mx{stalled-cycles}，而不是内联本身。
\end{enumerate}

\section{LTO：跨翻译单元的决策怎么记账}
\label{sec:pm10-lto}

\subsection{为什么模块收益不能相加}

\Tp{-flto} 把优化边界从单个源文件推到整个程序。它带来的典型收益来自三件单 TU 做不到的事：跨文件内联、
常量与偏移表的跨 TU 传播、死代码与未用虚函数分支的消除。问题是这三件事\textbf{互相抵消}：一次跨文件
内联可能同时增加代码体积，把原本局部的热路径挤散。于是``\Tp{-flto} 让 A 模块快若干个百分点、B 模块
也快若干个百分点''这类按模块测的数，在全程序上既不会相加，也不保证同号。唯一可靠的口径是：
\textbf{只在完整链接出来的可执行文件或共享库上测 LTO 与否}。

\subsection{编译期、链接期与构建时间也进账}

LTO 把一部分工作从编译期推到链接期，链接阶段常常变成构建里最慢的一步，而且并行度往往不足。这会从两个
方向伤到测量工作本身：迭代变慢让人更倾向于``少跑几轮''，直接违反 第~\ref{ch:pm03-stats} 章 的重复要求。
所以做 LTO 决策时，要把构建时间当作\textbf{结论的一部分}写下来，而不是当作工程抱怨。

GCC/Clang 侧常见的两个维度（名称与可用性随版本变，动手前用编译器的帮助输出或文档确认）：

\begin{itemize}
  \item \Tp{-flto=auto}（或显式给分片数）：让链接阶段并行分片，减少峰值内存与耗时。分片数影响
        \textbf{决策的确定性}，因此影响可复现性——同一份源码在 \Tp{auto} 下于不同核数的机器上可能
        得到不同产物。
  \item Full LTO 与 Thin LTO：前者把中间表示全部合并，跨模块决策更完整、代价更高；后者保留模块边界、
        只导出索引，链接快得多，但可做的决策集合不同（Clang 的 \Tp{-flto=thin} 与 GCC 的实现路线
        不可互换比较）。MSVC 侧对应 \Tp{/GL} 加 \Tp{/LTCG}，语义与粒度都不同。
\end{itemize}

\begin{compare}[title={要不要上 LTO：三笔账一起记}]
\begin{tabular}{>{\raggedright\arraybackslash}p{0.16\textwidth} >{\raggedright\arraybackslash}p{0.26\textwidth} >{\raggedright\arraybackslash}p{0.30\textwidth}}
\toprule
维度 & 只看运行时会得出什么 & 加上记账后该看什么 \\
\midrule
运行时 & LTO 版本更快 & 快多少是否有区间与判据（第~\ref{ch:pm03-stats} 章） \\
构建 & 链接慢一点 & 迭代次数下降会不会漏掉回归 \\
可复现 & 本机产物稳定 & 换核数或换链接并行度后产物是否变化 \\
\bottomrule
\end{tabular}
\end{compare}

\subsection{LTO 与剖析器的相互干扰}

LTO 之后，采样剖析看到的可能是一团被彻底重排的代码：符号还在但函数体面目全非，栈展开更难。若此时同时
开了 \Tp{-flto} 与 \Tp{-g}，调试信息的体积与质量都会成为问题（调试信息不能很好地跨 LTO 边界合并时，
报告里会出现半截的栈）。可行的分工是把两件事分开：\textbf{需要干净归属时不链接 LTO}，\textbf{需要
发布结论时必须链接 LTO}，并且不要把前者的相对差直接当成后者的绝对收益。

\section{PGO 与 BOLT：收益方向由训练数据决定}
\label{sec:pm10-pgo}

\subsection{profile 代表什么，就优化什么}

PGO 的两轮流程（先采集 profile 数据，再带着它重新编译）本身没有魔法：它只是把``哪条分支常走、哪个
调用点热、哪个函数值得内联''从猜测变成观测。于是它的\textbf{全部风险都集中在代表性}上。用一段样例流量
跑出来的 profile，会把优化预算花在样例里存在、真实负载里不存在的路径上——这不是``收益变小''，而是
可能\textbf{为不存在的路径付出体积与布局代价，让真实热路径更慢}。常见的三类失真：

\begin{itemize}
  \item 覆盖不全：样例没走到的分支，编译器按未预测处理；真实负载里它很热。
  \item 分布倒挂：样例是冷启动或小输入，真实是稳态或大输入，内联阈值与循环展开选择整个翻盘。
  \item 环境污染：采集时带着剖析器或压测工具，分支预测行为与线上不同。
\end{itemize}

下面那条注记给一条纪律：\textbf{PGO 的 profile 必须与它要服务的产物同源同负载}，并且要有一份能读出
``这份 profile 是怎么来的''的归档。

\begin{note}[title={PGO 不改变正确性判据，但改变可比性}]
带 PGO 的产物与不带的产物是两个二进制，\textbf{不能}与历史基线混在同一条时间序列里比较；一旦把 PGO
纳入发布，CI 基线要在切换当天重建（口径见 第~\ref{ch:pm13-ci} 章）。同理，profile 文件本身要进构建产物归档，
否则下一次重新采集得到的二进制与这一次不可比。
\end{note}

\subsection{把采集与校验写成脚本骨架}

\begin{lstlisting}[style=shellstyle,title={GCC 风格 PGO 两轮流程骨架（选项名以你的版本为准）},label={lst:pm10-pgo}]
set -eu
STEP=${1:-collect}

case "$STEP" in
collect)
  cmake -B build-pgo -DCMAKE_BUILD_TYPE=Release \
    -DCMAKE_CXX_FLAGS="-fprofile-generate=$PWD/profile-data"
  cmake --build build-pgo --parallel
  mkdir -p profile-data
  # 关键：这一步要跑代表性负载，不是单元测试
  ./build-pgo/myapp --workload typical --input-corpus ./corpus
  ;;
apply)
  cmake -B build-pgo2 -DCMAKE_BUILD_TYPE=Release \
    -DCMAKE_CXX_FLAGS="-fprofile-use=$PWD/profile-data -fprofile-correction"
  cmake --build build-pgo2 --parallel
  ;;
verify)
  # 产物校验：确认真的用了 profile，且 profile 不是旧的
  readelf -p .comment ./build-pgo2/myapp | head -20 || true
  nm -C --size-sort ./build-pgo2/myapp | tail -20
  ;;
esac
\end{lstlisting}

Clang 一族的差异要点：先由 \Tp{-fprofile-generate} 产出原始 profile 文件，再用 \Fn{llvm-profdata}
把它合并成一份索引式数据，使用时让 \Tp{-fprofile-use=} 指向合并后的文件；未合并就直接用，常常导致
``没有 profile 数据''一类的诊断并悄悄退化成非 PGO 构建——\textbf{这类诊断必须当失败处理}，否则你会
拿一个自认为带 PGO 的二进制去测。较新的工具链还会用二进制里的名字段来匹配 profile，名字不匹配同样
会静默降级；统一的确认方法是把构建日志当作产物的一部分，在里面搜 profile 族的诊断。

BOLT 走的是另一层：它在链接后的二进制上重排代码与基本块，因此它的收益主要是 I-cache 与分支布局方向的，
与 PGO 的算法层决策不完全重叠也不互相保证。BOLT 的输入需要带符号的重定位信息与一份执行样本（通常用
\Fn{perf} 的采样模式取分支记录类数据），于是它把 第~\ref{ch:pm04-sampling} 章 的工具链变成了构建前置；这条
链上任何一环（符号丢失、样本代表性）出问题，收益方向都可能翻。对它的正确态度与 PGO 相同：\textbf{只报
差分区间与判据，不报单点}。

\section{自动向量化：报告是线索不是结论}
\label{sec:pm10-vec}

\subsection{怎么读到``为什么没向量化''}

各编译器都提供了``向量化报告''一族开关，\textbf{名称与可用性随版本变}，落笔或动手前先在\textbf{你自己
手上那一个版本}的文档或帮助输出里确认它认哪一个：

\begin{itemize}
  \item GCC：一对 \Tp{-fopt-info-vec-optimized} 与 \Tp{-fopt-info-vec-missed}，前者报成功、后者报
        原因；更宽的 \Tp{-fopt-info-vec} 两者都报。这一族在不同版本里还细分过把报告写到文件的用法。
  \item Clang：\Tp{-Rpass=loop-vectorize} 报成功，\Tp{-Rpass-missed=loop-vectorize} 报``没做成''，
        \Tp{-Rpass-analysis=loop-vectorize} 报更细的原因；三条是分层的，只开第一条时你看不到失败。
  \item MSVC：\Tp{/Qvec-report} 系列（历史上还分过若干级别），以及 \Tp{/Qvec} 前缀下的细项。它与
        \Tp{/arch:} 支持到哪个指令集是两回事，别把``能生成''当成``确实生成了''。
  \item Intel 编译器：\Tp{-vec-report} 一族（同样随版本更名，较新版本走另一套报告体系）。
\end{itemize}

\begin{warning}[title={把向量化报告当性能结论}]
向量化报告只能回答``编译器\textbf{有没有}做这件事''，不能回答``做了之后\textbf{快没快}''。原因有三层：
一是向量化可能只是把一段本来就不是瓶颈的循环换了实现；二是向量化引入了对齐、掩码、余数循环，尾部代价在
循环次数小的调用点上会盖过收益；三是报告说的是``这一次编译的这一次判定''，换个 \Tp{-march}、换一份
profile、换一个 LTO 分片就可能不同。所以流程必须是：报告\textbf{定位}到某个循环，然后用微基准或真实
负载\textbf{测}这个循环在开关前后各自的区间，再用 第~\ref{ch:pm06-micro} 章 的 \Mx{instructions}、\Mx{IPC}
与缓存事件确认执行形态真的变了。最终仍然要在发布产物上反汇编看有没有向量指令——报告说有，编译器版本
一升级就可能没有。
\end{warning}

一段报告输出长这样（见清单 \ref{lst:pm10-vec}，\textbf{只是格式示例，非本仓库实测值}，措辞随版本变化
很大）：

\begin{lstlisting}[style=outstyle,title={向量化报告形态示意（分层三行，格式示例非实测）},label={lst:pm10-vec}]
# 示意输出：格式示例，非本仓库实测；行号与措辞随编译器版本变
src/dsp.cpp:118:3: remark: vectorized loop (vector width: 4, interleaved count: 2)
src/dsp.cpp:131:3: remark: loop not vectorized: the execution budget is too small
src/dsp.cpp:146:3: remark: loop not vectorized: cannot determine dependence between *dst and *src
\end{lstlisting}

最后那行``依赖关系判不出来''是最常见的一挡，它指向别名问题：编译器无法证明两个指针不重叠，就只能保守
处理。出路各自的代价不同，别混用：把循环拆成编译器能证明的部分；用语言层的别名声明（\CL{} 侧的
\Tp{restrict}，\CPP{} 里则要靠规范化的指针约定或编译器扩展，可移植性有限）；或者把数据组织成不产生
可疑重叠的形状（布局与局部性的口径见 第~\ref{ch:pm06-micro} 章）。

\subsection{归约、浮点重组与 \Tp{-ffast-math} 的语义代价}

浮点加法不满足结合律，所以``把归约重排成向量顺序''会\textbf{改变结果}。默认级别下编译器不会这么做，
这正是很多浮点循环向量化失败的直接原因。打开放松开关之后它才允许做，而放松开关不是免费的：

\begin{itemize}
  \item GCC/Clang 的 \Tp{-ffast-math} 是一组开关的合集（有限性假设、结合性、带符号零等若干子项），
        它允许丢失特殊值语义、允许重排、允许用非标准格式的近似实现。后果是结果可能不同、溢出行为
        可能改变、原本正确的边界判断可能失效。
  \item Clang 有较细的单独开关，可以只打开重排而保留其它部分；MSVC 的 \Tp{/fp:fast} 与
        \Tp{/fp:precise} 是同一维度的两档。各家的粒度并不可移植，逐项确认。
\end{itemize}

正确的写法是把这件事记成\textbf{正确性变更}而不是优化：需要放松浮点语义时，先给出``结果差在允许误差
内''的数值验证（用 第~\ref{ch:pm03-stats} 章 的判据比较两组输出，而不是比较运行时间），再报速度收益，并在
结论里写明它改变了浮点语义。反过来的顺序——先看到明显变快，再补一句``顺便打开了 \Tp{-ffast-math}''——
是本册意义上的反模式（详见 第~\ref{ch:pm15-antipattern} 章）。

\section{可复现性：同机同码为什么给出不同二进制}
\label{sec:pm10-repro}

如果两次构建产出不同的字节，那么``同一份代码两次跑出不同数字''里，就有一部分不是噪声而是真实代码差异。
常见的注入点：

\begin{itemize}
  \item \textbf{时间与路径}：编译日期时间宏、源文件路径宏、构建目录绝对路径被编进调试信息与字符串表，
        进而改变段偏移与对齐。
  \item \textbf{编译器内部随机性}：若干哈希种子、符号表遍历顺序、并行编译与 LTO 分片的完成顺序，都可能
        让代码布局不同。GCC 侧有 \Tp{-frandom-seed=} 用来钉住这类种子；构建目录名也常常间接影响它，
        所以固定 \Tp{-frandom-seed} 时同时固定输出路径才有效。
  \item \textbf{代码布局本身}：函数在 \Tp{.text} 里的先后决定它们落到哪些缓存组与页。布局不同带来的
        差异足以混淆一次小改进的量级（取决于工作集与访问模式，要读者自己标定），而且它\textbf{每次
        构建都可能变}。这就是``只跑一次的 A/B 一定不可信''最硬的来源之一，判据见
        第~\ref{ch:pm03-stats} 章。
\end{itemize}

地址随机化影响的是运行期地址，与产物字节是否相同是两件事，别混起来：随机化会改变页与缓存索引，从而让
同一产物在不同进程实例间有可观测差；在\textbf{测量专用}的前提下关掉它可以降低这类方差（可用性与副作用
要自己确认，别忘了它同时改变安全属性），但更稳妥的做法是把它当随机效应处理并靠重复。

把编译器版本写进结论，最省事的办法是让它成为产物的属性而不是外部记忆：ELF 侧的注释段会留下编译器串
（清单 \ref{lst:pm10-inline-check} 的第 0 步就是这件事）；把编译数据库、编译器版本串、版本控制提交号
三件一起归档，就得到一份``这个二进制是怎么来的''的可核对描述。链接顺序同理：静态库的输入顺序改变会
决定包含哪些目标文件，从而改变代码；对动态库，符号解析顺序不同会让同一个调用绑到不同实现（这一类问题
在 第~\ref{ch:pm07-memory} 章 里以分配器替换的形式最常见）。

\begin{bestpractice}[title={把构建口径做成一段可粘贴的文字}]
一个能进代码评审的性能结论，开头应该有一段``被测对象''，包含：编译器与版本串、构建类型与实际生效的
\Tp{-O} 档、\Tp{NDEBUG} 状态、是否 \Tp{-flto}（full 还是 thin、并行度多少）、是否 PGO/BOLT 以及
profile 来源、\Tp{-march}/\Tp{-mtune} 或等价的 \Tp{/arch:} 设置、产物大小与 \Tp{.text} 段大小、核对
内联与向量化所用的反汇编命令。这段文字不需要读者相信任何口头承诺，每一条都可以被复核。
\end{bestpractice}

\section{排障与速查}
\label{sec:pm10-check}

\begin{guideline}[title={改一行代码前先确认改了二进制}]
顺序永远是：先确认构建确实被触发（目标文件时间戳、产物字节变了），再确认改动落在你想让它落的地方
（符号表或反汇编），最后才读数。跳过前两步得到的``没有差别''，含义是``改动没进产物''，而不是
``改动无效''。
\end{guideline}

\begin{longtable}{>{\raggedright\arraybackslash}p{0.13\textwidth} >{\raggedright\arraybackslash}p{0.29\textwidth} >{\raggedright\arraybackslash}p{0.38\textwidth}}
\caption{产物核对四步：从``我以为改了''到``我确认改了''}
\label{tab:pm10-check}\\
\toprule
步骤 & 看什么 & 命令与判读 \\
\midrule
\endfirsthead
\toprule
步骤 & 看什么 & 命令与判读 \\
\midrule
\endhead
1 符号表 & 目标函数是否还在、大小是否变 & \Fn{nm} 加按大小排序的选项；函数消失多半是被展开或消除 \\
2 反汇编 & 调用点是否消失、有无向量指令 & \Fn{objdump} 反汇编后在函数体内查调用；展开与否只能在这里定 \\
3 大小与段 & \Tp{.text} 与整体体积 & \Fn{size}、\Fn{readelf} 段表；\Tp{-O3}、LTO、BOLT 都可能只动这一项 \\
4 链接与版本 & 输入顺序、编译器串 & 注释段与链接命令行；同一提交不同产物多半在这里 \\
\bottomrule
\end{longtable}

三条按症状的处置，按``先能确认、再能解释''的顺序排列：

\begin{enumerate}
  \item 现象是``加了 \Tp{inline} 之后毫无变化''。这就是预期行为：裸 \Tp{inline} 并不承诺展开。要展开
        就改用 \Tp{always\_{\Bk}inline}，要保留独立函数体就改用 \Tp{noinline}；两种意图最后都要在发布
        产物上用反汇编确认一次，否则写不出可复述的结论。
  \item 现象是``打开 \Tp{-O3} 反而变慢''。不要直接退回低档位就完事：先看是不是体积或布局问题（比较
        \Tp{.text} 段大小，再看 \Mx{IPC} 与前端受限事件，入口 第~\ref{ch:pm05-events} 章），也可能是循环
        展开类开关把尾部代价放大了。逐条关掉开关做二分，而不是猜哪一个。
  \item 现象是``LTO 或 PGO 之后基准变差，而功能测试全部通过''。优先怀疑 profile 的代表性
        （第~\ref{sec:pm10-pgo} 节），其次怀疑被测循环自身被优化掉了（第~\ref{ch:pm15-antipattern} 章 的微基准
        一节）：优化级别越高，越容易把基准循环当成常量算掉。
\end{enumerate}

最后把本章的立场写清楚：编译器开关不是性能旋钮，而是\textbf{被测对象的定义}。\Tp{inline} 给你的是链接
语义，\Tp{-O3} 给你的是一组可能相互抵消的决策，向量化报告给你的是编译器的一次自述。它们都必须在发布
产物上、以可复述的动作核对过，才能进入结论；否则你测的是构建系统，不是代码。

\part{异构与 GPU：把计时口径搬到设备侧}

\chapter{GPU 计时口径：CUDA event、流同步与 kernel 时间}
\label{ch:pm11-gpu}

前面各章的工作假设是：CPU 调一个函数，函数就开始跑，跑完才返回。把工作交给 GPU 之后这个假设断了——
提交（launch）只是把一条命令放进队列，CPU 立刻拿到返回值，真正的执行发生在稍后、
发生在另一个时钟域、发生在可能被别人占着的硬件上。于是要问的第一个问题不再是``怎么让它变快''，
而是\emph{我这一次量到的到底是哪一段时间}。本章依次给出：异步为什么让计时变成协议问题
（第~\ref{sec:pm11-async}~节，含三种口径的选择表）、CUDA event 各调用的确切语义与边界
（第~\ref{sec:pm11-event}~节，含 GPU 上为什么更常取最小值）、一段可跑的测量骨架
（第~\ref{sec:pm11-skeleton}~节）、设备侧看 kernel 时间与工具扰动
（第~\ref{sec:pm11-device-side}~节）、两个时钟域怎么对齐（第~\ref{sec:pm11-domain}~节）、
口径搬到 ROCm/oneAPI/图形 API 时哪些东西搬不动（第~\ref{sec:pm11-portability}~节），
最后是计时前的六项确认（第~\ref{sec:pm11-trouble}~节）。
接口：时钟源归第~\ref{ch:pm02-clock} 章，本章只接它留下的 GPU 那一格；重复次数与判据归第~\ref{ch:pm03-stats} 章；
数据搬运、PCIe 与多卡的账在第~\ref{ch:pm12-bandwidth} 章；GPU 基准能不能挂进流水线在第~\ref{ch:pm13-ci} 章；
最容易犯的几个写法作为反例收在第~\ref{ch:pm15-antipattern} 章。本仓库的 CI runner 里没有 GPU 也没有 Nsight，
所以本章所有工具输出都只是\emph{形态示意}：字段布局与列顺序可信，数字全部写成占位，不得引用。

\section{异步执行把计时变成协议问题}
\label{sec:pm11-async}

\subsection{提交、执行、同步是三段不同的时间}
\label{sub:pm11-queue}

在流（stream）上写 \Fn{cudaMem{\Bk}cpy{\Bk}Async} 或 launch 一个 kernel，
语义是``把这件事挂到这条流的尾部''，函数返回时设备可能还没开始。让 CPU 等到设备做完的只有显式同步
（\Fn{cudaDevice{\Bk}Synchronize}、\Fn{cudaStream{\Bk}Synchronize}、
\Fn{cudaEvent{\Bk}Synchronize}）或一次强制串行的读回。
于是``在 CPU 侧包一段墙钟''测到的不是一个数，而是三项之和：

\begin{itemize}
  \item \textbf{排队与提交}：驱动写队列的成本、前序未完成任务占着设备的时间、
        同卡其它使用者（图形合成器、另一个 context、别的进程）带来的等待；
  \item \textbf{执行}：kernel 与拷贝真正占用设备的那段，是优化最常想改的一格；
  \item \textbf{同步}：CPU 从``等''里被唤醒的成本，含自旋还是睡眠（取决于同步方式与驱动）、
        线程调度、以及同步点之后那一次读回。
\end{itemize}

三段各自由不同的东西决定：第一段看队列深度与邻居，第二段看算法与硬件，第三段看你选的同步方式。
把它们混进一个数，得到的结论就只是``GPU 变慢了''，下一步该改哪个模块毫无指向。
这就是``计时是协议问题''的确切含义：\emph{先声明计时段的左右边界落在哪一层，再谈数字}。

\subsection{三种口径与它们各自能回答的问题}
\label{sub:pm11-three-scope}

同一份 GPU 代码，本章承认三种互不等价的口径。差别不在精度，而在\emph{包含关系}：
表~\ref{tab:pm11-scope} 按``谁包含谁''排，不是按优劣排。

\begin{longtable}{>{\raggedright\arraybackslash}p{0.17\textwidth}
                  >{\raggedright\arraybackslash}p{0.22\textwidth}
                  >{\raggedright\arraybackslash}p{0.22\textwidth}
                  >{\raggedright\arraybackslash}p{0.22\textwidth}}
\caption{GPU 侧三种计时口径的适用问题域（口径选择表，非数据表）}\label{tab:pm11-scope}\\
\toprule
口径 & 时间从哪来 & 包含什么 & 适合回答什么 \\
\midrule
\endfirsthead
\toprule
口径 & 时间从哪来 & 包含什么 & 适合回答什么 \\
\midrule
\endhead
\bottomrule
\endfoot
CPU 侧 \MWall{} 包住提交与同步 &
CPU 的单调时钟（第~\ref{ch:pm02-clock} 章） &
排队 $+$ 执行 $+$ 同步唤醒 &
``调用方等了多久''、端到端吞吐、上线约束；不能用来判断某个 kernel 有没有变快 \\
CUDA event 差 &
设备侧打点，两个 event 之间的流内间隔 &
同一条流上被打点包住的那段设备工作 &
``这段 GPU 工作本身多久''、kernel 级 A/B；\textbf{不含}提交前的排队 \\
剖析器的设备时间戳 &
剖析器读到的逐任务时间戳 &
单个 kernel、单次 memcpy/memset，逐条列出 &
把一段时间拆到具体任务、看重叠与空隙；扰动大小由工具决定（第~\ref{sub:pm11-replay}~节） \\
\end{longtable}

常见错误是把第二行当成第一行：``event 说 0.4\,ms，端到端墙钟说 2.1\,ms，所以补丁只快了 0.4\,ms''。
差值是排队、同步与主机侧准备，它可能才是用户感受到的那部分。反过来也错：端到端 \MWall{} 变好，
可能只是队列变空了，kernel 一行代码没改。\textbf{两个都报、各自说明含义}，是唯一打不穿的写法。

\begin{guideline}[title={要点：先选口径，再选工具}]
\label{box:pm11-guideline-scope}
定口径的顺序是：问题（谁在等这件事）$\rightarrow$ 边界（计时段左右各落在哪个调用上）
$\rightarrow$ 时钟（CPU 单调钟还是设备打点）$\rightarrow$ 统计量（min、\Mx{p50} 还是 \Mx{p99}）。
工具只是把边界落实成代码的手段；边界先定下来，用 CUDA event 还是用 \Fn{nsys} 的时间戳，
剩下的只是扰动大小与可见粒度的取舍。
\end{guideline}

\section{CUDA event 的语义与边界}
\label{sec:pm11-event}

\subsection{各调用做什么}
\label{sub:pm11-event-api}

CUDA event 是``能被打进流的时间戳''。一次标准测量涉及四个调用，
把它们分开列是因为每一步都各有一个坑：

\begin{itemize}
  \item \Fn{cudaEvent{\Bk}Create} 建打点对象；纯同步用途可改用
        \Fn{cudaEvent{\Bk}Create{\Bk}WithFlags} 带上 \Tp{cudaEvent{\Bk}Disable{\Bk}Timing}，
        省掉打点成本；
  \item \Fn{cudaEvent{\Bk}Record} 把打点这件动作\emph{排进指定的流}：它不等设备，
        只是往队列尾部追加一项，打点时刻由设备在流推进到那里时才记录；
  \item \Fn{cudaEvent{\Bk}Synchronize} 让 CPU 等到该 event 在设备上完成；
  \item \Fn{cudaEvent{\Bk}ElapsedTime} 给出两个 event 打点之间的毫秒差，返回类型是 \Tp{float}。
\end{itemize}

由此得出三条必须写进代码注释的纪律。

\textbf{一，读值前两个 event 都必须已完成。}
end 还挂着的时候读那个差值，结果不可用；标准写法是先
\Fn{cudaEvent{\Bk}Synchronize}(end)，或改用 \Fn{cudaEvent{\Bk}Query} 轮询到完成。
哪个错误码表示``还没完成''随版本变——做一个实验自证：刚 record 完立刻读值，
把返回码与返回的值一起记下来，别猜也别抄论坛。

\textbf{二，两个 event 必须同流、且顺序与你的假设一致。}
跨流相减得到的是两个设备计数器之差：同一时钟域内可比，但不构成任何一段有意义的工作时间，
因为两条流可以同时推进，差值取决于它们怎么重叠。

\textbf{三，\Fn{cudaEvent{\Bk}Record} 之前的 CPU 侧代码进不了这个数。}
beg 打点在设备推进到那里时才发生，所以 beg 的时间戳晚于你调用它的时刻；
也就是说 event 差\emph{自动排除}了排队等待。想测排队只能用第一行口径或剖析器的时间线。

\subsection{分辨率、精度与``能测多短''的自测法}
\label{sub:pm11-resolution}

常被引用的打点粒度是零点几微秒那个量级，但这是要读者自己标定的量，不是保证值：
粒度、舍入方式、是否受流调度影响，都随驱动与实现变。标定只需三行代码——
在同一条空流上连续 record 两个 event、同步、相减，重复几十次看分布。
它给出``这段空档有多大''的下限感知，也顺带告诉你 \Tp{float} 在读回前舍掉了多少位。
两条推论：不要用 event 给亚微秒级的东西下结论，那个尺度属于第~\ref{ch:pm05-events} 章的 \Mx{IPC}
与第~\ref{ch:pm06-micro} 章的指令级证据；也不要拿它报分钟级任务，毫秒 \Tp{float} 有效位数有限，
长任务用 CPU 侧 \MWall{}，事后再与 event 差对账。

\subsection{默认流的隐式同步}
\label{sub:pm11-default-stream}

不写流参数的调用走 legacy default stream，而它与所有 blocking 流之间存在隐式双向同步。
后果在测量上很具体：默认流上测出的 event 差可能包含别人在其它流上正跑的工作，
你的 kernel 也可能被那条流上先前提交的无关工作挡住——于是``这段代码变慢了''
其实只是``这条流上多排了一会儿队''。
处理办法是把被测段放进显式创建的流，必要时用 \Tp{cudaStream{\Bk}NonBlocking} 创建，
或使用线程私有默认流；这个选项是否存在、默认朝哪边随 CUDA 版本变，
以 \Fn{cudaDevice{\Bk}Get{\Bk}Attribute} 的返回与你手上头文件的注释为准。
判据一句话：只要程序里除了被测 kernel 还有别的流在跑（图形、解码、另一个线程），
默认流上的任何数都不可信。

\begin{note}[title={注记：event 复用与打点边界}]
\label{box:pm11-note-event}
一个 event 只携带一次``最近打点''：同一条流上对同一个 event 连续 \Fn{cudaEvent{\Bk}Record} 两次，
第一次的时间戳被覆盖，它不是环形缓冲。循环里复用 beg/end 是常见写法，
但要在上一轮 \Fn{cudaEvent{\Bk}ElapsedTime} 之后才重新 record，
否则你会拿这一轮的 beg 去配上一轮的 end。另外，launch 与 record 的先后只在同一条流内成立；
把 record 发到另一条流，顺序就交给队列了。
\end{note}

\subsection{为什么 GPU 上更常取最小值}
\label{sub:pm11-min}

CPU 微基准常报中位数，GPU 圈里更常报 min。机制在噪声的形状：GPU 样本的长尾几乎都是\emph{加法}噪声——
多排一会儿队、被邻居抢一次设备、时钟还在往上爬、CPU 唤醒晚了、页表刚建立。
这些只会把某次样本推长，不会让它比``设备真跑一遍''更短。所以在空闲、独占、时钟稳定的机器上，
最小值最接近``这件事本身要多久''；均值被环境事件抬高，且不随你的改动下降。
但 min 有三个使用条件，必须连同结论一起写下来：

\begin{itemize}
  \item \textbf{样本数固定}。极值天然随 $n$ 单调变小，$n=20$ 的 min 与 $n=200$ 的 min 不能比；
        $n$ 先按第~\ref{ch:pm03-stats} 章定好，A/B 两侧用同一个 $n$。
  \item \textbf{环境独占}。同卡上有多进程服务、图形合成器或另一个 context 时，
        min 是``最好运气''而不是承诺；这时要报的恰好是 \Mx{p50} 与 \Mx{p99}（那是第~\ref{ch:pm14-production} 章的线上口径）。
  \item \textbf{别用 min 判回归}。一次没跑好不会体现在 min 上；判定要用区间（第~\ref{ch:pm03-stats} 章）
        或连续两次确认（第~\ref{ch:pm13-ci} 章）。
\end{itemize}

一句话：min 回答``这段 kernel 的地板在哪''，不回答``用户会不会觉得慢''。

\section{一段能用的 GPU 测量骨架}
\label{sec:pm11-skeleton}

\subsection{warmup 要过四关}
\label{sub:pm11-warmup}

第一次 launch 的代价里混着四件事，每件只影响前几次：
\textbf{上下文与模块加载}——进程首次使用设备要建 context，产物若只带 PTX 而不含当前架构的 SASS
还要现场 JIT（一次性成本，只在冷启动出现；架构覆盖由 \Tp{-gencode} 一类编译选项决定，
测前先确认目标卡有实时码）；\textbf{库内部的算法选择}——cuBLAS/cuBLASLt、cuDNN 一类
在首次调用时按形状挑 kernel 并缓存选择，开了自动调优成本更大；\textbf{缓存与地址翻译}——
L2、常量缓存、TLB 与页表冷的时候都慢，反复读写同一块显存后才热；\textbf{时钟爬升}——
空闲时频率在低位，第一批 kernel 在低频下跑，何时拉高不由你的代码决定
（第~\ref{sub:pm11-clocklock}~节）。
warmup 次数不该拍脑袋：跑一段更长的序列，把累积均值按次数画出来，平在哪儿就在哪儿起测
（判据与画法见第~\ref{ch:pm03-stats} 章）；显式做法是丢掉前若干次，并把``丢了几次''写进复现步骤。

\subsection{把搬运与计算分账}
\label{sub:pm11-transfer}

H2D/D2H 要单独计。一是它们走另一条通路（带宽与链路宽度见第~\ref{ch:pm12-bandwidth} 章）；
二是只有\emph{页锁定}主机内存上的 \Fn{cudaMem{\Bk}cpy{\Bk}Async} 才真能异步，
普通内存上的异步拷贝会被驱动拆成同步分段，于是``计算段时间''里藏着一次同步。
小拷贝还有一层坑：启动成本与几微秒同量级时，字节数越小越被启动成本主导，
``带宽随尺寸变小而变差''其实是分母错了。量级由读者自己标定：扫一圈尺寸，
把每次的 event 差与尺寸成对存下来，曲线自己会说哪一段是启动开销区。

\begin{lstlisting}[style=cppstyle, caption={GPU event 计时骨架：口径边界、warmup 与错误检查的位置都写在段外（可编译形态，数值由读者产生）}, label={lst:pm11-harness}]
// Convention: the timed section contains device work only.
// Warmup, printing and error checks live outside it.
// Both scopes are collected side by side: device time and wall time.
struct GpuTimer {
  cudaEvent_t beg;
  cudaEvent_t end;
  GpuTimer()  { cudaEventCreate(&beg); cudaEventCreate(&end); }
  ~GpuTimer() { cudaEventDestroy(beg); cudaEventDestroy(end); }
  void start(cudaStream_t s) { cudaEventRecord(beg, s); }
  // milliseconds; caller guarantees the timed section is enqueued
  double stop(cudaStream_t s) {
    cudaEventRecord(end, s);
    cudaEventSynchronize(end);       // drop this line -> garbage reading
    float ms = 0.0f;
    cudaError_t e = cudaEventElapsedTime(&ms, beg, end);
    if (e != cudaSuccess) { return -1.0; }
    return static_cast<double>(ms);
  }
};

void bench(cudaStream_t s, float* d_in, float* d_out, size_t n) {
  const int kWarm = 50;              // discarded iterations, record this
  const int kRun  = 200;             // same n on both sides of an A/B
  for (int i = 0; i < kWarm; ++i) {  // 1) warmup, never timed
    launch_kernel(d_in, d_out, n, s);
  }
  cudaStreamSynchronize(s);          // drain warmup before we start
  (void)cudaGetLastError();          // async errors surface here first

  std::vector<double> dev(kRun);
  std::vector<double> wall(kRun);
  GpuTimer t;
  for (int i = 0; i < kRun; ++i) {
    auto c0 = std::chrono::steady_clock::now();   // host scope, left edge
    t.start(s);                                   // device scope, left edge
    launch_kernel(d_in, d_out, n, s);
    double dt = t.stop(s);                        // sync at the right edge
    auto c1 = std::chrono::steady_clock::now();
    dev[i]  = dt;
    wall[i] = std::chrono::duration<double,
              std::milli>(c1 - c0).count();
  }
  // 2) report after the loop: one print inside would enter the
  //    timed section and move the minimum
  report_min_median_p99(dev);
  report_min_median_p99(wall);
  // 3) one error check outside the loop, not inside it
  cudaError_t post = cudaGetLastError();
  if (post != cudaSuccess) {
    std::cout << cudaGetErrorString(post) << std::endl;
  }
}
\end{lstlisting}

\begin{guideline}[title={要点：骨架里三条不动的规矩}]
\label{box:pm11-guideline-skeleton}
一，\textbf{同步只出现在计时段的右边界}：别在段中间插 \Fn{cudaDevice{\Bk}Synchronize} 图省事，
它把``等整台设备''的成本算进了你的 kernel。
二，\textbf{错误检查在段外}：launch 的错误是延迟报出的，把 \Fn{cudaGetLast{\Bk}Error} 放进循环
等于给每个样本加一次驱动调用，落在计时段里就污染 min。
三，\textbf{两个口径同时采}（骨架里的 \Tp{dev} 与 \Tp{wall}）：
只有并列两列才看得出``排队在涨''这种形状，单列数在评审里一定被追问。
\end{guideline}

\section{从设备侧看 kernel 时间}
\label{sec:pm11-device-side}

\subsection{三代工具的分工}
\label{sub:pm11-tools}

\begin{itemize}
  \item \textbf{nvprof}（旧命令行剖析器）直接给每个 kernel 的平均/最小/最大时间与调用次数，
        但在较新的 CUDA 版本与架构上已不再提供：先确认它还在不在，不在就用下面两个，
        别照抄教程命令。
  \item \textbf{Nsight Systems}（\Fn{nsys}）看\emph{时间线}：kernel、拷贝、memset 与
        CPU 侧提交、同步的相对位置与重叠，读的是设备活动记录。它回答``这段时间里设备在做什么、
        空隙在哪''，扰动相对小，但逐任务时间不是高精度证据（工具自己也要打点）。
  \item \textbf{Nsight Compute}（\Fn{ncu}）看\emph{单个 kernel 的内部}：计数器、访存效率、
        occupancy、stall 分类。它回答``为什么慢''。
\end{itemize}

三者不是精度递增，是问题域不同。把 \Fn{ncu} 的时间当成``这个 kernel 多久''报出去，
是本章最容易犯、也最容易被一眼看出的错。

\subsection{重放出来的时间不是运行时序}
\label{sub:pm11-replay}

物理计数器个数有限，一个 kernel 跑一遍读不完想读的事件。\Fn{ncu} 的办法是\emph{重放}：
同一 kernel 执行多遍、每遍读一组计数器，必要时保存并恢复设备状态（显存写入要回滚，
否则下一遍的数据不同）。于是报告里的时间是``工具接管下、第 $k$ 遍、某个时钟状态下那一次''，
不是你程序里那一次。再叠两件事：工具通常会\emph{接管频率}（把 GPU 时钟钉到某个基准，
开关名以 \Fn{ncu --help} 为准，形态见清单~\ref{lst:pm11-cmds}）；为精确取样还可能把 kernel
串行化。所以：\textbf{可以}从 \Fn{ncu} 读事件之间的\emph{比例}、瓶颈归因、哪类 stall 占大头、
访存命中形态；\textbf{不可以}读绝对时长、A/B 的加速比与``上线能省多少''——
加速比请用 CUDA event（第~\ref{sec:pm11-event}~节）或 \Fn{nsys} 时间线，在同一时钟状态下取。
第~\ref{ch:pm05-events} 章讲过 CPU 侧 PMU 复用会把一次运行拆成多轮、报告给缩放比例；GPU 这边是同一件事换了工具名，
判据也一样：\emph{不同轮次取到的量不能相加}。

\subsection{occupancy 与 kernel 时间不是线性关系}
\label{sub:pm11-occupancy}

occupancy（活跃 warp 数与该架构每 SM 最大可驻留 warp 数之比）由三件事卡住：每线程寄存器数、
每 block 共享内存量、每 block 与每 SM 的线程和 block 上限。它衡量的是``设备有多少在飞的活
可以用来掩盖延迟''，不是一个直接的时间，所以画不成一条直线：

\begin{itemize}
  \item \textbf{掩盖延迟需要互相独立的 warp}。若每个 warp 都在等同一条长依赖链
        （串行归约、被分支串行化的路径），再多 warp 也只能排队，时间不降。
  \item \textbf{瓶颈可能在带宽}。DRAM 或 L2 已被现有 warp 打满时，提高 occupancy 只是增加争抢，
        甚至因为更多 warp 抢同一份缓存行而变慢。
  \item \textbf{低 occupancy 不等于慢}。计算密度高、寄存器压力大的 kernel 常年在低 occupancy 下
        接近峰值算力；小 kernel 可能被启动粒度限制，驻留量到不了上限，时间却由启动与尾效应主导。
  \item \textbf{同名不同义}。报告里的 occupancy 常常是理论上限，实际是否够用取决于访存能否
        被掩盖——这只有 stall 构成与实测时间能回答。
\end{itemize}

正确用法是做\emph{差分实验}而不是照着数字优化：固定算法与数据，只改一个变量（block 尺寸、
或每线程寄存器上限、或共享内存用量），在同一时钟状态下分别用 event 取时间、用 \Fn{ncu} 取
stall 构成，两条曲线一起看。时间随 occupancy 单调下降，说明这段在延迟受限区；时间几乎不动，
说明它不在 occupancy 的控制之下——这与``该不该继续调''是两个完全不同的结论。
work-group 尺寸与数据并行的变换本身不在本章，那是《并行算法》册的内容；
本章只关心一件事：\emph{别把 occupancy 当成时间报出去}。

\subsection{时钟锁定：先确认你的驱动与工具支持哪一种}
\label{sub:pm11-clocklock}

不锁时钟的测量，样本里混着 DVFS 的历史：前几次低频、后面高频，或者机房散热一变整段漂移
（同第~\ref{ch:pm02-clock} 章的``同一台机器今天不一样''）。三种处理，能不能用取决于驱动与工具版本：

\begin{enumerate}
  \item \textbf{驱动侧锁频}。NVIDIA 的持久模式与锁频开关、AMD 的性能等级设置，都要先看
        \Fn{nvidia-smi --help} 与 \Fn{rocm-smi --help} 里到底有没有该项，
        再看时钟查询报的支持范围；容器与托管实例里这类开关常被拒（权限或虚拟化限制），
        必须实测一次能否生效、生效后频率是否真按你锁的值走。
  \item \textbf{工具侧接管}。\Fn{ncu} 一类的 clock control：剖瓶颈时保持同一基准即可，
        但它与运行时不同频（第~\ref{sub:pm11-replay}~节）。
  \item \textbf{不锁，只采样}。每个样本旁边记一次 SM 频率与 \Tp{pstate}，
        把频率分布连同时间一起报。这条路最诚实，代价是多存一列。
\end{enumerate}

任何 GPU 结论都应能回答一句``当时 SM 频率处于什么状态''；答不上来，版本之间的差就分不清
是代码的还是时钟的。

\begin{lstlisting}[style=shellstyle, caption={GPU 计时前后的命令形态（字段与支持项以本机 --help 输出为准；本仓库无 GPU，此处仅示意）}, label={lst:pm11-cmds}]
# 1) What this card can report: supported clocks and current pstate
nvidia-smi -q -d CLOCK,PERFORMANCE | less
nvidia-smi --query-gpu=name,pstate,clocks.sm,clocks.max.sm --format=csv

# 2) Is a clock lock offered at all? Option names move with the
#    driver version, so read the help before copying a recipe
nvidia-smi --help | grep -i -e lock -e clock
nvidia-smi --query-supported-clocks=graphics --format=csv | head

# 3) Timeline: where kernels, copies and host sync calls sit
nsys profile -o run001 --force-overwrite true ./gemm_bench
nsys stats run001.nsys-rep

# 4) One kernel in depth: replays it and takes over the clocks
ncu --set full --kernel-name regex:gemm --launch-count 3 -o gemm001 ./gemm_bench
ncu --help | grep -i clock          # which switch your version uses

# 5) Per-sample frequency column, for the measure-but-not-lock scope
nvidia-smi --query-gpu=timestamp,clocks.sm,pstate --format=csv,noheader
\end{lstlisting}

\begin{lstlisting}[style=outstyle, caption={形态示意：时间线统计与逐 kernel 报告的字段布局（非本仓库实测值；数值全部写成占位，只表示列的形状与单位）}, label={lst:pm11-shape}]
# (1) timeline statistics: one row per name. Total is a sum of device
#     segments, not the end-to-end wall time of the process
   Time Range                Avg (ms)  Total (ms)  Instances  Name
   x.xxxxxx - x.xxxxxx       x.xxx     x.xxx       NN         gemm_kernel
   x.xxxxxx - x.xxxxxx       x.xxx     x.xxx       NN         cudaMemcpyAsync

# (2) per-kernel report fields, measured while the tool owns the clocks
  Kernel Name               : gemm_kernel
    Duration                : x.xxx us    <- one replay, not a runtime sample
    Grid Size / Block Size  : (x, x, x) / (x, x, x)
    Theoretical Occupancy   : xx.xx %     <- an upper bound, not a duration
    Measured Occupancy      : xx.xx %
    Memory Throughput       : xx.xx % of peak
    Compute (SM) Throughput : xx.xx % of peak
\end{lstlisting}

清单~\ref{lst:pm11-shape} 的每一列都在说同一件事：\textbf{同一屏里同时存在``设备时间''
``重放时间''与``理论上限''三种量}，报数时必须写明取的是哪一列。

\section{CPU 与 GPU 时钟域的对齐}
\label{sec:pm11-domain}

第~\ref{ch:pm02-clock} 章的结论在这里继续成立：\emph{只有同源的量才能相减}。GPU 打点用设备侧计数器，
CPU 侧 \Fn{clock{\Bk}\_gettime} 或 \Fn{Query{\Bk}Performance{\Bk}Counter} 用主机计数器，
两者不同源、起点不同、走的路径也不同。三条可操作的规矩。

\textbf{一，跨域只做分解，不做相减。}``端到端墙钟减 event 差等于 CPU 开销''这个式子看着漂亮，
实际含四项：提交、驱动排队、同步唤醒、以及两个域之间的相对漂移。要写成
``其它 $=$ 端到端 $-$ 设备段''，逐项说明可能来源，且只在两个数各自可复现时才报。

\textbf{二，单位与换算要自己确认。}CUDA event 直接给毫秒；设备侧的 \Fn{clock}
（以及 \Tp{clock64} 形态）给的是 SM 周期，频率一变换算跟着错，只能当``周期数''这个无量纲量用；
PTX 层的 \Tp{globaltimer} 是纳秒级域，但\emph{多卡之间不保证同步}，跨卡比较属于第~\ref{ch:pm12-bandwidth} 章。
图形 API 的时间戳是 tick 计数，换算要用设备报告的频率或周期参数
（第~\ref{sec:pm11-portability}~节）。凡``由设备告诉你的换算''，都要在环境指纹里记下当时那个值——
它随驱动与电源状态变。

\textbf{三，对齐误差要量，不是假设。}若确实要把 CPU 事件与设备事件画到同一条时间线上
（做重叠分析时就要），标定办法是跑一段空 trace、两侧同时读时间戳，
看差值的漂移与抖动随时间呈什么形状。剖析器提供的关联时间戳（CUPTI 一类的活动记录）
其单位与时钟源随版本与平台变，以工具文档加上你自己的漂移实验为准；
``看起来对齐了''的一张截图不是证据。

\section{ROCm、oneAPI 与图形侧：口径怎么搬}
\label{sec:pm11-portability}

搬的是\emph{口径}，不是算法：work-group 尺寸、数据并行变换与归约结构属于《并行算法》册，
这里只回答``前面那三种口径分别落在哪些调用上''。

\textbf{HIP（ROCm）。}\Fn{hipEvent{\Bk}Create}、\Fn{hipEvent{\Bk}Record}、
\Fn{hipEvent{\Bk}Synchronize}、\Fn{hipEvent{\Bk}ElapsedTime} 与 CUDA 一一对应、参数位置相同，
所以第~\ref{sec:pm11-skeleton}~节的骨架能逐字搬。搬不动的是三个前提：默认流语义
（HIP 同样区分 \Tp{hipStream{\Bk}Default}、\Tp{hipStream{\Bk}NonBlocking} 与
\Tp{hipStream{\Bk}PerThread}，哪几条流之间隐式同步要在本机复测）、打点粒度
（用第~\ref{sub:pm11-resolution}~节的空档实验重标一次）、以及计数器工具与锁频方式
（\Fn{rocprof} 一族与 \Fn{rocm-smi} 的开关名以 --help 为准）。
在 CUDA 上取到的 min 分布形状不能指望在另一家驱动栈重现——那是环境问题，不是算法问题。

\textbf{SYCL / oneAPI。}队列 \Fn{submit} 返回一个 event，它天生就是``这段工作的句柄''；
同步用 \Fn{wait}，要把异步异常拿到手就用 \Fn{wait{\Bk}\_and{\Bk}\_throw}。
建队列时带上 profiling 属性，同一个 event 就能给三个时间戳
（\Tp{command{\Bk}\_submit}、\Tp{command{\Bk}\_start}、\Tp{command{\Bk}\_end}，单位纳秒），
这比 CUDA 的双 event 更直接把``排队''与``执行''拆成两列：第~\ref{sub:pm11-queue}~节要的
三段分账，在这里是两个差值。两个前提要确认：profiling 开关本身可能有代价
（实现可能因此改提交路径，所以带与不带要各跑一次对照）；时间戳起点由实现自定义，
只有差值有意义。换后端时换的是接口，不是数。

\begin{lstlisting}[style=cppstyle, caption={HIP 与 SYCL 的计时边界对照片段：同一套协议，落点不同（形态片段，不含实测值）}, label={lst:pm11-port}]
// --- HIP: same call shape as CUDA, so the harness moves verbatim.
//     Granularity and default-stream behaviour must be re-measured.
hipEvent_t hb, he;
hipEventCreate(&hb);
hipEventCreate(&he);
hipEventRecord(hb, hstream);
launch_kernel(d_in, d_out, n, hstream);
hipEventRecord(he, hstream);
hipEventSynchronize(he);              // as little skippable as in CUDA
float hms = 0.0f;
hipEventElapsedTime(&hms, hb, he);

// --- SYCL: one submit, one handle; three timestamps split queueing
//     from execution, so the two scopes become two separate numbers.
sycl::queue q(sycl::property::queue::enable_profiling{});
sycl::event e = q.submit([&](sycl::handler& h) {
  // kernel body here; its parallel structure is another book's topic
  h.single_task<class pm11_probe>([=]() { /* body */ });
});
e.wait();                          // values are valid only after this
using PI = sycl::info::event_profiling;
auto t0 = e.get_profiling_info<PI::command_submit>();
auto t1 = e.get_profiling_info<PI::command_start>();
auto t2 = e.get_profiling_info<PI::command_end>();
double queued_us = (t1 - t0) / 1000.0;  // ns -> us, queue part
double device_us = (t2 - t1) / 1000.0;  // device segment alone
\end{lstlisting}

\textbf{图形 API 的 GPU 时间戳查询。}Vulkan 用 \Fn{vkCmd{\Bk}Write{\Bk}Timestamp}
把设备时间戳写进 query pool：队列要具备时间戳能力，结果只有在提交完成、缓冲区读回之后才有效，
tick 到纳秒的换算取物理设备报告的周期参数；D3D12 用 \Fn{Write{\Bk}Timestamp} 写 timestamp 查询堆，
频率靠设备侧的 timestamp frequency 查询。共同点决定了口径：\emph{读回 query 缓冲区是一次额外往返}，
它既不能放进被测段，也不能忽略——一帧里放太多打点本身会改变 GPU 的工作形状。
图形侧结论还要连着``GPU busy 时间不等于帧时间''来读：帧时间由呈现与垂直同步决定，
那是第~\ref{ch:pm01-why} 章与第~\ref{ch:pm03-stats} 章的问题域，不是本节这两个数能替代的。

\begin{compare}[title={对比：CUDA event 与通用 GPU 时间戳查询}]
\label{box:pm11-compare-ts}
\small
\begin{tabular}{>{\raggedright\arraybackslash}p{0.17\textwidth}
                >{\raggedright\arraybackslash}p{0.29\textwidth}
                >{\raggedright\arraybackslash}p{0.29\textwidth}}
\toprule
维度 & CUDA event & 通用时间戳查询（SYCL profiling / Vulkan / D3D12） \\
\midrule
时间从哪来 &
设备在流推进到打点处记录，直接给毫秒（\Tp{float}） &
设备计数器 tick 或纳秒，换算要用设备报告的频率或周期 \\
能否拆出排队 &
不能；beg 已把提交前的排队排除在外（第~\ref{sub:pm11-event-api}~节） &
能；三个时间戳给出排队与执行两列 \\
读回成本 &
一次 \Fn{cudaEvent{\Bk}Synchronize} 加一次查询，固定在段右边界 &
通常要等提交完成后读回缓冲区，往返更显式 \\
边界含义 &
``这一条流上这一段工作'' &
``命令缓冲区里两个打点之间''，跨队列跨提交要自己核对 \\
适用问题 &
kernel 级 A/B、min 分布 &
把多个绘制与计算任务画到同一条设备时间线上 \\
\bottomrule
\end{tabular}
\end{compare}

\section{排障与速查}
\label{sec:pm11-trouble}

\begin{warning}[title={常见陷阱：忘了一个同步，测出 0.001\,ms}]
\label{box:pm11-warning-nosync}
最典型的一次：\Fn{cudaEvent{\Bk}ElapsedTime} 返回一个几乎为零的数，或者干脆返回错误。
机制是 end 还没在设备上完成（甚至 beg 都还没打上），你读的是两个未完成打点之间的垃圾值；
异步执行让``忘了同步''在功能上看不出来、在时间上一眼看出——快得不合理。三步判定：
一，补上 \Fn{cudaEvent{\Bk}Synchronize}(end) 再测一次，数字跳回合理量级就是漏同步；
二，把 \Fn{cudaEvent{\Bk}ElapsedTime} 与 \Fn{cudaGetLast{\Bk}Error} 的返回码一起检查，别只看值；
三，与设备侧对账：样本数乘单段时长应当对得上 \Fn{nsys} 时间线上的总占用，
对不上就是边界画错了（计时段多算或少算了拷贝，或者流被默认流串了）。
近亲是``把打印放进计时段''：一次 \Fn{std::cout} 就可能比被测 kernel 更久，
那属于第~\ref{ch:pm15-antipattern} 章的微基准反模式，在 GPU 上只是更容易触发。
\end{warning}

\begin{longtable}{>{\raggedright\arraybackslash}p{0.15\textwidth}
                  >{\raggedright\arraybackslash}p{0.23\textwidth}
                  >{\raggedright\arraybackslash}p{0.24\textwidth}
                  >{\raggedright\arraybackslash}p{0.21\textwidth}}
\caption{GPU 计时前的六项确认（清单先于数字）}\label{tab:pm11-gpu-check}\\
\toprule
确认项 & 为什么要确认 & 怎么确认 & 跳过的典型后果 \\
\midrule
\endfirsthead
\toprule
确认项 & 为什么要确认 & 怎么确认 & 跳过的典型后果 \\
\midrule
\endhead
\bottomrule
\endfoot
口径是否声明 &
设备段、端到端、工具重放三者互不等价 &
表~\ref{tab:pm11-scope} 选一行，写进报告标题 &
拿 event 差声称端到端提速，或反过来 \\
右边界有没有同步 &
\Fn{cudaEvent{\Bk}ElapsedTime} 要求两个打点都已完成 &
段尾有 \Fn{cudaEvent{\Bk}Synchronize}，返回码一并看 &
读到 0.001\,ms 一类的垃圾值 \\
流的选择与默认流 &
legacy 默认流与 blocking 流互相隐式同步 &
被测段放进显式创建的流；同卡有别的流时尤其要 &
把邻居的排队算成自己的 kernel 时间 \\
warmup 是否够 &
context、JIT、库内选择、缓存、爬升都在前几次 &
累积均值按次数看是否平（第~\ref{ch:pm03-stats} 章），记下丢了几次 &
首次样本主导 min，版本对照失真 \\
时钟状态 &
DVFS 让同一段代码在不同频率下跑 &
锁频（先确认支持，清单~\ref{lst:pm11-cmds}）或每样本记一次频率 &
``变快了''其实只是爬频更早 \\
取哪个统计量 &
GPU 噪声是右偏的加法噪声；min 随样本数单调变小 &
固定 $n$，min 与 \Mx{p50}/\Mx{p99} 并列报 &
用 min 判回归，或用均值看不见地板 \\
\end{longtable}

\begin{bestpractice}[title={实践：一个能活下来的 GPU 结论长什么样}]
\label{box:pm11-bestpractice}
六件套：\textbf{口径}（设备段 / 端到端 / 工具重放，选一并给出另两个的数）；
\textbf{环境指纹}（GPU 型号与架构、驱动与工具包版本、是否独占、\Tp{pstate} 与当时 SM 频率、
是否锁频以及锁到什么值）；\textbf{协议}（用哪条流、样本数 $n$、warmup 次数、取 min 还是分位数）；
\textbf{数据}（两个口径并列的两列，不是截图）；\textbf{边界}（这个数在什么条件下不再成立：
并发邻居、别的频率、别的输入尺寸）；\textbf{复现步骤}（一条命令加上它需要的环境设置）。
第~\ref{ch:pm13-ci} 章会把这六件套变成流水线里能检查的字段——那里还要记住一件与本仓库有关的事实：
\emph{没有 GPU 也没有 Nsight 的 runner 跑不了本章任何一条命令}，所以 GPU 基准要么挂在有卡的
专用机器上，要么就从门禁里摘出去；不要为了让它``通过''而把数字写死成常量。
\end{bestpractice}

本章收束成一句与全书同构的话：CPU 上说``测量先于优化''，到了设备侧要改成\textbf{``协议先于测量''}——
先约定谁在等、边界落在哪个调用、时间取自哪个计数器，再谈 min、\Mx{p50} 与加速比。
下一站是第~\ref{ch:pm12-bandwidth} 章，处理这套协议里最贵的一格：数据在哪里、搬它要花多久。

\chapter{带宽与拓扑：显存、PCIe、NUMA 与多卡的数据搬运}
\label{ch:pm12-bandwidth}

第~\ref{ch:pm11-gpu} 章把``时间''搬到了设备侧：什么时候开始算、算到什么时候、谁的时间算数。
本章接手它剩下的那一半问题——\textbf{数据移动本身怎么变成一个可辩护的数}。一条搬运路径上有
四种互相独立的量：链路的能力、平台实际能跑到的能力、你的程序真的搬了多少字节、以及这些字节走了
哪条物理路径（同一桥片还是跨了互连）。前三种决定``带宽够不够''，第四种决定``为什么今天不够''。
接口很清楚：计时协议取自第~\ref{ch:pm11-gpu} 章的 event 语义与第~\ref{ch:pm02-clock} 章的时钟域，
局部性成因归第~\ref{ch:pm06-micro} 章，占用与页驻留归第~\ref{ch:pm07-memory} 章，跨线程共享归
第~\ref{ch:pm09-concurrency} 章，事件可用性归第~\ref{ch:pm05-events} 章，三平台工具与权限并列
放在附录~\ref{ch:pmapA-platform}。与第~\ref{ch:pm09-concurrency} 章的分工要背下来：\textbf{那边
讲同一份数据被多个核抢，这边讲同一份数据离谁远}——两者会同时出现，但判据与修法完全不同。

本章一条硬规矩：所有带宽数字都只写成``规格上限 vs 实测口径''的机制说明，并给出读者在自己机器上
量出这个数的方法。清单里出现的输出一律标注为示意输出。

\section{带宽的三种数：规格上限、平台实测、应用实际}
\label{sec:pm12-three-numbers}

\subsection*{三种数分别是什么}
\textbf{第一种：规格上限}，来自器件手册与 SKU 表，是电气与时序意义上的最好情况。显存的规格
上限是两个量的乘积：位宽（bit）与等效时钟频率。换算成每秒字节数只需一步：位宽除以 8 换成字节，
再乘频率。这里最容易犯的口径错误是\textbf{把预取倍数再乘一遍}——DDR、GDDR、HBM 手册给的频率
通常已经是等效频率（预取与串行化已经算进去了）。主机内存还要再乘一个隐含因子：通道数，它不写在
CPU 型号里，要从平台手册确认再与实测对照。

PCIe 的口径另有一套算法：上限等于每 lane 的速率乘 lane 数，再从 bit 换成字节。两个陷阱要自己
确认而不是背数字。其一，每 lane 速率是``编码前''还是``编码后''取决于该代际的线路编码，各家写法
也不统一，落笔前回到该代规范的定义。其二，lane 数是\textbf{插槽能力}与\textbf{实际协商结果}
两回事：共享 lane 的槽位、转接卡、BIOS 设置都可能让它从 16 降到 8 甚至 4，协商值必须用 \Tp{lspci}
与 sysfs 现读（清单~\ref{lst:pm12-topo}）。另外 PCIe 是全双工链路，收与发各自有能力，``双向相加''
的口径只用来对齐厂商标称，绝不用来解释你的一次单向拷贝——那是典型的口径偷换，
第~\ref{ch:pm15-antipattern} 章把它列成反模式。

\textbf{第二种：平台实测上限}，是这台机器、这个构建、这个线程数下能跑到的最大流带宽。它一定低于
规格上限，而且低多少不是常数。三个前提决定它跑不跑得起来：工作集要\textbf{远大于最后一级缓存}
（否则你量的是末级缓存带宽）；线程数要\textbf{扫一遍}（太少填不满通道，太多在互连上互相踩）；
写路径要\textbf{避免读改写}（普通写先把整行读进来再改，白白吃掉一半读带宽——这正是 STREAM 一族
用函数级封装或临时存储指令的原因）。清单~\ref{lst:pm12-stream} 给一个最小可跑的 triad 探针：
它不给你结论，只给你``我自己机器上的上限''这个输入。

\begin{lstlisting}[style=cppstyle,caption={主机侧流带宽探针（triad）：数组大小、线程数与重复次数都由命令行给，数值由运行者得到},label={lst:pm12-stream}]
// Triad probe: a[i] = b[i] + s * c[i], reported in decimal GB/s.
// Build: g++ -O3 -fopenmp -std=c++17 stream_probe.cpp -o stream_probe
#include <chrono>
#include <cstdio>
#include <cstdlib>
#include <vector>
#include <omp.h>

static double triad(std::vector<float>& a, const std::vector<float>& b,
                    const std::vector<float>& c, float s) {
  const std::size_t n = a.size();
  auto t0 = std::chrono::steady_clock::now();
  #pragma omp parallel for schedule(static)
  for (std::size_t i = 0; i < n; ++i) a[i] = b[i] + s * c[i];
  auto t1 = std::chrono::steady_clock::now();
  double sec = std::chrono::duration<double>(t1 - t0).count();
  // Three arrays are touched: 2 reads + 1 write, i.e. 12 bytes per element.
  return 12.0 * static_cast<double>(n) / sec / 1e9;   // GB/s = 1e9 bytes/s
}

int main(int argc, char** argv) {
  std::size_t n = argc > 1 ? std::strtoull(argv[1], nullptr, 10) : (1u << 26);
  int thr = argc > 2 ? std::atoi(argv[2]) : omp_get_max_threads();
  omp_set_num_threads(thr);
  std::vector<float> a(n), b(n), c(n);
  for (std::size_t i = 0; i < n; ++i) { b[i] = 0.5f * i; c[i] = 2.0f * i; }
  triad(a, b, c, 3.0f);                  // warm-up: page faults + clock ramp
  for (int r = 0; r < 7; ++r)
    std::printf("thr=%d n=%zu run%d %.3f GB/s\n",
                thr, n, r, triad(a, b, c, 3.0f));
  return 0;
}
\end{lstlisting}

\textbf{第三种：应用实际}，是``这次运行真的搬了多少字节除以这段真花了多少时间''。分子要从算法
定义推（这个 kernel 读哪几个数组、写哪几个、每个多大、遍历几次），不要拿``进程驻留量''或
``张量总大小''冒充：那会把没参与这段计算的字节也算进去。分母用第~\ref{ch:pm11-gpu} 章的 event
口径，不要用 \MCpu{}。还有第四种分子可选：剖析器上报的``计数器口径字节数''。它与你自己数的字节
不等价，差别来自访问合并与写回粒度，因此\textbf{逻辑字节数往往低估真流量}。两个都记，差值本身
就是``访问模式有多坏''的证据。

\subsection*{为什么百分之六十是常态而不是失败}
三种数排成一个严格递减的梯队，每一级都在合理地掉，掉的原因不同，能动的部分也不同。
\begin{itemize}
\item 规格到实测：通道交织没打满、刷新与行命中开销、读写方向切换、时序余量、器件温度与频率。
  这一段属于平台，不由你的代码决定。
\item 实测到应用：启动与排队开销没被摊薄（本章第三节的``小拷贝''问题）、访问步长与未合并、
  地址翻译开销、计时区里混进了同步、时钟还没爬满（第~\ref{ch:pm02-clock} 章的调频行为）。
  这一段才是你的代码能动的那部分。
\end{itemize}
所以``利用率百分之六十''本身不是结论。有意义的表述是：应用实际相对于\textbf{你自己机器上的实测
上限}是多少，剩下那部分能不能被上面第二条里的某个具体机制解释。判据的显著性口径回指
第~\ref{ch:pm03-stats} 章：先定``要检出多大的差异''，再决定跑几次。

\subsection*{roofline 怎么读，怎么用}
roofline 是以算术强度为横轴、以可达性能为纵轴（两侧常取对数）的图，只有两条线：斜率等于实测带宽
的斜坡，高度等于实测算力的天花板；交点叫机器平衡点，落在斜坡左边的 kernel 是访存受限，落在右边的
\textbf{不能靠少搬数据变快}。纪律是\textbf{别拿别人的图读自己的代码}：两条线要用你自己的实测带宽
与实测算力画；先做数量级判断，只要差一个量级，结论就不依赖精确图。

\begin{longtable}{>{\raggedright\arraybackslash}p{0.16\textwidth} >{\raggedright\arraybackslash}p{0.22\textwidth} >{\raggedright\arraybackslash}p{0.22\textwidth} >{\raggedright\arraybackslash}p{0.20\textwidth}}
\caption{带宽的几种说法、来源与最常见的偷换}\label{tab:pm12-caliber}\\
\toprule
数是谁的 & 从哪里来 & 什么时候可信 & 最常见的偷换 \\
\midrule
\endfirsthead
\toprule
数是谁的 & 从哪里来 & 什么时候可信 & 最常见的偷换 \\
\midrule
\endhead
\midrule
\bottomrule
\endlastfoot
规格上限 & 器件手册、SKU 表、能力寄存器 & 只做``不可能超过''的上界校验；用于选硬件 & 当成``应该达到''；频率再乘一次预取倍数 \\
平台实测上限 & STREAM 类探针，扫线程与工作集 & 同机同构建同线程数可复现；是利用率的分母 & 换平台换插槽后继续沿用旧值 \\
应用实际 & 自己数的字节除以设备时间 & 唯一能进结论的数，要连字节口径一起报 & 用驻留量或张量大小冒充搬过的字节 \\
计数器口径字节 & 剖析器的流量统计 & 判断访问模式好坏（与逻辑字节对比） & 当成语义字节反过来推算法 \\
显存占用百分比 & \Tp{nvidia-smi} 容量字段 & 判断装不装得下 & 当成带宽利用率 \\
访存忙碌度 & \Tp{nvidia-smi} 利用率类字段 & 判断有没有在访存 & 当成带宽比例，它是时间占比不是速率 \\
多卡聚合带宽 & 每卡各自实测值之和 & 只在没有共享上游链路时近似成立 & 当成单卡的 $N$ 倍，忽略桥片与内存通道共享 \\
\end{longtable}

\begin{warning}[title={把``利用率''当带宽}]
\Tp{nvidia-smi} 与它的 \Tp{dmon} 模式给的访存利用率类字段，口径是\textbf{采样窗口里存在访存 kernel
在跑的时间占比}，与``跑到峰值带宽的百分之几''是两件事；容量那栏既不是带宽也不是延迟。要带宽只有
两条路：自己算字节除以时间，或者读剖析器的设备侧流量计数器并声明它是计数器口径。拿忙碌度宣称
``带宽还有余量''，等于拿``CPU 占用高''宣称``程序很快''。
\end{warning}

\section{算术强度：这个 kernel 该去看带宽还是去看算力}
\label{sec:pm12-intensity}

算术强度是一个比值，写成一句话：\textbf{这个 kernel 每搬运一个字节，做了多少次浮点运算}。它决定
优化方向——低强度的收益来自少搬，高强度的收益来自少算或算得更宽。判据也写成一句话：把算术强度
乘上实测带宽，得到这个强度下最多能有多少运算速率；拿它与实测算力比，小的一方就是瓶颈。这个乘积
不要背，每次现算，因为两个输入都得在你自己机器上取。

\begin{lstlisting}[style=cppstyle,caption={同一段计时同时报出带宽、算力与算术强度；靠改变 work 做差分，看时间听谁的},label={lst:pm12-intensity}]
// One timing region, three reported numbers. The point is the sweep:
// double `work` (more flops, same traffic) and see whether time moves.
#include <chrono>
#include <cstdio>
#include <vector>

struct Probe { double sec; double gbs; double gflops; double ai; };
using Clock = std::chrono::steady_clock;

static Probe run(const std::vector<float>& in, std::vector<float>& out,
                 const std::vector<float>& coef, int work) {
  const std::size_t n = in.size(), t = coef.size();
  double bytes = 8.0 * static_cast<double>(n);          // 1 read + 1 write
  double flops = 2.0 * work * t * static_cast<double>(n - 2 * t);
  auto t0 = Clock::now();
  for (std::size_t i = t; i + t < n; ++i) {
    float acc = 0.0f;
    for (int w = 0; w < work; ++w)
      for (std::size_t k = 0; k < t; ++k) acc += in[i + k] * coef[k];
    out[i] = acc;                    // traffic stays fixed while work grows
  }
  double sec = std::chrono::duration<double>(Clock::now() - t0).count();
  return Probe{sec, bytes / sec / 1e9, flops / sec / 1e9, flops / bytes};
}

int main() {
  const std::size_t n = 1u << 24;
  std::vector<float> in(n, 1.0f), out(n, 0.0f), coef{1.f, 2.f, 3.f, 4.f};
  for (int work : {1, 2, 4, 8, 16}) {
    Probe p = run(in, out, coef, work);
    std::printf("work=%2d ai=%6.2f flop/B  %7.2f GB/s  %7.2f GFLOP/s\n",
                work, p.ai, p.gbs, p.gflops);
  }
}
\end{lstlisting}

\subsection*{判据的三种测法}
\begin{enumerate}
\item \textbf{静态推算}：从循环体数出来。读几个数组、写几个、每次遍历多少元素，除以循环里的
  加减乘除次数（融合乘加按两次运算算）。这法最快，也最容易在``有没有片上复用''上骗自己。
\item \textbf{差分实验}（推荐，因为它不需要任何工具权限）：流量不变把运算量翻倍，或运算量不变
  把流量翻倍，看时间动谁。时间几乎不动的那一侧就是有富余的一侧。清单~\ref{lst:pm12-intensity}
  就是这个实验的形状。两个注意：编译器可能把多余运算折叠掉（要留住不可折叠的依赖）；翻倍运算量
  可能顺带改变向量化与寄存器压力，让实验不再是单因子变化。
\item \textbf{计数器口径}：读设备侧的流量与运算计数，直接得到实际强度。可用事件随平台与权限变，
  发现方法见第~\ref{ch:pm05-events} 章；拿不到是常态（包括本仓库的持续集成环境），不是异常。
\end{enumerate}

\subsection*{融合为什么在带宽受限下收益大}
融合改变的是\textbf{中间结果的搬运}：原本``写出去再读回来''的字节，现在留在寄存器或片上存储里。
于是在斜坡段上字节数直接少一截，时间近似按字节比例下降；在天花板段上你省的是本来就不缺的东西，
收益趋零甚至为负——融合通常\textbf{提高}算术强度，同时带来寄存器压力上升、占用率下降、循环结构
变复杂，这些代价在计算受限的一侧会直接落到时间上。占用率与 kernel 时间不是线性关系
（第~\ref{ch:pm11-gpu} 章的口径），别用``占用率低了''单独解释变慢。

正确的验证是一次\textbf{三格差分}：融合前、融合后、以及融合后但把中间数组的读写重新暴露出来
（模拟未融合的流量）。三格同机同构建交替跑（第~\ref{ch:pm03-stats} 章的 A/B/A/B）。若第三格与
第一格接近、第二格明显更快，收益确实是搬运省下来的；若第二格没差别，你原来的 kernel 本来就不在
斜坡上。

\begin{guideline}[title={动布局之前先定侧}]
\begin{itemize}
\item 先量算术强度与两个实测上限，再决定改数据布局（少搬、改步长、改分块）还是改计算（少算、
  向量化、换算法）。方向错了，改动越大越慢。
\item 分块的收益只在斜坡段成立：它减少的是重复搬运，不减少运算量。块大小要按你的片上存储容量
  与工作集推，不要抄别人的常数。
\item 主机侧与设备侧的两个上限不同源：同一个 kernel 在 CPU 上是访存受限，搬到加速卡上可能越过
  平衡点。换了执行单元就重算强度，结论不许跟着搬。
\end{itemize}
\end{guideline}

\section{PCIe 与主存—显存搬运}
\label{sec:pm12-pcie}

\subsection*{一次拷贝里到底有几段}
一条主机到设备的路径至少四段：应用侧调用与参数校验、驱动把请求排进队列、直接存储访问引擎从物理
页搬、完成通知回到计时体系。\Fn{cudaMemcpy} 的同步语义把四段全算进计时区；\Fn{cudaMemcpy{\Bk}Async}
只有在\textbf{源是页锁定内存}、并且入队后不立刻同步的前提下，才真的能把 CPU 侧那段与搬运重叠。
源不是页锁定时，驱动要先做一次临时的锁定拷贝，异步接口于是多出一段你没打算测的 CPU 侧搬运——
这就是``我用了异步为什么没变快''的头号原因。

页锁定用 \Fn{cudaHostAlloc}（或 \Fn{cudaMallocHost}）拿；写合并类标志改变的是\textbf{写}方向的
行为，两者要分开测。它为什么快：物理地址稳定，省掉驱动侧一次暂存拷贝，并且允许引擎连续读。它不是
白拿的：锁页内存不可换出，分配本身远慢于常规分配，过量锁定还会伤整台机器（占用口径，回指
第~\ref{ch:pm07-memory} 章）。\textbf{分配不许进计时循环}，把它放在 warmup 里，计时区只包搬运。

\begin{lstlisting}[style=cppstyle,caption={单段主机到设备拷贝的带宽测量骨架：计时边界与 event 语义取自第~\ref{ch:pm11-gpu} 章},label={lst:pm12-h2d}]
// Measures one host-to-device leg of a pinned buffer over a size sweep.
// Event semantics (record / synchronize / elapsed) are covered in ch 11.
#include <cstdio>
#include <cuda_runtime.h>

#define CHECK(call) do {                                          \
    cudaError_t e_ = (call);                                      \
    if (e_ != cudaSuccess) {                                      \
      std::printf("cuda error: %s\n", cudaGetErrorString(e_));    \
      return 0.0;                                                 \
    }                                                             \
  } while (0)

static double h2d_gb(void* dst, void* src, size_t bytes,
                     cudaStream_t s, cudaEvent_t e0, cudaEvent_t e1) {
  CHECK(cudaMemsetAsync(dst, 0, bytes, s));      // touch pages before timing
  CHECK(cudaEventRecord(e0, s));
  CHECK(cudaMemcpyAsync(dst, src, bytes, cudaMemcpyHostToDevice, s));
  CHECK(cudaEventRecord(e1, s));
  CHECK(cudaEventSynchronize(e1));
  float ms = 0.0f;
  CHECK(cudaEventElapsedTime(&ms, e0, e1));
  if (ms <= 0.0f) return 0.0;                    // never divide by a zero
  return (double)bytes / (ms * 1e-3) / 1e9;      // GB/s
}

int main(void) {
  void* dst = NULL;  CHECK(cudaMalloc(&dst, 1u << 28));     // 256 MiB
  cudaStream_t s;  CHECK(cudaStreamCreate(&s));
  cudaEvent_t e0, e1;  CHECK(cudaEventCreate(&e0)); CHECK(cudaEventCreate(&e1));
  for (size_t kb = 4; kb <= (1u << 16); kb *= 2) {
    size_t bytes = kb * 1024;
    void* pin = NULL;
    CHECK(cudaHostAlloc(&pin, bytes, cudaHostAllocDefault));
    double best = 0.0;
    for (int r = 0; r < 10; ++r) {
      double v = h2d_gb(dst, pin, bytes, s, e0, e1);
      if (v > best) best = v;              // max bandwidth = min time view
    }
    std::printf("size %8zu KiB  best %.3f GB/s\n", kb, best);
    CHECK(cudaFreeHost(pin));              // alloc cost stays outside
  }
  CHECK(cudaFree(dst));  CHECK(cudaStreamDestroy(s));
  return 0;
}
\end{lstlisting}

\subsection*{小拷贝是启动问题，不是带宽问题}
把 GB/s 对拷贝尺寸作图（两侧取对数）会得到一个稳定形状：小尺寸处近似直线上升（时间几乎与字节
无关，排队与启动开销主导），中段弯曲，大尺寸处平台化（撞上平台实测上限）。这条曲线给你的不是
``带宽多少''，而是\textbf{拐点位置}：小于拐点的多次独立拷贝先谈合并提交与批量，大于拐点的才谈
带宽与拓扑。取``最快的一次''而不是平均值是设备侧的常见口径（原因见第~\ref{ch:pm11-gpu} 章：噪声
主要来自调度），但报数必须写清统计量，并按第~\ref{ch:pm03-stats} 章定重复次数。

单向与双向分开报：同一方向的多路拷贝会争同一批引擎，两个方向同时跑得到的和通常不等于单向的两倍；
跨节点时还叠加第~\ref{sec:pm12-numa} 节的距离惩罚。

\subsection*{拓扑与链路宽度要自己读}
\begin{lstlisting}[style=shellstyle,caption={搬运路径的拓扑确认：链路协商值、设备亲和性、节点距离与容器实际可用节点},label={lst:pm12-topo}]
# BDF = PCI address of the device, e.g. as printed by nvidia-smi
sudo lspci -vv -s "${BDF}" | grep -E 'LnkCap:|LnkSta:'   # capability vs status
cat /sys/bus/pci/devices/${BDF}/current\_link\_speed
cat /sys/bus/pci/devices/${BDF}/current\_link\_width
cat /sys/bus/pci/devices/${BDF}/numa\_node               # -1 often means not exposed

nvidia-smi topo -m                          # affinity matrix, see the sample below
nvidia-smi --query-gpu=name,pci.bus\_id,memory.total --format=csv
nvidia-smi dmon -s u -s m -c 20             # utilization fields, NOT GB/s

numactl --hardware                          # nodes, sizes, distance matrix
lstopo --of txt topo.txt                    # full enumerated hardware tree
lscpu | grep -E 'NUMA|Socket|Model name'
cat /sys/fs/cgroup/cpuset.cpus.effective    # what this container is allowed
\end{lstlisting}

\begin{lstlisting}[style=outstyle,caption={示意输出（格式示例，非本仓库实测值）：设备亲和性矩阵的形状与字段读法},label={lst:pm12-out-topo}]
# illustrative shape only, not a measurement from this repository
        GPU0  GPU1  GPU2  GPU3  CPU Affinity    NUMA Affinity
GPU0     X    NV12   SYS   SYS   0-15            0
GPU1    NV12    X    SYS   SYS   0-15            0
GPU2     SYS   SYS    X    NV4   16-31           1
GPU3     SYS   SYS   NV4    X    16-31           1

Legend:  NV#  = number of inter-GPU links
         SYS  = traffic goes through system memory
         PHB  = through the host bridge of the root complex
         PIX  = a single PCIe bridge,  PXB = multiple PCIe bridges
\end{lstlisting}

矩阵怎么读：\textbf{同一 root complex 之内}（PIX、PXB、PHB）意味着两次搬运可以不上到主机内存；
\textbf{SYS} 意味着流量要穿过主机内存与互连，这一步同时吃掉带宽预算和距离预算。同一列还告诉你
``该把缓冲区分配在哪儿''——与设备同一节点的 CPU 核才是页锁定缓冲的落点。\Tp{lstopo} 的树更完整，
但它画的是\textbf{枚举到的}拓扑，协商降级仍要回 sysfs 对。访问控制重定向与输入输出内存管理会把
设备间直连改成绕经桥片的通路：矩阵名义可达，实测只有远端水平，本章第五节展开。虚拟化与直通模式下
你测到的是虚拟拓扑给的能力，不是物理插槽的能力。

\begin{warning}[title={用 CPU 时间证明``拷贝变快了''}]
异步接口下 CPU 只做入队，\MCpu{} 变短是必然的，与搬运快慢无关。要判断变没变快，计时区必须
\textbf{包住最终同步}（或用设备侧时间戳），并且把``CPU 侧花了多少''与``设备队列排了多久''分成
两笔记。否则你测到的是交付被推迟，不是吞吐被提高——这是第~\ref{ch:pm09-concurrency} 章
``延迟被搬进队列''在搬运上的版本。
\end{warning}

\section{NUMA：距离矩阵、绑定策略与``绑了核没绑内存''}
\label{sec:pm12-numa}

\subsection*{距离矩阵怎么读}
\Tp{numactl{\Bk} --hardware} 给每个节点的 CPU 列表、内存大小与一个距离矩阵；\Tp{lstopo} 画同样的
结构。读法三条纪律。第一，\textbf{距离是相对量纲}，是平台给的排序与代价刻度，不是纳秒也不是倍数；
本地与远端的真实倍数要用下面的四格实验自己量。第二，矩阵的\textbf{非对称}要照实记录：有些平台从
$A$ 到 $B$ 与从 $B$ 到 $A$ 不等，结论要按方向分开。第三，\textbf{设备也是节点邻居}：GPU 与网卡
挂在某个节点上（sysfs 的 \Tp{numa{\Bk}\_node} 字段），跨节点访问设备的惩罚是两层叠加。

\subsection*{为什么绑了核还是慢}
默认内存策略是``在哪个 CPU 上落地就尽量分配到本地节点''，但\textbf{落地时刻}与\textbf{落地线程}
未必是你以为的那个。三种典型错位：
\begin{itemize}
\item \textbf{首次触碰效应}：大数组的页在第一次被写时才真正分配，谁先碰页就归谁的节点。如果
  初始化线程与计算线程绑在不同节点，计算就在访问远端页。
\item \textbf{本地耗尽}：本地节点内存水位已紧，策略允许回落到其他节点，页散落远端。慢的是逐页
  而不是整体，所以看起来``时快时慢''。
\item \textbf{父进程先分配}：绑定发生在创建线程之后，而页早已在父进程所在节点落地，子线程绑得
  再准也是在读远端。
\end{itemize}
三种都能确认而不是猜：\Tp{/proc/{\Bk}pid{\Bk}/numa{\Bk}\_maps} 逐区域列出页的节点分布；
\Fn{move{\Bk}\_pages}（或 \Tp{migratepages} 工具）能把已知区域迁走再做\textbf{迁移前后差分}；
\Fn{get{\Bk}\_mempolicy} 读当前线程的策略。伪共享是同一缓存行上的跨线程流量，与距离无关：
\Tp{numa{\Bk}\_maps} 里全是远端页就先修距离，行级共享明显就先按第~\ref{ch:pm09-concurrency} 章
修共享。

策略入口分三层：进程级 \Tp{numactl}（\Opt{cpunodebind}、\Opt{membind}、\Opt{interleave}），
线程与区域级 \Fn{set{\Bk}\_mempolicy} 与 \Fn{mbind}（绑定、优先、交织三种语义各不相同），容器层
\Tp{cpuset.cpus} 与 \Tp{cpuset.mems}（决定你真的看得见几个节点）。交织策略提高的是\textbf{聚合
带宽}、代价是平均延迟，只在工作集巨大且访问均匀时划算，不要当默认值用。

\begin{lstlisting}[style=shellstyle,caption={四格对照实验：同一二进制在两种 CPU 绑定与两种内存落点下各跑一遍；远端代价是比值不是绝对值},label={lst:pm12-numa}]
BIN=./stream_probe
N=$((1 << 26)); T=8
for cpu in 0 1; do
  for mem in 0 1; do
    out="bw-c${cpu}m${mem}.txt"
    numactl --cpunodebind=${cpu} --membind=${mem} ${BIN} ${N} ${T} > "${out}"
  done
done
# c0m0 vs c0m1 : the memory-distance penalty seen from node 0
# c0m0 vs c1m0 : the cost of the CPU crossing the interconnect
# c1m1         : what you get when nobody configured anything
grep -h 'GB/s' bw-c*m*.txt      # report medians; the test itself is ch 3
\end{lstlisting}

\begin{note}[title={设备侧的 NUMA 账要一起记}]
页锁定缓冲落在远端节点时，一次主机到设备的搬运要先在 CPU 互连上走一段再上设备链路。实测表现是
``同样尺寸，带宽低一档''，而且降低幅度与当时跨互连的流量相关，不是一步台阶。控制办法是把
\textbf{分配线程也绑到设备同侧节点}，而不是只绑计算线程；然后在四格之外再加一格``同一设备、
缓冲换节点''的对照，把这一项单独量出来。
\end{note}

\section{多卡：可达性、争用，以及吞吐为什么不等于单卡的 $N$ 倍}
\label{sec:pm12-multigpu}

\subsection*{可达性是问出来的，不是推出来的}
\Fn{cuda{\Bk}Device{\Bk}CanAccessPeer} 回答``这两个设备之间有没有直接通路''，由硬件与拓扑决定；
返回不可达就只能经主机中转，那是两跳搬运、两次占用带宽。可达不等于启用：
\Fn{cuda{\Bk}Device{\Bk}Enable{\Bk}PeerAccess} 是\textbf{按对、按方向}的显式动作，还要求两端
上下文已经建立（口径见第~\ref{ch:pm11-gpu} 章）。矩阵里显示直连的平台也可能因为访问控制重定向或
内存管理配置而实际绕远。所以顺序永远是：读矩阵、查可达性、\textbf{实测一次设备间拷贝}、再与
``经主机中转的同一尺寸''对照——三个数一起报，才知道直连值不值。

\subsection*{吞吐掉下来的三种分解}
\begin{itemize}
\item \textbf{共享上游}：聚合带宽的上限不是单卡上限乘 $N$，因为桥片、内存通道、交换结构的容量
  是共享的，卡越多每张分到的越少。判据：从单卡开始逐次加卡，看聚合曲线的\textbf{增量}是否递减，
  同时记录每卡各自的实现带宽。
\item \textbf{同步}：一轮要等所有卡，于是每轮耗时由最慢那张决定。总时间看不出来，必须保留
  \textbf{每卡各自的时间序列}，分别报``每卡的中位数与尾部''和``整轮的尾部''：三者口径不同，不能
  互相替代（分位数口径见第~\ref{ch:pm03-stats} 章）。
\item \textbf{负载不均}：切分不对称、显存容量不同导致可用工作集不同、拓扑不对称导致只有某些设备
  对能直连。判据是把每张卡的字节数与时间并列画出来，而不是只看总吞吐。
\end{itemize}
还有常被忽略的第四项：\textbf{通信与计算争同一份资源}。同卡的拷贝引擎与计算单元争访存带宽，
集合通信的流量与 kernel 争链路和互连，表现是``没改 kernel，kernel 时间自己变长了''。归因办法是把
通信单独计时（独立流、独立 event），与 kernel 计时分成两笔账，再看两者的重叠区间；只报一个总时间
就等于把这条线索彻底埋掉。

\begin{compare}[title={多卡常见的四种错误归因与对应证据}]
\begin{tabular}{>{\raggedright\arraybackslash}p{0.20\textwidth} >{\raggedright\arraybackslash}p{0.24\textwidth} >{\raggedright\arraybackslash}p{0.30\textwidth}}
\toprule
看到的现象 & 别急着说 & 需要的证据 \\
\midrule
加卡后每卡带宽降 & ``卡在抢处理器'' & 聚合曲线增量，加上游链路的能力与实测值 \\
某轮突然变慢 & ``有卡坏了'' & 每卡时间序列的尾部，加上该轮的通信计时 \\
直连已启用却不见快 & ``互联没用'' & 直连实测对同尺寸经主机中转实测的差值 \\
kernel 时间漂长 & ``代码退化了'' & 通信与 kernel 的重叠区间，两笔分账 \\
\bottomrule
\end{tabular}
\end{compare}

\section{非 GPU 的异构对照：加速器、SoC 与搬运的计时边界}
\label{sec:pm12-hetero}

不是只有 GPU 有``另一侧''。专用加速器、可编程逻辑卡、嵌入式 SoC 的固定功能单元都有各自的搬运语义
与计时源，而\textbf{跨时钟域比较只在同一时间戳体系内成立}这条纪律（第~\ref{ch:pm02-clock} 章与
第~\ref{ch:pm11-gpu} 章）一字不改地适用：设备侧时间戳自洽，主机侧的``提交返回''不含执行、
``等待返回''含排队与调度，能跨域比较的只有同一体系内两个时间戳之差。

嵌入式 SoC 的区别更根本：CPU 与加速器\textbf{共用同一份内存和同一个内存控制器}，没有独立的
``显存带宽''这一层。于是带宽账只有一份，加速器跑满就是在抢 CPU 的带宽，两个``实测上限''必须同时量
才解释得清，``少搬数据''的收益不再来自链路而来自内存控制器。缓存一致性边界决定搬运的真实代价：
共享缓冲区究竟是一致性的（软件不必显式维护）还是非一致性的（要做清理与无效化，或者用非缓存映射），
差别直接体现为搬运字节数与 CPU 侧时间。这一层的机制、保留内存区与描述符队列属于《硬件访问与内核》
分册；本章只负责把它们变成可测的对象：\textbf{搬运的起止边界画在哪儿、用谁的时间戳、搬了多少字节}。

\begin{compare}[title={四种``搬运''的计时边界该怎么画}]
\begin{tabular}{>{\raggedright\arraybackslash}p{0.18\textwidth} >{\raggedright\arraybackslash}p{0.26\textwidth} >{\raggedright\arraybackslash}p{0.30\textwidth}}
\toprule
搬运类型 & 边界画在哪 & 为什么不能画在别处 \\
\midrule
主机到设备 & 设备侧流内两个 event 之间，字节按请求尺寸计 & 包住调用只量到 CPU 入队；漏掉最终同步就量到``推迟'' \\
跨节点访问同一页 & 没有``一次搬运''这个事件：按页驻留记账，比较两种绑定配置的耗时 & 内存访问不返回时间戳，只能做配置差分，并用 \Tp{numa{\Bk}\_maps} 证明页在哪 \\
设备间直连 & 两端各取一对 event，只比同一次传输的两侧 & 只在一端计时会把另一端的排队算进这一端的链路 \\
SoC 共享内存 & 显式维护与提交各自计时，再看完成信号 & 一致性区域内看不见的搬运会把边界抹平，归因随之失败 \\
\bottomrule
\end{tabular}
\end{compare}

\begin{guideline}[title={搬运预算的记账法}]
\begin{itemize}
\item 把每条路径拆成\textbf{链路段}（有速率上限）与\textbf{排队段}（有启动开销），两段分别计时；
  小尺寸的问题归排队段，大尺寸的问题归链路段。
\item 链路段的分母用\textbf{你自己机器上的实测上限}；规格上限只做上界校验，不用来判定``差多少
  就是缺陷''。
\item 聚合口径（多卡、多方向、多进程）要注明是否共享同一上游资源；不写这一句，数字就没有可比性。
\end{itemize}
\end{guideline}

\section{排障与速查}
\label{sec:pm12-trouble}

\begin{longtable}{>{\raggedright\arraybackslash}p{0.20\textwidth} >{\raggedright\arraybackslash}p{0.19\textwidth} >{\raggedright\arraybackslash}p{0.20\textwidth} >{\raggedright\arraybackslash}p{0.19\textwidth}}
\caption{搬运场景排障表：先看什么、常见误判、一句话处置}\label{tab:pm12-transfer}\\
\toprule
搬运场景 & 先看哪个指标 & 常见误判 & 一句话处置 \\
\midrule
\endfirsthead
\toprule
搬运场景 & 先看哪个指标 & 常见误判 & 一句话处置 \\
\midrule
\endhead
\midrule
\bottomrule
\endlastfoot
主机到设备带宽低 & 协商速率与 lane 数、是否页锁定 & 说``显存慢''，把链路问题当器件问题 & 先确认路径能力，再确认源内存是锁定的 \\
用了异步没变快 & CPU 入队时间与队列深度分开看 & 以为异步自动重叠 & 源改成页锁定，计时区包住最终同步 \\
小拷贝带宽极低 & GB/s 对尺寸的曲线与拐点 & 拿小拷贝算利用率 & 合并提交；小尺寸只报延迟不报带宽 \\
利用率报不出数 & 忙碌度字段的口径 & 把忙碌度当带宽 & 自己算字节除以时间，或读计数器口径 \\
想减少搬运 & 算术强度与实测带宽的乘积 & 见慢就融合 & 先定瓶颈侧，再用三格差分验收益 \\
多卡吞吐不升 & 聚合曲线增量与每卡时间序列 & 归因给``驱动不行'' & 查共享上游、同步等待、负载不均三处分解 \\
直连可达却不见快 & 拓扑矩阵与直连对中转的实测差 & 认为可达即高速 & 查访问控制与内存管理配置，按实测差决定用不用 \\
绑核后仍慢 & 页的节点分布 & 以为绑核等于绑内存 & 用内存绑定策略把页落对，再跑四格对照 \\
同机两次差很多 & 频率、调频策略、内存水位、邻居负载 & 归因给代码 & 记环境指纹（第~\ref{ch:pm02-clock} 章）后再判，必要时重建基线 \\
换机器结论就变 & lane 数、节点数、设备亲和性 & 沿用旧的``上限'' & 上限属于机器：换平台就重测，结论里写清是哪台 \\
跨设备比时间 & 两端时间戳是否同源 & CPU 时间减设备时间 & 只在同一时间戳体系内做差（第~\ref{ch:pm11-gpu} 章） \\
\end{longtable}

\begin{bestpractice}[title={把``数据在哪''写进结论}]
一个能活的搬运结论必须带五件东西，缺一件就退化成口号：
\begin{enumerate}
\item \textbf{口径}：这条时间是 \MWall{} 还是设备侧时间，字节是逻辑字节还是计数器口径字节；
\item \textbf{路径}：数据从哪个节点的哪种内存出发，经过哪几段链路，落在哪里——一句话拓扑描述，
  例如``远端节点的页锁定内存，经主机桥到设备''；
\item \textbf{三个数各报哪个}：规格上限只作上界，实测上限作分母，应用实际作结论；
\item \textbf{环境指纹}：协商速率与 lane 数、节点与距离矩阵、亲和性矩阵、设备型号，全部现读，
  不抄别人的；
\item \textbf{复现步骤}：清单~\ref{lst:pm12-stream}、~\ref{lst:pm12-h2d}、~\ref{lst:pm12-numa}
  里哪几条命令、哪个绑定配置、重复几次。
\end{enumerate}
``带宽够了''不是结论；``在某节点绑定下，这段搬运的实测带宽达到本机 triad 上限的某个比例，剩余
差距由未合并访问解释，证据是计数器字节比逻辑字节多出某个倍数''才是结论。数字由你测，句式照这个写。
\end{bestpractice}

\begin{warning}[title={``拓扑我不知道，先按理想算''}]
本章所有结论都挂在路径上，而路径\textbf{读得到}：协商速率、lane 数、节点距离、设备亲和性、可达
性矩阵都能在 sysfs、\Tp{lspci}、\Tp{numactl}、\Tp{lstopo}、\Tp{nvidia-smi} 里现读。按理想拓扑
算出的``可用带宽''一旦与实测差出量级，你会把配置问题（共享 lane、转接卡、容器只给了一个节点、
重定向）误诊成代码问题，然后花两周去改一个本来就不在瓶颈上的 kernel。开测第一步是把这些读数抄进
环境指纹，而不是先跑基准。
\end{warning}

\begin{note}[title={本章的数字从哪来}]
本章一个实测数都没给。清单~\ref{lst:pm12-stream}、~\ref{lst:pm12-intensity}、~\ref{lst:pm12-h2d}
只定义测量形状，数值由运行者得到；清单~\ref{lst:pm12-out-topo} 是\textbf{示意输出}，只教你怎么读列；
量级关系一律写成``谁高于谁''与``为什么掉一档''；三平台权限对照回指附录~\ref{ch:pmapA-platform}。
\end{note}

\part{持续性能：从 CI 门禁到生产采样}

\chapter{连续性能：CI 基线、回归检测与阈值设定}
\label{ch:pm13-ci}

前面十二章解决的是\textbf{怎么拿到一个数}：选哪个时钟、重复多少次、用什么事件、怎么把指标送回代码位置。从本章开始，问题换成另一个形状——\textbf{这个数怎么活下来}。它要每天在别人的机器上被重算一遍，要在半小时内的流水线里给出结论，要在合并请求上说出``可以''或``不行''，还要在三个月后换了一台机器时不至于把整个基线拖崩。

本章解决三件事：其一，CI 环境给出的性能数天然缺失哪些东西，为什么这些缺失决定了 CI \textbf{只能用来拦回归、不能用来报绝对值}；其二，判据怎么从``别人仓库里的那个百分之五''变成``你自己这条流水线跑出来的分布''；其三，门禁挂在流水线的哪一步、失败时留什么证据。与 第~\ref{ch:pm03-stats} 章 的接口是统计判据本身（区间、bootstrap、重复次数下限），本章只讲它在 CI 约束下的退化形态与工程封装；与 第~\ref{ch:pm14-production} 章 的接口是``拿不到 PMU 的那一侧''该怎么办；与 第~\ref{ch:pm15-antipattern} 章 的接口是``一个错误的结论怎么被自动化地复制一千次''。

口径先钉住（见第~\ref{ch:pm01-why} 章）：本章所谓``回归''默认指 \MWall{} 或 \MThr{} 上的分布位移，不是 \MCpu{}，更不是 \Mx{IPC}。共享 runner 上恰恰是后两者最不可信。

\section{CI 上的性能数为什么天然不可信}
\label{sec:pm13-untrust}

一台流水线里的机器和一个工程师自己的工作站，差别不在``快慢''，而在\textbf{可归因性}。CI 环境通常同时缺下面五样东西。

\textbf{其一，独占的计算资源。} 托管 runner 是虚拟机或容器，vCPU 由宿主调度，邻居负载不可见也不可拒绝。一个基准在\textbf{同一份代码、同一个镜像}上拿到两个相差一到两倍的 \MWall{}，在共享实例上是常态而不是异常。抢占时间、CPU 信用额度（burstable 实例积分耗尽后被限速）、宿主的实时迁移，都会把\textbf{时间}搬进你的测量段，而你的程序一行也没多跑。

\textbf{其二，频率与温度的可控性。} 第~\ref{ch:pm02-clock} 章 里那套``把调频策略设成 \Tp{performance}、关掉 \Tp{no{\Bk}\_turbo}、等热状态稳定''的手法，在虚拟机里通常\textbf{没有权限}，甚至连 \Tp{/sys/devices/{\Bk}system/cpu/} 下的 \Tp{cpufreq} 目录都不存在。你能观察到的只有``这台机今天比昨天慢''，无法施加任何影响。

\textbf{其三，PMU 权限。} 这是本章最重要的一条。共享 runner 上 \Tp{perf{\Bk}\_event{\Bk}\_paranoid} 通常被设成 $2$ 或更高，容器运行时还叠了 seccomp 与 cgroup 策略；虚拟化层往往根本不向 guest 暴露性能计数器。后果是：硬件事件（\Tp{cycles}、\Tp{instructions}、\Tp{cache{\Bk}-misses}）返回 \Tp{<not counted>} 或 \Tp{<not supported>}，\Fn{perf} 本身在很多 runner 镜像里\textbf{没有安装}。也就是说 第~\ref{ch:pm05-events} 章 里那一整套微架构分层，在 CI 上\textbf{只剩软件事件可用}（\Tp{cpu-clock}、\Tp{task{\Bk}-clock}、\Tp{page{\Bk}-faults}、\Tp{context{\Bk}-switches}）。

\textbf{其四，构建缓存的冷热不一致。} 第一次跑是全新编译，之后是增量；\Tp{ccache} 命中与否改变\textbf{编译时间}，链接顺序与 \Tp{PGO}/\Tp{BOLT} profile 是否新鲜改变\textbf{代码布局}。同一个 commit 在两个 job 里可能得到两个不同的二进制，而布局差异足以让 \MWall{} 动几个百分点（细节口径见第~\ref{ch:pm10-compiler} 章）。

\textbf{其五，时间基准本身。} 虚拟机里的 \Fn{rdtsc} 需要你自己确认 \Tp{constant{\Bk}\_tsc} 与跨核同步性，\Tp{/proc/cpuinfo} 的标志位在 guest 里未必可信；拿到\textbf{负耗时}或``一次跑三纳秒''的样本，通常是时钟源选错而不是程序真的那么快。

\begin{guideline}[title={CI 性能数的正确用法}]
\label{box:pm13-usage}
\begin{itemize}
\item \textbf{适合}：在同一 job、同一 runner、同一镜像内，对 base 与 head \textbf{两个二进制做相对比较}；对同一指标做\textbf{长序列的趋势观察}；拦住``慢了一个台阶''级别的粗回归。
\item \textbf{不适合}：报绝对值（``我们的解析器每秒若干兆字节''）、比较两个不同指标的绝对量级、做微架构归因、给一次优化``验效''。这些要在可控机器上做（第~\ref{sec:pm13-env} 节）。
\item \textbf{阈值来源}：只能来自你自己流水线上的\textbf{空跑分布}（见 第~\ref{sec:pm13-verdict} 节 的标定函数）。抄别人仓库的百分比，等于把你的噪声当成别人的噪声。
\end{itemize}
\end{guideline}

本仓库的 GitHub Actions runner 就属于上面这一类：没有 PMU 权限，通常也没有 \Fn{perf}。因此\textbf{本册所有清单里的工具输出都只是形态示意}，不是本仓库的实测记录。下面这段就是``没有权限''时报告长什么样，读它的方式是\textbf{看形状、不看数值}。

\begin{lstlisting}[style=outstyle, caption={形态示意（非本仓库实测值）：无 PMU 权限环境下
硬件事件整列为 not counted，只剩时间行可读；数值一律留白}, label={lst:pm13-notcounted}]
$ perf stat -e cycles,instructions,cache-misses ./build/release/bench
 Performance counter stats for './build/release/bench':

     <not counted>      cycles
     <not counted>      instructions
     <not counted>      cache-misses

     <value omitted>    seconds time elapsed
     <value omitted>    seconds sys

# 若 perf 根本没装，症状更简单：command not found。
# 不要把 not counted 读成"这个程序没有缓存未命中"：那是权限缺失，不是测量结果。
\end{lstlisting}

还有一个容易忽略的失真源：\textbf{剖析器自身的开销}。托管 runner 上跑 \Fn{perf} \Fn{record} 会把 CPU 税压在本已拥挤的实例上，采到的热点是失真的（第~\ref{ch:pm04-sampling} 章）。CI 里如果需要栈，优先用应用内自采或事件时间线，把火焰图留给可控机器与 第~\ref{ch:pm14-production} 章。

\section{基线的四种存法与各自的失效模式}
\label{sec:pm13-baseline}

``基线''不是一个数，是一个\textbf{存储策略}。四种常见存法的代价与偏差方向差别很大，选错会让门禁系统性地往一边漏。

\textbf{（一）历史序列。} 每次流水线跑完把样本追加进时序库或 CSV，判据用最近 $N$ 次滚动窗口（或按发布版分段）。优点是自动跟随硬件与编译器的漂移；缺点是\textbf{漂移会被吸收成基线}——性能每天慢千分之一，一个月后基线已经慢了三成，门禁依然全绿。这是历史序列最典型的漏报方向。

\textbf{（二）固定阈值。} 一个写进仓库的 JSON：每个用例一个上限时间。优点是判据可审计、不随环境漂；缺点是\textbf{阈值一开始就是猜的}，换一台机器就全体失真。固定阈值必须绑定环境指纹（第~\ref{sec:pm13-env} 节），否则它只是在等一次误报潮。

\textbf{（三）上一发布版。} 基线取某个发布标签那次构建的结果，每个 PR 与之比。优点是有``相对于产品''的语义，能挡住慢性漂移；缺点是发布版结果本身可能来自另一台机器，\textbf{不可比}，而且每个 PR 都跟同一个点值比，等于把一次测量的方差摊到一千次判断上。

\textbf{（四）同机自比。} 在同一个 job 里同时构建 base（目标分支）与 head，交替执行 A/B/A/B，只比这次运行内部的两组样本。优点是把机器、温度、邻居负载都消成公共项，是共享 runner 上\textbf{唯一站得住}的做法；代价是每个 PR 跑两倍时间，且``相对发布版慢了多少''这类绝对叙述无从谈起。

\begin{longtable}{|>{\raggedright\arraybackslash}p{0.15\textwidth}|>{\raggedright\arraybackslash}p{0.16\textwidth}|>{\raggedright\arraybackslash}p{0.25\textwidth}|>{\raggedright\arraybackslash}p{0.22\textwidth}|}
\caption{基线存储策略对照：成本、漏报方向（该红的不红）与误报方向（不该红的乱红）}
\label{tab:pm13-baseline}\\
\hline
\textbf{策略} & \textbf{成本} & \textbf{漏报方向} & \textbf{误报方向}\\
\hline
\endfirsthead
\hline
\textbf{策略} & \textbf{成本} & \textbf{漏报方向} & \textbf{误报方向}\\
\hline
\endhead
历史序列（滚动窗口） & 每次跑完自动更新，边际成本近零；需要存储 &
慢性漂移被吸收进基线，单次小回归始终在窗口内 &
一次异常慢的邻居负载污染窗口，之后几次都被比成``变快''\\
\hline
固定阈值 & 便宜，判定是纯函数；需要人工维护 &
换硬件或换编译器后阈值偏松，回归直接过关 &
冷缓存的首次运行、编译时间混进测量段\\
\hline
上一发布版 & 中等；发布流程要产出可复用的基线产物 &
发布版本身带回归时，之后所有人都``比它快'' &
发布版跑在另一台机器或另一镜像，系统性差异被当成回归\\
\hline
同机自比（A/B/A/B） & 最高：每个 PR 两倍执行时间与两倍构建 &
只挡这次运行内的相对变差，跨版本慢性漂移看不见 &
两组样本量不对称，或交替顺序被打断（base 阶段被抢占）\\
\hline
\end{longtable}

实践中建议\textbf{两轨并存}：PR 门禁走同机自比（拦突变），夜间流水线走历史序列加固定阈值（拦慢性漂移并给出绝对趋势）。两条轨道的判据与告警等级分开，不要合成一个红灯。

\begin{note}[title={为什么``每个 PR 都跑全量基准''会退化成噪声计数器}]
\label{box:pm13-noise-counter}
全量基准意味着几十个指标同时进入检验，而每个指标又只拿三五次重复去判——这是 第~\ref{ch:pm03-stats} 章 里多重比较的教科书场景：单个指标假阳率百分之五时，五十个指标里\textbf{至少一个报警}的概率接近九成。后果不是``偶尔红一次''，而是团队开始把红灯当天气，门禁失去意义。可用的收敛手法按优先级：只对少数关键指标做门禁，其余只记录；按指标族做一次合并判定；要求连续两次确认才阻断（第~\ref{sec:pm13-verdict} 节）；把全量基准放夜间而 PR 只跑冒烟集合。
\end{note}

\section{回归检测的判据：从分布到三档结论}
\label{sec:pm13-verdict}

判据不能写成``head 比 base 慢超过百分之五就失败''。点值比点值，等于把两台机器的方差塞进一个减法里。CI 上可用的判据有五条，都是\textbf{可执行规则}而不是推导。

\textbf{（一）先标定，再定阈值。} 在你要真正使用的那类 runner 上，对\textbf{同一个 commit、同一个二进制}重复跑 $N$ 次（下限取二十次起步，指标本身抖动大就取更多），每次内部再重复若干次采样。把这 $N$ 个代表值（建议用每次样本的中位数）的相对波动算出来，阈值取``波动若干倍''与``业务最小关心差''二者的大者。阈值于是有了出处：它记录的是这台机器今天的噪声水平，而不是一句口号。

\textbf{（二）用区间而不是点值。} 对两组样本做 bootstrap 中位数差，看置信区间的\textbf{位置}而不是中心：区间整体落在噪声门限之上才叫变差，跨过零就叫无结论。实现见清单 \ref{lst:pm13-verdict}，方法与样本量下限回指 第~\ref{ch:pm03-stats} 章。

\textbf{（三）单边。} 门禁只拦``变差''。``变快''应当只是记录——它既诱使你调松阈值，也会在指标方向语义相反时（\MThr{} 与 \MWall{} 一个越大越好一个越小越好）制造混乱。每个指标必须显式声明方向。

\textbf{（四）连续两次确认。} 第一次 \Tp{REGRESSED} 只标记，第二次仍红才阻断合并。它把共享 runner 上的假阳率按近似平方级压下去，代价是坏改动多活一次提交；在无法独占机器的现实下，这个交换通常值得。

\textbf{（五）阶跃偏置用交替跑排除。} 如果 base 全部先跑、head 全部后跑，任何发生在两个阶段之间的环境变化（邻居起来、宿主限流、缓存被清）都会写成回归。交替执行 A/B/A/B 并保留每次的顺序记录，才能把这类阶跃压成公共噪声。

\begin{lstlisting}[style=pythonstyle, caption={可跑：空跑标定 + bootstrap 区间 + 三档结论 +
两次确认。只用标准库；判据参数必须由调用方从标定结果传入}, label={lst:pm13-verdict}]
# tools/perf_gate.py  (Python 3.8+, 无第三方依赖)
import random
import statistics

MIN_SAMPLES = 5


def bootstrap_median_diff_ci(base, head, rounds=5000, alpha=0.05, seed=0):
    """返回 head - base 中位数差的近似置信区间 (lo, hi)。
    样本不足直接返回 nan, 让调用方给出 INCONCLUSIVE, 不要猜。"""
    if len(base) < MIN_SAMPLES or len(head) < MIN_SAMPLES:
        return (float('nan'), float('nan'))
    rng = random.Random(seed)
    diffs = []
    nb, nh = len(base), len(head)
    for _ in range(rounds):
        mb = statistics.median([base[rng.randrange(nb)] for _ in range(nb)])
        mh = statistics.median([head[rng.randrange(nh)] for _ in range(nh)])
        diffs.append(mh - mb)
    diffs.sort()
    lo = diffs[max(0, int(round((alpha / 2.0) * rounds)) - 1)]
    hi = diffs[min(rounds - 1, int(round((1.0 - alpha / 2.0) * rounds)) - 1)]
    return (lo, hi)


def calibrate_threshold(runs_of_same_commit, business_floor=0.03):
    """runs_of_same_commit: 同一 commit 重复 N 次, 每次一个代表值(该次样本的中位数)。
    建议阈值 = max(波动放大, 全距折算, 业务最小关心差)。
    N 小于 20 直接抛错: 阈值没有出处, 就不要写进流水线。"""
    reps = [float(v) for v in runs_of_same_commit]
    if len(reps) < 20:
        raise ValueError('need >= 20 empty runs to calibrate, got %d' % len(reps))
    med = statistics.median(reps)
    if med <= 0.0:
        raise ValueError('non-positive median, fix the metric first')
    mad = statistics.median([abs(v - med) for v in reps]) / med
    rng_rel = (max(reps) - min(reps)) / med
    return {'mad_rel': mad, 'range_rel': rng_rel,
            'suggested_min_rel': max(3.0 * mad, 0.5 * rng_rel, business_floor)}


def verdict(base, head, min_rel, higher_is_better=False):
    """三档结论: REGRESSED / NOISE / IMPROVED, 外加 INCONCLUSIVE。
    min_rel 必须由 calibrate_threshold() 给出, 不要手填百分比。
    head - base 为正表示耗时变长; 越大越好的指标要传 higher_is_better=True。"""
    lo, hi = bootstrap_median_diff_ci(base, head)
    med_b = statistics.median(base) if base else float('nan')
    detail = {'ci_abs': (lo, hi), 'min_rel': min_rel, 'median_base': med_b}
    if lo != lo or hi != hi or med_b <= 0.0:
        detail['reason'] = 'insufficient samples'
        return 'INCONCLUSIVE', detail
    rel_lo, rel_hi = lo / med_b, hi / med_b
    if higher_is_better:
        rel_lo, rel_hi = -rel_hi, -rel_lo      # 统一成: 正数表示变差
    detail['ci_rel'] = (rel_lo, rel_hi)
    if rel_lo > min_rel:
        return 'REGRESSED', detail
    if rel_hi < -min_rel:
        return 'IMPROVED', detail
    return 'NOISE', detail


def should_block(status_history, confirm=2):
    """只有连续 confirm 次 REGRESSED 才阻断合并; NOISE 与 INCONCLUSIVE 都不阻断。"""
    recent = status_history[-confirm:]
    return len(recent) == confirm and all(s == 'REGRESSED' for s in recent)
\end{lstlisting}

清单里三处口径值得在评审时逐条问。\Tp{calibrate{\Bk}\_threshold} 里的二十次是\textbf{下限}而不是建议值；\Tp{verdict} 用区间下界与阈值比，因此\textbf{默认不报警}（宁可漏一次小回归，不可误报一百次）；\Tp{should{\Bk}\_block} 的两次确认会把阻断推迟一次提交，需要在 PR 描述里明写``当前状态：待确认''，否则作者以为已经过关。

\begin{compare}[title={两种判据在同一批数据上的行为差别}]
\label{box:pm13-two-rules}
\begin{tabular}{>{\raggedright\arraybackslash}p{0.24\textwidth}%
>{\raggedright\arraybackslash}p{0.24\textwidth}%
>{\raggedright\arraybackslash}p{0.24\textwidth}}
\hline
\textbf{维度} & \textbf{点值比阈值} & \textbf{区间比阈值（推荐）}\\
\hline
需要多少样本 & 一组一个数，看着最省 & 每组要能撑起一次 bootstrap\\
\hline
被一次离群值影响 & 直接翻转结论 & 中位数与分位数天然钝化\\
\hline
输出形态 & 红或绿两态 & 红、无结论、绿三档\\
\hline
共享 runner 上的假阳 & 高（阈值被迫放大） & 可控（放大的是区间而不是阈值）\\
\hline
\end{tabular}
\end{compare}

对 \Mx{p99} 这类尾延迟指标要单独提醒：CI 里想拿尾延迟，样本量必须明显大于 第~\ref{ch:pm03-stats} 章 给出的下限；拿五十次样本算 \Mx{p99} 再进门禁，等于把排序最后一个元素当结论。这种情况下正确的做法是\textbf{把判据换成 \Mx{p95} 或``最大值与均值的比''}，并把这一改动写进指标定义，而不是偷偷换。

\section{让基准跑在可控环境：独占、绑定与指纹入库}
\label{sec:pm13-env}

上一节所有判据的前提是``同一次运行内可比''。要报绝对值、要做归因、要给发布结论，就得把环境真正管起来。下面这些手法在自建 Linux 机器上大多可行，在托管 runner 上大多\textbf{没有权限}，所以每条都附``怎么确认自己有没有''。

\begin{lstlisting}[style=shellstyle, caption={环境指纹采集：输出单行 JSON 与基准结果一起入库；
指纹不同的两次结果不得进同一条基线序列（托管 runner 上多数行会因权限而空）},
label={lst:pm13-fingerprint}]
#!/usr/bin/env bash
# tools/ci_env_fingerprint.sh -- 单行 JSON, 键里不含秘密
set -uo pipefail

read_or_na () { [ -r "$1" ] && tr -d '[:space:]' < "$1" || echo NA; }
clean () { printf '%s' "$1" | tr -d '"'; }   # 去掉双引号, 免得把 JSON 打坏

cpu_model=$(lscpu 2>/dev/null | awk -F: '/Model name/{gsub(/^ +/,"",$2); print $2; exit}')
cores=$(nproc 2>/dev/null || echo NA)
gov=$(read_or_na /sys/devices/system/cpu/cpu0/cpufreq/scaling_governor)
cur_hz=$(read_or_na /sys/devices/system/cpu/cpu0/cpufreq/scaling_cur_freq)
paranoid=$(read_or_na /proc/sys/kernel/perf_event_paranoid)
virt=$(systemd-detect-virt 2>/dev/null || echo NA)
kernel=$(uname -r)
compiler=$(${CXX:-g++} --version 2>/dev/null | head -n 1)
image=${CI_CONTAINER_IMAGE:-local}
runner=${RUNNER_NAME:-local}
load1=$(cut -d' ' -f1 /proc/loadavg)

printf '{"cpu_model":"%s","cores":"%s","governor":"%s","cur_hz":"%s","paranoid":"%s",' \
       "$(clean "$cpu_model")" "$(clean "$cores")" "$(clean "$gov")" \
       "$(clean "$cur_hz")" "$(clean "$paranoid")"
printf '"virt":"%s","kernel":"%s","compiler":"%s","image":"%s","runner":"%s","load1":"%s"}\n' \
       "$(clean "$virt")" "$(clean "$kernel")" "$(clean "$compiler")" \
       "$(clean "$image")" "$(clean "$runner")" "$(clean "$load1")"
\end{lstlisting}

清单里 \Tp{paranoid} 与 \Tp{virt} 两行是本章反复用到的字段：前者决定你能不能拿到硬件事件（第~\ref{ch:pm05-events} 章），后者决定``同机自比''里那个``同机''到底同不同。\Tp{load1} 是廉价的一次体检：它不为零时，这次结果\textbf{应该丢掉}，而不是交给判据去消化。丢掉一次的代价是一次重跑，判据消化它的代价是一条错基线。

除采集之外，可控环境还要落实六件\textbf{动作}。

\begin{itemize}
\item \textbf{独占}：自托管 runner，一次只跑一个基准 job；同机其他 CI 任务、容器、桌面会话一律排除。做不到独占就退回同机自比，放弃绝对值。
\item \textbf{绑核}：\Tp{taskset} \Opt{-c} 指定核，\Tp{numactl} 指定节点。只绑 CPU 不绑内存在双路上仍会慢（距离的账见 第~\ref{ch:pm12-bandwidth} 章，跨线程共享见 第~\ref{ch:pm09-concurrency} 章）。绑核要写进结论，否则别人复现不出来。
\item \textbf{频率}：能改就设 \Tp{performance}、关 \Tp{no{\Bk}\_turbo}；改不了就\textbf{记录}当前频率，并把频率漂移当成已知混淆项写进判据说明。虚机上这些文件通常不可写。
\item \textbf{预热与丢弃首轮}：首轮含页缓存、动态链接、分支预测器状态与频率爬升。CI 里常见的做法是丢弃第一次运行的全部样本，然后固定重复次数。
\item \textbf{固定迭代次数而不是固定时长}：Google Benchmark 的 \Opt{benchmark\_min\_time} 一类默认值会随机器负载自适应地改变实际执行时长，让\textbf{样本量}在两次运行之间不一致。门禁里宁可显式指定重复次数，把样本量当成协议的一部分（默认值不是结论，见 第~\ref{ch:pm03-stats} 章）。
\item \textbf{环境变了就重建基线}：把指纹哈希做成基线的键的一部分——镜像升级、宿主换型、编译器小版本都算变化。键不同就不许和旧序列比。这条纪律比任何判据都省事。
\end{itemize}

基准侧要按\textbf{样本}而不是按\textbf{均值}输出，否则判据无从下手。下面这段是最小可跑形态：每次迭代一行 JSON，均值只在最后打印，交给门禁的是逐行样本。

\begin{lstlisting}[style=cppstyle, caption={CI 友好的基准输出骨架：每个迭代单独成行、
显式重复次数、防死代码消除。数值不填，由读者在自己的机器上产出}, label={lst:pm13-sample}]
// tools/bench_emit.cpp -- 输出 NDJSON, 每行一个迭代样本
#include <chrono>
#include <cstdio>
#include <cstdlib>
#include <vector>

static void keep_alive(void* p) noexcept {
  // 反死代码消除: 编译器无法证明被测循环没有副作用
  asm volatile("" : : "r"(p) : "memory");
}

static double one_iteration(int n) {
  auto t0 = std::chrono::steady_clock::now();
  std::vector<double> buf(static_cast<size_t>(n), 1.0);
  double acc = 0.0;
  for (int i = 0; i < n; ++i) acc += buf[static_cast<size_t>(i)];
  keep_alive(&acc);
  auto t1 = std::chrono::steady_clock::now();
  return std::chrono::duration<double, std::milli>(t1 - t0).count();
}

int main(int argc, char** argv) {
  const int n = (argc > 1) ? std::atoi(argv[1]) : 1000000;
  const int iters = (argc > 2) ? std::atoi(argv[2]) : 30;  // 样本量由协议给定
  for (int i = 0; i < 3; ++i) (void)one_iteration(n);      // 预热, 结果丢弃
  double sum = 0.0;
  for (int i = 0; i < iters; ++i) {
    double ms = one_iteration(n);
    sum += ms;
    std::printf("{\"case\":\"accum\",\"iter\":%d,\"wall_ms\":%.6f}\n", i, ms);
  }
  std::fprintf(stderr, "mean_ms=%.6f n=%d iters=%d\n", sum / iters, n, iters);
  return 0;
}
\end{lstlisting}

两个提醒。其一，\Fn{steady{\Bk}\_clock} 在虚机里的分辨率与推进源由宿主决定，``两次跑出三成差''时先按 第~\ref{ch:pm02-clock} 章 检查时钟源，再怀疑代码。其二，这段代码\textbf{测的是 \MWall}；想同时拿 \MCpu{} 需要在同一进程内读 CPU 时间源，两个口径的分歧见 第~\ref{ch:pm01-why} 章。

\section{性能门禁与功能测试的关系}
\label{sec:pm13-gate}

性能门禁是\textbf{流水线的消费者}，不是独立系统。它挂载的位置、失败时的证据留存、``谁有权改阈值''三件事决定它能不能长期活下来。

\textbf{位置。} 放在构建之后并\textbf{复用同一二进制}（用产物传递，不在基准 job 里重新编译）。重新编译意味着你在比两个不同产物，那已经不是性能比较了（第~\ref{ch:pm10-compiler} 章）。执行顺序上，功能测试先跑：功能不正确的程序跑得再快也要先失败，而且功能测试失败时的日志价值高于任何基准数。快速集合挂 PR，全量集合挂夜间。

\textbf{判定用退出码，不用日志文本。} 本仓库门禁的做法可以直接对齐：判定只看退出码，需要人读的信息放两处——日志尾部若干行进 job 摘要，完整原始样本作为产物上传。

\begin{lstlisting}[style=shellstyle, caption={门禁封装：退出码单独判失败、日志尾部取证、
产物截断检查（样本行数不足即 INCONCLUSIVE，绝不判回归）}, label={lst:pm13-gate}]
#!/usr/bin/env bash
# tools/perf_gate.sh <bench-bin> <base-json> <head-json> <min-rel>
set -uo pipefail
BENCH=$1; BASE=$2; HEAD=$3; MIN_REL=$4
LOG=perf-gate.log
SUMMARY=${GITHUB_STEP_SUMMARY:-/dev/stdout}
MIN_ROWS=${MIN_ROWS:-20}

# 1) 采样: base 与 head 交替各跑一次, 顺序写进产物, 便于事后查阶跃偏置
rc_head=0; "${BENCH}" 1000000 30 > "$HEAD" 2>> "$LOG" || rc_head=$?
rc_base=0; "${BENCH}" 1000000 30 > "$BASE" 2>> "$LOG" || rc_base=$?
if [ "$rc_head" -ne 0 ] || [ "$rc_base" -ne 0 ]; then
  printf 'GATE=ERROR rc_head=%s rc_base=%s\n' "$rc_head" "$rc_base" >> "$SUMMARY"
  tail -n 40 "$LOG" >> "$SUMMARY"; exit 1
fi

# 2) 产物截断检查: 行数不足说明被超时或被抢占, 结论无效, 而不是"过关"
rows_head=$(grep -c . "$HEAD" || true)
rows_base=$(grep -c . "$BASE" || true)
if [ "$rows_head" -lt "$MIN_ROWS" ] || [ "$rows_base" -lt "$MIN_ROWS" ]; then
  printf 'GATE=INCONCLUSIVE rows_base=%s rows_head=%s need=%s\n' \
         "$rows_base" "$rows_head" "$MIN_ROWS" >> "$SUMMARY"
  exit 0        # 不阻断合并, 但把这次运行标黄并留档
fi

# 3) 判据交给 perf_gate.py（清单 lst:pm13-verdict），它只打印一个状态词
status=$(python3 tools/perf_gate.py "$BASE" "$HEAD" "$MIN_REL")
if [ -z "$status" ]; then
  echo 'GATE=ERROR: judge emitted nothing' >> "$SUMMARY"; exit 1
fi
printf 'GATE=%s base_rows=%s head_rows=%s min_rel=%s\n' \
       "$status" "$rows_base" "$rows_head" "$MIN_REL" >> "$SUMMARY"
tail -n 20 "$LOG" >> "$SUMMARY"
# 4) 只有 REGRESSED 让 job 失败; INCONCLUSIVE 由上层按重试策略处理
[ "$status" = REGRESSED ] && exit 1
exit 0
\end{lstlisting}

\textbf{证据留存。} 一次失败要能事后复盘，至少要留四样：两个二进制的路径与 \Tp{build-id}、完整样本文件（NDJSON，不要只留均值）、环境指纹（清单 \ref{lst:pm13-fingerprint} 的输出）、判据参数（阈值、样本量、以及阈值哪天用哪批空跑数据标定的）。缺任何一样，三个月后没人能判断那次红是不是真的。产物要检查\textbf{截断}：上传前核对行数与文件大小下限，被截断的样本比没有样本更坏，因为它看起来像有。

\begin{bestpractice}[title={门禁输出的三件套}]
\label{box:pm13-three}
任何一次基准 job 的输出都必须包含三件：(1) \textbf{结论}，即 \Tp{PASS}、\Tp{REGRESSED}、\Tp{INCONCLUSIVE}、\Tp{ERROR} 四态，不要压成两态；(2) \textbf{判据}，即阈值、样本量、确认次数与阈值的标定出处；(3) \textbf{原始数据位置}。四态里 \Tp{INCONCLUSIVE} 最关键，它把``测不出来''与``没变差''分开；否则噪声会伪装成过关。
\end{bestpractice}

``谁改阈值''要有制度答案：改阈值必须附一次新的标定数据（PR 里带上新的空跑分布），并让基线序列在新阈值下重放一遍，看历史会不会翻红。缺这一步，阈值会变成每次红灯后的逃生口（见 第~\ref{sec:pm13-trouble} 节）。

\section{对照侧：托管 runner、自建机与三个平台}
\label{sec:pm13-platforms}

\textbf{托管与自建。} 托管 runner 的 CPU 型号、宿主、邻居都不可选，同一仓库的两个 job 可能落在不同代际的机器上；优点是免运维、有弹性。自托管把可控性买回来，代价是你要负责清洁（重启、清缓存、单一负载）与容量。可行的折中：PR 用托管做同机自比，夜间与发布验证用自托管做绝对值与归因。

\textbf{Linux runner。} 唯一能提供完整工具面的一档，但\textbf{权限取决于平台方}。下面这段自证脚本可以在任何 Linux runner 上跑，它的输出决定你能用本册哪几章的手法。

\begin{lstlisting}[style=shellstyle, caption={CI 上的能力自证（Linux）：逐条探测, 而不是引用
别人的结论。每一行对应本册某一章的前置条件}, label={lst:pm13-probe}]
#!/usr/bin/env bash
# tools/ci_capability_probe.sh -- 打印本 runner 到底能测什么
p=/proc/sys/kernel/perf_event_paranoid
echo "paranoid=$([ -r "$p" ] && cat "$p" || echo unreadable)"
command -v perf >/dev/null 2>&1 && echo perf=present || echo perf=absent
# 硬件事件是否真有计数: 对自己跑一个极短命令, 看是不是 not counted
if command -v perf >/dev/null 2>&1; then
  perf stat -e cycles true 2>&1 | sed -n 's/^ *//p' | head -n 4
fi
for f in /sys/devices/system/cpu/cpu0/cpufreq/scaling_governor \
         /sys/devices/system/cpu/cpu0/cpufreq/scaling_cur_freq; do
  echo "$f=$([ -w "$f" ] && echo writable || echo readonly_or_absent)"
done
quota=NA
[ -r /sys/fs/cgroup/cpu.max ] && quota=$(xargs < /sys/fs/cgroup/cpu.max)
echo "nproc=$(nproc) cpu.max=$quota"
echo "virt=$(systemd-detect-virt 2>/dev/null || echo NA)"
grep -q constant_tsc /proc/cpuinfo && echo tsc=constant || echo tsc=unverified
\end{lstlisting}

三行最容易被忽略：\Tp{cpu.max} 不是 \Tp{max} 意味着基准被\textbf{限流}，长基准会拿到一个被 cgroup 时间片切碎的数；\Tp{perf{\Bk}\_event{\Bk}\_paranoid} 大于 $2$ 时连自身进程的采样都受限；\Tp{tsc=unverified} 时任何用周期计数换算出来的时间都要重算。

\textbf{Windows runner。} 采样类工具依赖 ETW 会话权限，通常需要管理员身份或 \Tp{Performance Log Users} 组成员；托管 runner 上一般拿不到，\Fn{xperf} 与 \Fn{wpr} 因此不可用，PMU 级工具同理。可行的替代是纯 \MWall{} 计时加进程内 CPU 时间读取，把剖析留到自建机。权限到底有没有，用``先试着开会话、失败就退回计时''的策略实测，而不是照文档假设。

\textbf{macOS runner。} \Fn{sample} 与 Instruments/\Fn{xctrace} 面向自身进程多数可用，但系统级剖析受完整性保护与权限模型约束，\Fn{dtrace} 需要额外授权，\Fn{powermetrics} 需要 \Fn{sudo}。Apple Silicon 上没有 x86 意义上的那份事件名体系，第~\ref{ch:pm05-events} 章 的事件表整张不适用。三平台权限与事件的完整并列见 附录~\ref{ch:pmapA-platform}。

三平台共同的一点：\textbf{CI 里没有免费的绝对值}。任何要报出去的数字都应在可控机器上重算一次，并把那台机器的指纹附在结论里（结论的写法模板见 第~\ref{ch:pm15-antipattern} 章）。

\section{排障与速查}
\label{sec:pm13-trouble}

\begin{warning}[title={CI 红三次就调阈值}]
\label{box:pm13-warning-threshold}
连续红灯之后最容易做的动作是把 \Tp{min\_rel} 放大，这是\textbf{最错}的一步：阈值是这条流水线噪声水平的度量，噪声不会因为你放大阈值而消失，只会转移成漏报。正确的顺序是：先查环境（有没有换镜像或换宿主、\Tp{load1} 与 \Tp{paranoid} 是否变了、构建缓存是不是冷了、是不是只有某一个用例红），再查代码（把两次红之间的提交做二分），最后才动阈值——而动阈值必须重做一次空跑标定并在 PR 里附上分布。判据错了调判据，测不到调测法，\textbf{不要调数字}。
\end{warning}

\begin{longtable}{|>{\raggedright\arraybackslash}p{0.10\textwidth}|>{\raggedright\arraybackslash}p{0.24\textwidth}|>{\raggedright\arraybackslash}p{0.22\textwidth}|>{\raggedright\arraybackslash}p{0.20\textwidth}|}
\caption{门禁假阳与假阴速查：先定类型，再按``先查''列动手}
\label{tab:pm13-trouble}\\
\hline
\textbf{类型} & \textbf{症状} & \textbf{先查} & \textbf{处置}\\
\hline
\endfirsthead
\hline
\textbf{类型} & \textbf{症状} & \textbf{先查} & \textbf{处置}\\
\hline
\endhead
假阳 & 同一 commit 重跑一次就绿 & 这次 job 的 \Tp{load1}、宿主是否换了、
base 与 head 的交替是否被打断 & 按 \Tp{INCONCLUSIVE} 重跑，不改判据\\
\hline
假阳 & 只有夜间流水线红，PR 全绿 & 夜间是否并发跑了别的 job、备份或索引任务 &
基准 job 加独占锁；夜间与 PR 用各自的阈值出处\\
\hline
假阳 & 冷构建那次必红，之后稳定 & 构建缓存是否命中、是否有预热阶段 &
丢弃首轮；把冷构建单列一类基线\\
\hline
假阴 & 明显改了热路径但门禁全绿 & 这个用例是否真走在被改的代码路径上、阈值是否过大 &
补关键用例；把阈值放大历史贴出来复审\\
\hline
假阴 & 发布版比三个版本前慢了两成而序列没红 & 是否用滚动窗口把漂移吸收了、
换机器时是否换了基线键 & 加``相对最近发布版''的第二轨判据\\
\hline
假阴 & 尾延迟恶化而均值不变 & 门禁指标是否只有均值、样本量够不够算分位数 &
换成有出处的分位数或最大值指标，样本量按 第~\ref{ch:pm03-stats} 章\\
\hline
\end{longtable}

最后三句速记。第一句：CI 给的数只能回答``和刚才比有没有变差''，回答不了``到底多快''。第二句：阈值必须有出处，出处是一次可复跑的空跑标定，不是一个凭空的百分比。第三句：红灯要留证据（产物、指纹、判据参数），不留证据的门禁等于在训练团队忽略它。再往前一步——绕开 CI，在生产环境里拿一份不用重跑的 profile——是第~\ref{ch:pm14-production} 章的事。
\chapter{生产环境的低开销采样与剖析预算}
\label{ch:pm14-production}

前一章把性能结论钉在流水线里，靠的是``每次构建都在同一台机器上跑同一批基准''；那一套前提在生产环境里一条都不成立：负载形状随用户变、机器是共享的、二进制是三天前发布的，而且你不能为了拿一个数把它停下来。本章解决生产侧的问题：\textbf{当观测本身要花钱、而花钱的对象又是正在挣钱的那台机器时，怎么把``看清楚''这件事设计成一个有预算、有上限、可回滚的东西}。

本章的三样核心代价是\textbf{采样率}、\textbf{停机时间}（剖析器占用被测程序的时间窗，以及它可能造成的服务停顿）与\textbf{符号化成本}；三者互相顶：栈拿得越全，单个样本越贵；样本越密，中断越频繁；符号化放在进程内做，就把离线成本搬进了热路径。排障速查在第~\ref{sec:pm14-trouble} 节，三平台的完整能力并列留给附录~\ref{ch:pmapA-platform}。

与前后章的接口：口径与判据回看第~\ref{ch:pm01-why} 章和第~\ref{ch:pm03-stats} 章；采样剖析本身的机制（三种栈展开、符号与折叠）在第~\ref{ch:pm04-sampling} 章，本章只谈它在生产约束下怎么取舍；权限模型与能不能拿到 PMU 的判定回看第~\ref{ch:pm05-events} 章；剖析器改变被测行为的机制在第~\ref{ch:pm15-antipattern} 章展开，本章只把它记进预算。

\section{观测预算是一等指标}
\label{sec:pm14-budget}

生产剖析与实验室剖析的差别不在工具，在约束。三条硬约束：

\begin{itemize}
  \item \textbf{不能停}。服务有可用性目标，任何``先把进程暂停再取栈''的手段（\Mx{stop-the-world} 快照、调试器 attach、部分运行时自带的 safepoint 采样）都要按``最坏情况停顿多少''来评估，而不是按``平均多少''。停顿落在请求上的后果是尾延迟与超时，见第~\ref{sec:pm14-overhead} 节。
  \item \textbf{不能重编}。观测对象必须是发布产物本身，否则你测的是另一个二进制。这条直接推出一条纪律：\textbf{发布构建就得为可观测性留下余地}，而不是出事了再用带调试信息的构建重跑一遍。要不要在发布里保留帧指针、留多少符号、build-id 怎么归档，是发版前的决定，构建口径在第~\ref{ch:pm10-compiler} 章。
  \item \textbf{不能占资源}。CPU 时间是容量，内存是容量，磁盘与网络也是容量；在容器里还有 quota，在笔记本与嵌入式设备上还有电与热。观测吃掉的每一个百分点都得从服务账上扣。
\end{itemize}

于是本章把\textbf{观测预算}（observation budget）当成一等指标：它是``为了看见，你愿意并且能够支付多少额外成本''的一个\emph{事先写下来的数}，而不是事后从日志里发现剖析器把机器打满。预算至少要记四笔：

\begin{enumerate}
  \item \textbf{CPU 税}：剖析器占用的 CPU 时间增量。口径要明确——是相对同一时刻另一组实例的 \MCpu{}（CPU 时间口径），还是相对 \MWall{}（墙钟时间口径）。两者在有并行度的服务上不等价，原因见第~\ref{ch:pm09-concurrency} 章。
  \item \textbf{内存税}：环形缓冲、每 CPU 样本页、未导出样本的驻留量（\Mx{RSS} 与峰值工作集）。它常表现为容量下降而不是变慢，因而最容易被漏掉。
  \item \textbf{符号化与存储税}：样本落地要变成人看得懂的东西才叫观测。这部分成本可以放在进程内、放在旁路采集器、也可以放在离线集群，三处的价格差得很远。
  \item \textbf{风险税}：权限提升、内核侧代码（驱动、eBPF 程序、ETW 会话）出问题导致整机重启或不可逆的性能退化。这一笔最难量化，却是``能不能在高峰时段开采样''的真正决定因素。
\end{enumerate}

预算的百分比从哪来？\textbf{不要抄别人的百分比}。可用的来源有三种：一是可用性与延迟目标留下的裕度（目标 \Mx{p99} 与当前 \Mx{p99} 之间那一段，你愿意花多少去买可见性）；二是容量水位（峰值利用率越接近上限，同样的采样率越危险，因为采样中断落在关键路径上的概率与代价同时上升）；三是历史故障（一次剖析器引发的抖动值不值得，取决于你上次出事的样子）。三种来源都指向同一个动作：\textbf{先定义预算，再选手段，最后测实际开销}，顺序反了就是把工具默认值当结论。

\begin{guideline}[title={观测预算的四要素模板}]
\label{box:pm14-budget-tpl}
写预算时用这四行，缺一行就会在事后争论：
\begin{itemize}
  \item \textbf{窗口}：一次采样持续多久（秒级），窗口之间静默多久。窗口与静默共同决定占空比（\Mx{duty cycle}）。
  \item \textbf{上限}：CPU 税、内存税、单窗口磁盘写入各自的上限值，以及超限后的动作（停止采样、丢样本，还是只告警）。
  \item \textbf{口径}：用 \MCpu{}、\MThr{}（吞吐）还是 \Mx{p99} 的增量来判断超预算；只写均值，等于默认接受尾部不管。
  \item \textbf{回滚}：不发版能不能关掉；关不掉的手段（编译进去的探针、常驻代理）要单独评审。
\end{itemize}
\end{guideline}

\section{低开销手段的层次}
\label{sec:pm14-methods}

手段的开销与能力大致成一条阶梯，越往上越贵，也越接近能直接给出原因。下面先给分层依据，再给一张对照表。

\subsection{定时采样：最低门槛的一层}
\label{sub:pm14-timer}

Linux 上 \Fn{perf} record 用 \Opt{freq} 请求每秒样本数，或用 \Opt{count} 配合 \Tp{-c} 请求每 N 个事件一个样本。两者的取舍不同：

\begin{itemize}
  \item \Opt{freq}（例如 \Tp{-F 99}）保证\emph{时间}均匀，负载变重时每个样本覆盖的指令数变少，而样本密度不变。适合回答``这段时间 CPU 花在哪''。
  \item \Opt{count} 保证\emph{事件}均匀，负载变重时样本变稀。适合跨版本比较同一段代码的相对占比，因为周期与实际工作绑定；但它在空闲时段几乎不出样本，容易被误读成``没工作''。
\end{itemize}

选频率没有通用答案，只有判据：\textbf{要看清的最短函数或最小占比，决定最低频率}；\textbf{中断能接受的额外开销，决定最高频率}。低频采样对短函数完全瞎——占比很低的函数需要成比例更长的窗口，这条统计关系在第~\ref{ch:pm04-sampling} 章讲过。频率的绝对开销请自己测（第~\ref{sec:pm14-overhead} 节给方法），不要用``别人说某个频率太贵''当结论，尤其不要把工具默认频率当推荐值：默认是为了在开发机上出数快而设的，不是为了在满载机器上稳而设的。

系统级（\Opt{system}）与按进程（\Opt{target}）的口径也不同：前者能看见内核与邻居，噪声与数据量都大得多；在容器里按 PID 采样还会漏掉不属于该 PID 的内核工作，容易被误读成``内核时间凭空消失''，机制见第~\ref{ch:pm05-events} 章。

\subsection{硬件精确采样：PEBS / IBS / SPE}
\label{sub:pm14-precise}

定时采样的样本落在中断返回处，与真正的事件发生点有偏差；硬件精确采样（Intel 的 PEBS、AMD 的 IBS、ARM 生态的 SPE 一族）把样本点拉到事件本身，还能带上访存地址。收益是归因质量（第~\ref{ch:pm06-micro} 章用得上），代价是三样：\textbf{权限更高}（很多发行版把精确事件与带地址采样挡在 \Tp{perf{\Bk}\_event{\Bk}\_paranoid} 与容器 capability 之外）、\textbf{每样本成本更高}（硬件写缓冲区、内核搬运）、\textbf{事件名随 CPU 型号与内核版本变}。最后一条决定了：任何抄来的事件表只能当线索，用 \Fn{perf} list 在你自己的机器上确认（附录~\ref{ch:pmapA-platform} 给三平台列出可用能力的方法）。

生产上的建议口径是：\textbf{精确采样属于排障窗口，不属于常驻观测}。它更贵、更依赖权限，而且一旦带上地址数据，样本体积与隐私风险同时上升。要用就写进变更单：什么时候开、开多大、谁来关、关不掉怎么办。

\subsection{内核内聚合：eBPF 栈与直方图}
\label{sub:pm14-ebpf}

把``每个样本都搬到用户态再符号化''改成``在内核里先把计数累加好，只把聚合结果搬出来''，是近年生产观测最重要的一个成本结构变化。典型做法：挂一个 \Tp{perf{\Bk}\_event} 程序，在回调里用栈回溯 helper 取栈、在内核 map 里按栈累加计数，用户态定期读走 map 的快照。开销结构因此变成：\textbf{每样本一个内核内成本}（取栈、算哈希、更新 map），而不再是\textbf{每样本一次数据搬运加一次用户态符号化}。

它的取舍必须写清：
\begin{itemize}
  \item \textbf{权限}：需要 \Tp{CAP{\Bk}\_BPF} 与 \Tp{CAP{\Bk}\_PERFMON}（老内核是 \Tp{CAP{\Bk}\_SYS{\Bk}\_ADMIN}），再加内核参数放宽；容器与 seccomp 可能连带挡掉，判定方法见附录~\ref{ch:pmapA-platform}。
  \item \textbf{聚合位置的不可逆性}：在内核里聚合意味着\emph{你事先决定了按什么分组}。分组选错（例如没按线程或没按版本分组），线上跑完就再也问不出第二个问题。这是低开销的代价：可查询性变窄。
  \item \textbf{栈深度与展开方式}：内核里能取的栈深度有上限，取的是帧指针链还是内核侧展开器，直接决定符号化后的质量。
  \item \textbf{内核侧代码本身就是风险税}：程序加载失败可能丢样本，验证器拒绝可能让采集静默为零。把``聚合数突然为零''当成必须告警的事件（第~\ref{sec:pm14-trouble} 节）。
\end{itemize}

\subsection{静态探针与分布式追踪：另一条轴}
\label{sub:pm14-usdt}

USDT（\Tp{probe} 一族的静态定义点，其他平台有对应机制）不属于采样，它属于``程序员预先在代码里画好的观测点''。开销特征是：\textbf{没有 attach 时接近一条 nop 加一点语义标记}，attach 后每个命中点变成一次回调。这条``接近零''是量级口径不是保证：编译器可能因为探针标记改变内联决策，机制与确认方法在第~\ref{ch:pm10-compiler} 章。USDT 的价值在于稳定契约——探针名可以随版本承诺兼容，栈采样给不出这种承诺。

把轴分清楚很关键：\textbf{采样}回答时间或事件花在哪，不需要预先知道代码结构；\textbf{静态探针}回答我自己定义的量，需要预先埋点；\textbf{分布式追踪}（trace 与 span）回答这一个请求走了哪些服务，是单请求路径。三者不能互相当证据，第~\ref{sec:pm14-linkage} 节展开。

\subsection{应用内自采：分位数摘要与环形缓冲}
\label{sub:pm14-inproc}

最后一层是进程自己做观测：在业务代码或运行时里维护\textbf{固定内存}的环形缓冲（记录最近若干秒的栈或事件），并维护\textbf{固定大小}的分位数摘要（桶式直方图、t-digest、DDSketch 这类可合并表示）。三条设计约束：

\begin{itemize}
  \item \textbf{内存必须事先封顶}。缓冲条数、栈深度上限、桶数量都应是编译期或启动期常量；任何按样本量增长的缓冲都不能进生产。
  \item \textbf{写路径必须极窄}。一次 \Fn{clock{\Bk}\_gettime} 加一次带放松序的原子自增，通常是可以接受的方向，但\emph{能不能接受由你测}：先按第~\ref{ch:pm02-clock} 章的方法测这个累加本身在目标机器上的成本，再和预算比。
  \item \textbf{导出留给异常时刻}。环形缓冲在崩溃或超时那一刻导出，是最经济的现场获取方式：平时零搬运成本，出事时才付符号化的钱。在信号处理程序里取栈要注意异步信号安全——很多展开库与分配器在其中并不安全，这条必须在测试里验证，不能在文档里假设。
\end{itemize}

\subsection{手段对照}
\label{sub:pm14-method-table}

\begin{longtable}{>{\raggedright\arraybackslash}p{0.17\textwidth}>{\raggedright\arraybackslash}p{0.15\textwidth}>{\raggedright\arraybackslash}p{0.16\textwidth}>{\raggedright\arraybackslash}p{0.19\textwidth}>{\raggedright\arraybackslash}p{0.16\textwidth}}
\caption{低开销手段对照：权限、开销类型、能拿到什么、风险（三平台完整并列见附录~\ref{ch:pmapA-platform}）}\label{tab:pm14-method}\\
\toprule
手段 & 需要的权限 & 开销类型 & 能拿到什么 & 主要风险 \\
\midrule
\endfirsthead
\toprule
手段 & 需要的权限 & 开销类型 & 能拿到什么 & 主要风险 \\
\midrule
\endhead
\bottomrule
\endfoot
\bottomrule
\endlastfoot
定时软件采样（\Opt{freq}） & 普通用户即可，取决于内核的采样权限参数 & 每样本一次中断，外加缓冲写入 & 按时间均匀的热点占比；配合栈展开得到调用关系 & 频率选高直接变成中断风暴；短函数欠采样 \\
事件周期采样（\Opt{count} 配 \Tp{-c}） & 同上，硬件事件常需放宽 & 同上，另加事件配置 & 与工作量绑定的占比，可跨版本比 & 空闲时段无样本；周期与实际成本不同源 \\
精确硬件采样（PEBS、IBS、SPE） & 高：精确事件常需驱动与特权 & 每样本更贵，可带地址与数据 & 事件发生点的栈、访存地址与大小 & 权限、隐私、样本体积；可用性随型号变 \\
内核内聚合（eBPF） & \Tp{CAP{\Bk}\_BPF}、\Tp{CAP{\Bk}\_PERFMON} & 内核回调，另加读 map 快照 & 预先分组好的直方图与按栈聚合计数 & 分组定了就不能改；程序被拒或静默为零 \\
静态探针（USDT 一类） & 读探针元数据即可，采集时视工具 & 未 attach 时近零，attach 后按命中计 & 自定义量、可承诺的稳定名称 & 埋点覆盖有限；标记可能影响优化决策 \\
应用内自采（摘要与环形缓冲） & 无需额外权限 & 热路径累加，外加固定内存 & 进程内视角的分位数与崩溃前现场 & 探针自身进热路径；实现不安全会崩 \\
停顿式快照（attach 或 safepoint） & 进程控制权 & 一次性长停顿 & 完整对象图与全部线程栈 & 与``不能停''冲突；只适合低峰或摘流实例 \\
全量追踪（每个调用都记） & 视工具而定 & 与调用次数成正比 & 单请求路径与完整顺序 & 不属低开销手段，只用于抽样后的深挖 \\
\end{longtable}

\begin{warning}[title={把低开销当成手段的固有属性}]
\label{box:pm14-warn-lowcost}
没有任何手段天生低开销。同一条 \Fn{perf} record 命令，在四核开发机上和在一台满载、开着超线程、容器 CPU quota 已用到上限的生产机上，代价差一个量级完全正常——因为采样中断抢的是\emph{剩余}的 CPU 时间，配额与邻居决定了它是零和还是溢出。判断标准只有一个：\textbf{在你的真实负载形状上，用第~\ref{sec:pm14-overhead} 节的方法测一次开与关的对照}。在测试环境测一次就把数当成生产预算，是本章最常见的错误。
\end{warning}

\section{常驻剖析器与连续 profile}
\label{sec:pm14-continuous}

\subsection{为什么总是有一份比出事再抓有效}
\label{sub:pm14-always-on}

生产故障的证据有三类会过期：现场（进程状态、缓存热度、邻居负载）、版本（等你装上剖析器，热修复已经发出去了）、以及时序（很多抖动只在特定流量形状下出现）。``出事再抓''的隐含假设是你能重现问题，而这个假设在生产里通常不成立。连续 profile 的价值不是更聪明的采样，而是\textbf{把观测变成一个已经存在的、带时间戳的历史序列}，出事时只需要查询，不需要部署。

实现方式是占空比：\textbf{短窗口采样加长静默}的循环，让长时间平均开销很低，而单个窗口内的密度足够回答这段时间花在哪。占空比的两个数分别控制两件事：窗口长度决定能看清多短的突发，静默长度决定平均税多少。两者都应当由预算反推，而不是照抄示例。

一个必须提前想清楚的副作用：占空比采样让 profile 在时间上\emph{不均匀}。当某个窗口正好覆盖一次抖动，聚合图上会出现一个尖峰，而它可能只是因为那几秒被采了。判据：\textbf{任何跨窗口比较都要知道每窗口的样本数与权重}，样本少的窗口不能和样本多的窗口直接叠。

\subsection{常驻剖析器自身的工程要求}
\label{sub:pm14-daemon}

常驻的东西必须有边界，否则它就是一次等待发生的故障：

\begin{itemize}
  \item \textbf{熔断}：CPU、内存、磁盘三个上限，超限后先丢样本再停采样；停止动作要幂等，且不发版也能触发。
  \item \textbf{不能把失败变成业务失败}：采集器写不进磁盘、拿不到符号，应该降级成没有 profile，绝不该让请求失败。任何采样器挂了业务也挂了的设计，都要在评审里被问住。
  \item \textbf{自观测}：剖析器自己要有指标——窗口时长、样本数、丢弃数、符号解析失败率。丢弃率突然上升是最早的异常信号，也是第~\ref{sec:pm14-trouble} 节多条症状的入口。
  \item \textbf{与摘流配合}：需要停顿式快照时，先在负载均衡上摘掉这个实例再抓，别在服务路径上做长停顿。
\end{itemize}

\subsection{符号、版本与归档}
\label{sub:pm14-symbols}

连续 profile 只有在``三年后的这份 profile 还能看懂''的前提下才有价值。三平台各有一套对应物：Linux 是 ELF 的 build-id 与分离调试信息；Windows 是 PDB 及其 GUID 与 age 的匹配；macOS 是 dSYM 与 UUID。共同纪律是同一条：\textbf{profile 必须和产生它的二进制一起归档，靠标识而不是靠版本字符串}。版本字符串会重复、会被重新打包覆盖；build-id 由内容与链接决定，用它做键。

Linux 侧的落地形态如下（命令可跑，数值请自己测；本仓库的 CI runner 里既没有 perf 也没有 PMU 权限）：

\begin{lstlisting}[style=shellstyle, caption={线上符号归档与低占空比采样骨架（形态示意，不含任何实测值）}, label={lst:pm14-duty}]
#!/usr/bin/env bash
# 三个常量就是预算的三个旋钮
WIN=5          # 单窗口秒数
IDLE=55        # 静默秒数
FREQ=99        # 采样频率, 合理值由读者自己测后决定
OUT=/var/tmp/pm14-profiles
BIN=/usr/local/bin/myservice
mkdir -p $OUT/data $OUT/sym

# 用 build-id 做符号目录名: 同版本不同构建也不会互相覆盖
bid=$(readelf -n $BIN | awk '/build id/{print $NF}')
mkdir -p $OUT/sym/$bid && cp -a $BIN $OUT/sym/$bid/

while :; do
  tag=$(date +%s)
  perf record --pid $MYPID -F $FREQ --call-graph fp \
    -o $OUT/data/win-$tag.data -- sleep $WIN || exit 1
  # 报表先落地, 原始样本压缩后按保留期删除
  perf report -i $OUT/data/win-$tag.data --stdio --percent-limit 1 \
    > $OUT/data/win-$tag.txt || true
  gzip -f $OUT/data/win-$tag.data
  find $OUT/data -name '*.gz' -mmin +2880 -delete
  sleep $IDLE
done
\end{lstlisting}

\begin{note}[title={帧指针这件事要在发版前决定}]
\label{box:pm14-fp}
生产上取栈的成本几乎全部由展开方式决定：帧指针链便宜、能在信号上下文里用，但要求编译时带上 \Tp{-fno-{\Bk}omit-{\Bk}frame-{\Bk}pointer}；DWARF 展开不要求重编，可每个样本贵得多，且缓冲有上限；LBR 便宜却只给最近几层，还依赖具体 PMU。三者对比细节在第~\ref{ch:pm04-sampling} 章。\textbf{结论性的纪律}：如果这个服务需要生产采样，就在发布构建里保留帧指针，并把它当产品需求写进构建配置。等出事后再补，就意味着拿一次带调试信息的重建二进制测出的 profile，去解释线上那个二进制的行为——那是第~\ref{ch:pm15-antipattern} 章的口径偷换。
\end{note}

\subsection{存储与隐私：只落栈，不落业务数据}
\label{sub:pm14-privacy}

样本可能顺带落进业务数据的途径，不止参数值这一种。要逐项判断：访存地址（精确采样带出的数据地址可反推布局）、字符串常量附近的指令指针、内联进来的第三方代码路径，以及解释型或 JIT 型运行时映射出的符号名（其中可能含用户提供的路由或表名）。可执行的纪律有三条：

\begin{enumerate}
  \item \textbf{默认不采数据}：只在有明确问题、有评审、有保留期时才带上地址与数据级别。
  \item \textbf{符号在本地解决，标识在入库前脱敏}：需要按用户或租户分组时，用不可逆摘要而不是原始 ID。
  \item \textbf{归档要有保留期与删除路径}：清单~\ref{lst:pm14-duty} 里那句按修改时间删除就是这条纪律的最小实现；删除路径要能在合规评审里演示，而不是脚本里有就算。
\end{enumerate}

\begin{bestpractice}[title={上线前一次性准备：符号、标识与开关}]
\label{box:pm14-prep}
发版流水线里多做四件事，能省下未来每一次紧急排障的时间：
（1）发布产物连同 build-id、PDB（GUID 加 age）、dSYM（UUID）归档到符号存储，与镜像摘要一起记录；
（2）在发布配置里就决定帧指针与调试信息的取舍，把它当产品需求而不是编译器随手参数；
（3）采样开关做成运行期可改（配置、环境变量或控制面），并默认在高峰时段关闭；
（4）把\emph{关}状态也当作需要验证的路径——定期验证关掉剖析器后服务仍然正常，否则你的回滚开关从来没被测过。
\end{bestpractice}

\section{开销怎么量：三笔税与一个占空比}
\label{sec:pm14-overhead}

\subsection{唯一可信的口径：同发布、同负载、开与关}
\label{sub:pm14-ab}

生产开销的测法与第~\ref{ch:pm03-stats} 章的比较方法同源，只是对象从两个版本换成同一个版本的两种观测状态：

\begin{itemize}
  \item \textbf{对照组必须是同一个二进制}。版本之间的差异远大于采样开销，任何拿上个版本没采样时的数来比较的做法都不成立。
  \item \textbf{两组要同时、同负载形状}。跨时间比较会把流量形状、缓存冷热、季节波动一起记进剖析器的账；阶跃偏置的处理在第~\ref{ch:pm03-stats} 章。
  \item \textbf{做不到分组就配对交替}。同一实例上开与关交替的多个窗口两两配对，窗口长度取到足够让流量形状相似。
  \item \textbf{报区间不报点值}。开销应当带一个区间以及组数，判据与最小可检出差异回看第~\ref{ch:pm03-stats} 章。区间跨过零时，正确结论是测不出来，而不是没有开销。
\end{itemize}

三笔税各自的测法：

\begin{enumerate}
  \item \textbf{CPU 税}：同一时间窗内处理组与对照组的每核 \MCpu{} 速率差；或者同一实例上关掉剖析器再测一遍总 CPU 时间（第~\ref{ch:pm04-sampling} 章给出的关一次再测口径）。注意后者测不到剖析器对\emph{别人}的干扰。
  \item \textbf{内存税}：\Mx{RSS} 与峰值驻留之差。Linux 读 \Tp{/proc/{\Bk}self/{\Bk}statm} 或 \Tp{smaps}，Windows 看进程工作集，macOS 用 \Fn{vm{\Bk}\_summarize} 一类。要连续取几个窗口而不是开局一次，因为缓冲是逐步填满的。
  \item \textbf{尾延迟税}：\Mx{p99} 及更靠后分位数之差，必须带样本量。小样本上的分位数完全不可信（第~\ref{ch:pm03-stats} 章），而在生产上更容易被忽略，因为排障窗口往往正好样本很少。
\end{enumerate}

\subsection{为什么必须看尾部而不是均值}
\label{sub:pm14-tail}

采样中断的代价\emph{不均}：落在等待时间的样本几乎免费，落在请求关键路径上的样本把中断处理、栈展开、缓冲写入的全部成本加进那一个请求。结果是一个非常典型的形态——\textbf{吞吐与均值几乎不动，高分位数明显移动}。只看均值的开销报告会得出开销可以忽略的结论，而真实用户看到的是超时。所以开销上报至少要同时给三个数：\MCpu{} 增量、\MThr{} 增量、\Mx{p99}（或更靠后的分位数）增量，并写清各自的区间。

\begin{compare}[title={同一次实验，两种标题造成的两种决策}]
\label{box:pm14-overhead-report}
\begin{tabular}{p{0.20\textwidth}>{\raggedright\arraybackslash}p{0.26\textwidth}>{\raggedright\arraybackslash}p{0.26\textwidth}}
\toprule
上报口径 & 排出的句子 & 导致的决定 \\
\midrule
只报均值 & ``开启采样后平均响应时间基本没变'' & 直接常开，频率继续拉高 \\
同时报尾部与容量 & ``均值变化落在噪声区间内；\Mx{p99} 在若干窗口内上移，吞吐余量下降'' & 压低频率或缩窄窗口，把尾部写进告警 \\
\bottomrule
\end{tabular}
\end{compare}

\subsection{停机时间：从占空比到最坏停顿}
\label{sub:pm14-downtime}

停机时间在生产语境里指两个完全不同的东西，混用会得出危险结论：

\begin{enumerate}
  \item \textbf{观测窗口本身占用的时间}（占空比）。它是统计量，摊在时间轴上，用第~\ref{sub:pm14-ab} 节的差分能测出来。
  \item \textbf{剖析器造成的服务停顿}。它是极值量，必须按最坏情况评估：缓冲满时的强制导出、符号解析持锁、attach 时暂停全部线程、采集进程被杀时留下的半状态。这些都不是平均一下就行的量。
\end{enumerate}

第二条的正确问法是：在满载实例上触发一次导出，最坏会让单个请求慢多少？测法：在摘流的实例上主动触发导出，同时用另一路已定义口径的延迟测量记录分布（时钟选择见第~\ref{ch:pm02-clock} 章），并把触发时间点对齐到延迟序列上。任何只测过平均值的停顿数据，都不能用来支持可以在线上触发这个结论。

\subsection{符号化成本放在哪}
\label{sub:pm14-symbolize}

符号化（把指令指针变成函数名与源位置）是三笔税里最容易被忽略的一笔，因为它常常不在被测机器上发生。三种位置，三种价格：

\begin{itemize}
  \item \textbf{进程内即时}：最省磁盘（可以直接输出聚合结果），但把解析成本放进被测进程，还要求符号表就在机器上，体积与启动时间一起上升。
  \item \textbf{同机旁路}：独立采集器读原始样本再符号化，成本从业务进程转移到可控进程；代价是要能访问被采进程的内存映射，权限模型立刻变复杂。
  \item \textbf{离线集中}：只上传原始样本与 build-id，符号在集群解决。生产进程最轻，但前提是第~\ref{sub:pm14-symbols} 节的归档纪律真的可靠，还要先评估样本的隐私面（第~\ref{sub:pm14-privacy} 节）。
\end{itemize}

判据是一条硬纪律：\textbf{热路径上不做任何需要打开文件、遍历符号表或分配增长型内存的事}。它与语言、平台无关；Windows 上按符号路径拉 PDB、macOS 上把 trace 交给工具后处理，本质上都是同一笔账的搬运。

\begin{lstlisting}[style=cppstyle, caption={应用内自采的最小骨架：固定内存环形缓冲与回调自计时（不含实测值；信号上下文的安全性与展开库实现有关，必须实测）}, label={lst:pm14-inproc}]
#include <atomic>
#include <chrono>
#include <cstdint>
#include <cstddef>
#include <vector>

// 编译期封顶的栈深度: 缓冲大小不随样本变化
constexpr std::uint32_t kMaxDepth = 24;

// 展开入口: 是否能在信号上下文里使用取决于实现, 需自行验证
std::uint32_t capture_frames(std::uint64_t* out, std::uint32_t max);

// 固定结构: 无动态分配, 一条样本一次写
struct Sample {
  std::uint64_t monoNs{0};        // 采样时刻, 时钟选择见第 2 章
  std::uint32_t tid{0};
  std::uint32_t depth{0};
  std::uint64_t ip[kMaxDepth];    // 深度封顶
};

class Ring {
public:
  explicit Ring(std::uint32_t cap) : cap_{cap}, slots_{cap} {}

  // 满时覆盖最旧并累加丢弃计数; 丢弃率本身就是指标
  void push(const Sample& s) {
    const std::uint32_t i =
        head_.fetch_add(1u, std::memory_order_relaxed) % cap_;
    slots_[i] = s;                       // 固定大小赋值, 不再分配
    if (count_.load(std::memory_order_relaxed) < cap_) {
      count_.fetch_add(1u, std::memory_order_relaxed);
    } else {
      dropped_.fetch_add(1u, std::memory_order_relaxed);
    }
    writes_.fetch_add(1u, std::memory_order_relaxed);  // 探针自身成本
  }

  // 只在导出路径上调用: 平时零搬运成本
  template <class F> void forEach(F&& f) const {
    const std::uint32_t n = count_.load(std::memory_order_relaxed);
    const std::uint32_t h = head_.load(std::memory_order_relaxed);
    for (std::uint32_t k = 0; k < n; ++k) {
      f(slots_[(h + cap_ - n + k) % cap_]);
    }
  }

  std::uint64_t probes() const { return writes_.load(std::memory_order_relaxed); }
  std::uint64_t drops() const { return dropped_.load(std::memory_order_relaxed); }

private:
  std::uint32_t cap_{0};
  std::vector<Sample> slots_;
  std::atomic<std::uint32_t> head_{0};
  std::atomic<std::uint32_t> count_{0};
  std::atomic<std::uint64_t> dropped_{0};
  std::atomic<std::uint64_t> writes_{0};
};

// 只量探针自己花了多少: 两次读时钟夹住回调体
// 它测不到中断对被中断代码的干扰, 那部分要用第 4.1 节的开关对照
struct ProbeCost {
  std::atomic<std::uint64_t> totalNs{0};
  std::atomic<std::uint64_t> calls{0};
};

void record_once(ProbeCost& cost, Ring& ring) {
  using clk = std::chrono::steady_clock;
  const auto t0 = clk::now();

  Sample s{};
  s.monoNs = static_cast<std::uint64_t>(t0.time_since_epoch().count());
  s.depth = capture_frames(s.ip, kMaxDepth);
  ring.push(s);

  const auto ns = std::chrono::duration_cast<std::chrono::nanoseconds>(
                      clk::now() - t0).count();
  cost.totalNs.fetch_add(static_cast<std::uint64_t>(ns),
                         std::memory_order_relaxed);
  cost.calls.fetch_add(1u, std::memory_order_relaxed);
}
\end{lstlisting}

\begin{warning}[title={回调自计时会系统性低估开销}]
\label{box:pm14-selftime}
清单~\ref{lst:pm14-inproc} 那种在回调里读两次时钟的做法只能测到回调体。它测不到三件事：中断进入与退出的成本、对被打断那条执行路径的流水线与缓存扰动、以及缓冲写入引起的内存序与跨线程伪共享（后者口径见第~\ref{ch:pm09-concurrency} 章）。因此\textbf{探针自计时只用来比较实现方案的相对好坏，永远不用来上报开销}；上报必须用第~\ref{sub:pm14-ab} 节的开与关对照，两个数不一致时以对照为准。
\end{warning}

\section{线上采样与告警联动，以及三平台对照}
\label{sec:pm14-linkage}

\subsection{把分布摘要送进指标系统}
\label{sub:pm14-metrics}

profile 与指标系统解决两个不同问题：指标回答现在有没有不对，profile 回答不对发生在哪个位置。要联动，就得把\emph{分布}而不是平均值放进指标：延迟与关键路径量的分位数摘要（固定桶直方图或可合并 sketch），按服务、版本、机型分组。三条纪律：

\begin{itemize}
  \item \textbf{桶边界事先定}。事后按观测到的分布挑边界，等于把噪声固化成告警线；边界与判据的选择回看第~\ref{ch:pm03-stats} 章。
  \item \textbf{跨实例合并要用可合并的表示}。原始样本不能直接合成分位数，就像跨线程的占用率不能直接相加（第~\ref{ch:pm09-concurrency} 章的歧义在跨机合并里同样出现）。
  \item \textbf{告警指向采样动作，不指向自动结论}。告警触发一次窗口采样是合理的；触发某个函数就是原因不是——那条推理属于第~\ref{ch:pm06-micro} 章的归因工作。
\end{itemize}

trace 与 profile 的接口：单请求路径告诉你一个请求经过哪些位置，profile 告诉你整体上时间被谁占。两者只能对齐\emph{时间与版本}，不能对齐\emph{样本与请求}——除非你愿意每次采样都写请求标识，而那会同时引入隐私面与选择偏差（见下面的 warning）。

\begin{warning}[title={只在慢请求上采样会引入幸存者偏差}]
\label{box:pm14-survivor}
``只在超过阈值的请求上取 profile''看起来省预算，实际破坏统计推断：你拿到的是已知慢的那一部分的栈分布，无法回答慢请求占多大比例、正常请求的分布长什么样、阈值是不是定错了。更糟的是，任何按 profile 占比估计的优化收益都变成以慢样本为分母，系统性偏大。正确做法是\textbf{两层采样并存}：全量低概率采样提供分母与基线，慢样本高概率采样提供细节，并且\emph{两个概率都作为元数据随 profile 上报}，让下游能还原权重。
\end{warning}

\subsection{Linux 侧}
\label{sub:pm14-linux}

生态最厚的一层：\Fn{perf} 一族（record、report、script、annotate）、\Tp{bpftrace} 与 eBPF 聚合、JIT 与解释型运行时的符号补充（\Tp{perf{\Bk}\_map} 与 jitdump 一类）、火焰图流水线（第~\ref{ch:pm04-sampling} 章）。生产上要额外确认三件事：
（1）\Tp{perf{\Bk}\_event{\Bk}\_paranoid} 与内核符号可见性参数的当前值，决定你能看到多少内核侧内容；
（2）容器有没有被 seccomp 与 capability 裁掉 \Tp{perf{\Bk}\_event{\Bk}\_open}；
（3）高频下缓冲区溢出的提示——报告里会明说丢了多少样本，丢了样本就不能算占比。
判定这台机器能做到哪一步的方法在附录~\ref{ch:pmapA-platform}，不要照抄别人的结论。

\subsection{Windows 侧：ETW 实时会话与 XPerf}
\label{sub:pm14-windows}

Windows 的生产采样主要走 ETW（Event Tracing for Windows）：内核 Profile 提供程序按间隔出样本，会话把事件流写进 ETL 文件，事后用 Windows Performance Analyzer 一类工具结合 PDB 解决符号。命令行入口通常是 \Fn{logman} 建会话、\Fn{xperf}（属于 \Tp{XPerf{\Bk}Kit}，需自行安装）直接开停、\Fn{wpr} 用预置 profile。低开销的旋钮同样是三个：采样间隔、开哪个提供程序、会话是否常驻。

\begin{lstlisting}[style=shellstyle, caption={Windows 低占空比采样骨架（形态示意：需管理员权限与会话权限；提供程序名与可用性请在自己的系统上确认）}, label={lst:pm14-win}]
# pwsh; 先把符号路径指到本地缓存与符号服务器
set _NT_SYMBOL_PATH=srv*D:/pm/syms*https://contoso/symsrv

# 短时窗口: 只开 CPU 采样提供程序, 不叠加其他提供程序
xperf -on PROFILE -f D:/pm/win-start.etl
Start-Sleep -Seconds 5
xperf -stop -d D:/pm/win-001.etl

# 想要栈: 展开方式与符号时机在会话侧决定, 代价要另外测
xperf -on PROFILE+CSWITCH -stackwalk Profile
Start-Sleep -Seconds 5
xperf -stop -d D:/pm/win-stack.etl

# logman 建的具名会话便于运行期开关, 也便于确认没有别人已开同类会话
logman create trace Pm14Short -p 'Windows Kernel Trace' profile ^
  -o D:/pm/svc.etl -ets
logman stop Pm14Short -ets

# 用 ETL 的字节数增长量自己估磁盘税: 不要引用别人的速率
(Get-Item D:/pm/win-001.etl).Length
\end{lstlisting}

三个 Windows 特有的口径问题：一是\textbf{符号匹配靠 PDB 的 GUID 与 age，不靠版本号}，构建产物与 profile 必须成对归档，等价于 Linux 的 build-id；二是\textbf{栈展开依赖映像本身}，优化后的二进制与协程、异步续延会让调用链断在框架层，看不穿异步是平台现象而不是你的代码写错（对照讨论见第~\ref{ch:pm04-sampling} 章）；三是\textbf{ETW 会话本身是共享资源}，多个代理同时开会话时开销不是相加而是互相干扰，常驻部署前要先枚举现有会话。计时侧 QPC 与 TSC 的关系见第~\ref{ch:pm02-clock} 章；采样间隔与 QPC 频率不同源这一点，在评估尾部时必须记住。

\subsection{macOS 侧：Instruments 模板、sample 与 sysdiagnose}
\label{sub:pm14-macos}

macOS 生产侧的可用面比前两者窄，但绝不是不能测。可用的四类：\textbf{Time Profiler}（可脚本化的 \Fn{xctrace}，或较老的 \Fn{instruments} 命令；指定模板、时间上限、附加 PID）适合有界窗口的 CPU 采样；\Fn{sample} 适合秒级的一次性调用栈快照，本质就是短时采样；\Fn{sysdiagnose} 是整机诊断包，属于事故现场级手段，产出大、耗时长、需要特权，绝不能当常驻机制；\Fn{powermetrics} 给的是能耗与频率视角，正好补上电与热也是预算这一笔（与硬实时口径的接口见《嵌入式与硬实时》册）。符号侧的对应物是 dSYM 与 UUID，归档纪律与第~\ref{sub:pm14-symbols} 节同一条；权限模型（含系统完整性保护）随版本变，执行前先确认。

\begin{lstlisting}[style=shellstyle, caption={macOS 有界窗口采样与导出骨架（形态示意：权限要求请自行确认；样本里不落业务内容）}, label={lst:pm14-mac}]
#!/usr/bin/env bash
OUT=/var/tmp/pm14-mac
mkdir -p $OUT

# 有界窗口: time-limit 就是停机时间的上界
xctrace record --template 'Time Profiler' ^
  --time-limit 5 s --pid 4242 --output $OUT/win.trace

# 无图形界面时把 trace 导出为可聚合的结构化文本
xctrace export --input $OUT/win.trace --xpath '/*/process-list' > $OUT/procs.xml

# 一次性栈快照: 适合确认现在卡在哪, 不适合做占比统计
sample 4242 3 > $OUT/sample.txt

# 能耗与频率侧的预算视角(通常需要特权, 自行确认)
sudo powermetrics --samplers cpu_power -i 1000 -t 5 > $OUT/power.txt
\end{lstlisting}

三平台在预算这件事上完全同构：都要选采样间隔、都要决定展开方式、都要把符号化搬到热路径之外、都受权限模型限制。差别在权限门槛与生态厚度，不在能不能测。任何这个平台做不到、只有 Linux 能测的结论，通常是还没找到对应旋钮：写结论前先把附录~\ref{ch:pmapA-platform} 的对照逐项找齐。

\begin{lstlisting}[style=outstyle, caption={开销上报的形态示例（格式示例，数值位置以尖括号占位，非本仓库实测值）}, label={lst:pm14-report}]
# 上报模板: 三笔税 + 区间 + 权重, 缺一项就退回重测
service            : ingest
binary build-id    : <id>                 (与 profile 同归档)
window             : 5 s on / 55 s idle, freq 99
groups             : on n=<n1>, off n=<n2> (同发布同负载配对)
cpu tax            : <median> [<lo>, <hi>]
throughput delta   : <median> [<lo>, <hi>]
p99 delta          : <median> [<lo>, <hi>] (小样本要标 n)
rss delta          : <median>              (连续窗口, 非开局单点)
symbolize location : offline (原始样本 + build-id 归档)
dropped samples    : <dropped> / <total>   (占比先于结论)
decision           : 保留 / 降频 / 缩窗口 / 摘流后触发
\end{lstlisting}

\section{排障与速查}
\label{sec:pm14-trouble}

\begin{bestpractice}[title={开线上采样的六项确认}]
\label{box:pm14-checklist-box}
上线前按顺序过六项，任何一项不合格先补齐再开：符号与产物成对归档、权限与内核参数现场确认、采样率与占空比写进配置、熔断与保留期落地、开与关对照的口径提前定义、运行期可关。逐项做法见表~\ref{tab:pm14-checklist}。
\end{bestpractice}

\begin{longtable}{>{\raggedright\arraybackslash}p{0.12\textwidth}>{\raggedright\arraybackslash}p{0.19\textwidth}>{\raggedright\arraybackslash}p{0.28\textwidth}>{\raggedright\arraybackslash}p{0.24\textwidth}}
\caption{开线上采样前的六项确认}\label{tab:pm14-checklist}\\
\toprule
项 & 确认什么 & 怎么确认 & 不合格会怎样 \\
\midrule
\endfirsthead
\toprule
项 & 确认什么 & 怎么确认 & 不合格会怎样 \\
\midrule
\endhead
\bottomrule
\endfoot
\bottomrule
\endlastfoot
符号 & 每个发布产物都能解析 & 列出符号归档目录，按 build-id、PDB 的 GUID 加 age、dSYM 的 UUID 反查得到对应符号文件 & profile 满屏 \Tp{[unknown]}，占比全部不可信 \\
权限 & 这台机器与这个容器能做到哪一步 & 读采样相关内核参数，先跑一个极短窗口并确认样本数不为零 & 采到空数据还以为一切正常 \\
采样率 & 频率与栈深度是否匹配目标函数 & 用两个不同频率各跑一个窗口，看目标函数占比随频率怎么变 & 欠采样把长函数误判成热点，过采样吃掉预算 \\
占空比 & 窗口与静默的比例有上限并被执行 & 配置里能读到这两个数并有校验；导出触发的最坏停顿已在摘流实例上测过 & 高峰时段连续采样导致尾延迟恶化 \\
熔断 & 超限之后是否真的停下来 & 人为把磁盘或内存限制压到超限，观察是否自动停并只降级为无 profile & 采集器把业务实例一起拖垮 \\
对照 & 开销与回滚路径可测可回 & 预先定义开与关的分组方式和上报模板（清单~\ref{lst:pm14-report}），不发版就能关 & 事后无法判断开销来自剖析器还是负载 \\
\end{longtable}

按症状的处置顺序：

\begin{itemize}
  \item \textbf{栈大量 \Tp{[unknown]} 或只剩一层}：先分辨是符号丢了还是展开断了。前者查 build-id 与符号路径；后者查发布构建是否保留帧指针，以及目标代码是不是尾调用或协程边界。两种修法完全不同，别混。
  \item \textbf{样本数或聚合数突然为零}：这不是没有热点，通常是权限被收紧、内核升级、eBPF 程序被拒或容器策略变了。把零样本本身做成告警，比任何指标都灵敏。
  \item \textbf{开启采样后 \Mx{p99} 变差而吞吐没变}：典型的观测者效应组合——中断落在关键路径上、缓冲写入造成跨线程干扰、或导出路径持锁。先缩窗口与降频，再把触发时间点对齐到延迟序列上（第~\ref{sub:pm14-downtime} 节）。
  \item \textbf{样本文件涨得比预算快}：单样本体积由栈深度、是否带地址数据、是否带额外时间戳决定。先量每窗口字节数与样本数（清单~\ref{lst:pm14-win} 最后一行就是最小测法），再决定砍深度还是砍频率。
  \item \textbf{某个函数占比很高但改它没效果}：占比是样本归属的结果，不是因果；采样点偏差（中断落在返回处或循环尾部）会人为放大某些位置。归因路径见第~\ref{ch:pm06-micro} 章，占比的读法见第~\ref{ch:pm04-sampling} 章。
  \item \textbf{同一份 profile 两次分析结论不同}：多半是折叠方式与过滤阈值不同（\Opt{percent-limit}、是否合并内核帧）。把生成报告的完整命令与参数一起归档，否则历史 profile 不可比。
  \item \textbf{只有排障窗口有数据，平时没有}：连续 profile 的占空比拉得太长，或者只有告警触发的采样在工作。回到第~\ref{sub:pm14-always-on} 节：常态窗口必须独立于告警存在。
\end{itemize}

最后把本章的判据收成三句：\textbf{预算先于手段}（第~\ref{sec:pm14-budget} 节）；\textbf{开销只能由同发布同负载的开与关对照给出，并且要报区间与尾部}（第~\ref{sec:pm14-overhead} 节）；\textbf{一份 profile 的价值取决于它能否和产生它的那个二进制一起被读懂}（第~\ref{sub:pm14-symbols} 节）。这三句在任何平台上都成立；各平台把旋钮找齐的具体做法在附录~\ref{ch:pmapA-platform} 与附录~\ref{ch:pmapB-params} 并列给出，按症状的总排障入口在附录~\ref{ch:pmapC-glossary}。

\part{反模式与收束}

\chapter{测量反模式：不可复现的结论是怎么产生的}
\label{ch:pm15-antipattern}

本章不引入任何新工具，也不讲新机制。它做一件相反的事：把前十四章正面讲过的纪律倒过来写一遍。
前文每一条纪律，都有一个对应的、会在真实项目里出现的坏样子——一段被随手贴进正文的 \Tp{\textbackslash begin\{strip\}}、
一句 ``比上个版本快了 30\%''、一个 ``加了关键字就快'' 的补丁说明、一次只在开发机上跑过一遍就进主干的决策。
本章要具体的就是这些 ``样子''：你能在它发生的那一刻认出它。全书地基那三句话（测量先于优化、
口径先于数字、归因先于改动，见 第~\ref{ch:pm01-why} 章）在这里变成一份否定式清单。收束章的用法只有一个：
在你把一个 ``更快'' 写进结论之前，拿本章的每一节对一遍。

这一章也是给前十四章准备的 ``反向目录''。每一节里的反模式都指向某一章的正面讲法：读者不需要重读
全书，只要认出自己犯的是本章哪一条，就能顺着章引用回到那一章去把正面方法重新做一遍。

\section{微基准反模式}

微基准（microbenchmark）就是 ``把一个函数或一小段代码单独圈出来计时''。它最容易出成绩，
也最容易出假成绩，因为编译器会极其主动地把你要测的东西消灭掉。第~\ref{ch:pm02-clock} 章 已经正面讲过
时钟选取、把返回值喂进 sink、反死代码消除；第~\ref{ch:pm03-stats} 章 讲过 warmup 与重复次数的判据；
第~\ref{ch:pm10-compiler} 章 讲过优化级别怎样改变可见行为。本节只把这些正面解法的反面样子列出来。

\subsection{被测工作被优化没了}

最常见的坏样子：你写了一段看起来在测 X 的代码，但 X 的结果没有被任何地方消费，编译器判定它是
死代码，整段删掉；或者你的输入是编译期常量，循环被常量折叠成一条赋值，跑的次数变成一次都没跑。
你测到的其实是空函数的调用开销，然后你以为你在测那个函数。症状很好认：数字 ``小得离谱''，
低到跟你写的那几行代码根本对不上——如果你写了 4 亿次循环，测出来跟什么都没做差不多，那多半就是
这段被优化没了。

下面这段是典型的反例。它名义上在测整数按位求和，实际上一次都没跑：

\begin{lstlisting}[style=cppstyle, caption={反例：输入是常量、结果没被消费，循环可能被整体删除（示意代码，非本仓库实测）}, label={lst:pm15-badloop}]
// BAD: constant input + unused result -> may be folded away
static long sum_of_digits(long n) {
    long s = 0;
    while (n) { s += n % 10; n /= 10; }
    return s;
}

void bench_bad() {
    auto t0 = std::chrono::steady_clock::now();
    for (int i = 0; i < 100000000; ++i) {
        sum_of_digits(1234567890L);   // same input every time
    }
    auto t1 = std::chrono::steady_clock::now();
    // return value is discarded -> call may be eliminated entirely
}
\end{lstlisting}

三件事叠在一起：\Tp{1234567890L} 是常量入参，编译器可以只算一次；返回值没进 sink，
于是那一次也不必算；\Fn{steady\_clock::now} 之间根本没有被测工作残留。修正就是三件事：
输入运行时才知道、结果必须被消费、计时段里只放被测工作。

\begin{lstlisting}[style=cppstyle, caption={修正版：来源可变、结果入 sink、时钟选取与调用开销见 ch02（示意代码，数值由读者自己标定）}, label={lst:pm15-goodloop}]
// GOOD: opaque input, result consumed, sink defeats DCE
int64_t sink;   // module-level: not optimized away

void bench_good(int argc, char** argv) {
    long v = std::atol(argv[1]);          // unknown until run time
    for (long k = 0; k < 16; ++k) {       // vary input: defeat folding
        v = sum_of_digits(v + k) ^ (k * 2654435761u);
    }
    sink = v;                             // consumed -> loop survives
}
\end{lstlisting}

这里的 \Tp{sink} 是最小可用的反死代码消除手段。更严谨的写法（\Fn{DoNotOptimize} 一类，
让优化器 ``以为'' 值被用掉了）在 第~\ref{ch:pm02-clock} 章 里讲过怎么用。核心判据只有一条：
如果你不能向自己解释 ``编译器为什么删不掉这段''，那这段的计时数字就没有意义。

\subsection{把取时钟、打日志算进被测工作}

第二个坏样子：把 \Fn{printf}、日志、或者取时钟本身放进计时段里。一次 \Fn{clock\_gettime}
通常在几十纳秒量级（具体取决于是否走 vDSO 与内核版本，见 第~\ref{ch:pm02-clock} 章）；一次带缓冲冲刷的
\Fn{printf} 可以到微秒级、并触发 I/O。你想测的那个 ``单函数'' 如果本来就在这个量级，
计时与打印会把它彻底淹没。正面口径留给 第~\ref{ch:pm02-clock} 章：只把被测工作夹在两次取时钟之间，
时钟调用与输出一律在计时段外，或者用 Google Benchmark 一类的注册式框架来保证顺序。

\subsection{``单函数多少纳秒'' 为什么推不出架构结论}

假设你真的量出了 ``这个函数 $2$ 纳秒''，然后据此下结论说 ``这块逻辑在热路径上，
值得为它重写数据结构''。错在哪？第~\ref{ch:pm01-why} 章 算过这笔账：收益 = 单次节省 × 调用次数 ×
是否在关键路径上。一个 $2$ 纳秒的函数，如果每秒只被调几千次，或者它根本不在关键路径上，
那它省下来的时间在整段 \MWall{} 里就是噪声。微基准只给你 ``单次节省'' 这一项，而且是在最理想
条件下（数据在 L1、分支全预测对）给——这两项恰恰是 第~\ref{ch:pm06-micro} 章 讲的缓存与分支前提。
所以 ``单函数多少纳秒'' 永远是必要条件而非充分条件：它不能单独变成架构结论，必须放回真实调用
频率与关键路径里再判一次。

\begin{guideline}[title={把 ``数是不是真的'' 这个问题前置}]
一个微基准只有在你同时能回答这三问时才可信：(1) 编译器为什么删不掉被测工作（sink、不透明输入）？
(2) 计时段里除了被测工作还有没有取时钟、日志、分配（见 第~\ref{ch:pm02-clock} 章）？
(3) 这个数放回调用频率与关键路径还剩多少（见 第~\ref{ch:pm01-why} 章）？三问有一个答不上来，
先别把它写进结论。
\end{guideline}

\section{观测者效应}

剖析器不是一面旁观的镜子，它本身就是一只伸进系统里的手。你每加一种观测，都在改变被观测对象的行为。
第~\ref{ch:pm04-sampling} 章 教过怎么量 ``采样自身开销''（关掉剖析器再测一次），第~\ref{ch:pm14-production} 章
把 ``观测预算'' 当一等指标；本节讲这只手具体改了什么、以及最常见的两种错觉。

\subsection{剖析器改变行为的四条通道}

第一，页缓存。剖析往往让程序跑得更慢、更久，于是同样的工作被重复做时命中缓存的概率变了，
冷/热差异被观测本身放大——第~\ref{ch:pm08-io} 章 专门讲过 ``第二次跑当然快'' 这个陷阱，加上剖析器
只会让它更明显。第二，内联决策。为了拿到完整调用栈你会用 \Opt{no-omit-frame-pointer} 重编，
而帧指针的存在会改变编译器的内联与寄存器分配，你测的已经不是发布那个二进制
（第~\ref{ch:pm10-compiler} 章）。第三，锁的时序。采样中断线程会拉长临界区，把一个本来不明显的锁竞争
放大成假热点，或者反过来把竞争打散（第~\ref{ch:pm09-concurrency} 章）。第四，频率爬升。
剖析带来的额外负载与停顿会影响 DVFS 决策，短基准尤其敏感（第~\ref{ch:pm02-clock} 章）。

四条通道有一个共同后果：你测的不再是 ``那个发布二进制的真实表现''，而是 ``发布二进制 + 剖析器''
的合成物。合成本身不一定坏，坏的是你把它当成前者去下结论。

\subsection{为什么加了采样后 \Mx{p99} 变差而吞吐没动}

这是最典型、也最容易被误读成 ``性能回归'' 的形态。定时采样在每次中断里要展开栈、记样本，
这笔开销摊到少数被中断打中的请求上，就是尾延迟：\Mx{p99} 里那批恰好 ``带了一份剖析器成本''
的请求变慢了。而 \MThr{} 看的是总体完成量——一次中断几十微秒量级，在总 CPU 工作量里占比极小，
吞吐基本不动。于是你看到 ``\Mx{p99} 涨了、\MThr{} 没变''，这不是被测代码变慢了，
是你把观测开销加进了尾部分布。

正确的口径是 第~\ref{ch:pm14-production} 章 的 ``开/关剖析器同机同负载各测一次再相减''，
把观测者的账记在观测者头上。相减之前还得过 第~\ref{ch:pm03-stats} 章 那关：差值要大于 run-to-run
变异才算 ``看得见''。

\begin{lstlisting}[style=outstyle, caption={形态示意（非本仓库实测值）：开/关剖析器对照的形状，重点在 p99 动而吞吐不动；尖括号是占位符}, label={lst:pm15-observer}]
# SHAPE EXAMPLE ONLY -- not measured on this repo's CI runner (no PMU access)
# Fill in the placeholders yourself on a machine where sampling is allowed.
run                 p99_lat     throughput
profiler OFF        <p99_off>   <thr_off>
profiler ON         <p99_on>    <thr_on>
delta (ON - OFF)    +obs tax    ~flat
# read this shape as: p99 gap is largely the sampler's own cost;
# subtract ON minus OFF on same host + same load before blaming the patch
\end{lstlisting}

清单里的 \Tp{<...>} 是占位符，不是数据。本仓库 CI runner 上没有 PMU 权限（第~\ref{ch:pm13-ci} 章
明说过），这类表只能在读者自己那台能采样的机器上填。判据不变：先证明差值大于 run-to-run 变异
（第~\ref{ch:pm03-stats} 章），再谈它是不是被测代码造成的。

\section{口径偷换}
\label{sec:pm15-antipattern-caliber}

口径先于数字，这是 第~\ref{ch:pm01-why} 章 立册的第二句话。反模式不是有人算错，而是标题换了口径、
数字还是那个数字，读者脑补出的结论却完全不同。四类最常见的偷换各自回指它被正面讲过的章：

把 \MCpu{} 当 \MWall{}：多线程程序 \MCpu{} 可以远大于 \MWall{}（第~\ref{ch:pm02-clock} 章），
``总 CPU 时间降了'' 不等于 ``用户等得短了''，尤其在被调度的进程里两个口径的的相关性很弱。
把 \MThr{} 当延迟：吞吐高可能只是缓冲足够大、把延迟搬进了队列（第~\ref{ch:pm08-io} 章、
第~\ref{ch:pm09-concurrency} 章），一次请求照样可能很慢——第~\ref{ch:pm09-concurrency} 章 里那张
``无界队列变快的假象'' 对照就是这个偷换的标准形态。把峰值当容量：带宽的理论峰值、STREAM 式实测、
应用实际是三种数（第~\ref{ch:pm12-bandwidth} 章），拿峰值报 ``我们有余量'' 会在真实混合负载上翻车。
把 ``跑得快的那次'' 当代表：只报最好的一次，等于把分布的左尾当成整体（第~\ref{ch:pm03-stats} 章）。

\begin{compare}[title={同一次实验，四种标题写法 = 四种决策}]
\begin{tabular}{@{}>{\raggedright\arraybackslash}p{0.30\textwidth}>{\raggedright\arraybackslash}p{0.46\textwidth}@{}}
\toprule
标题写成 & 读者会做出的决策 \\
\midrule
{``\MCpu{} 降 8\%''} & ``省了 CPU，云成本会降，合入。'' 但没人问 \MWall{} 有没有变。 \\
{``\MThr{} 升 8\%''} & ``能扛更多了，扩容推迟。'' 但 \Mx{p99} 可能同时变差。 \\
{``\Mx{p99} 降 8\%''} & ``尾延迟改善，用户体验更好。'' 但样本量不足时这个 p99 不可信（第~\ref{ch:pm03-stats} 章）。 \\
{``最好一次降 8\%''} & ``确实更快了。'' 而这可能只是热缓存运气，重跑就没了。 \\
\bottomrule
\end{tabular}
\end{compare}

同一个 8\%，四种口径就四种行动，其中三种可能都是错的——这就是口径偷换最贵的一面。
收束方式很简单：结论里第一个词必须是口径，并且这个口径要和它服务的真实提问匹配
（第~\ref{ch:pm01-why} 章 那张 ``三种真实提问'' 表就是拿来对齐这个的）。

\section{归因偷换}
\label{sec:pm15-antipattern-cause}

指标把你送到代码附近，但 ``在这附近改了它、于是变快了'' 需要另一套证据。归因偷换就是位置对了、
原因错了。这一节是最贵的一节：口径偷换骗的是读者，归因偷换骗的是你自己，而且骗完之后你会带着
错误的因果去指导下一批决策。第~\ref{ch:pm06-micro} 章 与 第~\ref{ch:pm05-events} 章 给的是正面归因链
（从 \Mx{cache-miss} 到布局、从 \Mx{branch-miss} 到数据依赖分支、TopDown 的四桶），
本节讲这条链断掉的三条常见路径与一个最顽固的具体反模式。

第一，同时改了两件事。补丁里既有算法改动又顺手升级了编译器或开了 \Tp{-O3}，变快的是后者，
你归功于前者。判据：任何 ``两件事一起改'' 的补丁都必须在能拆开的粒度上分别测一次，
否则你只知道 ``这一堆改动'' 有效，至于哪一半有效不知道。第二，只测了一个负载。你手上的样例
恰好命中了新代码的快路径，换一个分布就退化，而 ``只测一个负载'' 正是 第~\ref{ch:pm03-stats} 章
反复警告的单样本问题。第三，环境漂移。测试前后机器频率、缓存占用、同宿主机邻居变了，
这是 第~\ref{ch:pm03-stats} 章 里的阶跃偏置，在 CI 里则是 第~\ref{ch:pm13-ci} 章 讲过的基线失效——
基线失效在报回归时最容易伪装成 ``补丁变慢了''。

\subsection{内联不是提速开关}

把 ``我加了 \Tp{inline}，所以这段快了'' 当成因果，是归因偷换里最顽固的一种。这里的口径必须与
《代码风格参考手册》完全一致：\Tp{inline} 首先改变的是 ODR 与链接语义——它是给编译器的一个提示，
不是 ``必须展开'' 的指令；实际是否展开由内联启发式、优化级别、LTO、PGO、函数体积与调用点热度
共同决定（见 第~\ref{ch:pm10-compiler} 章）。真正硬保证的是 \Tp{noinline}，它保证不展开；
而 \Tp{always\_inline}、\Tp{\_\_forceinline} 只是尽力而为——GCC 在约束无法满足时可能报错，
Clang 可能只给警告，MSVC 的文档也明言不保证。

于是任何 ``加 \Tp{inline} 就快'' 的写法都是错的：加了 \Tp{inline} 变快，
要么根本没展开、变快的是别的改动（更可能是常量折叠或循环不变量外提）；要么真展开了，
但 ``展开'' 与 ``变快'' 之间没有必然关系（更大的函数体可能挤爆 I-cache，见 第~\ref{ch:pm06-micro} 章）。
唯一的确认方式是在实际发布配置的二进制上反汇编、数那个调用还在不在，探针构建上看到的数不作数
（第~\ref{ch:pm10-compiler} 章）。

\begin{lstlisting}[style=shellstyle, caption={在发布产物上确认展开是否发生（示意命令，非本仓库实测；开关名与工具随版本变，先查自己那套 objdump/nm 支持哪些）}, label={lst:pm15-inline-check}]
# Confirm: did the call actually disappear in the release binary?
# (1) keep the exact release binary, do NOT rebuild a probe
objdump -d ./release_binary | grep -E 'call.*\bhot_fn\b'
nm -C ./release_binary | grep -E '\bhot_fn\b'
# (2) expected: no 'call ... hot_fn' for the inlined hot path;
#     if the call survives, your 'inline' hint was ignored.
# (3) for MSVC build, use dumpbin /disasm instead; for macOS, use
#     otool -tV. Tool availability depends on your SDK version -- check.
\end{lstlisting}

这个清单本身不告诉你 ``快了多少''，它只把 ``到底有没有展开'' 变成可核对的事实。展开之后是不是更快，
还得按 第~\ref{ch:pm10-compiler} 章 的产物核对四步与 第~\ref{ch:pm03-stats} 章 的判据再走一遍。

\begin{warning}[title={补丁有效，因为回滚更慢}]
``我把改动 revert 回去它就又慢了，所以是我的改动带来了加速。''——这条推理常常错，
因为回滚的那次运行本身也在漂移：热缓存已经形成、频率已经爬上去、编译器可能因源文件顺序变化
给出不同的内联决策。回滚不等于回到起点，它是在一个新环境里跑旧代码。正确做法是 A/B/A/B 交替跑，
让漂移被对称地摊到两侧（第~\ref{ch:pm03-stats} 章），而不是先测新、再测旧。若一次交替跑不出稳定的
差值，那 ``差值'' 本来就是漂移给你的错觉，而不是补丁的功劳。
\end{warning}

\section{不可复现的结论长什么样}
\label{sec:pm15-antipattern-unrepro}

一个不会活过下周的结论，通常在这几处缺料：没记环境指纹（机器、内核、CPU 型号、频率 governor，见
第~\ref{ch:pm02-clock} 章）、没记编译器与构建口径（\Tp{-O} 级别、是否 LTO、\Tp{-DNDEBUG}，
见 第~\ref{ch:pm10-compiler} 章）、只有一张截图、只有均值没有分布、只有最好的一次。截图尤其危险：
它没有原始序列、没有重复次数、没有单位，别人既无法复现也无法证伪。

正面解法在 第~\ref{ch:pm02-clock} 章 给过一段采集环境指纹的清单，这里给它的否定式对应物——一份
``结论体检表''，每一项都问一句、并规定不合格时结论该怎么改写。拿它逐行审你即将发出去的那段话。

\begin{longtable}{@{}>{\raggedright\arraybackslash}p{0.20\textwidth}>{\raggedright\arraybackslash}p{0.32\textwidth}>{\raggedright\arraybackslash}p{0.30\textwidth}@{}}
\caption{结论体检表：一项一问，不合格就按右列改写}\label{tab:pm15-audit}\\
\toprule
检查项 & 合格长什么样 & 不合格怎么改写 \\
\midrule
\endfirsthead
\toprule
检查项 & 合格长什么样 & 不合格怎么改写 \\
\midrule
\endhead
\bottomrule
\endfoot
口径 & 写明是 \MWall{}、\MCpu{}、\MThr{} 还是 \Mx{p99}，与真实提问匹配 & 把 ``变快了'' 改成 ``在 \Mx{p99} 上快了'' \\
环境指纹 & CPU 型号、内核、governor、频率、并发度、内存配置都记录 & 补一行 ``某型号 CPU / 某内核 / performance governor'' \\
构建口径 & 编译器版本、\Tp{-O} 级别、LTO/PGO、\Tp{-DNDEBUG} 都写明 & 标注 ``用 \Tp{-O2} 无 LTO，发布配置另测'' \\
判据 & 重复次数、统计量、最小可检出差异、置信区间都写出 & 补 ``重复 N 次，差值大于 run-to-run 变异'' \\
数据来源 & 原始序列与分布随结论留存，不是截图 & 有截图处补上可复跑的日志或数据文件 \\
样本形状 & 报中位数与分位数，不只均值、不只最好一次 & 把 ``最好 8\%'' 改成 ``中位数 8\%，p90 另列'' \\
结论边界 & 写明在哪个负载、并发度、硬件上成立 & 加 ``仅在单线程冷缓存成立，多核待测'' \\
复现步骤 & 一段可粘贴的命令行，别人能跑出同形数据 & 补齐采样命令、参数、warmup 与丢弃规则 \\
\end{longtable}

体检表的每一行都对应本章前面某一节的正面口径。它的作用不是让结论更漂亮，而是让结论可以被证伪——
一份能被证伪的结论才是能活的结论。

\begin{lstlisting}[style=shellstyle, caption={可复现结论的最小环境指纹采集骨架（不含实测值，读者在发布机上运行）}, label={lst:pm15-fingerprint}]
# capture once, paste alongside every perf claim
uname -a                                        >  env.txt
lscpu | grep -E 'Model name|Mhz|Flags'        >> env.txt
cat /sys/devices/system/cpu/cpu0/cpufreq/scaling_governor >> env.txt
${CC} --version                               >> env.txt
git rev-parse HEAD                            >> env.txt
sha256sum ./release_binary                    >> env.txt
# env.txt is now part of the claim, not an afterthought
\end{lstlisting}

这份清单的每一项对应体检表里的一行。它的用途是把 ``我测的到底是什么'' 钉死成文字，
这样 第~\ref{ch:pm13-ci} 章 里说的 ``环境变了就重建基线'' 才有可比对象——没有指纹，你连 ``变了''
都判定不了。

\section{怎么写一个能活的结论}

把前面几节的否定式倒过来，就是一个能活过评审、活过下周、活过换人的结论。它的骨架是一条固定顺序：
\textbf{口径 $\rightarrow$ 环境指纹 $\rightarrow$ 判据 $\rightarrow$ 数据 $\rightarrow$
结论边界 $\rightarrow$ 复现步骤}。注意每一步都在收前面一章的口：口径来自 第~\ref{ch:pm01-why} 章，
环境指纹来自 第~\ref{ch:pm02-clock} 章，判据来自 第~\ref{ch:pm03-stats} 章，数据的形状纪律来自
第~\ref{ch:pm04-sampling} 章 与 第~\ref{ch:pm05-events} 章，归因边界来自 第~\ref{ch:pm06-micro} 章，
构建口径来自 第~\ref{ch:pm10-compiler} 章，长期存活来自 第~\ref{ch:pm13-ci} 章 与 第~\ref{ch:pm14-production} 章。

一次评审里，你至少该问三个问题，它们能挡掉本章大多数反模式：

\begin{itemize}
  \item ``这个数是什么口径，和它要回答的那个提问是同一个吗？''——挡口径偷换（\S\ref{sec:pm15-antipattern-caliber}）。
  \item ``同一份代码换台机器、换个负载还成立吗？你怎么知道的？''——挡归因偷换与单负载
        （\S\ref{sec:pm15-antipattern-cause}）。
  \item ``给我复现步骤，我跑出同形的数再说。''——挡不可复现：截图、只有均值、最好一次
        （\S\ref{sec:pm15-antipattern-unrepro}）。
\end{itemize}

\begin{bestpractice}[title={能活的结论模板}]
按六段写死：(1) 口径——一句话说明测的是 \MWall{}、\MCpu{}、\MThr{} 还是 \Mx{p99}，
以及它回答的问题；(2) 环境指纹——CPU、内核、governor、编译器版本、二进制哈希
（用 \ref{lst:pm15-fingerprint} 采集）；(3) 判据——warmup 规则、重复次数、统计量、
最小可检出差异，先注册再跑；(4) 数据——原始序列或分位数摘要，不用截图；(5) 结论边界——
``仅在 X 负载 / Y 并发 / Z 硬件成立''；(6) 复现步骤——可粘贴命令。缺任何一段，
把结论降级成 ``待确认''，不要发。
\end{bestpractice}

六个顺序不能换。口径先于数字，判据先于跑，边界先于发。这三条本来就是全书前十三章反复讲的东西，
本节只是把它们排成一条能贴到 PR 描述里的模板。

\section{与前面各章的收束：按症状索引全书}

本章是全册的否定式目录。反过来，它也是一张 ``按症状进来、去对应章拿正面解法'' 的索引。
下表把读者最常撞上的现象，映射回正面讲过它的章，并给第一步该做什么。章号全部用章级标签引用，
目标全部是注册表里那 18 个 label。

\begin{longtable}{@{}>{\raggedright\arraybackslash}p{0.32\textwidth}>{\raggedright\arraybackslash}p{0.13\textwidth}>{\raggedright\arraybackslash}p{0.36\textwidth}@{}}
\caption{按症状回指全书：症状 $\rightarrow$ 正面章 $\rightarrow$ 第一步}\label{tab:pm15-index}\\
\toprule
你看到的症状 & 回指章 & 第一步做什么 \\
\midrule
\endfirsthead
\toprule
你看到的症状 & 回指章 & 第一步做什么 \\
\midrule
\endhead
\bottomrule
\endfoot
数字小得离谱、像是没干活 & 第~\ref{ch:pm02-clock} 章 & 检查被测工作是否被优化没（sink、不透明输入） \\
两次跑出百分之几十的差 & 第~\ref{ch:pm03-stats} 章 & 先量 run-to-run 变异，别急着比点值 \\
热点是 \Tp{memcpy\_erms} 之类 & 第~\ref{ch:pm04-sampling} 章 & 先确认能不能拿到完整栈，再谈热点 \\
高 \Mx{IPC} 但就是不快 & 第~\ref{ch:pm05-events} 章 & 走 TopDown 分桶，别盯单事件 \\
换个遍历顺序就差很多 & 第~\ref{ch:pm06-micro} 章 & 查局部性与布局，再看分支可预测性 \\
内存一直涨就是不敢删 & 第~\ref{ch:pm07-memory} 章 & 增长不等于泄漏，用三步判据先定性 \\
第二次跑当然快 & 第~\ref{ch:pm08-io} 章 & 区分冷/热运行，问缓存有没有参与 \\
加了线程反而慢 & 第~\ref{ch:pm09-concurrency} 章 & 按顺序查串行比例、伪共享、调度迁移 \\
改了代码二进制没变 & 第~\ref{ch:pm10-compiler} 章 & 产物核对四步：符号 / 反汇编 / 大小 / 链接顺序 \\
GPU 段测出接近 $0$ & 第~\ref{ch:pm11-gpu} 章 & 检查是不是漏了同步，只测到排队 \\
带宽利用率只有六成 & 第~\ref{ch:pm12-bandwidth} 章 & 对齐理论峰值 / 实测 / 应用实际三种数 \\
CI 红三次就想调阈值 & 第~\ref{ch:pm13-ci} 章 & 先查环境漂移，不是改判据 \\
线上开采样后 \Mx{p99} 变差 & 第~\ref{ch:pm14-production} 章 & 开关剖析器同机同负载各测一次再相减 \\
结论只有截图或只有均值 & 第~\ref{ch:pm15-antipattern} 章 & 用结论体检表逐行改写，缺料就降级 \\
需要三平台工具与权限对照 & 附录~\ref{ch:pmapA-platform} & 去附录 A 找该平台能用的采样手段 \\
指标 / 事件 / 参数口径拿不准 & 附录~\ref{ch:pmapB-params} & 附录 B 的指标字典与参数分组 \\
不认识缩写或要按症状排障 & 附录~\ref{ch:pmapC-glossary} & 附录 C 术语表与按症状排障手册 \\
\end{longtable}

这张表就是本章存在的意义：当你怀疑某个 ``更快'' 站不住时，先在这里找到它的症状，
再回到那一章用正面的方法重测。测得准，从来不是靠更好的直觉，而是靠更少的自欺。

\begin{note}[title={三平台与 sanitizer 的口径收束}]
本章的反模式在三平台上都成立，具体手段随平台换：Linux 走 perf 一族（第~\ref{ch:pm04-sampling} 章、
第~\ref{ch:pm05-events} 章），Windows 走 XPerf/ETW/VTune（第~\ref{ch:pm14-production} 章 有对照小节），
macOS 走 Instruments 与 \Tp{powermetrics}。完整的三平台并列留给附录~\ref{ch:pmapA-platform}，
不要写 ``只有 Linux 能测''。sanitizer 相关的用法一律回指《内存泄漏检测与防御》册与
《LSP 与代码质量工具》第 12 章；本册只讲它们作为测量对象的那一面——开 sanitizer 跑出来的数
绝不是性能数。
\end{note}

\appendix
\chapter{三平台工具、权限与事件对照速查}
\label{ch:pmapA-platform}

正文各章以 Linux + \Tp{perf} 为主线，Windows 与 macOS 只在关键章各给一节对照；本附录把它们并到
同一张桌面上。这里不重复机制讲解——时钟语义在第~\ref{ch:pm02-clock} 章，采样与栈展开在
第~\ref{ch:pm04-sampling} 章，事件与计数器复用逻辑在第~\ref{ch:pm05-events} 章。本附录只回答四个
操作问题：这台机器上有什么工具、要什么权限、怎么现场确认、输出丢了什么。

三条纪律贯穿全表，写作与查阅时都不要破：

\begin{itemize}
  \item 工具名与开关随发行版、SDK、Xcode 版本变。本附录刻意不写版本号，因为写了就会过期；
        每一格都给``自己列一遍''的命令。
  \item 事件表一律不抄。硬件事件名由内核 PMU 驱动与 CPU 型号决定，同一代处理器换个内核就换
        写法；能背的只有``怎么发现''，不是``叫什么''。
  \item 权限结论必须自己跑一遍自查命令才知道。别人说``这台机器能采内核''不算证据，
        \Tp{perf stat -e cycles true} 的输出才算。
\end{itemize}

\section{能力总对照：同一件事在三个平台上叫什么}

表~\ref{tab:pmapA-tool} 按``能力''而不是按``工具''组织：左列是你想问的问题，三列给出该平台上
最常用的手段、权限前置与输出形态。每格的读法是``手段；权限；输出''。权限栏只写通常要求什么，
具体到某台机器仍以第~\ref{sec:pmapA-perm} 节的自查为准。

\begin{longtable}{>{\raggedright\arraybackslash}p{0.16\textwidth}
                  >{\raggedright\arraybackslash}p{0.22\textwidth}
                  >{\raggedright\arraybackslash}p{0.22\textwidth}
                  >{\raggedright\arraybackslash}p{0.20\textwidth}}
\caption{\label{tab:pmapA-tool}同一测量能力在 Linux / Windows / macOS 上的工具、权限与输出形态并列。
所有工具名与开关的可用性都要现场确认。}\\
能力 & Linux & Windows & macOS \\
\hline
\endfirsthead
\hline
能力 & Linux & Windows & macOS \\
\hline
\endhead
CPU 定时采样 & \Tp{perf record}；用户态可自采，采内核要放开档；二进制流经 \Tp{perf report} 出文本 &
\Tp{wpr} 与 \Tp{xperf} 的 CPU profile；内核会话要管理员；产出 \Tp{.etl} 在 WPA 里看 &
\Tp{xctrace} 的 Time Profiler 或 \Tp{sample}；对别人进程要提权；产出 \Tp{.trace} 包 \\
\hline
调用栈展开 & 帧指针 / DWARF 解卷 / LBR 三条路；DWARF 要 \Tp{-g}，LBR 要 PMU 权限 &
帧链依赖保留帧指针的构建；异步与协程栈需专门事件；\Tp{.etl} 内的 Stack 表 &
\Tp{xctrace} 记用户栈需未 strip；内核栈要提权；\Tp{.trace} 内含 Call Tree \\
\hline
硬件事件计数 & \Tp{perf stat}；虚拟化与容器里常只剩软件事件；文本摘要 &
VTune 与 \Tp{xperf} 的 PMU profile；要管理员且宿主暴露计数器；报表或 CSV &
Instruments 的 Counters 类模板；Apple 芯片与 Intel 事件集完全不同；图形报表 \\
\hline
分层归因 & 事件组或 \Tp{perf} 的指标组（名字随版本变，现场 \Tp{perf list}） &
厂商分析器的 TopDown 类分析类型 & 无直接等价物，靠分桶事件自己拼 \\
\hline
分配与堆 & heaptrack、\Tp{malloc\_stats} 族、Callgrind 族；开销大，见第~\ref{ch:pm07-memory} 章 &
ETW 堆事件加 WPA 的 Heap 视图 & Allocations 与 Leaks 模板，配环境变量打开栈记录 \\
\hline
磁盘与文件 I/O & \Tp{iostat}、\Tp{blktrace}、块层 tracepoint &
\Tp{xperf} 的 disk 与 file profile，WPA 的 Storage 视图 & File Activity 模板、\Tp{fs\_usage} \\
\hline
网络与管道 & \Tp{ss}、\Tp{nstat}、\Tp{iperf} 族 & ETW 网络 provider & Networks 模板、\Tp{nettop} \\
\hline
调度与等待 & \Tp{perf sched} 与 \Tp{/proc} 一族 & 就绪与切换事件，Thread Time 视图 &
System Trace 模板的 run queue \\
\hline
锁与竞争 & \Tp{perf c2c} 类（冲突事件名要现场确认） & WPA 的 Contention 视图 &
无 PMU 层等价物，靠系统调用与锁模板侧面看 \\
\hline
频率与功耗 & \Tp{turbostat}、\Tp{cpufreq} 目录 & ETW 的 Power profile & \Tp{powermetrics}，要 root \\
\hline
GPU 时间线 & Nsight 族 & 厂商分析器的 GPU 分析类型 & Metal 与 GPU Driver 类模板 \\
\hline
火焰图导出 & \Tp{perf script} 接折叠脚本，纯文本流水线 &
需先把 \Tp{.etl} 转成栈列表，官方无一等支持 & 导出 XML 再转换，或直接看 Call Tree \\
\hline
\end{longtable}

最后一列不是``macOS 更弱''的意思：Instruments 在等待原因、帧率与 GPU 时间线上的信息密度常超过
另两家，差别在可脚本化程度与权限模型，不在能不能测。

\section{权限模型：三平台各自怎么卡人}
\label{sec:pmapA-perm}

三个平台的卡点位于不同层，所以自查命令完全不同：Linux 卡在一个可调整的整数加一组 capability；
Windows 卡在你是否被提升、是否属于日志用户组，并且内核会话是独占资源；macOS 卡在 SIP 与设备
结点的访问控制。

\subsection{Linux：paranoid 档位、capability 与容器连带}

核心开关是 \Tp{/proc/sys/kernel/perf\_event{\Bk}\_paranoid}。取值从 $-1$ 到 3（较新内核还有更细的
数据地址采样限制档），每一档的准确语义以你本机内核的 \Tp{perf\_event{\Bk}\_open(2)} 手册页为准。
实践上的阶梯是：档位较高时只能测自己启动进程的用户态事件且采不到内核符号；常见发行版默认档
允许用户态剖析、禁止内核剖析；档位放开后 raw tracepoint 与 kprobe/uprobe 还需要
\Tp{CAP\_PERFMON}（内核 5.8 起从 \Tp{CAP\_SYS\_ADMIN} 拆出的能力位）或 \Tp{CAP\_SYS\_ADMIN}。
另外，附加到\textbf{别人的}进程受 ptrace 语义约束（如 \Tp{yama/ptrace{\Bk}\_scope}），两个条件
任一不满足都会失败，而报错常常只提第一个。内核符号地址另有独立开关 \Tp{kptr\_restrict}，它
不影响``能不能采''，只影响报告里内核帧``有没有名字''——这两件事混着排查会绕远路。

\begin{lstlisting}[style=shellstyle, caption={Linux 权限自查：逐条可复制执行，输出决定后面能做什么}, label={lst:pmapA-linux-perm}]
# 1) 当前档位。注意：它不是容器命名空间的值，容器里读到的就是宿主机的
cat /proc/sys/kernel/perf_event_paranoid

# 2) 内核符号地址可见性（决定报告里内核帧有没有名字）
cat /proc/sys/kernel/kptr_restrict
# 3) 附加到别人进程是否被 ptrace 层先拦掉
cat /proc/sys/kernel/yama/ptrace_scope 2>/dev/null

# 4) 本机到底有没有 perf，以及它自己怎么说
perf --version
perf list 2>&1 | head -n 20
# 5) 决定性一步：硬件事件到底可用不可用
perf stat -e cycles,instructions true

# 6) 对照：软件事件应当在任何档位下都能用
perf stat -e task-clock,context-switches,cpu-clock,major-faults true
# 7) 容器里的能力位；没有 capsh 就直接看掩码
grep -E 'CapEff|CapBnd' /proc/self/status
capsh --decode=<上一行的十六进制> 2>/dev/null
\end{lstlisting}

第 5 与第 6 条要一起看：``只有软件事件可用''是虚拟化与容器环境的常态，不是配置错误。

\begin{warning}[title={改全局档位不是解决办法，而且常常改不动}]
\begin{itemize}
  \item \Tp{kernel.perf\_event{\Bk}\_paranoid} 不是命名空间化的 sysctl。容器里 \Tp{sysctl -w}
        要么失败，要么改到\textbf{宿主机全局}，影响同机所有租户。共享机器上这不是你可以动的旋钮。
  \item 宿主机放开了档位，也挡不住 seccomp 或默认能力集把 \Fn{perf\_event{\Bk}\_open} 直接拒掉。
        为测一把就把生产容器改成特权模式，是把测量成本转嫁成安全成本。
  \item 正确顺序：先确认软件事件够不够用（第~\ref{ch:pm05-events} 章），再决定要不要申请一台
        放开权限的专用机，而不是反过来。
\end{itemize}
\end{warning}

\subsection{Windows：管理员、日志用户组与内核会话的独占}

\begin{itemize}
  \item 内核会话（\Tp{xperf -on} 与 \Tp{wpr -start} 的大多数 profile）要求管理员，而且是
        被 UAC 提升后的进程。
  \item 同一时刻系统里只能有\textbf{一个}内核会话。安全软件或诊断收集器占着时起不来，报错
        不会告诉你会话在谁手里——先看 \Tp{wpr -status} 与 \Tp{logman query -ets}。
  \item 不少用户态会话与计划跟踪只要求属于 Performance Log Users 组，不必是完整管理员。
  \item provider 的 verbose 档位可能带业务数据，部分 provider 在非提升进程里根本开不到该级别。
  \item 硬件计数器（厂商分析器与 \Tp{xperf} 的 PMU profile）要求管理员，且在多数云虚拟机与
        启用了虚拟化安全特性的宿主上不可用。
\end{itemize}

\begin{lstlisting}[style=shellstyle, caption={Windows 权限与可用性自查（工具本身要先装，可用性随 SDK 组件变）}, label={lst:pmapA-win-perm}]
:: 1) 是不是管理员、有没有被提升。组 SID 尾号 544 是 Administrators，559 是 Performance Log Users
whoami /groups | findstr /I S-1-5-32
net session >nul 2>&1 && echo ELEVATED || echo NOT-ELEVATED
:: 2) 内核会话是否已被别人占着——起 trace 之前必看
wpr -status
logman query -ets

:: 3) 现场列出可用 provider 与 profile，别背清单
logman query providers > providers.txt
wpr -profiles
xperf -providers
:: 4) 决定性一步：起一个最小内核会话并立刻收尾
xperf -on base
xperf -d trace.etl
:: 跑不通即权限不够；跑通了务必确认 -d 已经把会话关掉
\end{lstlisting}

\subsection{macOS：SIP、dtrace helper 与提权边界}

卡点分三层：SIP 决定能不能对受保护的二进制做内核级观测；\Tp{/dev/dtracehelper} 与 dtrace 策略
决定普通用户能不能用 dtrace；能不能观测某个进程则取决于进程属主与调试权限。

\begin{itemize}
  \item \Tp{csrutil status} 显示启用时，\Tp{dtrace} 对受保护进程与多数内核 probe 会失败。
        关闭 SIP 是文档里存在但日常开发不该走的路；换一台不敏感的机器更合适。
  \item \Fn{DevTools{\Bk}Security} 是更温和的官方放法：启用后把 dtrace 访问权授给开发者组成员，
        先查它的状态再谈别的。
  \item \Tp{powermetrics}、内核级 \Tp{xctrace} 模板与 \Tp{spindump} 需要 root；\Tp{sample} 与
        \Tp{vmmap} 对自己的进程通常不需要，对别人的进程需要。
\end{itemize}

\begin{lstlisting}[style=shellstyle, caption={macOS 权限与可用性自查}, label={lst:pmapA-mac-perm}]
# 1) SIP 状态
csrutil status
# 2) dtrace 的策略与设备结点
ls -l /dev/dtracehelper
sudo /usr/bin/DevToolsSecurity -status

# 3) xctrace 侧到底有哪些模板（模板名随 Xcode 变）
xcrun xctrace list templates
xcrun xctrace list devices
# 4) probe 现场列，别抄别人给的 provider 名
dtrace -l -n 'profile*' | head -n 10
dtrace -l -P sdt | head -n 20

# 5) 两条路各试一次：不需要 root 的自采样，与需要 root 的那一档
sample $$ 1 -file /tmp/self.txt
sudo powermetrics -n 1 -i 100 2>&1 | head -n 20
\end{lstlisting}

\begin{compare}[title={同一需求在三平台上的最低门槛}]
\begin{tabular}{>{\raggedright\arraybackslash}p{0.20\textwidth}
                >{\raggedright\arraybackslash}p{0.19\textwidth}
                >{\raggedright\arraybackslash}p{0.19\textwidth}
                >{\raggedright\arraybackslash}p{0.19\textwidth}}
需求 & Linux & Windows & macOS \\
\hline
采自己的用户态 & 多数默认档即可 & 用户态会话，通常够用 & \Tp{sample} 级别即可 \\
采内核态 & 需放开档并具备相应能力位 & 必须提升，且会话独占 & 受 SIP 约束，多需 root \\
采别人进程 & 还受 ptrace 语义约束 & 提升或同权 & 需提权或属主一致 \\
硬件计数器 & 虚拟化里常只剩软件事件 & 云虚拟机上常直接不可用 & 芯片族决定事件集 \\
\hline
\end{tabular}
\end{compare}

\section{时钟与延迟分账的三平台差别}

\Tp{perf} 报告里的时间、WPA 时间线的横轴、Instruments 的区间，背后是三种不同的时钟契约。机制
讨论在第~\ref{ch:pm02-clock} 章，这里只给对照与``睡眠之后会不会说谎''这一条。

\begin{itemize}
  \item Linux：\Fn{clock{\Bk}\_gettime} 的各 clock id 语义不同（是否含挂起、是否随时间同步步进、
        是否被频率调整影响），选错的表现是``耗时为负''或``挂起一夜后数字对不上''。\Fn{rdtsc}
        的恒定频率、非停、跨核同步三件事要分别确认，虚拟化再加一层缩放。
  \item Windows：契约计数源是 \Fn{Query{\Bk}Performance{\Bk}Counter}，频率由
        \Fn{Query{\Bk}Performance{\Bk}Frequency} 给出。这个频率\textbf{不是}常量：有的平台上报
        缩放后的标称频率，有的直接是硬件频率，所以必须现场打印。把 QPC 差值当``周期数''用是错的，
        它只是某个固定频率下的计数；睡眠与待机下它是否推进也要自己验。
  \item macOS：\Fn{mach{\Bk}\_absolute{\Bk}\_time} 返回的是 tick，必须用
        \Fn{mach{\Bk}\_timebase{\Bk}\_info} 的分子与分母现场换算，别假设系数；
        \Fn{clock{\Bk}\_gettime} 也在，但各 clock id 在休眠期间是否推进与 Linux 的直觉不一致，
        同样要验。
\end{itemize}

\begin{lstlisting}[style=cppstyle, caption={睡眠或挂起前后各取一次快照，验证所选时钟是否代表真实经过时间（骨架，数值自己填）}, label={lst:pmapA-clock}]
#include <ctime>
#include <iostream>
#if defined(_WIN32)
#include <windows.h>
#endif

struct Snap { double mono; double raw; double extra; };

static Snap now() {
  Snap s{0.0, 0.0, 0.0};
#if defined(_WIN32)
  LARGE_INTEGER c, f;
  QueryPerformanceCounter(&c);
  QueryPerformanceFrequency(&f);            // 频率现场取，不要假设
  const double freq = static_cast<double>(f.QuadPart);
  s.mono = static_cast<double>(c.QuadPart) / freq;
  s.raw = s.mono;
  s.extra = 0.0;                            // 自己接一路确认过语义的对照时钟
#else
  timespec a{};
  clock_gettime(CLOCK_MONOTONIC, &a);
  s.mono = a.tv_sec + 1e-9 * a.tv_nsec;
#if defined(CLOCK_MONOTONIC_RAW)
  clock_gettime(CLOCK_MONOTONIC_RAW, &a);
  s.raw = a.tv_sec + 1e-9 * a.tv_nsec;
#endif
#if defined(CLOCK_BOOTTIME)                 // Linux：含挂起
  clock_gettime(CLOCK_BOOTTIME, &a);
  s.extra = a.tv_sec + 1e-9 * a.tv_nsec;
#elif defined(CLOCK_UPTIME_RAW)             // macOS：不含挂起
  clock_gettime(CLOCK_UPTIME_RAW, &a);
  s.extra = a.tv_sec + 1e-9 * a.tv_nsec;
#endif
#endif
  return s;
}

int main() {
  const Snap a = now();
  std::cout << a.mono << ' ' << a.raw << ' ' << a.extra << std::endl;
  std::cin.get();                           // 机器睡眠数分钟后按回车
  const Snap b = now();
  std::cout << b.mono - a.mono << ' ' << b.raw - a.raw
            << ' ' << b.extra - a.extra << std::endl;
  // 三个差值里哪个等于你记下的真实经过时间，哪个才是你能信的时钟；
  // 其余两个不是坏了，而是它们的语义本来就不含挂起那一段。
  return 0;
}
\end{lstlisting}

延迟分账的差别比时钟本身更容易踩坑：同一次卡顿，Linux 上能把它拆成用户态、内核、中断与队列
四段；Windows 上等待原因来自 ETW 的线程就绪事件族，要在分析器里看等待分类；macOS 上对应的是
run queue 与调度事件，但受权限限制常常只能看到用户态一侧。三平台共同的纪律是：只看工具给的
时间线而不知道它取自哪个时钟，等于没有口径。

\section{事件名与可用性：为什么这里不给事件表}

\begin{warning}[title={抄来的事件名是最常见的假结论来源}]
一张从文章或旧书抄来的事件清单，在你机器上通常有三种结局：名字不存在（直接报错）、名字存在但
被别名改写（测到的不是你以为的东西）、名字存在却报不可用（此时计数常显示 0 或空白，被误读成
``这个开销是零''）。第三种最危险，症状与第~\ref{ch:pm05-events} 章排障节那条警告一致：
不可用不等于零。
\end{warning}

正确动作只有三个：列出来、逐个试、把可用清单连同内核与 CPU 型号一起归档。

\begin{lstlisting}[style=shellstyle, caption={三平台列出可用事件与模板的各自入口}, label={lst:pmapA-events}]
# Linux：分族列出，硬件、软件、指标组要分开看
perf list hw sw metricgroup
perf list --details
lscpu | head -n 20                 # 事件集受 CPU 型号与内核 PMU 驱动约束
# Windows：provider 与 profile 都是现场列
logman query providers
wpr -profiles
xperf -providers
# macOS：dtrace probe 一族，加 xctrace 模板
dtrace -l | wc -l
dtrace -l -n 'profile*'
xcrun xctrace list templates
\end{lstlisting}

\begin{lstlisting}[style=outstyle, caption={示意输出（格式示例，非本仓库实测值）：硬件事件不可用而软件事件可用时的典型形态}, label={lst:pmapA-out}]
$ perf stat -e cycles,instructions,task-clock,context-switches ./bench

 Performance counter stats for './bench':

     <not supported>      cycles
     <not supported>      instructions
              12.345 msec task-clock
                   98      context-switches

       0.014321 sec time elapsed
\end{lstlisting}

这段形态在 CI、容器与云虚拟机上都会遇到。它的正确读法是``这台机器只能给出软件事件层的结论''，
不是``这个程序没有周期消耗''。缺事件时口径怎么退而求其次，见第~\ref{ch:pm03-stats} 章与
附录~\ref{ch:pmapB-params}。

\section{符号与构建产物：build-id、PDB、dSYM}

三平台的符号体系是同一件事的三种实现：产物里留一枚``我是哪个二进制版本''的指纹，外部留一份
能把地址翻译回名字的文件；分析器只有两边指纹对得上才给名字。Linux 用 ELF 的 build-id，调试信息
可以在二进制内也可以剥离到独立文件，\Tp{perf} 按 build-id 命中缓存；Windows 用独立的 \Tp{.pdb}，
二进制里记录 PDB 标识（GUID 与 age），三者必须同时对上——只有一份同名的旧 PDB 会让栈
``看起来有名字但全是错的''，这比没名字更难发现；macOS 用 dSYM bundle，UUID 必须与二进制一致，
分析器按 UUID 找符号。

下面三条命令不是完整流程，而是三张``指纹在不在''的探针：每条只回答一个二值问题——产物里有没有
build-id 或 PDB 标识、外部符号文件能不能与之对上、缓存目录里躺着的是哪一枚。先把这三问回答完，
再谈符号解析失败的原因；把它们贴进 CI 的产物校验步骤，比剖析报告出来以后回头查更省事。

\begin{lstlisting}[style=shellstyle, caption={三平台的产物与符号一致性核对}, label={lst:pmapA-symbol}]
# Linux：build-id 与调试段还在不在；报告里大量裸地址先怀疑这里，再怀疑采样本身
readelf -n ./build/release/app
file ./build/release/app
# Windows：exe 内记录的 PDB 标识与手上的 .pdb 是否同一枚；缓存目录与符号服务器自己填
dumpbin /headers app.exe | findstr /I RSDS
symchk /r app.exe /s srv*<local-cache>*<symbol-server>
# macOS：二进制与 dSYM 的 UUID 必须相同
dwarfdump --uuid ./build/release/app
dwarfdump --uuid ./build/release/app.dSYM
\end{lstlisting}

\begin{note}[title={符号丢了会怎样：三种不同的``看起来还好''}]
\begin{itemize}
  \item 报告里全是裸地址或 \Tp{[unknown]}：你不会得到错误结论，但归因没了，火焰图缩成少数
        几帧，占比全部错位。
  \item 符号部分对部分错（内联帧、模板实例、重载函数最容易）：这是最危险的一档，图看着完整，
        函数名却指错了地方。
  \item 用发布产物的符号去解读一个重新编译过的二进制：可复现性直接归零。
        第~\ref{ch:pm10-compiler} 章``改一行代码前先确认改了二进制''那条纪律在这里同样适用。
\end{itemize}
\end{note}

三平台共同的实践：把 build-id、PDB 标识、dSYM 的 UUID 连同产物一起归档，profile 文件与它们
放在一起。线上采样（第~\ref{ch:pm14-production} 章）没有这份归档就只是噪声。

\section{CI、容器与``该用哪个工具''}

\begin{bestpractice}[title={三平台 runner 的通用事实：先划边界，再定基线}]
\begin{itemize}
  \item 云端 runner 是共享、突发、通常没有硬件事件权限的环境：Linux 上 \Tp{perf} 常常根本没装，
        Windows 上拿不到提升，macOS 上 SIP 默认启用。
  \item 因此本册正文与本附录里的所有输出清单都只是\textbf{形态示意}，不是本仓库的实测数据。
        任何要把数字写进基线的动作，都得发生在你能控的那台机器上（第~\ref{ch:pm13-ci} 章）。
  \item CI 上仍然能做的：软件事件、时间戳序列、产物指纹核对，以及``这台机器今天能不能测''
        的自检脚本。
  \item 写结论时把平台指纹一起写进去：CPU 型号、内核或系统版本、编译器版本、调频策略或电源
        计划、当前频率。缺一项，结论在别人机器上就不可复现。
\end{itemize}
\end{bestpractice}

表~\ref{tab:pmapA-decision} 是全附录的收束：按你想回答的问题给出第一步与前置条件。它与
附录~\ref{ch:pmapC-glossary} 的按症状排障表互补——那张按``你看到了什么''进，这张按
``你要问什么''进。

\begin{longtable}{>{\raggedright\arraybackslash}p{0.21\textwidth}
                  >{\raggedright\arraybackslash}p{0.24\textwidth}
                  >{\raggedright\arraybackslash}p{0.16\textwidth}
                  >{\raggedright\arraybackslash}p{0.13\textwidth}}
\caption{\label{tab:pmapA-decision}该用哪个工具：按问题入口的三平台并列。每行命令都只是第一步，
不是完整流程；回指章是机制出处。}\\
你想回答的问题 & 第一步（三平台各一条） & 权限前置 & 回指章 \\
\hline
\endfirsthead
\hline
你想回答的问题 & 第一步（三平台各一条） & 权限前置 & 回指章 \\
\hline
\endhead
慢在哪个函数 & \Tp{perf record}；\Tp{wpr -start CPU}；Time Profiler & 用户态即可起步 &
  第~\ref{ch:pm04-sampling} 章 \\
\hline
访存受限还是计算受限 & \Tp{perf stat} 最小事件集；VTune 或 Instruments 计数模板 &
  硬件事件要放行 & 第~\ref{ch:pm05-events} 章 \\
\hline
这个补丁真的变快了吗 & 交替 A/B/A/B 加区间判断，平台只是取数来源 & 无特殊要求 &
  第~\ref{ch:pm03-stats} 章 \\
\hline
栈为什么是断的 & \Tp{perf record -g} 报的展开方式；Stack 表；\Tp{sample} 输出 &
  多是构建问题而非权限 & 第~\ref{ch:pm04-sampling} 章 \\
\hline
是谁在等谁 & \Tp{perf sched}；Thread Time 视图；run queue & 内核侧要提权 &
  第~\ref{ch:pm09-concurrency} 章 \\
\hline
是不是伪共享 & \Tp{perf c2c} 或加 padding 做差分；Contention 视图 & 硬件事件与内核档 &
  第~\ref{ch:pm09-concurrency} 章 \\
\hline
内存是占用还是带宽 & \Tp{/proc} 与 smaps；ETW 堆事件；Allocations 模板 & 低开销档可自采 &
  第~\ref{ch:pm07-memory} 章 \\
\hline
慢在缓存还是设备上 & \Tp{iostat -x} 配 fio；disk profile；File Activity & 块层事件常需提权 &
  第~\ref{ch:pm08-io} 章 \\
\hline
GPU 那段时间归谁 & Nsight 族；GPU 分析类型；GPU 模板 & 驱动与权限双门槛 &
  第~\ref{ch:pm11-gpu} 章 \\
\hline
线上能不能一直采 & 低频 \Tp{perf} 或 eBPF；ETW 实时会话；系统诊断包 & 受治理策略约束 &
  第~\ref{ch:pm14-production} 章 \\
\hline
这个数能不能进 PR 描述 & 先跑清单~\ref{lst:pmapA-linux-perm} 那一族自检再取数 & 取数环境要可控 &
  第~\ref{ch:pm15-antipattern} 章 \\
\hline
\end{longtable}

\begin{guideline}[title={跨平台通用的起手式}]
\begin{enumerate}
  \item 先跑对应平台的自检清单（清单~\ref{lst:pmapA-linux-perm}、\ref{lst:pmapA-win-perm}、
        \ref{lst:pmapA-mac-perm} 之一），确认可用事件与可用档位。
  \item 再定口径：\MWall、\MCpu、\Mx{IPC}、\Mx{p99} 只能选一个当主指标，其余记为辅助。
  \item 取数前先确认符号与产物一致（清单~\ref{lst:pmapA-symbol}），否则后面每一步都在猜。
  \item 归档环境指纹与命令本身，让第二个人能原样重跑。
  \item 拿到数之后再谈归因：指标只把你送到代码位置附近，见第~\ref{ch:pm06-micro} 章。
\end{enumerate}
\end{guideline}

最后一句口径提醒：三平台能测的东西重叠度很高，差别在门槛与形态，不在可能性。``只有 Linux
能测''、``Windows 拿不到栈''、``macOS 做不了剖析''这类说法都可以用本节的自查命令证伪——跑得通
就能测；跑不通，是权限或构建的问题，不是平台的问题。指标与参数的逐条解释在
附录~\ref{ch:pmapB-params}，术语与按症状排障在附录~\ref{ch:pmapC-glossary}。

\chapter{指标、事件与命令行参数速查}
\label{ch:pmapB-params}

正文各章解释机制，本附录只回答三类查表式的问题：这个名字算的是什么、这个开关实际改变了什么、这行输出里的这一列能不能当结论用。它不重复推导——判定标准在第~\ref{ch:pm03-stats} 章，平台与权限的并列对照在附录 A，缩写与按症状排查在附录 C。

本附录有一条组织主线：任何参数只有两种作用，要么改变\textbf{你看到的是什么}（口径、归属、样本完整性），要么改变\textbf{被测对象自身的行为}（开销、侵入性、缓冲）。把这两件事混在一张表里，正是"我加了个开关就快了"这类结论的源头，所以第~\ref{sec:pmapB-params}~节按这两类分表。所有默认值一律按"默认是什么 + 为什么会误伤"来写：默认值不是推荐值，更不是结论；而工具版本会移动默认值，所以每条都附了确认方法。

\begin{guideline}[title={查一个参数只需要问三件事}]
第一，默认值是什么，在你的版本上确认过吗（\Tp{--help} 回显、\Tp{fio --enghelp}、\Tp{perf record -h}）。第二，它落在正确性侧还是开销侧——它改变的是口径还是侵入性。第三，生效后你打算用什么证据确认（反汇编、\Tp{perf list} 的事件回显、报告列名的变化）。三件事答不齐，就不要把这个参数写进结论。
\end{guideline}

\section{指标字典}
\label{sec:pmapB-metric}

这是全附录最长的一张表，用途是"看到别人报的指标先对一下口径"。第 4 列只列常见出处，不代表该工具只会报这个；最后一列写的是最典型的误读方式，不是该指标不可用。四个口径标签（\MWall{} 墙钟、\MCpu{} CPU 时间、\MLat{} 延迟、\MThr{} 吞吐）的含义见第~\ref{ch:pm01-why} 章，时钟源见第~\ref{ch:pm02-clock} 章。

{\footnotesize
\begin{longtable}{>{\raggedright\arraybackslash}p{0.155\textwidth} >{\raggedright\arraybackslash}p{0.26\textwidth} >{\raggedright\arraybackslash}p{0.075\textwidth} >{\raggedright\arraybackslash}p{0.16\textwidth} >{\raggedright\arraybackslash}p{0.17\textwidth}}
\caption{指标字典：口径、单位、常见出处与典型误读}\label{tab:pmapB-metric}\\
\toprule
指标 & 口径 & 单位 & 何处可见 & 常见误读 \\
\midrule
\endfirsthead
\toprule
指标 & 口径 & 单位 & 何处可见 & 常见误读 \\
\midrule
\endhead
\bottomrule
\endlastfoot
\MWall{} & 起点到终点经过的真实时间，含被抢占与等待 & 秒（常报 ms） & Google Benchmark、\Tp{hyperfine}、\Tp{perf stat} 的 elapsed & 把它当成程序干了这么多活 \\
\MCpu{} & 任务在各核上实际占用处理器的时间和 & 秒（常报 ms） & \Tp{perf stat}、\Fn{getrusage}、VTune & 多线程下拿 \MCpu{} 除核数去还原 \MWall{} \\
\Tp{user} / \Tp{sys} & \MCpu{} 按特权级的拆分 & 秒 & \Tp{perf stat}、\Tp{time}、\Tp{vmstat} & 把 \Tp{sys} 偏高直接归因于自己的代码 \\
\Mx{cycles} & 计数器口径下的时钟周期数，频率会变时不等于时间 & 周期 & \Tp{perf stat}、VTune & 跨机器比周期数当成比速度 \\
\Mx{instructions} & 退休指令数 & 条 & \Tp{perf stat} & 与 \Mx{cycles} 一样跨微架构比较 \\
\Mx{IPC} & 每周期退休指令数 & 无量纲 & \Tp{perf stat}、\Tp{perf stat} 的 metric 输出 & 高 IPC 等于快；访存受限时它本来就低 \\
\Mx{CPI} & \Mx{IPC} 的倒数 & 无量纲 & TopDown 一族工具 & 当成唯一能区分瓶颈的量 \\
\Mx{branch-miss} & 未命中分支占全部分支之比 & 百分比 & \Tp{perf stat}、VTune & 不看分支总数，只比百分比 \\
\Mx{LLC-miss} & 末级缓存未命中占访问之比 & 百分比 & \Tp{perf stat}、\Tp{perf evlist} & 把 miss 数当周期代价，不做层次换算 \\
\Mx{TLB-miss} & 地址翻译未命中之比 & 百分比 & \Tp{perf stat}（事件名随架构变） & 与缓存 miss 混为一谈 \\
\Mx{retiring} & TopDown 四桶之一：真正有用的槽位 & 百分比 & VTune、TopDown 事件组 & 桶名照抄却不了解槽位模型 \\
\Mx{bad-speculation} & TopDown：浪费在错误路径上的槽位 & 百分比 & 同上 & 看到就想去消分支 \\
\Mx{frontend-bound} & TopDown：取指与解码供给不上 & 百分比 & 同上 & 忽略代码布局与 I-cache 因素 \\
\Mx{backend-bound} & TopDown：执行资源或访存在等 & 百分比 & 同上 & 不细分 core 与 memory 就下结论 \\
\Mx{stalled-cycles} & 处于停顿的周期占比 & 百分比 & \Tp{perf stat} & 当成精确值，它其实是采样估计 \\
\Mx{p50} & 排序后一半样本不超过它 & 时间 & 各基准框架、\Tp{fio} & 把中位数当典型延迟报给业务方 \\
\Mx{p99} & 99\% 样本不超过它 & 时间 & \Tp{fio}、APM、自建直方图 & 样本不足时报 \Mx{p99}，实际是最大值 \\
\Mx{p999} & 更深的尾部分位数 & 时间 & 同上 & 与 \Mx{p99} 混用；两者噪声差一个量级 \\
\Mx{max} & 全程最差的一次观测 & 时间 & \Tp{hyperfine}、\Tp{fio} & 当成容量上限，忽略它对样本量的依赖 \\
\Mx{mean} & 算术平均 & 时间 & 几乎所有工具 & 长尾分布下用它代表典型值 \\
\Mx{stddev} & 样本标准差 & 时间 & \Tp{hyperfine}、\Tp{fio} & 多峰数据下它只描述离散度，不描述位置 \\
\Mx{IQR} & 上下四分位距 & 时间 & 自己统计、\Tp{fio} 的分布图 & 与标准差互相替换 \\
\Mx{MAD} & 中位数绝对偏差 & 时间 & 自己统计（第~\ref{ch:pm03-stats} 章） & 当成稳健性的免费午餐 \\
\Mx{CI} & 统计量的区间估计 & 与量同单位 & 自己算（bootstrap） & 读成"真值一定落在这个区间内" \\
\Mx{MDD} & 最小可检出差异 & 时间或百分比 & 判据脚本、第~\ref{ch:pm03-stats} 章 & 先跑实验再定阈值，事后凑显著 \\
\MThr{} & 单位时间完成量 & 次每秒、B/s & \Tp{fio}、Google Benchmark、服务侧 & 用吞吐变化解释单请求延迟变化 \\
\Mx{bandwidth} & 单位时间搬运字节 & B/s、GiB/s & \Tp{ncu}、STREAM 型测试、\Tp{iostat} & 用理论峰值当可达成上限 \\
\Mx{IOPS} & 单位时间完成的 I/O 数 & 次每秒 & \Tp{fio}、\Tp{iostat} & 丢掉块大小，与带宽互换着报 \\
\Mx{util} & 采样窗内设备有 I/O 在途的时间比 & 百分比 & \Tp{iostat -x} & 把到 $100$ 的 \Mx{util} 直接当饱和 \\
\Mx{await} & 请求在设备侧的平均停留时间 & ms & \Tp{iostat -x} & 当成 \Mx{p99}；均值会藏住尾部 \\
\Mx{aqu-sz} & 平均在途请求数 & 个 & \Tp{iostat -x} & 与 \Mx{util} 相加当成同一个概念 \\
\Mx{RSS} & 驻留物理内存 & KB、MB & \Tp{/{\Bk}proc/{\Bk}self/{\Bk}statm}、\Tp{top}、smaps & 当成虚拟地址空间占用 \\
\Mx{peak RSS} & 全程最大驻留 & KB & \Fn{getrusage}、\Tp{massif}、heaptrack & 当成稳态内存需求 \\
\Mx{minor fault} & 无需读盘的缺页 & 次 & \Tp{perf stat}、\Fn{getrusage} & 当成性能事故 \\
\Mx{major fault} & 需要读盘的缺页 & 次 & \Tp{perf stat} & 与 minor 合并计数后当成噪声 \\
\Mx{ctx-switch} & 主动加被动的上下文切换次数 & 次每秒 & \Tp{perf stat}、\Tp{vmstat}、\Tp{pidstat} & 只看绝对值判优劣，不看同负载基线 \\
\Mx{speedup} & 同一任务两配置的时间之比 & 无量纲 & 自己算 & 两配置口径不同仍然相除 \\
\Mx{efficiency} & 加速比除以并行度 & 百分比 & 第~\ref{ch:pm09-concurrency} 章的伸缩性曲线 & 超过 $100\%$ 一律当测量错误 \\
\Mx{occupancy} & 设备侧同时驻留线程与理论上限之比 & 百分比 & \Tp{ncu}、\Tp{nsys} & 当成与 kernel 时间成正比 \\
\Mx{AI} & 算术强度：每字节搬运对应的浮点操作数 & FLOP{\Bk}/byte & \Tp{ncu}、roofline 图 & 用单核标称算力当分母 \\
\Mx{profiler tax} & 开启剖析后同一口径的退化量 & 百分比、B、ms & 自己量（第~\ref{ch:pm14-production} 章） & 没关剖析器就报出自己的 CPU 税 \\
\end{longtable}
}

\section{命令行参数分组}
\label{sec:pmapB-params}

判别一个参数属于哪一侧的最快问法：把它换掉之后，\textbf{报告的列名或归属}变了，还是\textbf{程序跑得多慢}变了。前者属于正确性侧，它决定你能引用哪一句结论；后者属于开销侧，它决定这份数据还值不值得引用。

\begin{compare}[title={同一个开关可以同时落在两侧}]
\begin{tabular}{>{\raggedright\arraybackslash}p{0.34\textwidth} >{\raggedright\arraybackslash}p{0.42\textwidth}}
\toprule
\Opt{call-{\Bk}graph} \Tp{dwarf} & 栈变完整（口径），样本体积与缓冲压力变大（开销），缓冲不足还会丢样本，于是口径又被改坏 \\
\Opt{freq} 调高 & 样本变多让小函数终于可见（口径），中断更密让被测程序更慢（开销），尾延迟最先受冲击 \\
\Opt{direct} 打开 & 绕开页缓存后测到的才是设备（口径），同时失去缓存复用，数字必然变化（真实性能差） \\
\bottomrule
\end{tabular}
\end{compare}

\subsection{perf 的正确性侧参数}
\label{sub:pmapB-perf-correct}

\begin{longtable}{>{\raggedright\arraybackslash}p{0.19\textwidth} >{\raggedright\arraybackslash}p{0.10\textwidth} >{\raggedright\arraybackslash}p{0.155\textwidth} >{\raggedright\arraybackslash}p{0.35\textwidth}}
\caption{改变口径与归属的 \Tp{perf} 参数}\label{tab:pmapB-perf-correct}\\
\toprule
参数 & 子命令 & 默认是什么 & 改变了什么，写错会怎样 \\
\midrule
\endfirsthead
\toprule
参数 & 子命令 & 默认是什么 & 改变了什么，写错会怎样 \\
\midrule
\endhead
\bottomrule
\endlastfoot
\Opt{event} \Tp{-e} & \Tp{stat}、\Tp{record} & \Tp{record} 默认按周期事件做时间采样，取不到时退到软件事件；\Tp{stat} 默认一组常用硬件事件 & 决定你在数什么。事件名随内核与 CPU 型号变，写错会得到 \Tp{not supported} 与 $0$；把 $0$ 当成"没有开销"是最贵的一种误读 \\
\Opt{call-{\Bk}graph} \Tp{-g} & \Tp{record} & 默认展开方式是 \Tp{fp}（帧指针链），\Tp{-g} 只负责打开采集 & 决定栈是否完整。发布产物常省略帧指针，于是报告只剩叶子函数，热点位置被系统性挪到拷贝与内存操作一类 \\
\Tp{dwarf} 与 \Opt{max-stack} & \Tp{record} & \Opt{call-{\Bk}graph} \Tp{dwarf} 每次采样保存的栈快照大小有上限 & 跨优化帧的栈能恢复，代价是体积与丢失样本的风险；上限要和缓冲一起看，超出后在结尾以丢样本警告的形式出现 \\
\Tp{lbr} & \Tp{record} & 需要 CPU 与内核支持，且只覆盖最近若干次跳转 & 得到硬件级的轻量栈，代价是深栈与内核边界处截断；别把它当完整调用树 \\
\Opt{percent-{\Bk}limit} & \Tp{report} & 阈值默认压得很低（交互界面会折叠小占比条目，折叠线随版本变） & 过滤只影响显示，父节点的总占比不会重新归一。把过滤后的表当 $100\%$ 会得出"热点消失了" \\
\Opt{sort} & \Tp{report} & 默认按命令、DSO、符号三级聚合 & 换排序键会把同名线程合并、把跨 DSO 的同一函数拆成多行。数字没变，看上去的分摊全变 \\
\Opt{children} 与 \Opt{no-{\Bk}children} & \Tp{report} & 数据带调用图时，默认同时给总占比与自身占比 & 决定你引用哪一列。把自身占比当成"这个函数花掉的时间"，等于把它下游的成本记给了别人 \\
\Opt{symfs} & \Tp{report}、\Tp{annotate} & 默认不启用，按数据里记录的绝对路径找符号 & 路径与实际产物不一致就得到 \Tp{[unknown]}，此后所有符号级结论都失效。确认方法：\Tp{perf buildid-{\Bk}list} 与 \Tp{report -v} 的解析回显 \\
\Opt{buildid-{\Bk}all} & \Tp{record} & 默认只为需要的 DSO 记录 build-id & 全部记录后，归档能区分同路径不同构建的产物；不记录时，重新编译过但路径未变的二进制会被解析到旧符号上 \\
\Tp{-p}、\Tp{-t}、\Tp{-C} & \Tp{record}、\Tp{stat} & 默认只跟指定的进程、线程或 CPU & 作用域写错会采到邻居负载。绑核（\Tp{-C}）与只看单线程（\Tp{-t}）的归属含义不同，不能互换 \\
\Opt{system} \Tp{-a} & \Tp{record}、\Tp{stat} & 默认关闭 & 打开即全局计数，内核、中断与其他容器都进你的账。适合整机容量判断，不适合给某个函数定罪 \\
\Opt{event} 的花括号写法 & \Tp{stat} & 默认各事件独立计数 & 同一组事件在同一时刻读数，代价是必须放得进同一批硬件计数器；放不下就整组进复用，输出里的比例是缩放后的估计 \\
\Opt{interval-{\Bk}count} \Tp{-I} & \Tp{stat} & 默认只在结束时汇总一次 & 按固定间隔给时间序列，能把阶跃偏置（热、频率、回收）从均值里分离出来。注意 \Tp{-I} 在 \Tp{record} 下是另一件事（继承统计），别按同一个心智模型用 \\
\Tp{-u} 与 \Tp{-k} & \Tp{record} & 默认用户态与内核态都收 & 只收一侧就是把另一侧的成本从账上抹掉；系统调用密集的程序，内核侧样本往往比想象的多 \\
\end{longtable}

\subsection{perf 的开销侧参数}
\label{sub:pmapB-perf-cost}

\begin{longtable}{>{\raggedright\arraybackslash}p{0.19\textwidth} >{\raggedright\arraybackslash}p{0.10\textwidth} >{\raggedright\arraybackslash}p{0.155\textwidth} >{\raggedright\arraybackslash}p{0.35\textwidth}}
\caption{改变侵入性与开销的 \Tp{perf} 参数}\label{tab:pmapB-perf-cost}\\
\toprule
参数 & 子命令 & 默认是什么 & 改变了什么，写错会怎样 \\
\midrule
\endfirsthead
\toprule
参数 & 子命令 & 默认是什么 & 改变了什么，写错会怎样 \\
\midrule
\endhead
\bottomrule
\endlastfoot
\Opt{freq} \Tp{-F} & \Tp{record} & 常见版本默认 $4000$ Hz & 频率同时决定样本数与中断开销，两者都近似线性。默认值是为桌面交互取的，短小函数上它会把剖析开销算进被测工作；提高前先量一次开与关的差值 \\
\Opt{count} \Tp{-c} & \Tp{record} & 与 \Tp{-F} 互斥，默认走时间采样 & 按事件计数采样会让高频事件样本密度失控、低频事件几乎采不到；混用得到的不是一张覆盖均匀的热点图 \\
\Opt{mmap-{\Bk}pages} \Tp{-m} & \Tp{record} & 默认值随 perf 与内核版本变，须用 \Tp{record -v} 的回显确认 & 环形缓冲太小则产生者快于消费者，直接丢样本；DWARF 展开下尤其明显，丢失量在结尾警告里报出 \\
\Opt{data} \Tp{-d} & \Tp{record} & 默认不记录每条样本的地址与数据载荷 & 记录后 \Tp{perf mem} 与 \Tp{perf c2c} 才可用，代价是每次采样多写字段与额外开销；只开一次就长期留着，数据体积会失控 \\
\Opt{max-stack} & \Tp{record} & 默认对每条栈的保存深度有限制 & 调大能拿到更深的链，也放大每样本写入量；调小则把深栈截成两截，父帧归属错 \\
\Opt{aio} & \Tp{record} & 默认同步写数据文件 & 异步写降低写盘阻塞，但把成本移进页缓存与刷盘时机；采样率高时两侧都要试一次 \\
\Opt{switch-{\Bk}output} & \Tp{record} & 默认不切换文件 & 长时剖析按大小或信号切文件，避免单文件过大；不切又可能掩盖早期样本与时间跨度 \\
\Opt{overwrite} \Tp{-B} & \Tp{record} & 默认关闭（缓冲写满就覆盖旧样本） & 打开后写满即停，配合事件触发式采集；把它当常规剖析用会得到一份时间跨度不完整的数据 \\
\Opt{repeat} \Tp{-r} & \Tp{stat} & 默认只跑一次 & 一次没有离散度，无法区分回归与漂移；重复次数增加总时长，也增加热状态漂移的机会，两害相权要写进结论 \\
\Opt{dry-run} & \Tp{record} & 默认关闭 & 不真正跑负载就能估算给定频率下的采样开销比例，是"先量 CPU 税再开采集"的现成手段 \\
\Opt{output} \Tp{-o} 与 \Opt{append} & \Tp{record} & 默认写 \Tp{perf.data}，同名即覆盖 & 多轮数据混进一个文件后，报告里的百分比是两轮的平均，归属无法拆开 \\
\end{longtable}

\begin{warning}[title={六个默认值，六个方向不同的坑}]
\Opt{call-graph} 默认 \Tp{fp}，它假设产物保留帧指针，而 x86-64 上的常见优化配置默认省略帧指针，于是栈断在叶子。
\Fn{perf} \Tp{record} 的 \Tp{-F} 默认 $4000$ Hz，是给桌面交互用的采样率，不是给短函数用的。
Google Benchmark 的 \Opt{benchmark\_min\_time} 默认是一个几十毫秒量级的最小时间窗（不同版本给过 $0.01$ 与 $0.05$ 秒），在噪声机上足以把两个版本判成同级；\Opt{benchmark\_repetitions} 默认 $1$，意味着没有任何分布可以报。
\Tp{hyperfine} 默认不做预热运行，冷启动、页缓存与频率爬升全部计入第一轮。
\Tp{fio} 默认 \Opt{direct} 为关且 \Opt{iodepth} 为 $1$，于是"随机 4K IOPS"里既没有队列也没有设备压力，多半测的是页缓存。
CMake 的 \Tp{CMAKE\_BUILD\_TYPE} 默认为空，即不下发任何 \Tp{-O} 级别——它编译出来的是 \Tp{-O0} 的程序。
以上每一条都属于"默认值不是推荐值"。确认自己版本的默认，唯一可靠办法是让工具自己说：\Tp{perf record -h}、\Tp{./bench --help}、\Tp{hyperfine --help}、\Tp{fio --enghelp}、\Tp{cmake --help-variable CMAKE\_BUILD\_TYPE}。
\end{warning}

\subsection{Google Benchmark}
\label{sub:pmapB-gbench}

注册侧的宏决定基准存在与否，选项族决定报什么数，两者都要记进结论。

\begin{table}[htbp]
\centering
\small
\begin{tabularx}{\textwidth}{>{\raggedright\arraybackslash}p{0.24\textwidth} >{\raggedright\arraybackslash}p{0.17\textwidth} >{\raggedright\arraybackslash}X}
\toprule
注册写法 / 选项 & 默认 & 改变了什么，写错会怎样 \\
\midrule
\Tp{BENCHMARK} 与 \Tp{BENCHMARK\_{\Bk}REGISTER\_{\Bk}F} & 必须显式注册 & 注册名是 \Opt{benchmark\_{\Bk}filter} 的匹配对象；改名后基准静默不跑，报告少一行而退出码仍是成功 \\
\Tp{->Arg}、\Tp{->Args}、\Tp{->Range} & 无参数时只跑一点 & 参数矩阵决定你看到的是一条曲线还是一个点；把单点当趋势是这个框架最常见的滥用 \\
\Tp{->Unit} & 由框架按量级自动选 & 自动单位让相邻两行的单位不同，脚本按固定单位读列就会错；跨表比较前先把单位钉死 \\
\Tp{-{\Bk}>Iterations} 与 \Tp{->MinTime} & 由时间窗自动决定迭代数 & 钉死迭代数能提高跨版本可比性，代价是小函数上耗时不可控、大函数上可能超时 \\
\Tp{->Threads} 与线程族 & 单线程 & 并行用例的 \MCpu{} 与 \MWall{} 分家，聚合行的含义随之改变（见第~\ref{ch:pm09-concurrency} 章）；线程族写法在版本间有变动，以你集成到的头文件为准 \\
\Tp{-{\Bk}>Repetitions} & 默认不重复，等于 $1$ & 只有重复才产出 mean、median、stddev 一类的聚合行；没有重复就没有分布，任何分位数说法都不成立 \\
\Opt{benchmark\_{\Bk}min\_{\Bk}time} & 版本相关的最小时间窗 & 它是自动迭代次数的驱动量：太小则噪声吞掉差异，太大则总时长与热状态漂移。显式给值比依赖默认好 \\
\Opt{benchmark\_{\Bk}report\_{\Bk}aggregates\_{\Bk}only} & 默认逐次行与聚合行都打 & 打开后只留聚合行，逐次行消失；把聚合行当"一次运行的结果"就丢了分布 \\
\Opt{benchmark\_{\Bk}format} & 默认 console & console 文本的字段随版本改，脚本解析应选 json 或 csv，并把版本回显一起归档 \\
\Opt{benchmark\_{\Bk}filter} & 默认匹配全部 & 正则不匹配时不报错。CI 里要把实际运行的基准数与预期数量核对，否则"全绿"可能是一行都没跑 \\
\Opt{benchmark\_{\Bk}enable\_{\Bk}random\_{\Bk}interleaving} & 默认关闭 & 关闭时用例按注册顺序执行，先后顺序会混入热与频率的阶跃偏置；随机交错是配对交替的近似替代 \\
\Fn{DoNotOptimize} 与 \Fn{ClobberMemory} & 需要显式写 & 不写就可能测不出任何工作：结果被编译器丢掉，你测到的是循环骨架。这是正确性侧问题，不是开销侧 \\
\Fn{UseRealTime} & 默认按 \MCpu{} 报告 & 改变的是单位口径：多线程或含等待的用例下两者差数倍，切换前后的数字不可同表比较 \\
\end{tabularx}
\caption{Google Benchmark 的注册写法与选项族：默认值与作用}
\label{tab:pmapB-gbench}
\end{table}

\subsection{hyperfine}
\label{sub:pmapB-hyperfine}

\Tp{hyperfine} 计时的是完整的进程生命周期（fork、exec、动态加载、退出），与 Google Benchmark 的热循环口径不可互换比较——这句话本身就是它最主要的误用方向。

\begin{table}[htbp]
\centering
\small
\begin{tabularx}{\textwidth}{>{\raggedright\arraybackslash}p{0.24\textwidth} >{\raggedright\arraybackslash}p{0.17\textwidth} >{\raggedright\arraybackslash}X}
\toprule
选项 & 默认 & 改变了什么，写错会怎样 \\
\midrule
\Opt{warmup} & $0$ 轮，不预热 & 第一轮带着冷页缓存、冷分支预测与未爬升的频率。有文件与网络依赖的命令必须显式预热，否则除 \Mx{min} 之外的样本全被抬高 \\
\Opt{min-runs} 与 \Opt{max-runs} & 轮次自动决定（少则个位数，多则按时间预算收敛；上下界随版本变，用 \Tp{--help} 确认） & 重复次数直接决定能报哪些分位数。轮次少时 \Mx{max} 与 \Mx{p99} 只差运气；要分布就把下限钉高 \\
\Opt{prepare} & 无 & 每次正式计时前执行且不计入时间，用来重置外部状态；写成 \Opt{setup}（只执行一次）会让第二轮起的数据依赖第一轮的遗留 \\
\Opt{setup}、\Opt{cleanup}、\Opt{unset} & 无 & 进程级一次、不进计时。用它清状态会把清理成本留在时间窗内，用它建立状态则窗内看不见 \\
\Opt{shell} & 依平台取默认 shell & 管道、通配与退出码语义随之改变。比较两个命令时，shell 差异可以大于被测差异 \\
\Opt{command-{\Bk}name} 与 \Opt{parameter-{\Bk}sweep} & 原名 / 不做扫描 & 参数扫描让同一基准产出多条目，命名不清会把回归归因到别的参数上 \\
\Opt{time-unit} & 自动量级选择 & 与 Benchmark 的自动单位是同一问题：表内单位不统一，脚本与截图都会读错 \\
\Opt{export-{\Bk}json} 与 \Opt{export-{\Bk}markdown} & 不导出 & 导出的是聚合结果，逐次样本默认只留在终端。要事后重算统计就必须落 json，或改用能落盘的采集方式 \\
\Opt{show-{\Bk}output} & 隐藏被计时命令的输出 & 打开后把被计时命令的写盘与刷新计进时间；对高输出量的命令，这会显著抬高被测成本 \\
\end{tabularx}
\caption{\Tp{hyperfine} 关键选项：默认值与作用}
\label{tab:pmapB-hyperfine}
\end{table}

\subsection{fio 的关键项}
\label{sub:pmapB-fio}

\Tp{fio} 的每一项关键参数都在回答同一个问题：这次测的是缓存还是设备。完整口径见第~\ref{ch:pm08-io} 章。

\begin{table}[htbp]
\centering
\small
\begin{tabularx}{\textwidth}{>{\raggedright\arraybackslash}p{0.24\textwidth} >{\raggedright\arraybackslash}p{0.17\textwidth} >{\raggedright\arraybackslash}X}
\toprule
选项 & 默认 & 改变了什么，写错会怎样 \\
\midrule
\Opt{ioengine} & 同步引擎一族（可用引擎随平台与版本变，用 \Tp{fio -{\Bk}-{\Bk}enghelp} 列出） & 决定队列有没有意义：同步引擎下 \Opt{iodepth} 不起作用，你以为在压设备，其实在压单线程系统调用 \\
\Opt{iodepth} & $1$ & 在途请求数。默认值下设备并发为零；它与 \Opt{numjobs} 不等价，两者相乘才是总并发 \\
\Opt{direct} & 关，即缓冲 I/O & 决定页缓存是否参与。关着测出的高吞吐与设备无关；打开还要确认 \Opt{invalidate} 有没有清掉本轮之前的缓存 \\
\Opt{rw} 一族 & \Tp{read} & 模式名（\Tp{randread}、\Tp{readwrite}、\Tp{randrw} 一类）改变读写比与顺序性；名字相近的 \Tp{readwrite} 与 \Tp{randrw} 顺序性完全不同 \\
\Opt{size}、\Opt{runtime}、\Opt{time\_{\Bk}based} & 跑到 \Opt{size} 用完即停 & 默认样本量由数据量决定而非时长；只给 \Opt{runtime} 不给 \Opt{time\_{\Bk}based} 时作业可能提前结束 \\
\Opt{fsync}、\Opt{end\_{\Bk}group\_{\Bk}fsync}、\Opt{fdatasync} & 关 & 关着时"写完"只表示进了页缓存。数据库式负载要按同步点计时，否则带宽与设备落盘能力无关 \\
\Opt{randrepeat} & 开 & 每次跑同一随机序列，跨版本可比；关掉后各轮序列不同，两版本的差异里会混进序列差异 \\
\Opt{percentile\_{\Bk}list} & 一组常用分位（含 \Mx{p99} 与 \Mx{p999}，列表随版本变） & 你要的分位不在列表里就等于没有。报告里缺 \Mx{p999} 不代表尾部不存在 \\
\Opt{group\_{\Bk}reporting} & 关，按 job 报 & 打开后多 job 合并成一个分位数，跨 job 的双峰被抹平；不打开则容易把单个 job 的尾部当成整体 \\
\Opt{latency\_{\Bk}log} 与 \Opt{log\_{\Bk}avg\_{\Bk}msec} & 关 & 只有聚合值时拿不到时间序列，无法区分阶跃偏置与真实回归；按毫秒取平均还会抹掉尖峰 \\
\Opt{output-{\Bk}format} & 普通文本 & 脚本解析选 json；字段层级随版本变，归档时把版本回显一起入库 \\
\end{tabularx}
\caption{\Tp{fio} 关键项：默认值与作用}
\label{tab:pmapB-fio}
\end{table}

\subsection{numactl 与 taskset：绑了什么，没绑什么}
\label{sub:pmapB-numa}

\Tp{taskset} 只改 CPU 亲和性，\Tp{numactl} 才能同时改内存落点。"绑了核但没绑内存"仍然会让访问跨节点，代价由距离与带宽决定（机制在第~\ref{ch:pm12-bandwidth} 章，这里只给参数含义）。

\begin{itemize}
\item \Tp{taskset -c 0-3 ./bench}：把进程钉到给定 CPU 列表。默认继承父进程掩码，容器里还受 cgroup 配额约束，掩码与配额不一致时的表现是"绑了却仍在等"。确认：\Tp{taskset -p \$\$} 与 \Tp{/proc/self/status} 的 \Tp{Cpus\_allowed} 行。
\item \Tp{numactl --cpunodebind=0 --membind=0 ./bench}：同时指定计算节点与内存节点。只写前者时内存仍按首次接触落，远端访问会静默进账。确认：\Tp{numactl --hardware} 的距离矩阵与 \Tp{/proc/self/numa\_maps}。
\item \Tp{numactl --interleave=all}：跨节点交错分配，用来判断"慢是不是距离造成的"——交错之后更慢，瓶颈多半在带宽而不在距离。
\item \Tp{numactl --localalloc}：只影响当前线程的下一次分配，配合逐线程绑核使用；把它当全局策略用是常见写法错误。
\item 绑定带来的可复现性收益只对同机有效：换一台节点拓扑不同的机器，同一套绑定的含义完全不同，结论里必须写清拓扑。
\end{itemize}

\section{编译器开关速查}
\label{sec:pmapB-ccflag}

优化级别改变的是"你在测哪一段机器码"，不是"同一件事快了多少"，所以这一族参数几乎全部属于正确性侧，少数只影响构建时间的除外。

\begin{longtable}{>{\raggedright\arraybackslash}p{0.19\textwidth} >{\raggedright\arraybackslash}p{0.27\textwidth} >{\raggedright\arraybackslash}p{0.19\textwidth} >{\raggedright\arraybackslash}p{0.15\textwidth}}
\caption{编译器开关对测量的影响与确认方法}\label{tab:pmapB-ccflag}\\
\toprule
开关 & 改变了什么 & 对测量的影响 & 怎么确认生效 \\
\midrule
\endfirsthead
\toprule
开关 & 改变了什么 & 对测量的影响 & 怎么确认生效 \\
\midrule
\endhead
\bottomrule
\endlastfoot
\Tp{-O0}（以及 \Tp{-{\Bk}CMAKE\_{\Bk}BUILD\_{\Bk}TYPE} 为空） & 尽量保持源码结构的最低优化 & 内联与向量化基本不发生，测到的是抽象层自身的开销 & \Tp{gcc -{\Bk}Q -{\Bk}O0 -{\Bk}-{\Bk}help={\Bk}optimizers} \\
\Tp{-Og} & 面向调试体验的优化 & 可用于讨论结构成本，不可用于发布结论 & 同上 \\
\Tp{-O1} 与 \Tp{-O2} & 保守到常规的优化 & 差别主要在展开范围与向量化条件，两级的数字不可互换 & 产物反汇编、\Fn{nm} 符号数量 \\
\Tp{-O3} & 更激进的展开与向量化 & 不必然更快：I-cache 压力、寄存器溢出、超线程共享都会反向收费 & 同一程序对照 \Mx{instructions} 与耗时 \\
\Tp{-Os} 与 \Tp{-Oz}（后者以 Clang 一族为主） & 以体积为目标 & 体积换指令密度，代码布局型收益在这里出现，容易与算法收益混淆 & 产物大小与 \Tp{.text} 段尺寸 \\
\Tp{-Ofast} & 连同浮点语义一起放宽 & 结果可能改变，速度差不再只归因于代码质量 & 数值回归测试 \\
\Tp{-DNDEBUG} & 关掉 \Fn{assert} & 断言与它求值的表达式一起消失，被测工作本身变小 & \Fn{nm} 查是否还有断言相关符号 \\
\Tp{-g}、\Tp{-g0}、\Tp{-g1} & 调试信息的有无与级别 & 不改变生成的指令，但改变体积与 DWARF 展开的可用性 & \Tp{readelf -{\Bk}S} 查调试段是否存在 \\
\Tp{-{\Bk}fno-{\Bk}omit-{\Bk}frame-{\Bk}pointer} & 保留帧指针 & 采样栈从叶子变成完整调用链；代价是多占一个寄存器并可能少一次优化 & 反汇编看函数序言是否保存帧指针 \\
\Tp{-flto}、\Tp{-{\Bk}flto={\Bk}auto}、thin 与 full & 跨模块优化与内联 & 分模块测的收益不等于全程序收益；链接时间进构建账 \Mx{profiler tax} 之外的一笔 & 静态函数是否还有独立符号、链接耗时 \\
\Tp{-{\Bk}fprofile-{\Bk}generate} 与 \Tp{-{\Bk}fprofile-{\Bk}use} & 采集并使用训练数据 & 收益方向由采集负载的代表性决定；采集运行本身很慢，绝不能当性能数 & 有 profile 与无 profile 两个产物的差值 \\
\Tp{-{\Bk}march={\Bk}native} 与 \Tp{-march=} & 目标微架构与可用指令集 & 换机后可能非法指令；跨机器结果不可比 & 编译回显的 target 与 \Tp{lscpu} \\
\Tp{-{\Bk}funroll-{\Bk}loops} & 循环展开 & 默认并不随 \Tp{-O3} 打开；展开会同时改变体积与对齐 & 对比两个级别的 \Tp{-{\Bk}-{\Bk}help={\Bk}optimizers} 差集 \\
\Tp{-{\Bk}falign-{\Bk}functions={\Bk}} & 函数入口对齐 & 与 I-cache 命中相关，会同时改变噪声水平 & \Tp{objdump -{\Bk}d} 查地址尾数 \\
\Tp{-{\Bk}fno-{\Bk}exceptions}、\Tp{-fno-rtti} & 去掉运行时支持 & 体积与展开表变化，异常路径的成本从账上消失，与业务口径未必一致 & 产物段尺寸与符号表 \\
\Tp{-{\Bk}frandom-{\Bk}seed={\Bk}} 与 \Tp{-{\Bk}Wl,-{\Bk}-{\Bk}build-{\Bk}id} & 稳定产物标识 & 消除"同码不同二进制"造成的伪差异 & \Tp{readelf -{\Bk}n} 查 build-id \\
\end{longtable}

向量化报告开关这一族只回答"有没有发生"，不回答"快了多少"：GCC 一侧是 \Opt{fopt-info-vec} 与 \Opt{fopt-info-vec-optimized}、\Opt{fopt-info-vec-missed} 这一对；Clang 一侧是 \Tp{-Rpass=loop-vectorize}、\Tp{-Rpass-missed=loop-vectorize}、\Tp{-Rpass-analysis=loop-vectorize}；MSVC 用 \Tp{/Qvec-report} 的级别形式。这一族的名字与可用性随编译器版本变，先在 \Tp{--help} 或 \Tp{/Q} 列表里确认你手上这个版本认哪一个；报告说"已向量化"之后，仍要回到发布产物上反汇编看真实指令。

\begin{note}[title={内联相关：属性不是提速开关}]
\Tp{inline} 首先是 \CPP{} 的 ODR 与链接语义，不是提速开关；是否真的展开由编译器启发式、LTO 与 PGO 决定。属性族里 \Tp{noinline} 是硬保证，\Tp{always\_inline} 与 \Tp{\_\_forceinline\_\_} 只是尽力而为，失败时 GCC 可能报错、Clang 可能警告、MSVC 文档明言不保证。想知道"到底有没有展开"，只能在\textbf{实际发布配置}的产物上反汇编确认（\Tp{objdump} 查调用点、\Tp{nm} 查该函数是否还有独立符号），在探针构建上看没有意义。完整口径见第~\ref{ch:pm10-compiler} 章与《代码风格参考手册》，两处必须一致。
\end{note}

\begin{lstlisting}[style=shellstyle,caption={参数与默认值的自查命令：让工具自己说，别背},label={lst:pmapB-selfcheck}]
# 1) your compiler, not the documented one: what a level really enables
gcc -Q -O2 --help=optimizers > o2.txt
gcc -Q -O3 --help=optimizers > o3.txt
diff o2.txt o3.txt                       # the delta you are actually paying for

# 2) frame pointers present? this decides --call-graph fp
objdump -d --no-show-raw-insn ./bench | less

# 3) ask each tool for its own defaults instead of trusting a cheat sheet
perf record -h | grep -Ei 'call-graph|freq|mmap|max-stack'
perf list                                # events available on THIS kernel+CPU
fio --enghelp ; fio --version
hyperfine --help | grep -Ei 'warmup|min-runs|max-runs'
./bench --help | grep -Ei 'min_time|repetitions|format'
cmake --help-variable CMAKE_BUILD_TYPE   # empty value means no -O at all

# 4) binding: cpuset and memory placement are two different facts
numactl --hardware
taskset -p $$
grep Cpus_allowed /proc/self/status
\end{lstlisting}

\section{口径换算}
\label{sec:pmapB-convert}

换算本身只是算术，出错的地方是单位前缀与"谁除以谁"。本节把常用关系写成正句，不用公式块，因为这张表会被抄进 PR 描述。

{\footnotesize
\begin{longtable}{>{\raggedright\arraybackslash}p{0.20\textwidth} >{\raggedright\arraybackslash}p{0.42\textwidth} >{\raggedright\arraybackslash}p{0.18\textwidth}}
\caption{口径换算表：把关系写成一句可核对的话}\label{tab:pmapB-convert}\\
\toprule
换算 & 正句关系 & 最容易错的一处 \\
\midrule
\endfirsthead
\toprule
换算 & 正句关系 & 最容易错的一处 \\
\midrule
\endhead
\bottomrule
\endlastfoot
频率与周期 & 周期数等于时间乘以频率。在恰好 $1$ GHz 的恒定频率下，$1$ 个周期就是 $1$ ns；频率一变，同一个周期数对应的时间跟着变 & 睿频与降频之后仍用周期数比快慢 \\
时间与吞吐 & 吞吐等于完成数除以总时间；取单次耗时的倒数即每秒次数，$0.5$ ms 一次对应 $2000$ 次每秒，前提是单线程且不重叠 & 用并发下的 \MThr{} 反推单请求 \MLat{} \\
带宽与数据量 & 带宽等于字节数除以搬运时间；$1$ GiB 在 $10$ GB/s 下约需 $0.1$ 秒量级，这是名义值，还要减掉启动开销与协议开销 & GB 与 GiB 混用；小拷贝里启动开销占主导 \\
单位前缀 & 十进制 $1$ GB 等于 $10^{9}$ 字节，二进制 $1$ GiB 等于 $2^{30}$ 字节，两者相差约 $7\%$；$1$ ms 等于 $1000$ us，$1$ us 等于 $1000$ ns & 同一张表里 ms 与 us 混排时被按同一单位读 \\
队列与停留时间 & 稳态下平均在途请求数等于完成速率乘以平均停留时间；三者知道两个就能推第三个，这也是 \Mx{aqu-sz} 的来源 & 把窗口平均当瞬时值，尖峰被抹平 \\
miss 与周期 & 访存代价近似等于各层未命中数乘以该层相对上一层的额外周期数；层次越深权重越大，所以单看总命中率会把 L1 与 DRAM 混为一谈 & 拿别人的层次延迟当乘数 \\
IPC 与 CPI & \Mx{CPI} 是 \Mx{IPC} 的倒数，两者出自同一次计数；\Mx{IPC} 上升而耗时也上升，说明工作量变大而不是变好 & 跨微架构比较 \Mx{IPC} \\
分位与样本量 & 报 \Mx{p99} 至少需要百量级以上的样本，\Mx{p999} 需要更多；样本不足时所谓 \Mx{p99} 只是最大值附近的一个点 & 只跑 $5$ 次就报分位数 \\
CPU 与墙钟 & 单线程下两者同量级；多核下 \MCpu{} 可达 \MWall{} 的若干倍，倍数上限由核数决定，利用率歧义的判据见第~\ref{ch:pm09-concurrency} 章 & 用 \MCpu{} 除以核数去还原 \MWall{} \\
占比与样本 & 报告里的百分比是样本命中比；除以采样频率才得到估计的时间占比，样本少时百分比的区间宽到足以盖住结论 & 把 $2\%$ 的抖动当成回归 \\
串行比例与加速比 & 加速比的上限由串行部分占比决定，而串行部分要靠实测伸缩性曲线得到，不能由代码行数估计 & 用峰值并行度当效率的分母 \\
观测预算 & CPU 税等于开启剖析后的 \MCpu{} 除以关闭后的 \MCpu{} 再减一；尾延迟税要单独看 \Mx{p99} 的差值，不能用均值代替 & 拿生产负载下的差值当同机同负载的差值 \\
\end{longtable}
}

\section{常见输出的字段解释}
\label{sec:pmapB-fields}

下面所有数字都只是形态示意，用来对齐"哪一列是什么"。本仓库的编写与自检都在 Windows 加 MiKTeX 上进行，发布用的 CI runner 既没有 PMU 权限，通常也没有 \Tp{perf}，因此这里不存在任何可引用的实测值；量级一律写成区间或"取决于……"，测量方法由读者在自己机器上做。

\begin{longtable}{>{\raggedright\arraybackslash}p{0.14\textwidth} >{\raggedright\arraybackslash}p{0.17\textwidth} >{\raggedright\arraybackslash}p{0.29\textwidth} >{\raggedright\arraybackslash}p{0.21\textwidth}}
\caption{常见工具输出的字段口径}\label{tab:pmapB-field}\\
\toprule
出处 & 字段 & 口径 & 不看清楚会怎样 \\
\midrule
\endfirsthead
\toprule
出处 & 字段 & 口径 & 不看清楚会怎样 \\
\midrule
\endhead
\bottomrule
\endlastfoot
\Tp{perf report} & Overhead & 该行的样本命中比；给的是总占比还是自身占比取决于是否带调用图 & 与 Self 互换着引用，成本记错人 \\
\Tp{perf report} & Children & 含被调用子树的样本归属 & 把同一笔成本在链上每层各算一遍 \\
\Tp{perf report} & Self & 只算落在该符号自身的样本 & 看到 \Tp{memcpy\_{\Bk}erms} 排第一就下结论 \\
\Tp{perf report} & Command、Shared Object & 命令名与 DSO 归属，是默认排序键的前两级 & 聚合把不同线程或不同动态库的同类符号并成一行 \\
\Tp{perf report} & \Tp{[unknown]} & 符号解析失败，不是"这里没有代码" & 当成优化已完成，实际只是丢了符号表 \\
\Tp{perf script} & period & 该样本代表的事件计数（时间采样时是推算权重） & 把 period 当命中次数直接累加 \\
\Tp{perf script} & weight、addr & 事件附带的加权信息与数据地址，需要 \Opt{data} 才存在 & 没开 \Opt{data} 就去做访存归因 \\
折叠栈 & 链加合计 & 分号自左向右是从根到叶的调用链，末尾数字是这条链的样本合计 & 顺序读反，把叶子当入口 \\
火焰图 & 横向宽度 & 宽度是样本占比，横轴不是时间顺序 & 从左往右读成执行时间线 \\
火焰图 & 纵向高度 & 栈深度；平顶表示自身命中多 & 把宽的顶层框当成"最重要的函数" \\
火焰图 & 颜色 & 默认随机着色，只为区分相邻帧 & 以为暖色表示热点或异常 \\
\Tp{iostat -x} & \%util & 采样窗内设备有 I/O 在途的时间比 & 并行设备上 \Mx{util} 到 $100$ 也不代表饱和 \\
\Tp{iostat -x} & r\_await、w\_await & 该类请求的平均停留时间，含排队 & 当成 \Mx{p99}；均值藏住了尾部 \\
\Tp{iostat -x} & aqu-sz & 平均在途请求数，由窗口平均得出 & 与 \Mx{util} 相加当成同一个概念 \\
\Tp{iostat -x} & rareq-sz、rkB/s & 平均请求大小与吞吐 & 只报 IOPS 而丢块大小，跨配置无法比较 \\
\Tp{vmstat} & r、b & 可运行任务数与不可运行（多在等 I/O）任务数 & 用它们直接推断核数够不够 \\
\Tp{vmstat} & cs、in & 每秒上下文切换与中断次数 & 不看同负载基线就判"切换太多" \\
\Tp{vmstat} & wa、st & 等待 I/O 的空闲比、被宿主抢走的时间比 & 把 wa 计入自己程序的 \MCpu{}；把 st 当代码变慢 \\
\Tp{nvidia-{\Bk}smi} & GPU-Util & 有内核在执行的采样时间比 & 当成算力占用率；一个稀疏内核也能报很高的利用率 \\
\Tp{nvidia-{\Bk}smi} & Memory-Usage & 已被各上下文保有的显存 & 当成活跃数据量或泄漏证据 \\
\Tp{nvidia-{\Bk}smi} & Perf、Pwr & 性能状态与当前功耗读数 & 频率未锁定时把耗时差异当成回归 \\
VTune & Elapsed、CPU Time & 前者是墙钟，后者是占用时间，差值主要是等待 & 用 Elapsed 解释单函数收益 \\
VTune & Effective Time & 扣除已判定为浪费部分后的墙钟视图 & 把它的分层直接搬到别的工具上 \\
Instruments & Weight、Self Weight & 采样命中与自身命中 & 与 perf 的 Children、Self 互换术语时读错列 \\
Instruments & Running、Blocked、Preempted & 线程三态时长拆分 & 把被抢占的时间当成自己的代码时间 \\
Instruments & Persistent、Transient & 保留与临时分配字节 & 当成 RSS，忽略分配器自身占用 \\
\end{longtable}

\begin{table}[htbp]
\centering
\small
\begin{tabular}{>{\raggedright\arraybackslash}p{0.26\textwidth} >{\raggedright\arraybackslash}p{0.44\textwidth}}
\toprule
你要判断的事 & 该看哪一列，以及不该看哪一列 \\
\midrule
这个函数值不值得改 & Self 与调用链上的 Children 一起看；不要只看排序第一的行 \\
改动是否真的更快 & 同一口径下两次 \Tp{perf stat} 的 \Mx{instructions} 与周期；不要只比耗时 \\
慢在等还是慢在算 & \MCpu{} 与 \MWall{} 的比值、线程三态拆分；不要看单个百分比 \\
尾延迟是否变差 & 分位数行与 \Tp{iostat} 的停留时间分布；不要看均值与 \Mx{util} \\
\bottomrule
\end{tabular}
\caption{同一次输出里"该看哪一列"的最小对照}
\label{tab:pmapB-column-choice}
\end{table}

\begin{lstlisting}[style=outstyle,caption={形态示意：\Tp{perf report --stdio} 的列头与行形状（非本仓库实测值）},label={lst:pmapB-report}]
# Overhead  Children  Self  Command  Shared Object   Symbol
# ........  ........  ....  .......  .............   ..............
    12.40%    38.90%  2.10%  bench    libbench.so     [.] PackFrames
     8.70%     8.70%  8.70%  bench    libc.so.6       [.] memcpy_erms
     5.10%    21.30%  1.20%  bench    bench           [.] main
     3.90%     3.90%  3.90%  bench    [unknown]       [k] 0xffffffff8f12
#
# percentages are sample counts, not wall-clock shares
\end{lstlisting}

Self 与 Children 必须同时留在视野里：上例第二行自身占比最高，但它没有子调用；第一行的下游成本远大于自身。出现 \Tp{[unknown]} 时，先按附录 A 的符号与权限条目把符号补齐，否则这张表不能作为结论。

折叠栈（folded stack）是把一棵调用树压成一行的写法：链用分号从根连到叶，末尾以制表符分隔这条链的
采样总数。FlameGraph 一族脚本吃的就是这个格式，所以从 \Tp{perf script} 到折叠栈的转换脚本一旦把
制表符写成空格，图上的宽度就整体失真——火焰图照样画得出来，只是每一段的长短不再对应样本数。
设备侧同理：\Tp{iostat -x} 的列名与列序随版本变动，任何按列号取值的脚本都必须先把表头读进来核对。

\begin{lstlisting}[style=outstyle,caption={形态示意：折叠栈的一行与 \Tp{iostat -x} 的字段排布（非实测值）},label={lst:pmapB-folded}]
# one line per collapsed chain: root ... leaf <tab> sample total
main;PackFrames;CopyBlock 4182
main;Decode 1265
# iostat -x 1 : column order moves between versions, read the header
Device   r/s   w/s   rkB/s   wkB/s  r_await  w_await  aqu-sz  %util
nvme0n1  812.0 33.5  10312.0  402.0     0.42     1.05    0.35   61.0
# rkB/s divided by r/s gives the average request size
# aqu-sz and r/s together give the average time a request spent in flight
\end{lstlisting}

\begin{lstlisting}[style=outstyle,caption={形态示意：\Tp{nvidia-smi} 与 \Tp{vmstat} 的关键列（非实测值）},label={lst:pmapB-dev}]
# nvidia-smi: GPU-Util and Memory-Usage are not what people assume
|  0%  31C    P8    24W /  300W |   1503MiB /  24576MiB |      0%      Default |
#      temp  pstate  power draw / cap   mem used / total     util   pers-mode
# vmstat: r runnable, b blocked on I/O, st stolen by the hypervisor
procs -----------memory---------- ---swap-- -----io---- -system-- ------cpu-----
 r  b   swpd   free   buff  cache   si   so    bi    bo   in   cs us sy id wa st
 3  0      0 812340  30211 905332    0    0   120   344 2100 1850  8  4 84  2  4
\end{lstlisting}

\begin{bestpractice}[title={把参数记录成可复核的三行}]
每次基准在归档里留下三行就能自证：完整命令行（包括你以为会生效的每一项默认）、工具与编译器的版本回显、口径声明（\MWall{} 还是 \MCpu{}，单位是 ms 还是 us，重复次数与统计量取法）。这三行是第~\ref{ch:pm13-ci} 章的基线可比性与第~\ref{ch:pm15-antipattern} 章的结论体检表的最小输入；缺任何一行，结论的有效期就只能缩到"这台机器这一次"。
\end{bestpractice}

\section{排障与速查：先查哪一张表}
\label{sec:pmapB-trouble}

\begin{enumerate}
\item 两个数字对不上：先查第~\ref{sec:pmapB-convert}~节的单位前缀与"谁除以谁"，再确认口径标签是 \MWall{} 还是 \MCpu{}。
\item 热点表里满是 \Tp{[unknown]}：查 \Opt{symfs}、build-id 与符号可用性（附录 A），不要动 \Opt{sort}。
\item 栈只有一层：查 \Opt{call-graph} 的默认值与帧指针开关是否真的进了发布配置，见第~\ref{sub:pmapB-perf-correct}~节与第~\ref{sec:pmapB-ccflag}~节。
\item 报告里没有 \Mx{p99}：查该工具的默认分位列表与重复次数（第~\ref{sub:pmapB-gbench}~节、第~\ref{sub:pmapB-fio}~节），不是查你的程序。
\item 同一命令两次结果差很多：先确认默认时间窗、预热轮次与是否随机交错（第~\ref{sub:pmapB-hyperfine}~节），再谈统计显著性（第~\ref{ch:pm03-stats} 章）。
\item 加了剖析器就变慢：按 CPU 税与尾延迟税分别记，用 \Tp{record --dry-run} 先估开销比例，见第~\ref{sub:pmapB-perf-cost}~节。
\item 吞吐升了但尾延迟崩：查 \Mx{util}、\Mx{aqu-sz} 与 \Mx{await} 的关系（第~\ref{sec:pmapB-fields}~节），把排队与处理分开。
\item 不知道指标该叫什么：回到第~\ref{sec:pmapB-metric}~节对一遍口径，再决定引用哪一列。
\end{enumerate}

需要术语展开与按症状的完整排查表，走附录 C；需要三平台工具与权限的并列对照，走附录 A。

\chapter{术语、缩写与按症状排障手册}
\label{ch:pmapC-glossary}

前两本附录是把\emph{能用什么}并列出来：附录\ref{ch:pmapA-platform}摆三平台的工具、权限与事件面，
附录\ref{ch:pmapB-params}摆指标、参数与字段口径。这一本收的是另外两种检索方式：
一是\emph{词}——读到别人报告里 \Mx{IPC}、\Tp{nonstop\_{\Bk}tsc}、\Tp{perf\_event\_{\Bk}paranoid}
这些词时，一句话告诉你它指什么、它在本册哪章被展开；二是\emph{症状}——当你手上只有一个
``两次跑出 30\% 差''的困惑、还没有任何指标时，按症状给你一条排查顺序和该跳去的章。
本附录不重讲机制，只做索引与起手式；所有``为什么''都在正文章节里，
表中最右列一律用章级引用回指。

词表有两个诚实的限制，先说在前面。第一，事件名、开关名与工具子命令随内核版本、CPU 型号和编译器
版本变化，本表给的是概念而不是某台机器上的确切拼写；第二，量级一律写成区间或``取决于…''，
凡是看着像实测值的数字都不该出现在这里，也不该出现在你写给别人的一句话结论里
（判断依据与自检见第\ref{ch:pm15-antipattern}章）。

\section{术语与缩写}
\label{sec:pmapC-glossary}

表~\ref{tab:pmapC-glossary} 按``缩写/术语 —— 一句定义 —— 在哪章展开''三列排。定义只负责让词不再
陌生，不替代定义所在章的完整口径；同一个词在不同口径下含义不同（例如``时间''可能是 \MWall{}、
\MCpu{}、\MLat{} 三者之一），拿不准时先回到第\ref{ch:pm01-why}章的口径四分法。

\begin{longtable}{>{\raggedright\arraybackslash}p{0.19\textwidth}
                  >{\raggedright\arraybackslash}p{0.54\textwidth}
                  >{\raggedright\arraybackslash}p{0.08\textwidth}}
\caption{本册术语与缩写：定义与展开位置}\label{tab:pmapC-glossary}\\
\toprule
词 & 在本册里指什么 & 展开 \\
\midrule
\endfirsthead
\toprule
词 & 在本册里指什么 & 展开 \\
\midrule
\endhead
\bottomrule
\endfoot
\MWall & 墙钟时间：从起点到终点经过的真实时间，包含等待、缺页、调度与睡眠。 & 第\ref{ch:pm01-why}章 \\
\MCpu & CPU 时间：进程或线程实际占用处理器的累计时间，多线程下可以大于 \MWall。 & 第\ref{ch:pm01-why}章 \\
\MThr & 吞吐量：单位时间完成的工作量；它不告诉你是谁在等，见 \MLat 的对照。 & 第\ref{ch:pm01-why}章 \\
\MLat & 延迟：单个请求或单次操作从提交到完成的时间，常用分位数报告。 & 第\ref{ch:pm01-why}章 \\
口径 & 指标的测量定义：谁包住谁、在哪个窗口内统计、用什么计数器得到。 & 第\ref{ch:pm01-why}章 \\
\Mx{p99} & 分位数：把样本排序后取第 99 百分位；样本量不足时它不稳定而不是更严格。 & 第\ref{ch:pm03-stats}章 \\
尾延迟 & 分布右端（\Mx{p99}、\Mx{p999} 一类）的延迟，通常由调度、缺页、回收或队列造成。 & 第\ref{ch:pm03-stats}章 \\
\Mx{IPC} & 每周期指令数：指令计数除以周期计数；高 \Mx{IPC} 不等于快，见该章解释。 & 第\ref{ch:pm05-events}章 \\
\Mx{CPI} & \Mx{IPC} 的倒数：每条指令平均花多少周期，TopDown 分层时常被反过来读。 & 第\ref{ch:pm05-events}章 \\
\Mx{MAD} & 中位数绝对偏差：对离群值稳健的离散度度量，用来判断今天的机器适不适合测。 & 第\ref{ch:pm03-stats}章 \\
\Mx{IQR} & 四分位距：上四分位减下四分位，长尾数据里比标准差更可解释。 & 第\ref{ch:pm03-stats}章 \\
bootstrap & 对同一批样本重采样来估计统计量分布的方法，用来给出中位数差的置信区间。 & 第\ref{ch:pm03-stats}章 \\
置信区间 & 一个区间加一个置信水平；结论落在区间里还是落在点值上，差别见该章判据三档。 & 第\ref{ch:pm03-stats}章 \\
\Mx{MDE} & 最小可检出差异：在当前重复次数与噪声水平下能可靠分辨的最小变化量。 & 第\ref{ch:pm03-stats}章 \\
预热 & 让首次效应（页缓存、即时编译、时钟爬升、分配器初始化）先过去再计量的那几轮。 & 第\ref{ch:pm03-stats}章 \\
冷启动 & 缓存、页缓存、分支预测器与频率都还没进入稳态的那一段，通常显著更慢。 & 第\ref{ch:pm03-stats}章 \\
热启动 & 同一负载刚跑过一遍、各项局部性已经建立的状态；冷热要成对报告。 & 第\ref{ch:pm08-io}章 \\
阶跃偏置 & 两组数据整体错开一段的偏置，多由版本先后带来的频率、缓存、温度差异造成，用交替跑排除。 & 第\ref{ch:pm03-stats}章 \\
观测者效应 & 测量手段本身改变被测行为（插桩、采样、内联决策、锁时序）。 & 第\ref{ch:pm15-antipattern}章 \\
PMU & 性能监测单元：CPU 内可实现事件计数与采样的硬件计数器集合，数量与事件集随型号变。 & 第\ref{ch:pm05-events}章 \\
复用 & 事件多于物理计数器时的分时采集：报告给出缩放比例，同一轮拿几十个事件尤其不可信。 & 第\ref{ch:pm05-events}章 \\
TopDown & 把周期分成 \Mx{retiring}、\Mx{bad-speculation}、\Mx{frontend-bound}、\Mx{backend-bound} 四桶的分层读法。 & 第\ref{ch:pm05-events}章 \\
PEBS & 处理器精确采样：把采样记录指向事件发生的指令本身而不是中断点，需权限且随平台变。 & 第\ref{ch:pm05-events}章 \\
IBS & AMD 的指令级采样机制，与 PEBS 扮演同一角色，事件名与可用性另成一套。 & 第\ref{ch:pm05-events}章 \\
LBR & Last Branch Record：硬件保留的最近分支跳转对，可低成本给出浅栈，只到叶子附近。 & 第\ref{ch:pm04-sampling}章 \\
DWARF & 调试信息格式，用来在栈展开时从指令地址还原帧链；采样缓冲有上限，深栈可能被截。 & 第\ref{ch:pm04-sampling}章 \\
帧指针 & \Tp{rbp} 一类链指针；编译时保留它可让栈展开便宜可靠，代价是占一个寄存器。 & 第\ref{ch:pm04-sampling}章 \\
\Tp{rdtsc} & 读时间戳计数器的指令；是否可用取决于恒定频率、非停与跨核同步三个条件。 & 第\ref{ch:pm02-clock}章 \\
TSC & 时间戳计数器：随周期递增的计数，虚拟化与降频历史会让它与真实时间脱钩。 & 第\ref{ch:pm02-clock}章 \\
\Tp{constant\_{\Bk}tsc} 与 \Tp{nonstop\_{\Bk}tsc} & 内核在 \Tp{/proc/{\Bk}cpuinfo} 里声明的 TSC 属性：频率是否恒定、休眠与降频时是否继续走。 & 第\ref{ch:pm02-clock}章 \\
\Tp{art} & 平均时基比率校验，说明 TSC 是否与实际时间保持固定换算；拿不到就要谨慎换算。 & 第\ref{ch:pm02-clock}章 \\
QPC & Windows 的 \Fn{QueryPerformance{\Bk}Counter}：单调高分辨计时，底层可能是 TSC 也可能不是。 & 第\ref{ch:pm02-clock}章 \\
vDSO & 内核映射进用户态的轻量库，把 \Fn{clock\_{\Bk}gettime} 一类调用变成不进内核的指令序列。 & 第\ref{ch:pm02-clock}章 \\
\Tp{CLOCK\_{\Bk}MONOTONIC} & 单调但不跟随 NTP 步进的时钟 id；相对的 \Tp{CLOCK\_{\Bk}REALTIME} 会被调整，可能倒走。 & 第\ref{ch:pm02-clock}章 \\
\Tp{CLOCK\_{\Bk}BOOTTIME} & 计入挂起时间的时钟 id，跨休眠测量时才需要它。 & 第\ref{ch:pm02-clock}章 \\
\Tp{task{\Bk}-clock} & 软件事件形式报告的线程占用 CPU 时间，不依赖 PMU。 & 第\ref{ch:pm05-events}章 \\
DVFS & 动态电压频率调整：负载与温度改变运行频率，于是同一份代码在不同时刻跑出不同周期数。 & 第\ref{ch:pm02-clock}章 \\
SMT & 单核多线程（超线程一类）：逻辑核共享执行资源，占用率与伸缩曲线的拐点会提前。 & 第\ref{ch:pm09-concurrency}章 \\
缺页 & 访问尚未建立物理映射的页导致的内核介入；major 与 minor 的代价差一个量级。 & 第\ref{ch:pm02-clock}章 \\
\Mx{RSS} & 驻留集大小：当前真正占用物理内存的量，与虚拟地址空间大小是两件事。 & 第\ref{ch:pm07-memory}章 \\
\Mx{VSZ} & 虚拟内存大小：地址空间被映射的量，包含未驻留与被共享的部分。 & 第\ref{ch:pm07-memory}章 \\
\Mx{THP} & 透明大页：内核自动把连续页合成大页，能改变缺页率与延迟分布，也可能带来整理开销。 & 第\ref{ch:pm07-memory}章 \\
\Mx{smaps} & \Tp{/proc/self/{\Bk}smaps} 一族文件，逐区域给出驻留、脏页与回收统计。 & 第\ref{ch:pm07-memory}章 \\
NUMA & 非统一内存访问：本地与远端内存代价不同，距离矩阵与绑定策略都要自己确认。 & 第\ref{ch:pm12-bandwidth}章 \\
缓存行 & 一次搬运的最小粒度（大小随平台变）；跨线程写同一行会产生伪共享。 & 第\ref{ch:pm09-concurrency}章 \\
伪共享 & 两个线程各自写同一缓存行的不同字段，导致行在线程间来回失效。 & 第\ref{ch:pm09-concurrency}章 \\
冲突未命中 & 容量足够但映射路径相撞造成的失效；与局部性未命中是两种原因。 & 第\ref{ch:pm06-micro}章 \\
局部性 & 数据在时间与空间上的聚集程度，决定缓存与预取能不能帮上忙。 & 第\ref{ch:pm06-micro}章 \\
\Mx{branch-miss} & 分支预测错误：可预测分支与数据依赖分支的处理方式完全不同。 & 第\ref{ch:pm06-micro}章 \\
roofline & 用算术强度把性能上限画成访存与计算两条界的读法；量级用比值而不是绝对数。 & 第\ref{ch:pm12-bandwidth}章 \\
算术强度 & 每字节数据搬运做多少次运算，决定一段代码是带宽受限还是计算受限。 & 第\ref{ch:pm12-bandwidth}章 \\
HBM & 高带宽显存：容量与带宽由堆栈与总线宽度决定，和 GPU 计算峰值是两笔账。 & 第\ref{ch:pm12-bandwidth}章 \\
PCIe & 主机与设备之间的链路；带宽要按代次、通道数与是否双向分开算。 & 第\ref{ch:pm12-bandwidth}章 \\
ACS & 访问控制重定向相关能力，决定 PCIe 点对点路径是否真的能绕过主机桥。 & 第\ref{ch:pm12-bandwidth}章 \\
P2P & 设备之间直接搬运而不经过主机内存；可达性要逐对确认，不是有卡就能走。 & 第\ref{ch:pm12-bandwidth}章 \\
CUDA event & 入流的计时标记：\Fn{cudaEventElapsedTime} 给出两个标记之间设备侧经过的时间。 & 第\ref{ch:pm11-gpu}章 \\
kernel & 一次 GPU 执行的单位；它是被排队的，所以计时窗口与 CPU 侧调用不是同一件事。 & 第\ref{ch:pm11-gpu}章 \\
页锁定内存 & 允许异步拷贝与真正重叠的宿主缓冲区；普通 \Fn{malloc} 出来的内存会让异步退化。 & 第\ref{ch:pm11-gpu}章 \\
\Mx{occupancy} & GPU 上同时驻留的线程比例；它与 kernel 时间的关系不是线性的。 & 第\ref{ch:pm11-gpu}章 \\
重放 & 分析工具为了拿某一类计数把 kernel 重复执行；它给出的数不是运行时序。 & 第\ref{ch:pm11-gpu}章 \\
利用率歧义 & \MCpu{} 与 \MWall{} 在并发下分家的三种形态：更慢但更费 CPU、更快但烧核、吞吐升而尾延迟崩。 & 第\ref{ch:pm09-concurrency}章 \\
背压 & 下游跟不上时把等待传回上游的机制；它会让 \MThr{} 稳定而 \MLat{} 变长。 & 第\ref{ch:pm09-concurrency}章 \\
抖动 & 延迟在时间上的不稳定性，常表现为同一均值的两种完全不同的分布。 & 第\ref{ch:pm08-io}章 \\
\Mx{build-id} & 二进制内容哈希，用来把 profile 与产物对上；\Opt{buildid-all} 影响归档方式。 & 第\ref{ch:pm04-sampling}章 \\
\Tp{perf\_event\_{\Bk}paranoid} & 内核限制非特权用户性能事件的开关，取值决定你能拿到什么。 & 附录\ref{ch:pmapA-platform} \\
USDT & 用户态静态探针：程序里预先埋好的观测点，线上采样常与它配合。 & 第\ref{ch:pm14-production}章 \\
eBPF & 在内核里安全执行的小程序，可做直方图与栈采集，开销与权限都要自己确认。 & 第\ref{ch:pm14-production}章 \\
ETW & Windows 的事件跟踪基础设施，XPerf 与多数剖析工具在它之上。 & 附录\ref{ch:pmapA-platform} \\
dSYM & macOS 的调试符号 bundle，与 Linux 的 DWARF、Windows 的 PDB 对应。 & 附录\ref{ch:pmapA-platform} \\
PGO & 按训练运行得到的画像来指导优化；训练负载的代表性决定收益方向。 & 第\ref{ch:pm10-compiler}章 \\
BOLT & 在链接后的二进制上重排代码以改善取指局部性的工具族。 & 第\ref{ch:pm10-compiler}章 \\
LTO & 链接期跨模块优化；模块级收益之和不等于全程序收益，构建时间也要进账。 & 第\ref{ch:pm10-compiler}章 \\
火焰图 & 把栈按宽度画成占比的可视化，形态有 \Mx{icicle}、\Mx{sunburst}、\Mx{treemap} 三种。 & 第\ref{ch:pm04-sampling}章 \\
自身占比 & 只算栈顶那一层的占比；把它当``这个函数花掉的时间''是最常见的误读。 & 第\ref{ch:pm04-sampling}章 \\
\MWall{} 与 \MCpu{} 分家 & 并发测量里 CPU 时间轴与墙钟轴不再同步的现象，先定口径再解释。 & 第\ref{ch:pm09-concurrency}章 \\
\Tp{\_\_OPTIMIZE\_{\Bk}\_\_} & 判断当前产物是不是优化构建的线索之一，不是结论本身。 & 第\ref{ch:pm10-compiler}章 \\
\Tp{-DNDEBUG} & 关掉断言的构建开关；带着断言跑与不带断言跑是两个不同的被测程序。 & 第\ref{ch:pm10-compiler}章 \\
\end{longtable}

\begin{guideline}[title={读这张表的三个规矩}]
第一，表里凡带 \Tp{} 的拼写都要在你自己的机器上核对一遍：事件名、clock id、sysfs 路径随内核与
硬件变，抄表不等于可用。第二，``展开''列只指到章，不指到小节；进章后先看它的口径定义，再看命令。
第三，含义相近而不可互换的词要靠口径分开：\Mx{MAD}（离散度）与 \Mx{MDE}（可检出差异）不同，
\Mx{IPC} 与 \Mx{CPI} 互为倒数而不是同一个东西，\MWall{} 与 \MCpu{} 更不是一回事。
\end{guideline}

\section{按症状排障手册}
\label{sec:pmapC-trouble}

这一节的用法是：手上还没有指标、只有一个``不对劲''的时候，按左列找最像的一行，从左到右执行。
``先查''永远是便宜且能一票否决的动作（产物对不对、时钟对不对、有没有基线），
``若正确再查''才需要工具与更长的窗口。最右列是章级回指，进去后按该章的排障小节接着走。
表~\ref{tab:pmapC-trouble} 的行之间没有优先级差别，但同一台机器上同时出现两三行时，
先把环境漂移那一组处理干净再谈其余（环境指纹见第~\ref{sec:pmapC-checklist}节）。

\begin{longtable}{>{\raggedright\arraybackslash}p{0.20\textwidth}
                  >{\raggedright\arraybackslash}p{0.24\textwidth}
                  >{\raggedright\arraybackslash}p{0.24\textwidth}
                  >{\raggedright\arraybackslash}p{0.13\textwidth}}
\caption{按症状排障：症状、先查、若正确再查与对应章}\label{tab:pmapC-trouble}\\
\toprule
症状 & 先查 & 若正确再查 & 对应章 \\
\midrule
\endfirsthead
\toprule
症状 & 先查 & 若正确再查 & 对应章 \\
\midrule
\endhead
\bottomrule
\endfoot
两次跑出 30\% 的差 & 构建产物是否同一个；当前频率与 governor & 变异性是不是本来就很大（\Mx{IQR}、\Mx{MAD}）；交替跑 A/B/A/B & 第\ref{ch:pm03-stats}章 \\
\MCpu{} 降了而 \MWall{} 升了 & 先确认问的是哪个口径，别急着归因 & 并行度与背压；分配器与频率变化；锁等待 & 第\ref{ch:pm01-why}章 \\
\MCpu{} 大于 \MWall{} & 是否多线程（多核并行时这本来就该出现） & 单线程基线有没有；利用率歧义的三种形态 & 第\ref{ch:pm09-concurrency}章 \\
耗时算出负数 & 是否用了 \Tp{CLOCK\_{\Bk}REALTIME} 且期间被步进 & 跨挂起测量；Windows 上计时源切换 & 第\ref{ch:pm02-clock}章 \\
\Tp{rdtsc} 与 \Fn{perf} 的数对不上 & \Tp{/proc/{\Bk}cpuinfo} 里的 \Tp{constant\_{\Bk}tsc}、\Tp{nonstop\_{\Bk}tsc}、\Tp{art} & 是否在虚拟机里；改用 \Fn{clock\_{\Bk}gettime} 复核 & 第\ref{ch:pm02-clock}章 \\
火焰图大片 \Mx{[unknown]} & 是否带符号编译；\Opt{symfs} 与 \Mx{build-id} 是否对上 & 内核符号权限（\Tp{kallsyms}）；发布产物 strip 掉了什么 & 第\ref{ch:pm04-sampling}章 \\
栈只有一两层就断 & 用了哪种展开方式（帧指针、\Opt{call\_{\Bk}graph dwarf}、LBR） & 是否 \Tp{-fomit{\Bk}-frame-pointer}；缓冲上限截断；协程与异步回调 & 第\ref{ch:pm04-sampling}章 \\
\Mx{memcpy} 一类排在最前 & 自身占比还是总占比；过滤阈值是否吃掉了别的东西 & 谁在调用它、能不能少搬；这是线索不是结论 & 第\ref{ch:pm04-sampling}章 \\
\Fn{perf stat} 里有 not supported & 有没有 PMU 权限；是否在 CI runner 里 & 事件名在你的机器上叫什么（\Fn{perf} list）；能否用软件事件替代 & 第\ref{ch:pm05-events}章 \\
计数为 0 被当成 0 开销 & 该事件是不支持还是真的没发生 & 换事件或换机器做交叉核对 & 第\ref{ch:pm05-events}章 \\
几十个事件一轮采完 & 物理计数器数量与复用缩放比例 & 分组（\Opt{group}）或每轮只带少数几个事件 & 第\ref{ch:pm05-events}章 \\
高 \Mx{IPC} 却更慢 & 分子分母各是什么；是否在比不同负载 & \Mx{backend-bound} 与访存；\MWall{} 里的等待部分 & 第\ref{ch:pm05-events}章 \\
加采样后 \Mx{p99} 变差而 \MThr{} 没动 & 采样频率与栈展开方式带来的开销 & 缓冲落盘与信号处理；把它当观测者效应单独记账 & 第\ref{ch:pm15-antipattern}章 \\
第一次跑明显慢 & 冷启动（页缓存、分配器、即时编译、频率爬升） & 预热轮数是否够；累积均值什么时候才平 & 第\ref{ch:pm03-stats}章 \\
第二次跑明显快 & 页缓存与文件是否已在内存里 & 冷热分开报；需要冷时如何制造冷以及它自身的失真 & 第\ref{ch:pm08-io}章 \\
基准报 0 或 1 个纳秒 & 被测工作有没有被优化掉；结果有没有被使用 & 反死代码消除手段；常量折叠与循环不变量提出 & 第\ref{ch:pm15-antipattern}章 \\
改一行代码数字不动 & 二进制是否真的重建（大小、符号、时间戳） & 构建缓存；跑的是不是旧产物；内联决策本来就未变 & 第\ref{ch:pm10-compiler}章 \\
以为加了 \Tp{inline} 就快了 & \Tp{inline} 是 ODR 与链接语义，不是提速开关 & 在发布产物上反汇编确认到底有没有展开 & 第\ref{ch:pm10-compiler}章 \\
线程加到 8 倍吞吐只到 2 倍 & 串行比例从哪来；绑核是否干扰复现 & 伸缩曲线拐点与物理核数、SMT 邻居数；锁与伪共享 & 第\ref{ch:pm09-concurrency}章 \\
两个线程写同一行 & 缓存行大小与结构体字段排布 & \Tp{perf c2c} 一类工具；padding 是否只是掩盖算法层的共享 & 第\ref{ch:pm09-concurrency}章 \\
\Mx{RSS} 一直涨但在用量平稳 & 是否碎片或缓存未释放；驻留与在用量分别怎么算 & 区域数与 \Tp{smaps} 形态；什么场景下峰值内存才是约束 & 第\ref{ch:pm07-memory}章 \\
内存工具慢一个数量级 & 用的是不是插桩或采样型工具 & 那份数只能看形状不能当性能数；工具用法回看《内存泄漏检测与防御》册 & 第\ref{ch:pm07-memory}章 \\
磁盘吞吐远低于标称 & 队列深度、\Opt{iodepth} 与 \Opt{direct} 是否够 & 调度器与页缓存路径；设备侧上限与 RAID、分层 & 第\ref{ch:pm08-io}章 \\
I/O 延迟高但 CPU 空闲 & 四段分账落在哪一层（应用、内核、块层、设备） & 同步调用（\Fn{fsync}）时机；写放大；尾延迟来源 & 第\ref{ch:pm08-io}章 \\
GPU 一段计时只有极小的值 & 计时窗口里是否漏了同步 & 默认流的隐式同步；首个与末个 event 的边界 & 第\ref{ch:pm11-gpu}章 \\
加了同步后总时长暴涨 & 原本有多少工作其实根本没被执行 & 把 H2D 与 D2H 单独计；流水线是否真的重叠 & 第\ref{ch:pm11-gpu}章 \\
小拷贝的时间不按字节走 & 启动与排队开销在总量里的占比 & 页锁定内存是否到位；拓扑与链路宽度 & 第\ref{ch:pm12-bandwidth}章 \\
多卡吞吐不等于单卡乘 N & 每卡基线与同步开销 & 带宽争用、负载不均、跨卡通信路径与 P2P 可达性 & 第\ref{ch:pm12-bandwidth}章 \\
绑核之后更慢 & 绑到了什么（CPU 还是 NUMA 节点）与是否和邻居争用 & 内存没绑造成的远端访问；SMT 配对；隔离策略 & 第\ref{ch:pm12-bandwidth}章 \\
换一台机器结论反过来 & 两台机器的环境指纹差在哪 & 频率策略、缓存层次、内核与编译器版本；把结论限定在一台机器 & 第\ref{ch:pm13-ci}章 \\
CI 偶尔红一次 & 是否共享或突发资源被邻居挤占 & 判据是否用了单边检验与连续两次确认；要不要同机自比 & 第\ref{ch:pm13-ci}章 \\
阈值一调就绿 & 调的是阈值还是环境 & 回归判据与 \Mx{MDE}；别把噪声计数器当门禁 & 第\ref{ch:pm13-ci}章 \\
线上采样只看慢请求 & 采样条件是否引入了幸存者偏差 & 观测预算怎么定；CPU 税与尾延迟分开记 & 第\ref{ch:pm14-production}章 \\
只有一个均值加一张截图 & 分布、区间、环境指纹、复现步骤都在吗 & 按结论体检表逐项补齐，缺项就改成无结论 & 第\ref{ch:pm15-antipattern}章 \\
\end{longtable}

\begin{warning}[title={用这张排障表最容易踩的两件事}]
一是跳步：``先查''那一列之所以便宜，是因为它能一票否决。产物不对、时钟不对、没有基线时，
后面所有工具的读数都是在给噪声拍照。二是把表里的行序当因果：症状相同、机制不同的情况非常多
（``两次跑出 30\% 差''至少有构建差异、频率漂移、变异性过大、阶跃偏置四条路），
走完表只是缩小范围，把相关性变成因果还需要对应章自己的证据。
\end{warning}

\section{测量前自检清单}
\label{sec:pmapC-checklist}

下面这份清单是开工前的三十秒，每一条都能用一个动作回答；回答不了的那条就是本次测量的风险来源。
把它抄进你的基准脚本头部注释，或者贴在评审意见里。

\begin{itemize}
  \item 我要回答的是哪一种提问：绝对慢、比上一版慢、还是能不能扛住 N？三者的最小证据不同。
  \item 口径写死了吗：\MWall、\MCpu、\MThr、\Mx{p99} 里哪一个是我的结论用词？
  \item 被测二进制是哪一个？构建类型、优化级别、\Tp{-DNDEBUG}、链接顺序都记下来了吗？
  \item 这个二进制是我改的那行代码重建出来的吗？产物大小与符号表核对过吗？
  \item 计时用的时钟是哪个？它单调吗？会不会被步进、挂起或降频影响？
  \item 这台机器今天的频率策略是什么？有没有别人在跑？温度与限频看得见吗？
  \item 有没有 PMU 权限？拿不到时是否改用软件事件，而不是把 0 当结论？
  \item 预热怎么做的？轮数是看过累积均值之后定的，还是拍脑袋？
  \item 重复次数是多少？这个次数能检出我关心的最小差异吗？
  \item A 与 B 是交替跑的还是先后跑的？先后跑就要准备解释阶跃偏置。
  \item 我要报的是区间还是点值？多参数扫描有没有预先写好的判据？
  \item 环境指纹存在哪里？内核、CPU 型号与微码、内存配置、编译器版本、负载参数都要在。
  \item 原始数据落在哪？逐轮样本或分位数摘要能不能连同脚本一起归档？
  \item 复现步骤是别人可以照敲的吗？还是只有我机器上缓存热着才成立？
  \item 结论的边界写明了吗：哪台机器、哪个负载、哪个版本之外不成立？
\end{itemize}

清单第 12 条的环境指纹建议固化成一个可跑脚本，把输出和数据放在一起；下面是形态骨架，
字段名与文件路径随发行版不同，缺哪一条就从你的系统里找等价来源，找不到也要在文件里留下
一行说明这一项拿不到，不要静默跳过。

\begin{lstlisting}[style=shellstyle,caption={环境指纹采集骨架（可跑；字段随发行版与内核变，不含任何实测值）},label={lst:pmapC-fingerprint}]
#!/bin/sh
# 用法: sh pm-fingerprint.sh > pm-fingerprint.txt 2>&1
# 目的: 让同一份代码在不同时刻跑出不同数这件事可被解释
echo '== host =='
uname -a
cat /proc/sys/kernel/perf_event_paranoid 2>/dev/null
echo '== cpu =='
lscpu 2>/dev/null
grep -m1 -e 'model name' -e 'flags' /proc/cpuinfo
echo '== tsc flags =='
grep -e constant_tsc -e nonstop_tsc -e art /proc/cpuinfo
echo '== frequency =='
cat /sys/devices/system/cpu/cpu0/cpufreq/scaling_governor 2>/dev/null
cat /sys/devices/system/cpu/cpu0/cpufreq/scaling_cur_freq 2>/dev/null
echo '== memory and numa =='
grep -e MemTotal -e Hugepagesize /proc/meminfo
numactl --hardware 2>/dev/null
echo '== toolchain =='
cc --version
c++ --version
echo '== artifact =='
ls -l ./bench
sha256sum ./bench 2>/dev/null
readelf -n ./bench 2>/dev/null
echo '== load right now =='
uptime
\end{lstlisting}

Windows 与 macOS 上这些字段各有对应来源（QPC 与 TSC 的关系、ETW 会话权限、SIP 与符号 bundle），
并列对照留在附录\ref{ch:pmapA-platform}，这里不重复；要紧的是同一台机器同一份指纹这条纪律，
而不是具体敲了哪条命令。

\section{两套起手式配方}
\label{sec:pmapC-recipes}

\subsection{起手式一：我只有 30 分钟}
\label{sub:pmapC-thirtymin}

目标不是得出结论，而是知道该去哪一章。顺序是：交替跑两版各若干轮拿 \MWall；一轮 \Fn{perf stat}
看有没有微架构异常；一次采样先看栈能不能拿到、热点形状长什么样。三步里任何一步失败，
就先修工具与权限，不要急着解释数字。

\begin{lstlisting}[style=shellstyle,caption={30 分钟起手式（Linux 与 perf 形态；另两平台见附录 A）},label={lst:pmapC-thirtymin}]
#!/bin/sh
# 前提: 两个版本都带 -O2 -g -fno-omit-frame-pointer 且已重建
# 交替跑, 各 5 轮; 顺序固定, 不做挑选
set -e
for i in 1 2 3 4 5; do
  ./bench_old > out_old_$i.txt
  ./bench_new > out_new_$i.txt
done
# 一轮事件计数: 只用一定可用的最小事件集, 别一次塞几十个
perf stat -e task-clock,cycles,instructions,context-switches,page-faults ./bench_new
# 一次采样: 先确认能不能拿到栈, 再看热点在哪
perf record -F 99 --call-graph dwarf -o perf_old.data ./bench_old
perf record -F 99 --call-graph dwarf -o perf_new.data ./bench_new
perf report -i perf_new.data --stdio
\end{lstlisting}

\begin{compare}[title={30 分钟能得出什么、得不出什么}]
能得出：两版差异的方向是否稳定；有没有微架构层面的可疑信号（\Mx{IPC} 形态、缺页与切换计数）；
热点是否落在你改的那段代码附近；栈能不能展开（这决定后续用哪种剖析代价）。
得不出：任何可以写进发布说明的绝对值；\Mx{p99} 一类的尾延迟结论（样本量不够）；
``快了多少''的百分比（5 轮里 30\% 的差完全可能只是噪声）。
如果非要把这一步的数报出去，措辞只能是：观察到方向，还需要重复次数与区间。
\end{compare}

\subsection{起手式二：我要写一个能进 PR 描述的结论}
\label{sub:pmapC-pr}

这一段是模板而不是命令；把它贴进 PR 描述，缺哪一格就是还没测完。核心要求是每一格都能由别人复现，
并且结论只在自己的适用域内说话。

\begin{lstlisting}[style=shellstyle,caption={结论模板的填空骨架（六格顺序与结论体检表对应）},label={lst:pmapC-pr}]
#!/bin/sh
# 每格只允许填: 数字 + 口径 + 可复现来源
echo '== metric =='
echo 'wall-clock, per-run total, fixed workload, 9 rounds'
echo '== environment fingerprint =='
sh pm-fingerprint.sh
echo '== decision (written BEFORE running) =='
echo 'A/B alternated; bootstrap 95 percent interval of median difference'
echo 'interval excludes zero -> better or worse; otherwise -> no conclusion'
echo '== data =='
echo 'median A, median B, interval, IQR, dropped rounds and why'
echo '== scope of the claim =='
echo 'machine, governor, compiler version, workload parameters'
echo '== reproduce =='
echo 'commit, build command, artifact sha256, run script path, raw log path'
\end{lstlisting}

\begin{guideline}[title={写结论的六格顺序，缺一格就别写数字}]
口径 $\rightarrow$ 环境指纹 $\rightarrow$ 判据（跑之前写好）$\rightarrow$ 数据（区间而不是点值）
$\rightarrow$ 结论边界（哪台机器、哪个负载之外不成立）$\rightarrow$ 复现步骤（命令、产物哈希、原始日志）。
这个顺序本身就是防御：先定口径再取数，先写判据再跑实验，先记指纹再报差异。
判据三档（明显变差、无结论、明显变好）里的``无结论''是合法结论，不是失败。
\end{guideline}

\begin{bestpractice}[title={评审时该问的三个问题}]
一，这个数是什么口径，和标题里的用词一致吗？二，同一份代码重跑还会得到同样的方向和区间吗，
重了几次，变异性多大？三，如果把差异归因到某个改动，除了时间先后还有别的证据吗
（数据布局、依赖关系、分支形态、拓扑与距离）？三个问题里有任何一个答不上，
PR 里那句``性能提升''就还只是叙事。
\end{bestpractice}

\section{与其他分册的接口}
\label{sec:pmapC-crossref}

本册只管怎么证明，不管怎么实现。别的分册里的性能问题进本册时按下面这条路走；
反向的接口（本册的读数解释不了机制时该去查什么）也在同一条里。

\begin{itemize}
  \item 《硬件访问与内核》：那边问的是寄存器、中断、DMA 与设备时钟域。进本册一般走
        第\ref{ch:pm02-clock}章（计数器与换算）和第\ref{ch:pm05-events}章（事件与权限），
        拿到数之后仍要回去确认机制与硬件约束。
  \item 《并行算法》：那边问的是算法能不能并行、并行度怎么定。本册提供加速比曲线、拐点与
        利用率歧义的读法（第\ref{ch:pm09-concurrency}章），但串行比例从哪来是算法层的判断，
        本册替不了。
  \item 《音视频开发》与《嵌入式与硬实时》：那两册关心按时交付，指标是截止期、抖动与丢帧。
        本册的第\ref{ch:pm03-stats}章（分位数与变异性）和第\ref{ch:pm14-production}章
        （线上低开销采样）可以直接用；但硬实时的最坏情况口径不在本册，别把 \Mx{p99} 当它。
  \item 《代码风格参考手册》：\Tp{inline}、\Tp{noinline}、\Tp{always\_{\Bk}inline} 一族的语义以那边为准，
        本册第\ref{ch:pm10-compiler}章只补``怎么在发布产物上确认到底有没有展开''这一面。
  \item 《内存泄漏检测与防御》：ASan、TSan、UBSan 与 heaptrack 一类工具的用法与门禁挂法在那边；
        本册只谈它们作为测量对象时的开销口径（第\ref{ch:pm07-memory}章），
        以及为什么 sanitizer 跑出来的数不能当性能数。
  \item 《LSP 与代码质量工具》：CI 里 clang-tidy 与 sanitizers 的门禁结构参照那边；
        本册第\ref{ch:pm13-ci}章补的是性能门禁特有的部分：基线存法、假阳假阴与阈值纪律。
\end{itemize}

最后一条纪律：本册所有清单里的数字、火焰图形态与工具输出都只是形态示意。
写章与构建环境没有 PMU 权限，通常也没有 \Fn{perf}，任何看起来像实测值的东西都不该被引用；
真正可用的数只能来自你自己的机器与你自己的负载，并按第~\ref{sec:pmapC-recipes}节的模板归档。


\end{document}
